% Exploratory analyses X3-X7 and the readings beside the released boundary.
% Moved from sections/supplement-exploratory-x3-x7.tex of the rule-assembled
% attempt and revised during the journal revision (agent S, 2026-10-13):
% machine labels translated with the original label kept once, the X3
% confound, procedure-specific X4 size flags, both X6 readings, the X6
% precision sensitivity and the X7 features (from the declaration
% README-x3-x7.md and declaration-x3-x7.json) stated. The X5 paragraph moved
% here from the X1-X2 file; its numbers are transcribed from results/x5.json.
% Round 5 (2026-10-14): the free-share reading follows round5.md (one fit; share
% shift, trait weights and lag-three miss from tab:x3-fit); the nested fit is
% described from README-x3-x7.md; the X5 lag signature (j.x5.signature.*), the
% attrition literature and a pointer to the shared realistic channel (S8.4) were added.
% Round 5b: grid shift, a = 0.60 fits, free gaps by strength and the r + 0.01
% trial point (task T11 keys); per-cell stress table tab:x5-lags.
% Round 6: the third-decimal hook is a sensitivity statement with a comparison of
% scales (round6.md item 1); the lag-three miss is matching error (item 4).

\subsection{Free signal share (X3)}\label{sec:x3}\phantomsection\label{supp:x3-x7}

\paragraph{Design.}
X3 lets persistence model B choose its own signal share. In the free-share arm, B's frozen latent law passes through the observed channel at a fitted share $r_B$, matched together with B's two persistence parameters to the lag-one to lag-three moments. The declared fit is nested: at each trial share the frozen profile solver matches B's two persistence parameters to lags one and two, and an outer one-dimensional search, starting at the design share, moves the share until lag three is matched. Because lag three becomes a matched moment, both X3 contrasts train the LDA on the lag-four and lag-five features only, unlike the main discriminator, which uses lags three to five. The known-share contrast uses the frozen two-parameter B at the cell's own $r$ and the free-share contrast uses the fitted $r_B$, both on the same panels with influence; each cell reports the free-minus-known AUC difference with a paired bootstrap interval, read against the declared margin of $\pm\resultnum{v4x.x3.summary.margin}{0.02}$. The two contrasts therefore differ by one fit, which matches lag three by moving the share. The difference measures the effect of that fit for the lag-4--5 LDA and is not a statement about the confirmatory discriminator.
These intervals are conditional on the accepted fits and the trained discriminators and omit the uncertainty of the three-parameter matching; the contrasts do not estimate the frozen boundary.

\paragraph{Readings.}
Fitting B's signal share and matching lag three reduced separability for the specified lag-4--5 LDA in \resultnum{v4x.x3.count.supports-objection}{47} cells, whose paired AUC-difference interval lay entirely below zero (reading ``reduced''; machine label \texttt{supports-objection}; Table~\ref{tab:v4x-x3}).
In \resultnum{v4x.x3.count.against-objection}{3} cells the paired interval lay within the declared margin (``within margin''; \texttt{against-objection}). This is local equivalence for the specified contrast; it does not establish equivalence of mechanisms or of their full laws.
In \resultnum{v4x.x3.count.inconclusive}{4} cells neither criterion was met (``neither''; frozen label \texttt{inconclusive}); the intervals in Table~\ref{tab:v4x-x3} retain the uncertainty under the declared conditional analysis.
There were \resultnum{v4x.x3.count.gap}{3} known-share fit gaps (\texttt{gap}), in which both contrasts remain missing and the known-share attempts are retained, and \resultnum{v4x.x3.count.free-gap}{15} free-share fit gaps (\texttt{free-gap}; \resultnum{j.x3.free_gap.a0010.n}{2} at $a=0.001$, \resultnum{j.x3.free_gap.a0030.n}{1} at 0.003, \resultnum{j.x3.free_gap.a3000.n}{1} at 0.30, \resultnum{j.x3.free_gap.a6000.n}{5} at 0.60 and \resultnum{j.x3.free_gap.a9000.n}{6} at 0.90). Table~\ref{tab:v4x-x3-attempts} reports each recorded attempt, its matching size and the last status; the archived values catalog reports residuals and evaluations. Search-budget exhaustion, a numerical search without an accepted fit and a candidate rejected by validation remain distinct outcomes, and a family-limit or no-root search label alone does not demonstrate that the free-share family cannot reproduce the target moments.
In \resultnum{v4x.x3.count.lost_crossing}{18} cells, the known-share point AUC was at least \resultnum{v4.protocol.target}{0.70} while the free-share interval's upper end was below that target (``lost crossing''; \texttt{lost\_crossing}). This is the declared paired-cell flag, not a crossing of a frozen band or a confirmatory boundary claim. Every flagged cell lies at $a=0.10$--0.60: of the \resultnum{j.x3.lost_crossings.cells}{21} cells in that range in which both fits succeeded, the flag was raised in \resultnum{j.x3.lost_crossings.n}{18}. There, with B's signal share fitted and lag three matched, the lag-4--5 LDA stayed below the target although it reached the target with the share known. The other three cells are those at $a=0.10$ with $r=0.50$ and 0.55, where the known-share AUC was below 0.70, and the cell at $a=0.30$ with $r=0.90$, where the free-share interval reached 0.70. The count is descriptive: it is taken over the cells in which both fits succeeded, without correction across cells, and it is neither a whole-curve inference nor an estimate of the power an analyst would lose. Where the free-share fit succeeded, no free-share point AUC reached 0.70. At $a=0.60$, B's trait weight was exactly zero, the lower edge of its box, in \resultnum{j.x3.free.a6000.lambda_zero.n}{4} of \resultnum{j.x3.free.a6000.fits.n}{4} accepted fits, which reduced B to an autoregression without trait ($\phi=\resultnum{j.x3.free.a6000.phi.min}{0.931}$--\resultnum{j.x3.free.a6000.phi.max}{0.932}), and the free-share AUC was \resultnum{j.x3.free.a6000.auc.min}{0.602}--\resultnum{j.x3.free.a6000.auc.max}{0.649}; no free-share fit succeeded at $a=0.90$.
The archived values catalog holds, for every cell, the known-share and free-share AUCs with their intervals, the fitted parameters, the points of the share search and the mean lag correlations of the panels; Table~\ref{tab:v4x-x3} reports the paired difference and Table~\ref{tab:x3-fit} the fitted parameters at the first-boundary cells.

Read after the results became available: at the cell of the released first boundary against B, the free-share contrast reduced the AUC of the lag-4--5 LDA at all nine shares, by \resultnum{j.x3.auc_reduction.min}{0.126} at $r=\resultnum{j.x3.auc_reduction.min.r}{0.55}$ to \resultnum{j.x3.auc_reduction.max}{0.273} at $r=\resultnum{j.x3.auc_reduction.max.r}{0.50}$ (Table~\ref{tab:v4x-selected-comparisons}).
At those cells the lag-4--5 LDA reached AUCs from \resultnum{j.x3.known.auc.min}{0.668} (at $r=\resultnum{j.x3.known.auc.min.r}{0.55}$) to \resultnum{j.x3.known.auc.max}{0.858} (at $r=\resultnum{j.x3.known.auc.max.r}{0.50}$, where the boundary cell lies at $a=0.30$) with the share known, and only \resultnum{j.x3.free.auc.min}{0.534} (at $r=\resultnum{j.x3.free.auc.min.r}{0.65}$) to \resultnum{j.x3.free.auc.max}{0.585} (at $r=\resultnum{j.x3.free.auc.max.r}{0.50}$) with B's share fitted and lag three matched (Table~\ref{tab:x3-levels}); these are AUCs of the exploratory lag-4--5 LDA, not of the main classifier.
The fit did this by moving B's signal share very little. The fitted share lay below the design share at every boundary cell, by \resultnum{j.x3.free.r_shift.min}{0.0011} at $r=\resultnum{j.x3.free.r_shift.min.r}{0.55}$ to \resultnum{j.x3.free.r_shift.max}{0.0028} at $r=\resultnum{j.x3.free.r_shift.max.r}{0.50}$ (Tables~\ref{tab:x3-levels} and~\ref{tab:x3-fit}). Over the grid it lay below the design share in \resultnum{j.x3.free.r_shift.negative.n}{39} of \resultnum{j.x3.free.fits.n}{39} accepted fits at $a\ge0.01$, by \resultnum{j.x3.free.r_shift.a0100.min}{0.0001}--\resultnum{j.x3.free.r_shift.a0100.max}{0.0004} at $a=0.01$, \resultnum{j.x3.free.r_shift.a0300.min}{0.0003}--\resultnum{j.x3.free.r_shift.a0300.max}{0.0009} at 0.03, \resultnum{j.x3.free.r_shift.a1000.min}{0.0011}--\resultnum{j.x3.free.r_shift.a1000.max}{0.0022} at 0.10, \resultnum{j.x3.free.r_shift.a3000.min}{0.0028}--\resultnum{j.x3.free.r_shift.a3000.max}{0.0051} at 0.30 (\resultnum{j.x3.free.r_shift.a3000.n}{8} fits) and \resultnum{j.x3.free.r_shift.a6000.min}{0.0069}--\resultnum{j.x3.free.r_shift.a6000.max}{0.0100} at 0.60 (\resultnum{j.x3.free.r_shift.a6000.n}{4} fits); at $a\le0.003$ none of the \resultnum{j.x3.free.r_shift.low.n}{15} accepted fits moved it by more than \resultnum{j.x3.free.r_shift.low.max}{0.0002} at four decimals. A shift of one sign that grows with strength, as at $a\ge0.01$, is not rounding noise, although by itself it does not show what the shift represents. At the boundary cells this small shift more than halved B's trait weight, from \resultnum{j.x3.known.lambda.min}{0.139}--\resultnum{j.x3.known.lambda.max}{0.363} with the share known to \resultnum{j.x3.free.lambda.min}{0.061}--\resultnum{j.x3.free.lambda.max}{0.133} with it fitted. At the design share, B matched to lags one and two could not match lag three: its lag-three correlation exceeded the influence arm's by \resultnum{j.x3.l3_miss_se.abs.min}{7.1} to \resultnum{j.x3.l3_miss_se.abs.max}{13.4} standard errors of the matching fit, which averages over many simulated panels; in correlation units the miss is small against the spread of the lag-three correlation in a single panel. In the other direction, at the second point of the share search, \resultnum{j.x3.plus.offset}{0.01} above the design share, B's trait weight was \resultnum{j.x3.plus.lambda.min}{0.376}--\resultnum{j.x3.plus.lambda.max}{0.596} with lags one and two still matched, and its lag-three correlation exceeded the influence arm's by \resultnum{j.x3.plus.l3_miss.min}{0.0069} to \resultnum{j.x3.plus.l3_miss.max}{0.0081} in correlation units, on the matching panels. The share shift and the extra matched lag are one fit whose contributions were not disentangled.
With one answer per wave, a share lower by a small fraction is equivalent, for the continuous response, to the design share with a component carrying that fraction of the latent variance added to the state and renewed every wave, a temporary true component that quasi-simplex models also count as error \cite[p.~183]{saris2014}; the free-share fit therefore cannot tell a slightly misjudged reliability from B with a small short-lived component added to its state. Separation was thus sensitive to the assumed share in its third decimal; for a comparison of scales only, three-wave reliability estimates for the SHP trust item have standard errors of \resultnum{j.shp.heise.w2004.se}{0.011} to \resultnum{j.shp.heise.w2002.se}{0.014} (Section~\ref{sec:shp-results}).

% Generated by v4_exploratory_export; edited during the journal revision
% (agent S, 2026-10-13): captions, headers and categorical labels translated
% into plain English. Every number keeps its \resultnum binding.
\begin{table*}[htbp]
\setlength{\tabcolsep}{2pt}
\caption{Free signal share (X3): AUC with B fitted at its own signal share minus AUC with B at the known share, both for the LDA on lags 4 and 5 and on the same panels with influence, with paired bootstrap 95\% intervals. Reading under the declared rule: reduced, interval entirely below zero (machine label supports-objection); within margin, interval inside $\pm$0.02 (against-objection); neither, any other interval, including an interval entirely above zero (frozen label inconclusive); known-share or free-share fit gap, no contrast (gap, free-gap). Intervals condition on accepted fits and training and omit the uncertainty of the three-parameter matching. The free-share arm also matches lag 3 (lags 1--3 instead of 1--2), so the difference is the joint effect of the fitted share and the extra matched lag, whose separate contributions are unresolved; both arms use lags 4--5 rather than the main classifier's lags 3--5, and the contrast does not estimate the confirmatory boundary.}\label{tab:v4x-x3}
\centering\footnotesize
\begin{tabular}{@{}p{0.090\textwidth}p{0.090\textwidth}p{0.269\textwidth}p{0.169\textwidth}p{0.329\textwidth}@{}}
\toprule
\raggedright $r$ & \raggedright $a$ & \raggedright \mbox{Reading} & \raggedright $\Delta$AUC & \raggedright \mbox{Interval} \tabularnewline
\midrule
\raggedright \resultnum{v4x.x3.cells.r050.a0010.cell.r}{0.5} & \raggedright \resultnum{v4x.x3.cells.r050.a0010.cell.a}{0.001} & \raggedright within margin & \raggedright \resultnum{v4x.x3.cells.r050.a0010.delta.point}{0.00954406} & \raggedright [\resultnum{v4x.x3.cells.r050.a0010.delta.interval_95.i000}{0.0018215}, \resultnum{v4x.x3.cells.r050.a0010.delta.interval_95.i001}{0.0185533}] \tabularnewline
\raggedright \resultnum{v4x.x3.cells.r050.a0030.cell.r}{0.5} & \raggedright \resultnum{v4x.x3.cells.r050.a0030.cell.a}{0.003} & \raggedright within margin & \raggedright \resultnum{v4x.x3.cells.r050.a0030.delta.point}{0} & \raggedright [\resultnum{v4x.x3.cells.r050.a0030.delta.interval_95.i000}{0}, \resultnum{v4x.x3.cells.r050.a0030.delta.interval_95.i001}{0}] \tabularnewline
\raggedright \resultnum{v4x.x3.cells.r050.a0100.cell.r}{0.5} & \raggedright \resultnum{v4x.x3.cells.r050.a0100.cell.a}{0.01} & \raggedright reduced & \raggedright $\resultnum{v4x.x3.cells.r050.a0100.delta.point}{-0.0197241}$ & \raggedright [$\resultnum{v4x.x3.cells.r050.a0100.delta.interval_95.i000}{-0.0258999}$, $\resultnum{v4x.x3.cells.r050.a0100.delta.interval_95.i001}{-0.0136672}$] \tabularnewline
\raggedright \resultnum{v4x.x3.cells.r050.a0300.cell.r}{0.5} & \raggedright \resultnum{v4x.x3.cells.r050.a0300.cell.a}{0.03} & \raggedright reduced & \raggedright $\resultnum{v4x.x3.cells.r050.a0300.delta.point}{-0.0636435}$ & \raggedright [$\resultnum{v4x.x3.cells.r050.a0300.delta.interval_95.i000}{-0.0718572}$, $\resultnum{v4x.x3.cells.r050.a0300.delta.interval_95.i001}{-0.0553668}$] \tabularnewline
\raggedright \resultnum{v4x.x3.cells.r050.a1000.cell.r}{0.5} & \raggedright \resultnum{v4x.x3.cells.r050.a1000.cell.a}{0.1} & \raggedright reduced & \raggedright $\resultnum{v4x.x3.cells.r050.a1000.delta.point}{-0.11445}$ & \raggedright [$\resultnum{v4x.x3.cells.r050.a1000.delta.interval_95.i000}{-0.119276}$, $\resultnum{v4x.x3.cells.r050.a1000.delta.interval_95.i001}{-0.109896}$] \tabularnewline
\raggedright \resultnum{v4x.x3.cells.r050.a3000.cell.r}{0.5} & \raggedright \resultnum{v4x.x3.cells.r050.a3000.cell.a}{0.3} & \raggedright reduced & \raggedright $\resultnum{v4x.x3.cells.r050.a3000.delta.point}{-0.27282}$ & \raggedright [$\resultnum{v4x.x3.cells.r050.a3000.delta.interval_95.i000}{-0.280625}$, $\resultnum{v4x.x3.cells.r050.a3000.delta.interval_95.i001}{-0.264818}$] \tabularnewline
\raggedright \resultnum{v4x.x3.cells.r050.a6000.cell.r}{0.5} & \raggedright \resultnum{v4x.x3.cells.r050.a6000.cell.a}{0.6} & \raggedright free-share fit gap & \raggedright not estimable & \raggedright (fit gap) \tabularnewline
\raggedright \resultnum{v4x.x3.cells.r050.a9000.cell.r}{0.5} & \raggedright \resultnum{v4x.x3.cells.r050.a9000.cell.a}{0.9} & \raggedright known-share fit gap & \raggedright not estimable & \raggedright (fit gap) \tabularnewline
\raggedright \resultnum{v4x.x3.cells.r055.a0010.cell.r}{0.55} & \raggedright \resultnum{v4x.x3.cells.r055.a0010.cell.a}{0.001} & \raggedright free-share fit gap & \raggedright not estimable & \raggedright (fit gap) \tabularnewline
\raggedright \resultnum{v4x.x3.cells.r055.a0030.cell.r}{0.55} & \raggedright \resultnum{v4x.x3.cells.r055.a0030.cell.a}{0.003} & \raggedright free-share fit gap & \raggedright not estimable & \raggedright (fit gap) \tabularnewline
\raggedright \resultnum{v4x.x3.cells.r055.a0100.cell.r}{0.55} & \raggedright \resultnum{v4x.x3.cells.r055.a0100.cell.a}{0.01} & \raggedright reduced & \raggedright $\resultnum{v4x.x3.cells.r055.a0100.delta.point}{-0.0136261}$ & \raggedright [$\resultnum{v4x.x3.cells.r055.a0100.delta.interval_95.i000}{-0.019653}$, $\resultnum{v4x.x3.cells.r055.a0100.delta.interval_95.i001}{-0.00689803}$] \tabularnewline
\raggedright \resultnum{v4x.x3.cells.r055.a0300.cell.r}{0.55} & \raggedright \resultnum{v4x.x3.cells.r055.a0300.cell.a}{0.03} & \raggedright reduced & \raggedright $\resultnum{v4x.x3.cells.r055.a0300.delta.point}{-0.0340532}$ & \raggedright [$\resultnum{v4x.x3.cells.r055.a0300.delta.interval_95.i000}{-0.0418134}$, $\resultnum{v4x.x3.cells.r055.a0300.delta.interval_95.i001}{-0.0265757}$] \tabularnewline
\raggedright \resultnum{v4x.x3.cells.r055.a1000.cell.r}{0.55} & \raggedright \resultnum{v4x.x3.cells.r055.a1000.cell.a}{0.1} & \raggedright reduced & \raggedright $\resultnum{v4x.x3.cells.r055.a1000.delta.point}{-0.12614}$ & \raggedright [$\resultnum{v4x.x3.cells.r055.a1000.delta.interval_95.i000}{-0.134898}$, $\resultnum{v4x.x3.cells.r055.a1000.delta.interval_95.i001}{-0.117152}$] \tabularnewline
\raggedright \resultnum{v4x.x3.cells.r055.a3000.cell.r}{0.55} & \raggedright \resultnum{v4x.x3.cells.r055.a3000.cell.a}{0.3} & \raggedright free-share fit gap & \raggedright not estimable & \raggedright (fit gap) \tabularnewline
\raggedright \resultnum{v4x.x3.cells.r055.a6000.cell.r}{0.55} & \raggedright \resultnum{v4x.x3.cells.r055.a6000.cell.a}{0.6} & \raggedright free-share fit gap & \raggedright not estimable & \raggedright (fit gap) \tabularnewline
\raggedright \resultnum{v4x.x3.cells.r055.a9000.cell.r}{0.55} & \raggedright \resultnum{v4x.x3.cells.r055.a9000.cell.a}{0.9} & \raggedright known-share fit gap & \raggedright not estimable & \raggedright (fit gap) \tabularnewline
\raggedright \resultnum{v4x.x3.cells.r060.a0010.cell.r}{0.6} & \raggedright \resultnum{v4x.x3.cells.r060.a0010.cell.a}{0.001} & \raggedright free-share fit gap & \raggedright not estimable & \raggedright (fit gap) \tabularnewline
\raggedright \resultnum{v4x.x3.cells.r060.a0030.cell.r}{0.6} & \raggedright \resultnum{v4x.x3.cells.r060.a0030.cell.a}{0.003} & \raggedright reduced & \raggedright $\resultnum{v4x.x3.cells.r060.a0030.delta.point}{-0.00589237}$ & \raggedright [$\resultnum{v4x.x3.cells.r060.a0030.delta.interval_95.i000}{-0.00962166}$, $\resultnum{v4x.x3.cells.r060.a0030.delta.interval_95.i001}{-0.00252203}$] \tabularnewline
\raggedright \resultnum{v4x.x3.cells.r060.a0100.cell.r}{0.6} & \raggedright \resultnum{v4x.x3.cells.r060.a0100.cell.a}{0.01} & \raggedright reduced & \raggedright $\resultnum{v4x.x3.cells.r060.a0100.delta.point}{-0.0153227}$ & \raggedright [$\resultnum{v4x.x3.cells.r060.a0100.delta.interval_95.i000}{-0.0201444}$, $\resultnum{v4x.x3.cells.r060.a0100.delta.interval_95.i001}{-0.0109601}$] \tabularnewline
\raggedright \resultnum{v4x.x3.cells.r060.a0300.cell.r}{0.6} & \raggedright \resultnum{v4x.x3.cells.r060.a0300.cell.a}{0.03} & \raggedright reduced & \raggedright $\resultnum{v4x.x3.cells.r060.a0300.delta.point}{-0.0707461}$ & \raggedright [$\resultnum{v4x.x3.cells.r060.a0300.delta.interval_95.i000}{-0.078441}$, $\resultnum{v4x.x3.cells.r060.a0300.delta.interval_95.i001}{-0.0636868}$] \tabularnewline
\raggedright \resultnum{v4x.x3.cells.r060.a1000.cell.r}{0.6} & \raggedright \resultnum{v4x.x3.cells.r060.a1000.cell.a}{0.1} & \raggedright reduced & \raggedright $\resultnum{v4x.x3.cells.r060.a1000.delta.point}{-0.159901}$ & \raggedright [$\resultnum{v4x.x3.cells.r060.a1000.delta.interval_95.i000}{-0.165875}$, $\resultnum{v4x.x3.cells.r060.a1000.delta.interval_95.i001}{-0.153723}$] \tabularnewline
\raggedright \resultnum{v4x.x3.cells.r060.a3000.cell.r}{0.6} & \raggedright \resultnum{v4x.x3.cells.r060.a3000.cell.a}{0.3} & \raggedright reduced & \raggedright $\resultnum{v4x.x3.cells.r060.a3000.delta.point}{-0.325963}$ & \raggedright [$\resultnum{v4x.x3.cells.r060.a3000.delta.interval_95.i000}{-0.335442}$, $\resultnum{v4x.x3.cells.r060.a3000.delta.interval_95.i001}{-0.317165}$] \tabularnewline
\raggedright \resultnum{v4x.x3.cells.r060.a6000.cell.r}{0.6} & \raggedright \resultnum{v4x.x3.cells.r060.a6000.cell.a}{0.6} & \raggedright free-share fit gap & \raggedright not estimable & \raggedright (fit gap) \tabularnewline
\raggedright \resultnum{v4x.x3.cells.r060.a9000.cell.r}{0.6} & \raggedright \resultnum{v4x.x3.cells.r060.a9000.cell.a}{0.9} & \raggedright free-share fit gap & \raggedright not estimable & \raggedright (fit gap) \tabularnewline
\raggedright \resultnum{v4x.x3.cells.r065.a0010.cell.r}{0.65} & \raggedright \resultnum{v4x.x3.cells.r065.a0010.cell.a}{0.001} & \raggedright neither & \raggedright $\resultnum{v4x.x3.cells.r065.a0010.delta.point}{-0.00602244}$ & \raggedright [$\resultnum{v4x.x3.cells.r065.a0010.delta.interval_95.i000}{-0.0273767}$, \resultnum{v4x.x3.cells.r065.a0010.delta.interval_95.i001}{0.0146704}] \tabularnewline
\raggedright \resultnum{v4x.x3.cells.r065.a0030.cell.r}{0.65} & \raggedright \resultnum{v4x.x3.cells.r065.a0030.cell.a}{0.003} & \raggedright neither & \raggedright $\resultnum{v4x.x3.cells.r065.a0030.delta.point}{-0.0122469}$ & \raggedright [$\resultnum{v4x.x3.cells.r065.a0030.delta.interval_95.i000}{-0.0359192}$, \resultnum{v4x.x3.cells.r065.a0030.delta.interval_95.i001}{0.00987345}] \tabularnewline
\raggedright \resultnum{v4x.x3.cells.r065.a0100.cell.r}{0.65} & \raggedright \resultnum{v4x.x3.cells.r065.a0100.cell.a}{0.01} & \raggedright reduced & \raggedright $\resultnum{v4x.x3.cells.r065.a0100.delta.point}{-0.0337366}$ & \raggedright [$\resultnum{v4x.x3.cells.r065.a0100.delta.interval_95.i000}{-0.0589632}$, $\resultnum{v4x.x3.cells.r065.a0100.delta.interval_95.i001}{-0.0101436}$] \tabularnewline
\raggedright \resultnum{v4x.x3.cells.r065.a0300.cell.r}{0.65} & \raggedright \resultnum{v4x.x3.cells.r065.a0300.cell.a}{0.03} & \raggedright reduced & \raggedright $\resultnum{v4x.x3.cells.r065.a0300.delta.point}{-0.0741078}$ & \raggedright [$\resultnum{v4x.x3.cells.r065.a0300.delta.interval_95.i000}{-0.0767876}$, $\resultnum{v4x.x3.cells.r065.a0300.delta.interval_95.i001}{-0.0715649}$] \tabularnewline
\raggedright \resultnum{v4x.x3.cells.r065.a1000.cell.r}{0.65} & \raggedright \resultnum{v4x.x3.cells.r065.a1000.cell.a}{0.1} & \raggedright reduced & \raggedright $\resultnum{v4x.x3.cells.r065.a1000.delta.point}{-0.171633}$ & \raggedright [$\resultnum{v4x.x3.cells.r065.a1000.delta.interval_95.i000}{-0.177105}$, $\resultnum{v4x.x3.cells.r065.a1000.delta.interval_95.i001}{-0.166476}$] \tabularnewline
\raggedright \resultnum{v4x.x3.cells.r065.a3000.cell.r}{0.65} & \raggedright \resultnum{v4x.x3.cells.r065.a3000.cell.a}{0.3} & \raggedright reduced & \raggedright $\resultnum{v4x.x3.cells.r065.a3000.delta.point}{-0.299154}$ & \raggedright [$\resultnum{v4x.x3.cells.r065.a3000.delta.interval_95.i000}{-0.307842}$, $\resultnum{v4x.x3.cells.r065.a3000.delta.interval_95.i001}{-0.290462}$] \tabularnewline
\raggedright \resultnum{v4x.x3.cells.r065.a6000.cell.r}{0.65} & \raggedright \resultnum{v4x.x3.cells.r065.a6000.cell.a}{0.6} & \raggedright reduced & \raggedright $\resultnum{v4x.x3.cells.r065.a6000.delta.point}{-0.388224}$ & \raggedright [$\resultnum{v4x.x3.cells.r065.a6000.delta.interval_95.i000}{-0.400112}$, $\resultnum{v4x.x3.cells.r065.a6000.delta.interval_95.i001}{-0.376105}$] \tabularnewline
\raggedright \resultnum{v4x.x3.cells.r065.a9000.cell.r}{0.65} & \raggedright \resultnum{v4x.x3.cells.r065.a9000.cell.a}{0.9} & \raggedright free-share fit gap & \raggedright not estimable & \raggedright (fit gap) \tabularnewline
\raggedright \resultnum{v4x.x3.cells.r070.a0010.cell.r}{0.7} & \raggedright \resultnum{v4x.x3.cells.r070.a0010.cell.a}{0.001} & \raggedright reduced & \raggedright $\resultnum{v4x.x3.cells.r070.a0010.delta.point}{-0.0250951}$ & \raggedright [$\resultnum{v4x.x3.cells.r070.a0010.delta.interval_95.i000}{-0.030417}$, $\resultnum{v4x.x3.cells.r070.a0010.delta.interval_95.i001}{-0.0197927}$] \tabularnewline
\raggedright \resultnum{v4x.x3.cells.r070.a0030.cell.r}{0.7} & \raggedright \resultnum{v4x.x3.cells.r070.a0030.cell.a}{0.003} & \raggedright reduced & \raggedright $\resultnum{v4x.x3.cells.r070.a0030.delta.point}{-0.025215}$ & \raggedright [$\resultnum{v4x.x3.cells.r070.a0030.delta.interval_95.i000}{-0.0263931}$, $\resultnum{v4x.x3.cells.r070.a0030.delta.interval_95.i001}{-0.0238922}$] \tabularnewline
\raggedright \resultnum{v4x.x3.cells.r070.a0100.cell.r}{0.7} & \raggedright \resultnum{v4x.x3.cells.r070.a0100.cell.a}{0.01} & \raggedright reduced & \raggedright $\resultnum{v4x.x3.cells.r070.a0100.delta.point}{-0.0432194}$ & \raggedright [$\resultnum{v4x.x3.cells.r070.a0100.delta.interval_95.i000}{-0.0448747}$, $\resultnum{v4x.x3.cells.r070.a0100.delta.interval_95.i001}{-0.0416101}$] \tabularnewline
\raggedright \resultnum{v4x.x3.cells.r070.a0300.cell.r}{0.7} & \raggedright \resultnum{v4x.x3.cells.r070.a0300.cell.a}{0.03} & \raggedright reduced & \raggedright $\resultnum{v4x.x3.cells.r070.a0300.delta.point}{-0.0710723}$ & \raggedright [$\resultnum{v4x.x3.cells.r070.a0300.delta.interval_95.i000}{-0.0739473}$, $\resultnum{v4x.x3.cells.r070.a0300.delta.interval_95.i001}{-0.0683929}$] \tabularnewline
\raggedright \resultnum{v4x.x3.cells.r070.a1000.cell.r}{0.7} & \raggedright \resultnum{v4x.x3.cells.r070.a1000.cell.a}{0.1} & \raggedright reduced & \raggedright $\resultnum{v4x.x3.cells.r070.a1000.delta.point}{-0.178288}$ & \raggedright [$\resultnum{v4x.x3.cells.r070.a1000.delta.interval_95.i000}{-0.183328}$, $\resultnum{v4x.x3.cells.r070.a1000.delta.interval_95.i001}{-0.173138}$] \tabularnewline
\raggedright \resultnum{v4x.x3.cells.r070.a3000.cell.r}{0.7} & \raggedright \resultnum{v4x.x3.cells.r070.a3000.cell.a}{0.3} & \raggedright reduced & \raggedright $\resultnum{v4x.x3.cells.r070.a3000.delta.point}{-0.322506}$ & \raggedright [$\resultnum{v4x.x3.cells.r070.a3000.delta.interval_95.i000}{-0.331696}$, $\resultnum{v4x.x3.cells.r070.a3000.delta.interval_95.i001}{-0.312521}$] \tabularnewline
\raggedright \resultnum{v4x.x3.cells.r070.a6000.cell.r}{0.7} & \raggedright \resultnum{v4x.x3.cells.r070.a6000.cell.a}{0.6} & \raggedright free-share fit gap & \raggedright not estimable & \raggedright (fit gap) \tabularnewline
\raggedright \resultnum{v4x.x3.cells.r070.a9000.cell.r}{0.7} & \raggedright \resultnum{v4x.x3.cells.r070.a9000.cell.a}{0.9} & \raggedright known-share fit gap & \raggedright not estimable & \raggedright (fit gap) \tabularnewline
\raggedright \resultnum{v4x.x3.cells.r075.a0010.cell.r}{0.75} & \raggedright \resultnum{v4x.x3.cells.r075.a0010.cell.a}{0.001} & \raggedright reduced & \raggedright $\resultnum{v4x.x3.cells.r075.a0010.delta.point}{-0.00756713}$ & \raggedright [$\resultnum{v4x.x3.cells.r075.a0010.delta.interval_95.i000}{-0.00851919}$, $\resultnum{v4x.x3.cells.r075.a0010.delta.interval_95.i001}{-0.00654083}$] \tabularnewline
\raggedright \resultnum{v4x.x3.cells.r075.a0030.cell.r}{0.75} & \raggedright \resultnum{v4x.x3.cells.r075.a0030.cell.a}{0.003} & \raggedright reduced & \raggedright $\resultnum{v4x.x3.cells.r075.a0030.delta.point}{-0.00572125}$ & \raggedright [$\resultnum{v4x.x3.cells.r075.a0030.delta.interval_95.i000}{-0.00630034}$, $\resultnum{v4x.x3.cells.r075.a0030.delta.interval_95.i001}{-0.00509347}$] \tabularnewline
\raggedright \resultnum{v4x.x3.cells.r075.a0100.cell.r}{0.75} & \raggedright \resultnum{v4x.x3.cells.r075.a0100.cell.a}{0.01} & \raggedright reduced & \raggedright $\resultnum{v4x.x3.cells.r075.a0100.delta.point}{-0.040827}$ & \raggedright [$\resultnum{v4x.x3.cells.r075.a0100.delta.interval_95.i000}{-0.0422523}$, $\resultnum{v4x.x3.cells.r075.a0100.delta.interval_95.i001}{-0.0392359}$] \tabularnewline
\raggedright \resultnum{v4x.x3.cells.r075.a0300.cell.r}{0.75} & \raggedright \resultnum{v4x.x3.cells.r075.a0300.cell.a}{0.03} & \raggedright reduced & \raggedright $\resultnum{v4x.x3.cells.r075.a0300.delta.point}{-0.0729784}$ & \raggedright [$\resultnum{v4x.x3.cells.r075.a0300.delta.interval_95.i000}{-0.076011}$, $\resultnum{v4x.x3.cells.r075.a0300.delta.interval_95.i001}{-0.0699015}$] \tabularnewline
\raggedright \resultnum{v4x.x3.cells.r075.a1000.cell.r}{0.75} & \raggedright \resultnum{v4x.x3.cells.r075.a1000.cell.a}{0.1} & \raggedright reduced & \raggedright $\resultnum{v4x.x3.cells.r075.a1000.delta.point}{-0.172387}$ & \raggedright [$\resultnum{v4x.x3.cells.r075.a1000.delta.interval_95.i000}{-0.177462}$, $\resultnum{v4x.x3.cells.r075.a1000.delta.interval_95.i001}{-0.167671}$] \tabularnewline
\raggedright \resultnum{v4x.x3.cells.r075.a3000.cell.r}{0.75} & \raggedright \resultnum{v4x.x3.cells.r075.a3000.cell.a}{0.3} & \raggedright reduced & \raggedright $\resultnum{v4x.x3.cells.r075.a3000.delta.point}{-0.308323}$ & \raggedright [$\resultnum{v4x.x3.cells.r075.a3000.delta.interval_95.i000}{-0.317558}$, $\resultnum{v4x.x3.cells.r075.a3000.delta.interval_95.i001}{-0.299006}$] \tabularnewline
\raggedright \resultnum{v4x.x3.cells.r075.a6000.cell.r}{0.75} & \raggedright \resultnum{v4x.x3.cells.r075.a6000.cell.a}{0.6} & \raggedright reduced & \raggedright $\resultnum{v4x.x3.cells.r075.a6000.delta.point}{-0.369447}$ & \raggedright [$\resultnum{v4x.x3.cells.r075.a6000.delta.interval_95.i000}{-0.381423}$, $\resultnum{v4x.x3.cells.r075.a6000.delta.interval_95.i001}{-0.357629}$] \tabularnewline
\raggedright \resultnum{v4x.x3.cells.r075.a9000.cell.r}{0.75} & \raggedright \resultnum{v4x.x3.cells.r075.a9000.cell.a}{0.9} & \raggedright free-share fit gap & \raggedright not estimable & \raggedright (fit gap) \tabularnewline
\raggedright \resultnum{v4x.x3.cells.r080.a0010.cell.r}{0.8} & \raggedright \resultnum{v4x.x3.cells.r080.a0010.cell.a}{0.001} & \raggedright reduced & \raggedright $\resultnum{v4x.x3.cells.r080.a0010.delta.point}{-0.000953875}$ & \raggedright [$\resultnum{v4x.x3.cells.r080.a0010.delta.interval_95.i000}{-0.00144369}$, $\resultnum{v4x.x3.cells.r080.a0010.delta.interval_95.i001}{-0.000473233}$] \tabularnewline
\raggedright \resultnum{v4x.x3.cells.r080.a0030.cell.r}{0.8} & \raggedright \resultnum{v4x.x3.cells.r080.a0030.cell.a}{0.003} & \raggedright reduced & \raggedright $\resultnum{v4x.x3.cells.r080.a0030.delta.point}{-0.00365556}$ & \raggedright [$\resultnum{v4x.x3.cells.r080.a0030.delta.interval_95.i000}{-0.00402726}$, $\resultnum{v4x.x3.cells.r080.a0030.delta.interval_95.i001}{-0.00323865}$] \tabularnewline
\raggedright \resultnum{v4x.x3.cells.r080.a0100.cell.r}{0.8} & \raggedright \resultnum{v4x.x3.cells.r080.a0100.cell.a}{0.01} & \raggedright reduced & \raggedright $\resultnum{v4x.x3.cells.r080.a0100.delta.point}{-0.0207484}$ & \raggedright [$\resultnum{v4x.x3.cells.r080.a0100.delta.interval_95.i000}{-0.0217408}$, $\resultnum{v4x.x3.cells.r080.a0100.delta.interval_95.i001}{-0.0197012}$] \tabularnewline
\raggedright \resultnum{v4x.x3.cells.r080.a0300.cell.r}{0.8} & \raggedright \resultnum{v4x.x3.cells.r080.a0300.cell.a}{0.03} & \raggedright reduced & \raggedright $\resultnum{v4x.x3.cells.r080.a0300.delta.point}{-0.0801886}$ & \raggedright [$\resultnum{v4x.x3.cells.r080.a0300.delta.interval_95.i000}{-0.0823959}$, $\resultnum{v4x.x3.cells.r080.a0300.delta.interval_95.i001}{-0.0779875}$] \tabularnewline
\botrule
\end{tabular}
\end{table*}
\begin{table*}[htbp]
\setlength{\tabcolsep}{2pt}
\addtocounter{table}{-\csname @ne\endcsname}
\providecommand{\theHtable}{\thetable}
\renewcommand{\theHtable}{v4x-x3-part1}
\caption{Free signal share (X3) (continued, part 2 of 2)}\label{tab:v4x-x3-part1}
\centering\footnotesize
\begin{tabular}{@{}p{0.090\textwidth}p{0.090\textwidth}p{0.269\textwidth}p{0.169\textwidth}p{0.329\textwidth}@{}}
\toprule
\raggedright $r$ & \raggedright $a$ & \raggedright \mbox{Reading} & \raggedright $\Delta$AUC & \raggedright \mbox{Interval} \tabularnewline
\midrule
\raggedright \resultnum{v4x.x3.cells.r080.a1000.cell.r}{0.8} & \raggedright \resultnum{v4x.x3.cells.r080.a1000.cell.a}{0.1} & \raggedright reduced & \raggedright $\resultnum{v4x.x3.cells.r080.a1000.delta.point}{-0.173259}$ & \raggedright [$\resultnum{v4x.x3.cells.r080.a1000.delta.interval_95.i000}{-0.17741}$, $\resultnum{v4x.x3.cells.r080.a1000.delta.interval_95.i001}{-0.169218}$] \tabularnewline
\raggedright \resultnum{v4x.x3.cells.r080.a3000.cell.r}{0.8} & \raggedright \resultnum{v4x.x3.cells.r080.a3000.cell.a}{0.3} & \raggedright reduced & \raggedright $\resultnum{v4x.x3.cells.r080.a3000.delta.point}{-0.301138}$ & \raggedright [$\resultnum{v4x.x3.cells.r080.a3000.delta.interval_95.i000}{-0.310852}$, $\resultnum{v4x.x3.cells.r080.a3000.delta.interval_95.i001}{-0.292093}$] \tabularnewline
\raggedright \resultnum{v4x.x3.cells.r080.a6000.cell.r}{0.8} & \raggedright \resultnum{v4x.x3.cells.r080.a6000.cell.a}{0.6} & \raggedright reduced & \raggedright $\resultnum{v4x.x3.cells.r080.a6000.delta.point}{-0.370742}$ & \raggedright [$\resultnum{v4x.x3.cells.r080.a6000.delta.interval_95.i000}{-0.382766}$, $\resultnum{v4x.x3.cells.r080.a6000.delta.interval_95.i001}{-0.35886}$] \tabularnewline
\raggedright \resultnum{v4x.x3.cells.r080.a9000.cell.r}{0.8} & \raggedright \resultnum{v4x.x3.cells.r080.a9000.cell.a}{0.9} & \raggedright free-share fit gap & \raggedright not estimable & \raggedright (fit gap) \tabularnewline
\raggedright \resultnum{v4x.x3.cells.r085.a0010.cell.r}{0.85} & \raggedright \resultnum{v4x.x3.cells.r085.a0010.cell.a}{0.001} & \raggedright neither & \raggedright $\resultnum{v4x.x3.cells.r085.a0010.delta.point}{-0.00937225}$ & \raggedright [$\resultnum{v4x.x3.cells.r085.a0010.delta.interval_95.i000}{-0.0261625}$, \resultnum{v4x.x3.cells.r085.a0010.delta.interval_95.i001}{0.00806083}] \tabularnewline
\raggedright \resultnum{v4x.x3.cells.r085.a0030.cell.r}{0.85} & \raggedright \resultnum{v4x.x3.cells.r085.a0030.cell.a}{0.003} & \raggedright neither & \raggedright $\resultnum{v4x.x3.cells.r085.a0030.delta.point}{-0.0169097}$ & \raggedright [$\resultnum{v4x.x3.cells.r085.a0030.delta.interval_95.i000}{-0.0336468}$, \resultnum{v4x.x3.cells.r085.a0030.delta.interval_95.i001}{0.00116805}] \tabularnewline
\raggedright \resultnum{v4x.x3.cells.r085.a0100.cell.r}{0.85} & \raggedright \resultnum{v4x.x3.cells.r085.a0100.cell.a}{0.01} & \raggedright reduced & \raggedright $\resultnum{v4x.x3.cells.r085.a0100.delta.point}{-0.0416128}$ & \raggedright [$\resultnum{v4x.x3.cells.r085.a0100.delta.interval_95.i000}{-0.0453075}$, $\resultnum{v4x.x3.cells.r085.a0100.delta.interval_95.i001}{-0.0375678}$] \tabularnewline
\raggedright \resultnum{v4x.x3.cells.r085.a0300.cell.r}{0.85} & \raggedright \resultnum{v4x.x3.cells.r085.a0300.cell.a}{0.03} & \raggedright reduced & \raggedright $\resultnum{v4x.x3.cells.r085.a0300.delta.point}{-0.0816274}$ & \raggedright [$\resultnum{v4x.x3.cells.r085.a0300.delta.interval_95.i000}{-0.0838939}$, $\resultnum{v4x.x3.cells.r085.a0300.delta.interval_95.i001}{-0.0794306}$] \tabularnewline
\raggedright \resultnum{v4x.x3.cells.r085.a1000.cell.r}{0.85} & \raggedright \resultnum{v4x.x3.cells.r085.a1000.cell.a}{0.1} & \raggedright reduced & \raggedright $\resultnum{v4x.x3.cells.r085.a1000.delta.point}{-0.194977}$ & \raggedright [$\resultnum{v4x.x3.cells.r085.a1000.delta.interval_95.i000}{-0.200278}$, $\resultnum{v4x.x3.cells.r085.a1000.delta.interval_95.i001}{-0.189724}$] \tabularnewline
\raggedright \resultnum{v4x.x3.cells.r085.a3000.cell.r}{0.85} & \raggedright \resultnum{v4x.x3.cells.r085.a3000.cell.a}{0.3} & \raggedright reduced & \raggedright $\resultnum{v4x.x3.cells.r085.a3000.delta.point}{-0.296039}$ & \raggedright [$\resultnum{v4x.x3.cells.r085.a3000.delta.interval_95.i000}{-0.305678}$, $\resultnum{v4x.x3.cells.r085.a3000.delta.interval_95.i001}{-0.285938}$] \tabularnewline
\raggedright \resultnum{v4x.x3.cells.r085.a6000.cell.r}{0.85} & \raggedright \resultnum{v4x.x3.cells.r085.a6000.cell.a}{0.6} & \raggedright reduced & \raggedright $\resultnum{v4x.x3.cells.r085.a6000.delta.point}{-0.350819}$ & \raggedright [$\resultnum{v4x.x3.cells.r085.a6000.delta.interval_95.i000}{-0.362525}$, $\resultnum{v4x.x3.cells.r085.a6000.delta.interval_95.i001}{-0.338954}$] \tabularnewline
\raggedright \resultnum{v4x.x3.cells.r085.a9000.cell.r}{0.85} & \raggedright \resultnum{v4x.x3.cells.r085.a9000.cell.a}{0.9} & \raggedright free-share fit gap & \raggedright not estimable & \raggedright (fit gap) \tabularnewline
\raggedright \resultnum{v4x.x3.cells.r090.a0010.cell.r}{0.9} & \raggedright \resultnum{v4x.x3.cells.r090.a0010.cell.a}{0.001} & \raggedright within margin & \raggedright \resultnum{v4x.x3.cells.r090.a0010.delta.point}{0} & \raggedright [\resultnum{v4x.x3.cells.r090.a0010.delta.interval_95.i000}{0}, \resultnum{v4x.x3.cells.r090.a0010.delta.interval_95.i001}{0}] \tabularnewline
\raggedright \resultnum{v4x.x3.cells.r090.a0030.cell.r}{0.9} & \raggedright \resultnum{v4x.x3.cells.r090.a0030.cell.a}{0.003} & \raggedright reduced & \raggedright $\resultnum{v4x.x3.cells.r090.a0030.delta.point}{-0.00968181}$ & \raggedright [$\resultnum{v4x.x3.cells.r090.a0030.delta.interval_95.i000}{-0.0127321}$, $\resultnum{v4x.x3.cells.r090.a0030.delta.interval_95.i001}{-0.00684528}$] \tabularnewline
\raggedright \resultnum{v4x.x3.cells.r090.a0100.cell.r}{0.9} & \raggedright \resultnum{v4x.x3.cells.r090.a0100.cell.a}{0.01} & \raggedright reduced & \raggedright $\resultnum{v4x.x3.cells.r090.a0100.delta.point}{-0.0319049}$ & \raggedright [$\resultnum{v4x.x3.cells.r090.a0100.delta.interval_95.i000}{-0.0359587}$, $\resultnum{v4x.x3.cells.r090.a0100.delta.interval_95.i001}{-0.027546}$] \tabularnewline
\raggedright \resultnum{v4x.x3.cells.r090.a0300.cell.r}{0.9} & \raggedright \resultnum{v4x.x3.cells.r090.a0300.cell.a}{0.03} & \raggedright reduced & \raggedright $\resultnum{v4x.x3.cells.r090.a0300.delta.point}{-0.0849248}$ & \raggedright [$\resultnum{v4x.x3.cells.r090.a0300.delta.interval_95.i000}{-0.0896719}$, $\resultnum{v4x.x3.cells.r090.a0300.delta.interval_95.i001}{-0.0799967}$] \tabularnewline
\raggedright \resultnum{v4x.x3.cells.r090.a1000.cell.r}{0.9} & \raggedright \resultnum{v4x.x3.cells.r090.a1000.cell.a}{0.1} & \raggedright reduced & \raggedright $\resultnum{v4x.x3.cells.r090.a1000.delta.point}{-0.201556}$ & \raggedright [$\resultnum{v4x.x3.cells.r090.a1000.delta.interval_95.i000}{-0.207505}$, $\resultnum{v4x.x3.cells.r090.a1000.delta.interval_95.i001}{-0.194782}$] \tabularnewline
\raggedright \resultnum{v4x.x3.cells.r090.a3000.cell.r}{0.9} & \raggedright \resultnum{v4x.x3.cells.r090.a3000.cell.a}{0.3} & \raggedright reduced & \raggedright $\resultnum{v4x.x3.cells.r090.a3000.delta.point}{-0.271204}$ & \raggedright [$\resultnum{v4x.x3.cells.r090.a3000.delta.interval_95.i000}{-0.280612}$, $\resultnum{v4x.x3.cells.r090.a3000.delta.interval_95.i001}{-0.261419}$] \tabularnewline
\raggedright \resultnum{v4x.x3.cells.r090.a6000.cell.r}{0.9} & \raggedright \resultnum{v4x.x3.cells.r090.a6000.cell.a}{0.6} & \raggedright free-share fit gap & \raggedright not estimable & \raggedright (fit gap) \tabularnewline
\raggedright \resultnum{v4x.x3.cells.r090.a9000.cell.r}{0.9} & \raggedright \resultnum{v4x.x3.cells.r090.a9000.cell.a}{0.9} & \raggedright free-share fit gap & \raggedright not estimable & \raggedright (fit gap) \tabularnewline
\botrule
\end{tabular}
\end{table*}

\begin{table}[tp]
\caption{Known and fitted signal-share comparisons at the first boundary.}\label{tab:x3-levels}
\centering\scriptsize
\setlength{\tabcolsep}{2pt}
\begin{tabular}{l@{\quad}l@{\quad}>{\raggedright\arraybackslash}p{16mm}>{\raggedright\arraybackslash}p{16mm}>{\raggedright\arraybackslash}p{38mm}>{\raggedright\arraybackslash}p{15mm}>{\raggedright\arraybackslash}p{12mm}}
\toprule
Design $r$ & $a^*$ & Known-share AUC & Free-share AUC & Difference [interval] & $r_B-r$ & Lost crossing \\
\midrule
\resultnum{j.b.r050.r}{0.50} & \resultnum{j.b.r050.first.point}{0.30} & \resultnum{j.x3.r050.known.auc}{0.858} & \resultnum{j.x3.r050.free.auc}{0.585} & $\resultnum{j.x3.r050.delta.point}{-0.273}$ [$\resultnum{j.x3.r050.delta.lower}{-0.281}$, $\resultnum{j.x3.r050.delta.upper}{-0.265}$] & $\resultnum{j.x3.r050.free.r_shift}{-0.0028}$ & yes \\
\resultnum{j.b.r055.r}{0.55} & \resultnum{j.b.r055.first.point}{0.10} & \resultnum{j.x3.r055.known.auc}{0.668} & \resultnum{j.x3.r055.free.auc}{0.542} & $\resultnum{j.x3.r055.delta.point}{-0.126}$ [$\resultnum{j.x3.r055.delta.lower}{-0.135}$, $\resultnum{j.x3.r055.delta.upper}{-0.117}$] & $\resultnum{j.x3.r055.free.r_shift}{-0.0011}$ & no \\
\resultnum{j.b.r060.r}{0.60} & \resultnum{j.b.r060.first.point}{0.10} & \resultnum{j.x3.r060.known.auc}{0.703} & \resultnum{j.x3.r060.free.auc}{0.543} & $\resultnum{j.x3.r060.delta.point}{-0.160}$ [$\resultnum{j.x3.r060.delta.lower}{-0.166}$, $\resultnum{j.x3.r060.delta.upper}{-0.154}$] & $\resultnum{j.x3.r060.free.r_shift}{-0.0014}$ & yes \\
\resultnum{j.b.r065.r}{0.65} & \resultnum{j.b.r065.first.point}{0.10} & \resultnum{j.x3.r065.known.auc}{0.706} & \resultnum{j.x3.r065.free.auc}{0.534} & $\resultnum{j.x3.r065.delta.point}{-0.172}$ [$\resultnum{j.x3.r065.delta.lower}{-0.177}$, $\resultnum{j.x3.r065.delta.upper}{-0.166}$] & $\resultnum{j.x3.r065.free.r_shift}{-0.0015}$ & yes \\
\resultnum{j.b.r070.r}{0.70} & \resultnum{j.b.r070.first.point}{0.10} & \resultnum{j.x3.r070.known.auc}{0.730} & \resultnum{j.x3.r070.free.auc}{0.552} & $\resultnum{j.x3.r070.delta.point}{-0.178}$ [$\resultnum{j.x3.r070.delta.lower}{-0.183}$, $\resultnum{j.x3.r070.delta.upper}{-0.173}$] & $\resultnum{j.x3.r070.free.r_shift}{-0.0017}$ & yes \\
\resultnum{j.b.r075.r}{0.75} & \resultnum{j.b.r075.first.point}{0.10} & \resultnum{j.x3.r075.known.auc}{0.733} & \resultnum{j.x3.r075.free.auc}{0.560} & $\resultnum{j.x3.r075.delta.point}{-0.172}$ [$\resultnum{j.x3.r075.delta.lower}{-0.177}$, $\resultnum{j.x3.r075.delta.upper}{-0.168}$] & $\resultnum{j.x3.r075.free.r_shift}{-0.0017}$ & yes \\
\resultnum{j.b.r080.r}{0.80} & \resultnum{j.b.r080.first.point}{0.10} & \resultnum{j.x3.r080.known.auc}{0.735} & \resultnum{j.x3.r080.free.auc}{0.562} & $\resultnum{j.x3.r080.delta.point}{-0.173}$ [$\resultnum{j.x3.r080.delta.lower}{-0.177}$, $\resultnum{j.x3.r080.delta.upper}{-0.169}$] & $\resultnum{j.x3.r080.free.r_shift}{-0.0017}$ & yes \\
\resultnum{j.b.r085.r}{0.85} & \resultnum{j.b.r085.first.point}{0.10} & \resultnum{j.x3.r085.known.auc}{0.763} & \resultnum{j.x3.r085.free.auc}{0.568} & $\resultnum{j.x3.r085.delta.point}{-0.195}$ [$\resultnum{j.x3.r085.delta.lower}{-0.200}$, $\resultnum{j.x3.r085.delta.upper}{-0.190}$] & $\resultnum{j.x3.r085.free.r_shift}{-0.0019}$ & yes \\
\resultnum{j.b.r090.r}{0.90} & \resultnum{j.b.r090.first.point}{0.10} & \resultnum{j.x3.r090.known.auc}{0.783} & \resultnum{j.x3.r090.free.auc}{0.582} & $\resultnum{j.x3.r090.delta.point}{-0.202}$ [$\resultnum{j.x3.r090.delta.lower}{-0.208}$, $\resultnum{j.x3.r090.delta.upper}{-0.195}$] & $\resultnum{j.x3.r090.free.r_shift}{-0.0022}$ & yes \\
\botrule
\end{tabular}
\par\smallskip\parbox{\linewidth}{Each row selects the first boundary of the main classifier against trait-plus-autoregressive persistence (B). Both exploratory classifiers use lag four and lag five; their AUCs are distinct from the main three-correlation classifier. The known-share fit matches the first two lags; the free-share fit estimates B's signal share and also matches lag three. Differences are free minus known, with paired ninety-five percent bootstrap intervals conditional on the fitted classifiers; uncertainty from matching is not included. Lost crossing is the archived flag: known-share point AUC reaches \resultnum{j.target.r070}{0.70} while the free-share interval's upper end is below that target. $r_B-r$ is fitted minus design share; Table~\ref{tab:x3-fit} gives B's fitted parameters.}
\end{table}

\begin{table}[tp]
\caption{Known-share and free-share persistence fits at the first boundary.}\label{tab:x3-fit}
\centering\scriptsize
\setlength{\tabcolsep}{2pt}
\begin{tabular}{l@{\quad}l@{\quad}>{\raggedright\arraybackslash}p{12mm}>{\raggedright\arraybackslash}p{12mm}>{\raggedright\arraybackslash}p{12mm}>{\raggedright\arraybackslash}p{12mm}>{\raggedright\arraybackslash}p{15mm}>{\raggedright\arraybackslash}p{15mm}>{\raggedright\arraybackslash}p{18mm}}
\toprule
Design $r$ & $a^*$ & Known $\phi$ & Known $\lambda$ & Free $\phi$ & Free $\lambda$ & Fitted $r_B$ & $r_B-r$ & Lag-three miss (fit SEs) \\
\midrule
\resultnum{j.b.r050.r}{0.50} & \resultnum{j.b.r050.first.point}{0.30} & \resultnum{j.x3.r050.known.phi}{0.839} & \resultnum{j.x3.r050.known.lambda}{0.363} & \resultnum{j.x3.r050.free.phi}{0.888} & \resultnum{j.x3.r050.free.lambda}{0.133} & \resultnum{j.x3.r050.free.r}{0.4972} & $\resultnum{j.x3.r050.free.r_shift}{-0.0028}$ & $\resultnum{j.x3.r050.l3_miss_se}{-13.4}$ \\
\resultnum{j.b.r055.r}{0.55} & \resultnum{j.b.r055.first.point}{0.10} & \resultnum{j.x3.r055.known.phi}{0.855} & \resultnum{j.x3.r055.known.lambda}{0.139} & \resultnum{j.x3.r055.free.phi}{0.869} & \resultnum{j.x3.r055.free.lambda}{0.064} & \resultnum{j.x3.r055.free.r}{0.5489} & $\resultnum{j.x3.r055.free.r_shift}{-0.0011}$ & $\resultnum{j.x3.r055.l3_miss_se}{-7.1}$ \\
\resultnum{j.b.r060.r}{0.60} & \resultnum{j.b.r060.first.point}{0.10} & \resultnum{j.x3.r060.known.phi}{0.853} & \resultnum{j.x3.r060.known.lambda}{0.151} & \resultnum{j.x3.r060.free.phi}{0.869} & \resultnum{j.x3.r060.free.lambda}{0.062} & \resultnum{j.x3.r060.free.r}{0.5986} & $\resultnum{j.x3.r060.free.r_shift}{-0.0014}$ & $\resultnum{j.x3.r060.l3_miss_se}{-9.0}$ \\
\resultnum{j.b.r065.r}{0.65} & \resultnum{j.b.r065.first.point}{0.10} & \resultnum{j.x3.r065.known.phi}{0.852} & \resultnum{j.x3.r065.known.lambda}{0.150} & \resultnum{j.x3.r065.free.phi}{0.869} & \resultnum{j.x3.r065.free.lambda}{0.061} & \resultnum{j.x3.r065.free.r}{0.6485} & $\resultnum{j.x3.r065.free.r_shift}{-0.0015}$ & $\resultnum{j.x3.r065.l3_miss_se}{-9.4}$ \\
\resultnum{j.b.r070.r}{0.70} & \resultnum{j.b.r070.first.point}{0.10} & \resultnum{j.x3.r070.known.phi}{0.851} & \resultnum{j.x3.r070.known.lambda}{0.157} & \resultnum{j.x3.r070.free.phi}{0.867} & \resultnum{j.x3.r070.free.lambda}{0.069} & \resultnum{j.x3.r070.free.r}{0.6983} & $\resultnum{j.x3.r070.free.r_shift}{-0.0017}$ & $\resultnum{j.x3.r070.l3_miss_se}{-9.6}$ \\
\resultnum{j.b.r075.r}{0.75} & \resultnum{j.b.r075.first.point}{0.10} & \resultnum{j.x3.r075.known.phi}{0.850} & \resultnum{j.x3.r075.known.lambda}{0.155} & \resultnum{j.x3.r075.free.phi}{0.866} & \resultnum{j.x3.r075.free.lambda}{0.074} & \resultnum{j.x3.r075.free.r}{0.7483} & $\resultnum{j.x3.r075.free.r_shift}{-0.0017}$ & $\resultnum{j.x3.r075.l3_miss_se}{-10.1}$ \\
\resultnum{j.b.r080.r}{0.80} & \resultnum{j.b.r080.first.point}{0.10} & \resultnum{j.x3.r080.known.phi}{0.851} & \resultnum{j.x3.r080.known.lambda}{0.151} & \resultnum{j.x3.r080.free.phi}{0.865} & \resultnum{j.x3.r080.free.lambda}{0.075} & \resultnum{j.x3.r080.free.r}{0.7983} & $\resultnum{j.x3.r080.free.r_shift}{-0.0017}$ & $\resultnum{j.x3.r080.l3_miss_se}{-9.5}$ \\
\resultnum{j.b.r085.r}{0.85} & \resultnum{j.b.r085.first.point}{0.10} & \resultnum{j.x3.r085.known.phi}{0.850} & \resultnum{j.x3.r085.known.lambda}{0.156} & \resultnum{j.x3.r085.free.phi}{0.865} & \resultnum{j.x3.r085.free.lambda}{0.074} & \resultnum{j.x3.r085.free.r}{0.8481} & $\resultnum{j.x3.r085.free.r_shift}{-0.0019}$ & $\resultnum{j.x3.r085.l3_miss_se}{-11.8}$ \\
\resultnum{j.b.r090.r}{0.90} & \resultnum{j.b.r090.first.point}{0.10} & \resultnum{j.x3.r090.known.phi}{0.849} & \resultnum{j.x3.r090.known.lambda}{0.159} & \resultnum{j.x3.r090.free.phi}{0.865} & \resultnum{j.x3.r090.free.lambda}{0.075} & \resultnum{j.x3.r090.free.r}{0.8978} & $\resultnum{j.x3.r090.free.r_shift}{-0.0022}$ & $\resultnum{j.x3.r090.l3_miss_se}{-11.7}$ \\
\botrule
\end{tabular}
\par\smallskip\parbox{\linewidth}{The share shift is fitted minus design share. The signed lag-three miss is influence minus B at the true share, divided by the matching attempt's standard error. It uses matching panels; the mean lag differences in Table~\ref{tab:x3-lags} use check panels. The free-share fit also matches lag three; the contributions of the share shift and the extra matched lag were not disentangled.}
\end{table}


Read after the results became available: at every share's first-boundary cell, the influence arm's mean within-person correlations fell below those of known-share B at lags three to five, and the shortfall grew with the lag (Table~\ref{tab:x3-lags}; negative at all three lags at \resultnum{j.x3.lag_shortfall.negative.n}{9} of nine shares and growing with the lag at \resultnum{j.x3.lag_shortfall.increasing.n}{9}), the lag-five shortfall being \resultnum{j.x3.lag_shortfall.ratio_l5_l3.min}{5.8} to \resultnum{j.x3.lag_shortfall.ratio_l5_l3.max}{6.9} times the lag-three one. Against free-share B, which matches lag three through its share, the shortfalls at lags four and five were much smaller. These are means over the check panels of the free-share analysis, not the confirmatory cells.

\begin{table}[tp]
\caption{Mean lag-correlation differences at the first boundary: influence minus known-share B and influence minus free-share B. These are means over the check panels of the free-share analysis.}\label{tab:x3-lags}
\centering\scriptsize
\setlength{\tabcolsep}{2pt}
\begin{tabular}{l@{\quad}l@{\quad}>{\raggedright\arraybackslash}p{16mm}>{\raggedright\arraybackslash}p{16mm}>{\raggedright\arraybackslash}p{16mm}>{\raggedright\arraybackslash}p{16mm}>{\raggedright\arraybackslash}p{16mm}>{\raggedright\arraybackslash}p{16mm}}
\toprule
Design $r$ & $a^*$ & Known: lag three & Known: lag four & Known: lag five & Free: lag three & Free: lag four & Free: lag five \\
\midrule
\resultnum{j.b.r050.r}{0.50} & \resultnum{j.b.r050.first.point}{0.30} & $\resultnum{j.x3.lag_shortfall.r050.l3}{-0.0018}$ & $\resultnum{j.x3.lag_shortfall.r050.l4}{-0.0060}$ & $\resultnum{j.x3.lag_shortfall.r050.l5}{-0.0120}$ & $\resultnum{j.x3.lag_shortfall.r050.free.l3}{0.0002}$ & $\resultnum{j.x3.lag_shortfall.r050.free.l4}{-0.0005}$ & $\resultnum{j.x3.lag_shortfall.r050.free.l5}{-0.0021}$ \\
\resultnum{j.b.r055.r}{0.55} & \resultnum{j.b.r055.first.point}{0.10} & $\resultnum{j.x3.lag_shortfall.r055.l3}{-0.0007}$ & $\resultnum{j.x3.lag_shortfall.r055.l4}{-0.0024}$ & $\resultnum{j.x3.lag_shortfall.r055.l5}{-0.0048}$ & $\resultnum{j.x3.lag_shortfall.r055.free.l3}{0.0001}$ & $\resultnum{j.x3.lag_shortfall.r055.free.l4}{-0.0003}$ & $\resultnum{j.x3.lag_shortfall.r055.free.l5}{-0.0011}$ \\
\resultnum{j.b.r060.r}{0.60} & \resultnum{j.b.r060.first.point}{0.10} & $\resultnum{j.x3.lag_shortfall.r060.l3}{-0.0010}$ & $\resultnum{j.x3.lag_shortfall.r060.l4}{-0.0031}$ & $\resultnum{j.x3.lag_shortfall.r060.l5}{-0.0060}$ & $\resultnum{j.x3.lag_shortfall.r060.free.l3}{0.0000}$ & $\resultnum{j.x3.lag_shortfall.r060.free.l4}{-0.0004}$ & $\resultnum{j.x3.lag_shortfall.r060.free.l5}{-0.0010}$ \\
\resultnum{j.b.r065.r}{0.65} & \resultnum{j.b.r065.first.point}{0.10} & $\resultnum{j.x3.lag_shortfall.r065.l3}{-0.0011}$ & $\resultnum{j.x3.lag_shortfall.r065.l4}{-0.0032}$ & $\resultnum{j.x3.lag_shortfall.r065.l5}{-0.0061}$ & $\resultnum{j.x3.lag_shortfall.r065.free.l3}{0.0000}$ & $\resultnum{j.x3.lag_shortfall.r065.free.l4}{-0.0003}$ & $\resultnum{j.x3.lag_shortfall.r065.free.l5}{-0.0009}$ \\
\resultnum{j.b.r070.r}{0.70} & \resultnum{j.b.r070.first.point}{0.10} & $\resultnum{j.x3.lag_shortfall.r070.l3}{-0.0011}$ & $\resultnum{j.x3.lag_shortfall.r070.l4}{-0.0035}$ & $\resultnum{j.x3.lag_shortfall.r070.l5}{-0.0070}$ & $\resultnum{j.x3.lag_shortfall.r070.free.l3}{0.0001}$ & $\resultnum{j.x3.lag_shortfall.r070.free.l4}{-0.0003}$ & $\resultnum{j.x3.lag_shortfall.r070.free.l5}{-0.0013}$ \\
\resultnum{j.b.r075.r}{0.75} & \resultnum{j.b.r075.first.point}{0.10} & $\resultnum{j.x3.lag_shortfall.r075.l3}{-0.0012}$ & $\resultnum{j.x3.lag_shortfall.r075.l4}{-0.0036}$ & $\resultnum{j.x3.lag_shortfall.r075.l5}{-0.0072}$ & $\resultnum{j.x3.lag_shortfall.r075.free.l3}{-0.0001}$ & $\resultnum{j.x3.lag_shortfall.r075.free.l4}{-0.0004}$ & $\resultnum{j.x3.lag_shortfall.r075.free.l5}{-0.0015}$ \\
\resultnum{j.b.r080.r}{0.80} & \resultnum{j.b.r080.first.point}{0.10} & $\resultnum{j.x3.lag_shortfall.r080.l3}{-0.0012}$ & $\resultnum{j.x3.lag_shortfall.r080.l4}{-0.0037}$ & $\resultnum{j.x3.lag_shortfall.r080.l5}{-0.0074}$ & $\resultnum{j.x3.lag_shortfall.r080.free.l3}{-0.0001}$ & $\resultnum{j.x3.lag_shortfall.r080.free.l4}{-0.0006}$ & $\resultnum{j.x3.lag_shortfall.r080.free.l5}{-0.0017}$ \\
\resultnum{j.b.r085.r}{0.85} & \resultnum{j.b.r085.first.point}{0.10} & $\resultnum{j.x3.lag_shortfall.r085.l3}{-0.0013}$ & $\resultnum{j.x3.lag_shortfall.r085.l4}{-0.0044}$ & $\resultnum{j.x3.lag_shortfall.r085.l5}{-0.0084}$ & $\resultnum{j.x3.lag_shortfall.r085.free.l3}{0.0001}$ & $\resultnum{j.x3.lag_shortfall.r085.free.l4}{-0.0006}$ & $\resultnum{j.x3.lag_shortfall.r085.free.l5}{-0.0017}$ \\
\resultnum{j.b.r090.r}{0.90} & \resultnum{j.b.r090.first.point}{0.10} & $\resultnum{j.x3.lag_shortfall.r090.l3}{-0.0015}$ & $\resultnum{j.x3.lag_shortfall.r090.l4}{-0.0047}$ & $\resultnum{j.x3.lag_shortfall.r090.l5}{-0.0092}$ & $\resultnum{j.x3.lag_shortfall.r090.free.l3}{0.0000}$ & $\resultnum{j.x3.lag_shortfall.r090.free.l4}{-0.0006}$ & $\resultnum{j.x3.lag_shortfall.r090.free.l5}{-0.0018}$ \\
\botrule
\end{tabular}
\par\smallskip\parbox{\linewidth}{Each row selects the first boundary of the main classifier against trait-plus-autoregressive persistence (B). The differences use the mean correlations recorded in the free-share analysis. Known-share B matches lags one and two; free-share B also fits the signal share and matches lag three.}
\end{table}

% Generated by v4_exploratory_export; edited during the journal revision
% (agent S, 2026-10-13): captions, headers and categorical labels translated
% into plain English. Every number keeps its \resultnum binding.
\begin{table*}[htbp]
\setlength{\tabcolsep}{2pt}
\caption{Free signal share (X3): every recorded matching attempt, by cell and arm (B at the known share or at a fitted share), with the matching size (1,024 pairs, or 4,096 in the retry), the final solver status and the independent validation. Status: search budget exhausted (machine label budget), no accepted fit found by the inner solver (inner-failure), and the frozen statuses ill-conditioned and non-monotone profile (Table~\ref{tab:glossary-states}). Budget exhaustion, a search without an accepted fit and a candidate rejected by validation are distinct outcomes; a search label alone does not show that the free-share family cannot reproduce the target moments. Residuals and evaluation counts are in the archived values catalog.}\label{tab:v4x-x3-attempts}
\centering\footnotesize
\begin{tabular}{@{}p{0.170\textwidth}p{0.150\textwidth}p{0.150\textwidth}p{0.270\textwidth}p{0.190\textwidth}@{}}
\toprule
\raggedright $a,r$ & \raggedright \mbox{Arm} & \raggedright \mbox{Matching} \mbox{pairs} & \raggedright \mbox{Status} & \raggedright \mbox{Validation} \tabularnewline
\midrule
\raggedright \resultnum{v4x.x3.cells.r050.a0010.cell.a}{0.001}, \resultnum{v4x.x3.cells.r050.a0010.cell.r}{0.5} & \raggedright known share & \raggedright \resultnum{v4x.x3.cells.r050.a0010.known.attempts.i000.matching_pairs}{1024} & \raggedright interior (accepted) & \raggedright passed \tabularnewline
\raggedright \resultnum{v4x.x3.cells.r050.a0010.cell.a}{0.001}, \resultnum{v4x.x3.cells.r050.a0010.cell.r}{0.5} & \raggedright free share & \raggedright \resultnum{v4x.x3.cells.r050.a0010.free.attempts.i000.matching_pairs}{1024} & \raggedright interior (accepted) & \raggedright passed \tabularnewline
\raggedright \resultnum{v4x.x3.cells.r050.a0030.cell.a}{0.003}, \resultnum{v4x.x3.cells.r050.a0030.cell.r}{0.5} & \raggedright known share & \raggedright \resultnum{v4x.x3.cells.r050.a0030.known.attempts.i000.matching_pairs}{1024} & \raggedright interior (accepted) & \raggedright passed \tabularnewline
\raggedright \resultnum{v4x.x3.cells.r050.a0030.cell.a}{0.003}, \resultnum{v4x.x3.cells.r050.a0030.cell.r}{0.5} & \raggedright free share & \raggedright \resultnum{v4x.x3.cells.r050.a0030.free.attempts.i000.matching_pairs}{1024} & \raggedright interior (accepted) & \raggedright passed \tabularnewline
\raggedright \resultnum{v4x.x3.cells.r050.a0100.cell.a}{0.01}, \resultnum{v4x.x3.cells.r050.a0100.cell.r}{0.5} & \raggedright known share & \raggedright \resultnum{v4x.x3.cells.r050.a0100.known.attempts.i000.matching_pairs}{1024} & \raggedright interior (accepted) & \raggedright passed \tabularnewline
\raggedright \resultnum{v4x.x3.cells.r050.a0100.cell.a}{0.01}, \resultnum{v4x.x3.cells.r050.a0100.cell.r}{0.5} & \raggedright free share & \raggedright \resultnum{v4x.x3.cells.r050.a0100.free.attempts.i000.matching_pairs}{1024} & \raggedright interior (accepted) & \raggedright passed \tabularnewline
\raggedright \resultnum{v4x.x3.cells.r050.a0300.cell.a}{0.03}, \resultnum{v4x.x3.cells.r050.a0300.cell.r}{0.5} & \raggedright known share & \raggedright \resultnum{v4x.x3.cells.r050.a0300.known.attempts.i000.matching_pairs}{1024} & \raggedright interior (accepted) & \raggedright passed \tabularnewline
\raggedright \resultnum{v4x.x3.cells.r050.a0300.cell.a}{0.03}, \resultnum{v4x.x3.cells.r050.a0300.cell.r}{0.5} & \raggedright free share & \raggedright \resultnum{v4x.x3.cells.r050.a0300.free.attempts.i000.matching_pairs}{1024} & \raggedright interior (accepted) & \raggedright passed \tabularnewline
\raggedright \resultnum{v4x.x3.cells.r050.a1000.cell.a}{0.1}, \resultnum{v4x.x3.cells.r050.a1000.cell.r}{0.5} & \raggedright known share & \raggedright \resultnum{v4x.x3.cells.r050.a1000.known.attempts.i000.matching_pairs}{1024} & \raggedright interior (accepted) & \raggedright passed \tabularnewline
\raggedright \resultnum{v4x.x3.cells.r050.a1000.cell.a}{0.1}, \resultnum{v4x.x3.cells.r050.a1000.cell.r}{0.5} & \raggedright free share & \raggedright \resultnum{v4x.x3.cells.r050.a1000.free.attempts.i000.matching_pairs}{1024} & \raggedright interior (accepted) & \raggedright passed \tabularnewline
\raggedright \resultnum{v4x.x3.cells.r050.a3000.cell.a}{0.3}, \resultnum{v4x.x3.cells.r050.a3000.cell.r}{0.5} & \raggedright known share & \raggedright \resultnum{v4x.x3.cells.r050.a3000.known.attempts.i000.matching_pairs}{1024} & \raggedright interior (accepted) & \raggedright passed \tabularnewline
\raggedright \resultnum{v4x.x3.cells.r050.a3000.cell.a}{0.3}, \resultnum{v4x.x3.cells.r050.a3000.cell.r}{0.5} & \raggedright free share & \raggedright \resultnum{v4x.x3.cells.r050.a3000.free.attempts.i000.matching_pairs}{1024} & \raggedright interior (accepted) & \raggedright passed \tabularnewline
\raggedright \resultnum{v4x.x3.cells.r050.a6000.cell.a}{0.6}, \resultnum{v4x.x3.cells.r050.a6000.cell.r}{0.5} & \raggedright known share & \raggedright \resultnum{v4x.x3.cells.r050.a6000.known.attempts.i000.matching_pairs}{1024} & \raggedright interior (accepted) & \raggedright passed \tabularnewline
\raggedright \resultnum{v4x.x3.cells.r050.a6000.cell.a}{0.6}, \resultnum{v4x.x3.cells.r050.a6000.cell.r}{0.5} & \raggedright free share & \raggedright \resultnum{v4x.x3.cells.r050.a6000.free.attempts.i000.matching_pairs}{1024} & \raggedright search budget exhausted & \raggedright not run (no fit) \tabularnewline
\raggedright \resultnum{v4x.x3.cells.r050.a9000.cell.a}{0.9}, \resultnum{v4x.x3.cells.r050.a9000.cell.r}{0.5} & \raggedright known share & \raggedright \resultnum{v4x.x3.cells.r050.a9000.known.attempts.i000.matching_pairs}{1024} & \raggedright ill-conditioned & \raggedright not run (no fit) \tabularnewline
\raggedright \resultnum{v4x.x3.cells.r055.a0010.cell.a}{0.001}, \resultnum{v4x.x3.cells.r055.a0010.cell.r}{0.55} & \raggedright known share & \raggedright \resultnum{v4x.x3.cells.r055.a0010.known.attempts.i000.matching_pairs}{1024} & \raggedright interior (accepted) & \raggedright passed \tabularnewline
\raggedright \resultnum{v4x.x3.cells.r055.a0010.cell.a}{0.001}, \resultnum{v4x.x3.cells.r055.a0010.cell.r}{0.55} & \raggedright free share & \raggedright \resultnum{v4x.x3.cells.r055.a0010.free.attempts.i000.matching_pairs}{1024} & \raggedright no accepted fit found & \raggedright not run (no fit) \tabularnewline
\raggedright \resultnum{v4x.x3.cells.r055.a0030.cell.a}{0.003}, \resultnum{v4x.x3.cells.r055.a0030.cell.r}{0.55} & \raggedright known share & \raggedright \resultnum{v4x.x3.cells.r055.a0030.known.attempts.i000.matching_pairs}{1024} & \raggedright interior (accepted) & \raggedright passed \tabularnewline
\raggedright \resultnum{v4x.x3.cells.r055.a0030.cell.a}{0.003}, \resultnum{v4x.x3.cells.r055.a0030.cell.r}{0.55} & \raggedright free share & \raggedright \resultnum{v4x.x3.cells.r055.a0030.free.attempts.i000.matching_pairs}{1024} & \raggedright no accepted fit found & \raggedright not run (no fit) \tabularnewline
\raggedright \resultnum{v4x.x3.cells.r055.a0100.cell.a}{0.01}, \resultnum{v4x.x3.cells.r055.a0100.cell.r}{0.55} & \raggedright known share & \raggedright \resultnum{v4x.x3.cells.r055.a0100.known.attempts.i000.matching_pairs}{1024} & \raggedright interior (accepted) & \raggedright passed \tabularnewline
\raggedright \resultnum{v4x.x3.cells.r055.a0100.cell.a}{0.01}, \resultnum{v4x.x3.cells.r055.a0100.cell.r}{0.55} & \raggedright free share & \raggedright \resultnum{v4x.x3.cells.r055.a0100.free.attempts.i000.matching_pairs}{1024} & \raggedright interior (accepted) & \raggedright passed \tabularnewline
\raggedright \resultnum{v4x.x3.cells.r055.a0300.cell.a}{0.03}, \resultnum{v4x.x3.cells.r055.a0300.cell.r}{0.55} & \raggedright known share & \raggedright \resultnum{v4x.x3.cells.r055.a0300.known.attempts.i000.matching_pairs}{1024} & \raggedright interior (accepted) & \raggedright passed \tabularnewline
\raggedright \resultnum{v4x.x3.cells.r055.a0300.cell.a}{0.03}, \resultnum{v4x.x3.cells.r055.a0300.cell.r}{0.55} & \raggedright free share & \raggedright \resultnum{v4x.x3.cells.r055.a0300.free.attempts.i000.matching_pairs}{1024} & \raggedright interior (accepted) & \raggedright passed \tabularnewline
\raggedright \resultnum{v4x.x3.cells.r055.a1000.cell.a}{0.1}, \resultnum{v4x.x3.cells.r055.a1000.cell.r}{0.55} & \raggedright known share & \raggedright \resultnum{v4x.x3.cells.r055.a1000.known.attempts.i000.matching_pairs}{1024} & \raggedright interior (accepted) & \raggedright passed \tabularnewline
\raggedright \resultnum{v4x.x3.cells.r055.a1000.cell.a}{0.1}, \resultnum{v4x.x3.cells.r055.a1000.cell.r}{0.55} & \raggedright free share & \raggedright \resultnum{v4x.x3.cells.r055.a1000.free.attempts.i000.matching_pairs}{1024} & \raggedright interior (accepted) & \raggedright passed \tabularnewline
\raggedright \resultnum{v4x.x3.cells.r055.a3000.cell.a}{0.3}, \resultnum{v4x.x3.cells.r055.a3000.cell.r}{0.55} & \raggedright known share & \raggedright \resultnum{v4x.x3.cells.r055.a3000.known.attempts.i000.matching_pairs}{1024} & \raggedright interior (accepted) & \raggedright passed \tabularnewline
\raggedright \resultnum{v4x.x3.cells.r055.a3000.cell.a}{0.3}, \resultnum{v4x.x3.cells.r055.a3000.cell.r}{0.55} & \raggedright free share & \raggedright \resultnum{v4x.x3.cells.r055.a3000.free.attempts.i000.matching_pairs}{1024} & \raggedright search budget exhausted & \raggedright not run (no fit) \tabularnewline
\raggedright \resultnum{v4x.x3.cells.r055.a6000.cell.a}{0.6}, \resultnum{v4x.x3.cells.r055.a6000.cell.r}{0.55} & \raggedright known share & \raggedright \resultnum{v4x.x3.cells.r055.a6000.known.attempts.i000.matching_pairs}{1024} & \raggedright interior (accepted) & \raggedright passed \tabularnewline
\raggedright \resultnum{v4x.x3.cells.r055.a6000.cell.a}{0.6}, \resultnum{v4x.x3.cells.r055.a6000.cell.r}{0.55} & \raggedright free share & \raggedright \resultnum{v4x.x3.cells.r055.a6000.free.attempts.i000.matching_pairs}{1024} & \raggedright no accepted fit found & \raggedright not run (no fit) \tabularnewline
\raggedright \resultnum{v4x.x3.cells.r055.a9000.cell.a}{0.9}, \resultnum{v4x.x3.cells.r055.a9000.cell.r}{0.55} & \raggedright known share & \raggedright \resultnum{v4x.x3.cells.r055.a9000.known.attempts.i000.matching_pairs}{1024} & \raggedright ill-conditioned & \raggedright not run (no fit) \tabularnewline
\raggedright \resultnum{v4x.x3.cells.r060.a0010.cell.a}{0.001}, \resultnum{v4x.x3.cells.r060.a0010.cell.r}{0.6} & \raggedright known share & \raggedright \resultnum{v4x.x3.cells.r060.a0010.known.attempts.i000.matching_pairs}{1024} & \raggedright interior (accepted) & \raggedright passed \tabularnewline
\raggedright \resultnum{v4x.x3.cells.r060.a0010.cell.a}{0.001}, \resultnum{v4x.x3.cells.r060.a0010.cell.r}{0.6} & \raggedright free share & \raggedright \resultnum{v4x.x3.cells.r060.a0010.free.attempts.i000.matching_pairs}{1024} & \raggedright no accepted fit found & \raggedright not run (no fit) \tabularnewline
\raggedright \resultnum{v4x.x3.cells.r060.a0030.cell.a}{0.003}, \resultnum{v4x.x3.cells.r060.a0030.cell.r}{0.6} & \raggedright known share & \raggedright \resultnum{v4x.x3.cells.r060.a0030.known.attempts.i000.matching_pairs}{1024} & \raggedright interior (accepted) & \raggedright passed \tabularnewline
\raggedright \resultnum{v4x.x3.cells.r060.a0030.cell.a}{0.003}, \resultnum{v4x.x3.cells.r060.a0030.cell.r}{0.6} & \raggedright free share & \raggedright \resultnum{v4x.x3.cells.r060.a0030.free.attempts.i000.matching_pairs}{1024} & \raggedright interior (accepted) & \raggedright passed \tabularnewline
\raggedright \resultnum{v4x.x3.cells.r060.a0100.cell.a}{0.01}, \resultnum{v4x.x3.cells.r060.a0100.cell.r}{0.6} & \raggedright known share & \raggedright \resultnum{v4x.x3.cells.r060.a0100.known.attempts.i000.matching_pairs}{1024} & \raggedright interior (accepted) & \raggedright passed \tabularnewline
\raggedright \resultnum{v4x.x3.cells.r060.a0100.cell.a}{0.01}, \resultnum{v4x.x3.cells.r060.a0100.cell.r}{0.6} & \raggedright free share & \raggedright \resultnum{v4x.x3.cells.r060.a0100.free.attempts.i000.matching_pairs}{1024} & \raggedright interior (accepted) & \raggedright passed \tabularnewline
\raggedright \resultnum{v4x.x3.cells.r060.a0300.cell.a}{0.03}, \resultnum{v4x.x3.cells.r060.a0300.cell.r}{0.6} & \raggedright known share & \raggedright \resultnum{v4x.x3.cells.r060.a0300.known.attempts.i000.matching_pairs}{1024} & \raggedright interior (accepted) & \raggedright passed \tabularnewline
\raggedright \resultnum{v4x.x3.cells.r060.a0300.cell.a}{0.03}, \resultnum{v4x.x3.cells.r060.a0300.cell.r}{0.6} & \raggedright free share & \raggedright \resultnum{v4x.x3.cells.r060.a0300.free.attempts.i000.matching_pairs}{1024} & \raggedright interior (accepted) & \raggedright passed \tabularnewline
\raggedright \resultnum{v4x.x3.cells.r060.a1000.cell.a}{0.1}, \resultnum{v4x.x3.cells.r060.a1000.cell.r}{0.6} & \raggedright known share & \raggedright \resultnum{v4x.x3.cells.r060.a1000.known.attempts.i000.matching_pairs}{1024} & \raggedright interior (accepted) & \raggedright passed \tabularnewline
\raggedright \resultnum{v4x.x3.cells.r060.a1000.cell.a}{0.1}, \resultnum{v4x.x3.cells.r060.a1000.cell.r}{0.6} & \raggedright free share & \raggedright \resultnum{v4x.x3.cells.r060.a1000.free.attempts.i000.matching_pairs}{1024} & \raggedright interior (accepted) & \raggedright passed \tabularnewline
\raggedright \resultnum{v4x.x3.cells.r060.a3000.cell.a}{0.3}, \resultnum{v4x.x3.cells.r060.a3000.cell.r}{0.6} & \raggedright known share & \raggedright \resultnum{v4x.x3.cells.r060.a3000.known.attempts.i000.matching_pairs}{1024} & \raggedright interior (accepted) & \raggedright passed \tabularnewline
\raggedright \resultnum{v4x.x3.cells.r060.a3000.cell.a}{0.3}, \resultnum{v4x.x3.cells.r060.a3000.cell.r}{0.6} & \raggedright free share & \raggedright \resultnum{v4x.x3.cells.r060.a3000.free.attempts.i000.matching_pairs}{1024} & \raggedright interior (accepted) & \raggedright passed \tabularnewline
\raggedright \resultnum{v4x.x3.cells.r060.a6000.cell.a}{0.6}, \resultnum{v4x.x3.cells.r060.a6000.cell.r}{0.6} & \raggedright known share & \raggedright \resultnum{v4x.x3.cells.r060.a6000.known.attempts.i000.matching_pairs}{1024} & \raggedright interior (accepted) & \raggedright passed \tabularnewline
\raggedright \resultnum{v4x.x3.cells.r060.a6000.cell.a}{0.6}, \resultnum{v4x.x3.cells.r060.a6000.cell.r}{0.6} & \raggedright free share & \raggedright \resultnum{v4x.x3.cells.r060.a6000.free.attempts.i000.matching_pairs}{1024} & \raggedright search budget exhausted & \raggedright not run (no fit) \tabularnewline
\raggedright \resultnum{v4x.x3.cells.r060.a9000.cell.a}{0.9}, \resultnum{v4x.x3.cells.r060.a9000.cell.r}{0.6} & \raggedright known share & \raggedright \resultnum{v4x.x3.cells.r060.a9000.known.attempts.i000.matching_pairs}{1024} & \raggedright interior (accepted) & \raggedright passed \tabularnewline
\raggedright \resultnum{v4x.x3.cells.r060.a9000.cell.a}{0.9}, \resultnum{v4x.x3.cells.r060.a9000.cell.r}{0.6} & \raggedright free share & \raggedright \resultnum{v4x.x3.cells.r060.a9000.free.attempts.i000.matching_pairs}{1024} & \raggedright no accepted fit found & \raggedright not run (no fit) \tabularnewline
\raggedright \resultnum{v4x.x3.cells.r065.a0010.cell.a}{0.001}, \resultnum{v4x.x3.cells.r065.a0010.cell.r}{0.65} & \raggedright known share & \raggedright \resultnum{v4x.x3.cells.r065.a0010.known.attempts.i000.matching_pairs}{1024} & \raggedright interior (accepted) & \raggedright passed \tabularnewline
\raggedright \resultnum{v4x.x3.cells.r065.a0010.cell.a}{0.001}, \resultnum{v4x.x3.cells.r065.a0010.cell.r}{0.65} & \raggedright free share & \raggedright \resultnum{v4x.x3.cells.r065.a0010.free.attempts.i000.matching_pairs}{1024} & \raggedright interior (accepted) & \raggedright passed \tabularnewline
\raggedright \resultnum{v4x.x3.cells.r065.a0030.cell.a}{0.003}, \resultnum{v4x.x3.cells.r065.a0030.cell.r}{0.65} & \raggedright known share & \raggedright \resultnum{v4x.x3.cells.r065.a0030.known.attempts.i000.matching_pairs}{1024} & \raggedright interior (accepted) & \raggedright passed \tabularnewline
\raggedright \resultnum{v4x.x3.cells.r065.a0030.cell.a}{0.003}, \resultnum{v4x.x3.cells.r065.a0030.cell.r}{0.65} & \raggedright free share & \raggedright \resultnum{v4x.x3.cells.r065.a0030.free.attempts.i000.matching_pairs}{1024} & \raggedright interior (accepted) & \raggedright passed \tabularnewline
\raggedright \resultnum{v4x.x3.cells.r065.a0100.cell.a}{0.01}, \resultnum{v4x.x3.cells.r065.a0100.cell.r}{0.65} & \raggedright known share & \raggedright \resultnum{v4x.x3.cells.r065.a0100.known.attempts.i000.matching_pairs}{1024} & \raggedright interior (accepted) & \raggedright passed \tabularnewline
\raggedright \resultnum{v4x.x3.cells.r065.a0100.cell.a}{0.01}, \resultnum{v4x.x3.cells.r065.a0100.cell.r}{0.65} & \raggedright free share & \raggedright \resultnum{v4x.x3.cells.r065.a0100.free.attempts.i000.matching_pairs}{1024} & \raggedright interior (accepted) & \raggedright passed \tabularnewline
\botrule
\end{tabular}
\end{table*}
\begin{table*}[htbp]
\setlength{\tabcolsep}{2pt}
\addtocounter{table}{-\csname @ne\endcsname}
\providecommand{\theHtable}{\thetable}
\renewcommand{\theHtable}{v4x-x3-attempts-part1}
\caption{Free signal share (X3): matching attempts (continued, part 2 of 3)}\label{tab:v4x-x3-attempts-part1}
\centering\footnotesize
\begin{tabular}{@{}p{0.170\textwidth}p{0.150\textwidth}p{0.150\textwidth}p{0.270\textwidth}p{0.190\textwidth}@{}}
\toprule
\raggedright $a,r$ & \raggedright \mbox{Arm} & \raggedright \mbox{Matching} \mbox{pairs} & \raggedright \mbox{Status} & \raggedright \mbox{Validation} \tabularnewline
\midrule
\raggedright \resultnum{v4x.x3.cells.r065.a0300.cell.a}{0.03}, \resultnum{v4x.x3.cells.r065.a0300.cell.r}{0.65} & \raggedright known share & \raggedright \resultnum{v4x.x3.cells.r065.a0300.known.attempts.i000.matching_pairs}{1024} & \raggedright interior (accepted) & \raggedright passed \tabularnewline
\raggedright \resultnum{v4x.x3.cells.r065.a0300.cell.a}{0.03}, \resultnum{v4x.x3.cells.r065.a0300.cell.r}{0.65} & \raggedright free share & \raggedright \resultnum{v4x.x3.cells.r065.a0300.free.attempts.i000.matching_pairs}{1024} & \raggedright interior (accepted) & \raggedright passed \tabularnewline
\raggedright \resultnum{v4x.x3.cells.r065.a1000.cell.a}{0.1}, \resultnum{v4x.x3.cells.r065.a1000.cell.r}{0.65} & \raggedright known share & \raggedright \resultnum{v4x.x3.cells.r065.a1000.known.attempts.i000.matching_pairs}{1024} & \raggedright interior (accepted) & \raggedright passed \tabularnewline
\raggedright \resultnum{v4x.x3.cells.r065.a1000.cell.a}{0.1}, \resultnum{v4x.x3.cells.r065.a1000.cell.r}{0.65} & \raggedright free share & \raggedright \resultnum{v4x.x3.cells.r065.a1000.free.attempts.i000.matching_pairs}{1024} & \raggedright interior (accepted) & \raggedright passed \tabularnewline
\raggedright \resultnum{v4x.x3.cells.r065.a3000.cell.a}{0.3}, \resultnum{v4x.x3.cells.r065.a3000.cell.r}{0.65} & \raggedright known share & \raggedright \resultnum{v4x.x3.cells.r065.a3000.known.attempts.i000.matching_pairs}{1024} & \raggedright interior (accepted) & \raggedright passed \tabularnewline
\raggedright \resultnum{v4x.x3.cells.r065.a3000.cell.a}{0.3}, \resultnum{v4x.x3.cells.r065.a3000.cell.r}{0.65} & \raggedright free share & \raggedright \resultnum{v4x.x3.cells.r065.a3000.free.attempts.i000.matching_pairs}{1024} & \raggedright interior (accepted) & \raggedright passed \tabularnewline
\raggedright \resultnum{v4x.x3.cells.r065.a6000.cell.a}{0.6}, \resultnum{v4x.x3.cells.r065.a6000.cell.r}{0.65} & \raggedright known share & \raggedright \resultnum{v4x.x3.cells.r065.a6000.known.attempts.i000.matching_pairs}{1024} & \raggedright interior (accepted) & \raggedright passed \tabularnewline
\raggedright \resultnum{v4x.x3.cells.r065.a6000.cell.a}{0.6}, \resultnum{v4x.x3.cells.r065.a6000.cell.r}{0.65} & \raggedright free share & \raggedright \resultnum{v4x.x3.cells.r065.a6000.free.attempts.i000.matching_pairs}{1024} & \raggedright interior (accepted) & \raggedright passed \tabularnewline
\raggedright \resultnum{v4x.x3.cells.r065.a9000.cell.a}{0.9}, \resultnum{v4x.x3.cells.r065.a9000.cell.r}{0.65} & \raggedright known share & \raggedright \resultnum{v4x.x3.cells.r065.a9000.known.attempts.i000.matching_pairs}{1024} & \raggedright interior (accepted) & \raggedright passed \tabularnewline
\raggedright \resultnum{v4x.x3.cells.r065.a9000.cell.a}{0.9}, \resultnum{v4x.x3.cells.r065.a9000.cell.r}{0.65} & \raggedright free share & \raggedright \resultnum{v4x.x3.cells.r065.a9000.free.attempts.i000.matching_pairs}{1024} & \raggedright no accepted fit found & \raggedright not run (no fit) \tabularnewline
\raggedright \resultnum{v4x.x3.cells.r070.a0010.cell.a}{0.001}, \resultnum{v4x.x3.cells.r070.a0010.cell.r}{0.7} & \raggedright known share & \raggedright \resultnum{v4x.x3.cells.r070.a0010.known.attempts.i000.matching_pairs}{1024} & \raggedright interior (accepted) & \raggedright passed \tabularnewline
\raggedright \resultnum{v4x.x3.cells.r070.a0010.cell.a}{0.001}, \resultnum{v4x.x3.cells.r070.a0010.cell.r}{0.7} & \raggedright free share & \raggedright \resultnum{v4x.x3.cells.r070.a0010.free.attempts.i000.matching_pairs}{1024} & \raggedright interior (accepted) & \raggedright passed \tabularnewline
\raggedright \resultnum{v4x.x3.cells.r070.a0030.cell.a}{0.003}, \resultnum{v4x.x3.cells.r070.a0030.cell.r}{0.7} & \raggedright known share & \raggedright \resultnum{v4x.x3.cells.r070.a0030.known.attempts.i000.matching_pairs}{1024} & \raggedright interior (accepted) & \raggedright passed \tabularnewline
\raggedright \resultnum{v4x.x3.cells.r070.a0030.cell.a}{0.003}, \resultnum{v4x.x3.cells.r070.a0030.cell.r}{0.7} & \raggedright free share & \raggedright \resultnum{v4x.x3.cells.r070.a0030.free.attempts.i000.matching_pairs}{1024} & \raggedright interior (accepted) & \raggedright passed \tabularnewline
\raggedright \resultnum{v4x.x3.cells.r070.a0100.cell.a}{0.01}, \resultnum{v4x.x3.cells.r070.a0100.cell.r}{0.7} & \raggedright known share & \raggedright \resultnum{v4x.x3.cells.r070.a0100.known.attempts.i000.matching_pairs}{1024} & \raggedright interior (accepted) & \raggedright passed \tabularnewline
\raggedright \resultnum{v4x.x3.cells.r070.a0100.cell.a}{0.01}, \resultnum{v4x.x3.cells.r070.a0100.cell.r}{0.7} & \raggedright free share & \raggedright \resultnum{v4x.x3.cells.r070.a0100.free.attempts.i000.matching_pairs}{1024} & \raggedright interior (accepted) & \raggedright passed \tabularnewline
\raggedright \resultnum{v4x.x3.cells.r070.a0300.cell.a}{0.03}, \resultnum{v4x.x3.cells.r070.a0300.cell.r}{0.7} & \raggedright known share & \raggedright \resultnum{v4x.x3.cells.r070.a0300.known.attempts.i000.matching_pairs}{1024} & \raggedright interior (accepted) & \raggedright passed \tabularnewline
\raggedright \resultnum{v4x.x3.cells.r070.a0300.cell.a}{0.03}, \resultnum{v4x.x3.cells.r070.a0300.cell.r}{0.7} & \raggedright free share & \raggedright \resultnum{v4x.x3.cells.r070.a0300.free.attempts.i000.matching_pairs}{1024} & \raggedright interior (accepted) & \raggedright passed \tabularnewline
\raggedright \resultnum{v4x.x3.cells.r070.a1000.cell.a}{0.1}, \resultnum{v4x.x3.cells.r070.a1000.cell.r}{0.7} & \raggedright known share & \raggedright \resultnum{v4x.x3.cells.r070.a1000.known.attempts.i000.matching_pairs}{1024} & \raggedright interior (accepted) & \raggedright passed \tabularnewline
\raggedright \resultnum{v4x.x3.cells.r070.a1000.cell.a}{0.1}, \resultnum{v4x.x3.cells.r070.a1000.cell.r}{0.7} & \raggedright free share & \raggedright \resultnum{v4x.x3.cells.r070.a1000.free.attempts.i000.matching_pairs}{1024} & \raggedright interior (accepted) & \raggedright passed \tabularnewline
\raggedright \resultnum{v4x.x3.cells.r070.a3000.cell.a}{0.3}, \resultnum{v4x.x3.cells.r070.a3000.cell.r}{0.7} & \raggedright known share & \raggedright \resultnum{v4x.x3.cells.r070.a3000.known.attempts.i000.matching_pairs}{1024} & \raggedright interior (accepted) & \raggedright passed \tabularnewline
\raggedright \resultnum{v4x.x3.cells.r070.a3000.cell.a}{0.3}, \resultnum{v4x.x3.cells.r070.a3000.cell.r}{0.7} & \raggedright free share & \raggedright \resultnum{v4x.x3.cells.r070.a3000.free.attempts.i000.matching_pairs}{1024} & \raggedright interior (accepted) & \raggedright passed \tabularnewline
\raggedright \resultnum{v4x.x3.cells.r070.a6000.cell.a}{0.6}, \resultnum{v4x.x3.cells.r070.a6000.cell.r}{0.7} & \raggedright known share & \raggedright \resultnum{v4x.x3.cells.r070.a6000.known.attempts.i000.matching_pairs}{1024} & \raggedright interior (accepted) & \raggedright passed \tabularnewline
\raggedright \resultnum{v4x.x3.cells.r070.a6000.cell.a}{0.6}, \resultnum{v4x.x3.cells.r070.a6000.cell.r}{0.7} & \raggedright free share & \raggedright \resultnum{v4x.x3.cells.r070.a6000.free.attempts.i000.matching_pairs}{1024} & \raggedright no accepted fit found & \raggedright not run (no fit) \tabularnewline
\raggedright \resultnum{v4x.x3.cells.r070.a9000.cell.a}{0.9}, \resultnum{v4x.x3.cells.r070.a9000.cell.r}{0.7} & \raggedright known share & \raggedright \resultnum{v4x.x3.cells.r070.a9000.known.attempts.i000.matching_pairs}{1024} & \raggedright non-monotone profile & \raggedright not run (no fit) \tabularnewline
\raggedright \resultnum{v4x.x3.cells.r075.a0010.cell.a}{0.001}, \resultnum{v4x.x3.cells.r075.a0010.cell.r}{0.75} & \raggedright known share & \raggedright \resultnum{v4x.x3.cells.r075.a0010.known.attempts.i000.matching_pairs}{1024} & \raggedright interior (accepted) & \raggedright passed \tabularnewline
\raggedright \resultnum{v4x.x3.cells.r075.a0010.cell.a}{0.001}, \resultnum{v4x.x3.cells.r075.a0010.cell.r}{0.75} & \raggedright free share & \raggedright \resultnum{v4x.x3.cells.r075.a0010.free.attempts.i000.matching_pairs}{1024} & \raggedright interior (accepted) & \raggedright passed \tabularnewline
\raggedright \resultnum{v4x.x3.cells.r075.a0030.cell.a}{0.003}, \resultnum{v4x.x3.cells.r075.a0030.cell.r}{0.75} & \raggedright known share & \raggedright \resultnum{v4x.x3.cells.r075.a0030.known.attempts.i000.matching_pairs}{1024} & \raggedright interior (accepted) & \raggedright passed \tabularnewline
\raggedright \resultnum{v4x.x3.cells.r075.a0030.cell.a}{0.003}, \resultnum{v4x.x3.cells.r075.a0030.cell.r}{0.75} & \raggedright free share & \raggedright \resultnum{v4x.x3.cells.r075.a0030.free.attempts.i000.matching_pairs}{1024} & \raggedright interior (accepted) & \raggedright passed \tabularnewline
\raggedright \resultnum{v4x.x3.cells.r075.a0100.cell.a}{0.01}, \resultnum{v4x.x3.cells.r075.a0100.cell.r}{0.75} & \raggedright known share & \raggedright \resultnum{v4x.x3.cells.r075.a0100.known.attempts.i000.matching_pairs}{1024} & \raggedright interior (accepted) & \raggedright passed \tabularnewline
\raggedright \resultnum{v4x.x3.cells.r075.a0100.cell.a}{0.01}, \resultnum{v4x.x3.cells.r075.a0100.cell.r}{0.75} & \raggedright free share & \raggedright \resultnum{v4x.x3.cells.r075.a0100.free.attempts.i000.matching_pairs}{1024} & \raggedright interior (accepted) & \raggedright passed \tabularnewline
\raggedright \resultnum{v4x.x3.cells.r075.a0300.cell.a}{0.03}, \resultnum{v4x.x3.cells.r075.a0300.cell.r}{0.75} & \raggedright known share & \raggedright \resultnum{v4x.x3.cells.r075.a0300.known.attempts.i000.matching_pairs}{1024} & \raggedright interior (accepted) & \raggedright passed \tabularnewline
\raggedright \resultnum{v4x.x3.cells.r075.a0300.cell.a}{0.03}, \resultnum{v4x.x3.cells.r075.a0300.cell.r}{0.75} & \raggedright free share & \raggedright \resultnum{v4x.x3.cells.r075.a0300.free.attempts.i000.matching_pairs}{1024} & \raggedright interior (accepted) & \raggedright passed \tabularnewline
\raggedright \resultnum{v4x.x3.cells.r075.a1000.cell.a}{0.1}, \resultnum{v4x.x3.cells.r075.a1000.cell.r}{0.75} & \raggedright known share & \raggedright \resultnum{v4x.x3.cells.r075.a1000.known.attempts.i000.matching_pairs}{1024} & \raggedright interior (accepted) & \raggedright passed \tabularnewline
\raggedright \resultnum{v4x.x3.cells.r075.a1000.cell.a}{0.1}, \resultnum{v4x.x3.cells.r075.a1000.cell.r}{0.75} & \raggedright free share & \raggedright \resultnum{v4x.x3.cells.r075.a1000.free.attempts.i000.matching_pairs}{1024} & \raggedright interior (accepted) & \raggedright passed \tabularnewline
\raggedright \resultnum{v4x.x3.cells.r075.a3000.cell.a}{0.3}, \resultnum{v4x.x3.cells.r075.a3000.cell.r}{0.75} & \raggedright known share & \raggedright \resultnum{v4x.x3.cells.r075.a3000.known.attempts.i000.matching_pairs}{1024} & \raggedright interior (accepted) & \raggedright passed \tabularnewline
\raggedright \resultnum{v4x.x3.cells.r075.a3000.cell.a}{0.3}, \resultnum{v4x.x3.cells.r075.a3000.cell.r}{0.75} & \raggedright free share & \raggedright \resultnum{v4x.x3.cells.r075.a3000.free.attempts.i000.matching_pairs}{1024} & \raggedright interior (accepted) & \raggedright passed \tabularnewline
\raggedright \resultnum{v4x.x3.cells.r075.a6000.cell.a}{0.6}, \resultnum{v4x.x3.cells.r075.a6000.cell.r}{0.75} & \raggedright known share & \raggedright \resultnum{v4x.x3.cells.r075.a6000.known.attempts.i000.matching_pairs}{1024} & \raggedright interior (accepted) & \raggedright passed \tabularnewline
\raggedright \resultnum{v4x.x3.cells.r075.a6000.cell.a}{0.6}, \resultnum{v4x.x3.cells.r075.a6000.cell.r}{0.75} & \raggedright free share & \raggedright \resultnum{v4x.x3.cells.r075.a6000.free.attempts.i000.matching_pairs}{1024} & \raggedright interior (accepted) & \raggedright passed \tabularnewline
\raggedright \resultnum{v4x.x3.cells.r075.a9000.cell.a}{0.9}, \resultnum{v4x.x3.cells.r075.a9000.cell.r}{0.75} & \raggedright known share & \raggedright \resultnum{v4x.x3.cells.r075.a9000.known.attempts.i000.matching_pairs}{1024} & \raggedright interior (accepted) & \raggedright passed \tabularnewline
\raggedright \resultnum{v4x.x3.cells.r075.a9000.cell.a}{0.9}, \resultnum{v4x.x3.cells.r075.a9000.cell.r}{0.75} & \raggedright free share & \raggedright \resultnum{v4x.x3.cells.r075.a9000.free.attempts.i000.matching_pairs}{1024} & \raggedright no accepted fit found & \raggedright not run (no fit) \tabularnewline
\raggedright \resultnum{v4x.x3.cells.r080.a0010.cell.a}{0.001}, \resultnum{v4x.x3.cells.r080.a0010.cell.r}{0.8} & \raggedright known share & \raggedright \resultnum{v4x.x3.cells.r080.a0010.known.attempts.i000.matching_pairs}{1024} & \raggedright interior (accepted) & \raggedright passed \tabularnewline
\raggedright \resultnum{v4x.x3.cells.r080.a0010.cell.a}{0.001}, \resultnum{v4x.x3.cells.r080.a0010.cell.r}{0.8} & \raggedright free share & \raggedright \resultnum{v4x.x3.cells.r080.a0010.free.attempts.i000.matching_pairs}{1024} & \raggedright interior (accepted) & \raggedright passed \tabularnewline
\raggedright \resultnum{v4x.x3.cells.r080.a0030.cell.a}{0.003}, \resultnum{v4x.x3.cells.r080.a0030.cell.r}{0.8} & \raggedright known share & \raggedright \resultnum{v4x.x3.cells.r080.a0030.known.attempts.i000.matching_pairs}{1024} & \raggedright interior (accepted) & \raggedright passed \tabularnewline
\raggedright \resultnum{v4x.x3.cells.r080.a0030.cell.a}{0.003}, \resultnum{v4x.x3.cells.r080.a0030.cell.r}{0.8} & \raggedright free share & \raggedright \resultnum{v4x.x3.cells.r080.a0030.free.attempts.i000.matching_pairs}{1024} & \raggedright interior (accepted) & \raggedright passed \tabularnewline
\raggedright \resultnum{v4x.x3.cells.r080.a0100.cell.a}{0.01}, \resultnum{v4x.x3.cells.r080.a0100.cell.r}{0.8} & \raggedright known share & \raggedright \resultnum{v4x.x3.cells.r080.a0100.known.attempts.i000.matching_pairs}{1024} & \raggedright interior (accepted) & \raggedright passed \tabularnewline
\raggedright \resultnum{v4x.x3.cells.r080.a0100.cell.a}{0.01}, \resultnum{v4x.x3.cells.r080.a0100.cell.r}{0.8} & \raggedright free share & \raggedright \resultnum{v4x.x3.cells.r080.a0100.free.attempts.i000.matching_pairs}{1024} & \raggedright interior (accepted) & \raggedright passed \tabularnewline
\raggedright \resultnum{v4x.x3.cells.r080.a0300.cell.a}{0.03}, \resultnum{v4x.x3.cells.r080.a0300.cell.r}{0.8} & \raggedright known share & \raggedright \resultnum{v4x.x3.cells.r080.a0300.known.attempts.i000.matching_pairs}{1024} & \raggedright interior (accepted) & \raggedright passed \tabularnewline
\raggedright \resultnum{v4x.x3.cells.r080.a0300.cell.a}{0.03}, \resultnum{v4x.x3.cells.r080.a0300.cell.r}{0.8} & \raggedright free share & \raggedright \resultnum{v4x.x3.cells.r080.a0300.free.attempts.i000.matching_pairs}{1024} & \raggedright interior (accepted) & \raggedright passed \tabularnewline
\raggedright \resultnum{v4x.x3.cells.r080.a1000.cell.a}{0.1}, \resultnum{v4x.x3.cells.r080.a1000.cell.r}{0.8} & \raggedright known share & \raggedright \resultnum{v4x.x3.cells.r080.a1000.known.attempts.i000.matching_pairs}{1024} & \raggedright interior (accepted) & \raggedright passed \tabularnewline
\raggedright \resultnum{v4x.x3.cells.r080.a1000.cell.a}{0.1}, \resultnum{v4x.x3.cells.r080.a1000.cell.r}{0.8} & \raggedright free share & \raggedright \resultnum{v4x.x3.cells.r080.a1000.free.attempts.i000.matching_pairs}{1024} & \raggedright interior (accepted) & \raggedright passed \tabularnewline
\raggedright \resultnum{v4x.x3.cells.r080.a3000.cell.a}{0.3}, \resultnum{v4x.x3.cells.r080.a3000.cell.r}{0.8} & \raggedright known share & \raggedright \resultnum{v4x.x3.cells.r080.a3000.known.attempts.i000.matching_pairs}{1024} & \raggedright interior (accepted) & \raggedright passed \tabularnewline
\botrule
\end{tabular}
\end{table*}
\begin{table*}[htbp]
\setlength{\tabcolsep}{2pt}
\addtocounter{table}{-\csname @ne\endcsname}
\providecommand{\theHtable}{\thetable}
\renewcommand{\theHtable}{v4x-x3-attempts-part2}
\caption{Free signal share (X3): matching attempts (continued, part 3 of 3)}\label{tab:v4x-x3-attempts-part2}
\centering\footnotesize
\begin{tabular}{@{}p{0.170\textwidth}p{0.150\textwidth}p{0.150\textwidth}p{0.270\textwidth}p{0.190\textwidth}@{}}
\toprule
\raggedright $a,r$ & \raggedright \mbox{Arm} & \raggedright \mbox{Matching} \mbox{pairs} & \raggedright \mbox{Status} & \raggedright \mbox{Validation} \tabularnewline
\midrule
\raggedright \resultnum{v4x.x3.cells.r080.a3000.cell.a}{0.3}, \resultnum{v4x.x3.cells.r080.a3000.cell.r}{0.8} & \raggedright free share & \raggedright \resultnum{v4x.x3.cells.r080.a3000.free.attempts.i000.matching_pairs}{1024} & \raggedright interior (accepted) & \raggedright passed \tabularnewline
\raggedright \resultnum{v4x.x3.cells.r080.a6000.cell.a}{0.6}, \resultnum{v4x.x3.cells.r080.a6000.cell.r}{0.8} & \raggedright known share & \raggedright \resultnum{v4x.x3.cells.r080.a6000.known.attempts.i000.matching_pairs}{1024} & \raggedright interior (accepted) & \raggedright passed \tabularnewline
\raggedright \resultnum{v4x.x3.cells.r080.a6000.cell.a}{0.6}, \resultnum{v4x.x3.cells.r080.a6000.cell.r}{0.8} & \raggedright free share & \raggedright \resultnum{v4x.x3.cells.r080.a6000.free.attempts.i000.matching_pairs}{1024} & \raggedright interior (accepted) & \raggedright passed \tabularnewline
\raggedright \resultnum{v4x.x3.cells.r080.a9000.cell.a}{0.9}, \resultnum{v4x.x3.cells.r080.a9000.cell.r}{0.8} & \raggedright known share & \raggedright \resultnum{v4x.x3.cells.r080.a9000.known.attempts.i000.matching_pairs}{1024} & \raggedright interior (accepted) & \raggedright passed \tabularnewline
\raggedright \resultnum{v4x.x3.cells.r080.a9000.cell.a}{0.9}, \resultnum{v4x.x3.cells.r080.a9000.cell.r}{0.8} & \raggedright free share & \raggedright \resultnum{v4x.x3.cells.r080.a9000.free.attempts.i000.matching_pairs}{1024} & \raggedright no accepted fit found & \raggedright not run (no fit) \tabularnewline
\raggedright \resultnum{v4x.x3.cells.r085.a0010.cell.a}{0.001}, \resultnum{v4x.x3.cells.r085.a0010.cell.r}{0.85} & \raggedright known share & \raggedright \resultnum{v4x.x3.cells.r085.a0010.known.attempts.i000.matching_pairs}{1024} & \raggedright interior (accepted) & \raggedright passed \tabularnewline
\raggedright \resultnum{v4x.x3.cells.r085.a0010.cell.a}{0.001}, \resultnum{v4x.x3.cells.r085.a0010.cell.r}{0.85} & \raggedright free share & \raggedright \resultnum{v4x.x3.cells.r085.a0010.free.attempts.i000.matching_pairs}{1024} & \raggedright interior (accepted) & \raggedright passed \tabularnewline
\raggedright \resultnum{v4x.x3.cells.r085.a0030.cell.a}{0.003}, \resultnum{v4x.x3.cells.r085.a0030.cell.r}{0.85} & \raggedright known share & \raggedright \resultnum{v4x.x3.cells.r085.a0030.known.attempts.i000.matching_pairs}{1024} & \raggedright interior (accepted) & \raggedright passed \tabularnewline
\raggedright \resultnum{v4x.x3.cells.r085.a0030.cell.a}{0.003}, \resultnum{v4x.x3.cells.r085.a0030.cell.r}{0.85} & \raggedright free share & \raggedright \resultnum{v4x.x3.cells.r085.a0030.free.attempts.i000.matching_pairs}{1024} & \raggedright interior (accepted) & \raggedright passed \tabularnewline
\raggedright \resultnum{v4x.x3.cells.r085.a0100.cell.a}{0.01}, \resultnum{v4x.x3.cells.r085.a0100.cell.r}{0.85} & \raggedright known share & \raggedright \resultnum{v4x.x3.cells.r085.a0100.known.attempts.i000.matching_pairs}{1024} & \raggedright interior (accepted) & \raggedright passed \tabularnewline
\raggedright \resultnum{v4x.x3.cells.r085.a0100.cell.a}{0.01}, \resultnum{v4x.x3.cells.r085.a0100.cell.r}{0.85} & \raggedright free share & \raggedright \resultnum{v4x.x3.cells.r085.a0100.free.attempts.i000.matching_pairs}{1024} & \raggedright interior (accepted) & \raggedright passed \tabularnewline
\raggedright \resultnum{v4x.x3.cells.r085.a0300.cell.a}{0.03}, \resultnum{v4x.x3.cells.r085.a0300.cell.r}{0.85} & \raggedright known share & \raggedright \resultnum{v4x.x3.cells.r085.a0300.known.attempts.i000.matching_pairs}{1024} & \raggedright interior (accepted) & \raggedright passed \tabularnewline
\raggedright \resultnum{v4x.x3.cells.r085.a0300.cell.a}{0.03}, \resultnum{v4x.x3.cells.r085.a0300.cell.r}{0.85} & \raggedright free share & \raggedright \resultnum{v4x.x3.cells.r085.a0300.free.attempts.i000.matching_pairs}{1024} & \raggedright interior (accepted) & \raggedright passed \tabularnewline
\raggedright \resultnum{v4x.x3.cells.r085.a1000.cell.a}{0.1}, \resultnum{v4x.x3.cells.r085.a1000.cell.r}{0.85} & \raggedright known share & \raggedright \resultnum{v4x.x3.cells.r085.a1000.known.attempts.i000.matching_pairs}{1024} & \raggedright interior (accepted) & \raggedright passed \tabularnewline
\raggedright \resultnum{v4x.x3.cells.r085.a1000.cell.a}{0.1}, \resultnum{v4x.x3.cells.r085.a1000.cell.r}{0.85} & \raggedright free share & \raggedright \resultnum{v4x.x3.cells.r085.a1000.free.attempts.i000.matching_pairs}{1024} & \raggedright interior (accepted) & \raggedright passed \tabularnewline
\raggedright \resultnum{v4x.x3.cells.r085.a3000.cell.a}{0.3}, \resultnum{v4x.x3.cells.r085.a3000.cell.r}{0.85} & \raggedright known share & \raggedright \resultnum{v4x.x3.cells.r085.a3000.known.attempts.i000.matching_pairs}{1024} & \raggedright interior (accepted) & \raggedright passed \tabularnewline
\raggedright \resultnum{v4x.x3.cells.r085.a3000.cell.a}{0.3}, \resultnum{v4x.x3.cells.r085.a3000.cell.r}{0.85} & \raggedright free share & \raggedright \resultnum{v4x.x3.cells.r085.a3000.free.attempts.i000.matching_pairs}{1024} & \raggedright interior (accepted) & \raggedright passed \tabularnewline
\raggedright \resultnum{v4x.x3.cells.r085.a6000.cell.a}{0.6}, \resultnum{v4x.x3.cells.r085.a6000.cell.r}{0.85} & \raggedright known share & \raggedright \resultnum{v4x.x3.cells.r085.a6000.known.attempts.i000.matching_pairs}{1024} & \raggedright interior (accepted) & \raggedright passed \tabularnewline
\raggedright \resultnum{v4x.x3.cells.r085.a6000.cell.a}{0.6}, \resultnum{v4x.x3.cells.r085.a6000.cell.r}{0.85} & \raggedright free share & \raggedright \resultnum{v4x.x3.cells.r085.a6000.free.attempts.i000.matching_pairs}{1024} & \raggedright interior (accepted) & \raggedright passed \tabularnewline
\raggedright \resultnum{v4x.x3.cells.r085.a9000.cell.a}{0.9}, \resultnum{v4x.x3.cells.r085.a9000.cell.r}{0.85} & \raggedright known share & \raggedright \resultnum{v4x.x3.cells.r085.a9000.known.attempts.i000.matching_pairs}{1024} & \raggedright interior (accepted) & \raggedright passed \tabularnewline
\raggedright \resultnum{v4x.x3.cells.r085.a9000.cell.a}{0.9}, \resultnum{v4x.x3.cells.r085.a9000.cell.r}{0.85} & \raggedright free share & \raggedright \resultnum{v4x.x3.cells.r085.a9000.free.attempts.i000.matching_pairs}{1024} & \raggedright no accepted fit found & \raggedright not run (no fit) \tabularnewline
\raggedright \resultnum{v4x.x3.cells.r090.a0010.cell.a}{0.001}, \resultnum{v4x.x3.cells.r090.a0010.cell.r}{0.9} & \raggedright known share & \raggedright \resultnum{v4x.x3.cells.r090.a0010.known.attempts.i000.matching_pairs}{1024} & \raggedright interior (accepted) & \raggedright passed \tabularnewline
\raggedright \resultnum{v4x.x3.cells.r090.a0010.cell.a}{0.001}, \resultnum{v4x.x3.cells.r090.a0010.cell.r}{0.9} & \raggedright free share & \raggedright \resultnum{v4x.x3.cells.r090.a0010.free.attempts.i000.matching_pairs}{1024} & \raggedright interior (accepted) & \raggedright passed \tabularnewline
\raggedright \resultnum{v4x.x3.cells.r090.a0030.cell.a}{0.003}, \resultnum{v4x.x3.cells.r090.a0030.cell.r}{0.9} & \raggedright known share & \raggedright \resultnum{v4x.x3.cells.r090.a0030.known.attempts.i000.matching_pairs}{1024} & \raggedright interior (accepted) & \raggedright passed \tabularnewline
\raggedright \resultnum{v4x.x3.cells.r090.a0030.cell.a}{0.003}, \resultnum{v4x.x3.cells.r090.a0030.cell.r}{0.9} & \raggedright free share & \raggedright \resultnum{v4x.x3.cells.r090.a0030.free.attempts.i000.matching_pairs}{1024} & \raggedright interior (accepted) & \raggedright passed \tabularnewline
\raggedright \resultnum{v4x.x3.cells.r090.a0100.cell.a}{0.01}, \resultnum{v4x.x3.cells.r090.a0100.cell.r}{0.9} & \raggedright known share & \raggedright \resultnum{v4x.x3.cells.r090.a0100.known.attempts.i000.matching_pairs}{1024} & \raggedright interior (accepted) & \raggedright passed \tabularnewline
\raggedright \resultnum{v4x.x3.cells.r090.a0100.cell.a}{0.01}, \resultnum{v4x.x3.cells.r090.a0100.cell.r}{0.9} & \raggedright free share & \raggedright \resultnum{v4x.x3.cells.r090.a0100.free.attempts.i000.matching_pairs}{1024} & \raggedright interior (accepted) & \raggedright passed \tabularnewline
\raggedright \resultnum{v4x.x3.cells.r090.a0300.cell.a}{0.03}, \resultnum{v4x.x3.cells.r090.a0300.cell.r}{0.9} & \raggedright known share & \raggedright \resultnum{v4x.x3.cells.r090.a0300.known.attempts.i000.matching_pairs}{1024} & \raggedright interior (accepted) & \raggedright passed \tabularnewline
\raggedright \resultnum{v4x.x3.cells.r090.a0300.cell.a}{0.03}, \resultnum{v4x.x3.cells.r090.a0300.cell.r}{0.9} & \raggedright free share & \raggedright \resultnum{v4x.x3.cells.r090.a0300.free.attempts.i000.matching_pairs}{1024} & \raggedright interior (accepted) & \raggedright passed \tabularnewline
\raggedright \resultnum{v4x.x3.cells.r090.a1000.cell.a}{0.1}, \resultnum{v4x.x3.cells.r090.a1000.cell.r}{0.9} & \raggedright known share & \raggedright \resultnum{v4x.x3.cells.r090.a1000.known.attempts.i000.matching_pairs}{1024} & \raggedright interior (accepted) & \raggedright passed \tabularnewline
\raggedright \resultnum{v4x.x3.cells.r090.a1000.cell.a}{0.1}, \resultnum{v4x.x3.cells.r090.a1000.cell.r}{0.9} & \raggedright free share & \raggedright \resultnum{v4x.x3.cells.r090.a1000.free.attempts.i000.matching_pairs}{1024} & \raggedright interior (accepted) & \raggedright passed \tabularnewline
\raggedright \resultnum{v4x.x3.cells.r090.a3000.cell.a}{0.3}, \resultnum{v4x.x3.cells.r090.a3000.cell.r}{0.9} & \raggedright known share & \raggedright \resultnum{v4x.x3.cells.r090.a3000.known.attempts.i000.matching_pairs}{1024} & \raggedright interior (accepted) & \raggedright passed \tabularnewline
\raggedright \resultnum{v4x.x3.cells.r090.a3000.cell.a}{0.3}, \resultnum{v4x.x3.cells.r090.a3000.cell.r}{0.9} & \raggedright free share & \raggedright \resultnum{v4x.x3.cells.r090.a3000.free.attempts.i000.matching_pairs}{1024} & \raggedright interior (accepted) & \raggedright passed \tabularnewline
\raggedright \resultnum{v4x.x3.cells.r090.a6000.cell.a}{0.6}, \resultnum{v4x.x3.cells.r090.a6000.cell.r}{0.9} & \raggedright known share & \raggedright \resultnum{v4x.x3.cells.r090.a6000.known.attempts.i000.matching_pairs}{1024} & \raggedright interior (accepted) & \raggedright passed \tabularnewline
\raggedright \resultnum{v4x.x3.cells.r090.a6000.cell.a}{0.6}, \resultnum{v4x.x3.cells.r090.a6000.cell.r}{0.9} & \raggedright free share & \raggedright \resultnum{v4x.x3.cells.r090.a6000.free.attempts.i000.matching_pairs}{1024} & \raggedright interior (accepted) & \raggedright rejected \tabularnewline
\raggedright \resultnum{v4x.x3.cells.r090.a6000.cell.a}{0.6}, \resultnum{v4x.x3.cells.r090.a6000.cell.r}{0.9} & \raggedright free share & \raggedright \resultnum{v4x.x3.cells.r090.a6000.free.attempts.i001.matching_pairs}{4096} & \raggedright interior (accepted) & \raggedright rejected \tabularnewline
\raggedright \resultnum{v4x.x3.cells.r090.a9000.cell.a}{0.9}, \resultnum{v4x.x3.cells.r090.a9000.cell.r}{0.9} & \raggedright known share & \raggedright \resultnum{v4x.x3.cells.r090.a9000.known.attempts.i000.matching_pairs}{1024} & \raggedright interior (accepted) & \raggedright passed \tabularnewline
\raggedright \resultnum{v4x.x3.cells.r090.a9000.cell.a}{0.9}, \resultnum{v4x.x3.cells.r090.a9000.cell.r}{0.9} & \raggedright free share & \raggedright \resultnum{v4x.x3.cells.r090.a9000.free.attempts.i000.matching_pairs}{1024} & \raggedright no accepted fit found & \raggedright not run (no fit) \tabularnewline
\botrule
\end{tabular}
\end{table*}

\clearpage

\subsection{Single-panel tests (X4a, X4b)}\label{sec:x4}

\paragraph{Two procedures.}
X4 asks how often B would be rejected for a single observed panel with influence. Both procedures, the fitted one included, score each panel with the frozen LDA direction trained against the true alternative of that cell, a direction that an analyst would not have; neither is a procedure an analyst could apply to observed data. The oracle single-panel test (X4a) is a known-parameter procedure: its critical value is the 0.95 quantile of \resultnum{v4x.x4.sizes.check_prose}{4,000} scores of the population-matched B, and it was applied to 1,000 new panels in each of the 81 cells against B. The fitted single-panel test (X4b) is a fitted-panel procedure: B is fitted to the lag-1 and lag-2 correlations of each single panel, with \resultnum{v4x.x4.sizes.analyst}{256} simulated panels per fitting evaluation and no validation; \resultnum{v4x.x4.sizes.nulls}{199} fresh B panels at the fitted parameter give a parametric-bootstrap p-value, and B is rejected when $p\le0.05$. X4b was declared only at $r=0.50$, 0.70 and 0.90, on the first 200 test panels of each cell; values outside this domain are not missing in-domain estimates. A failed single-panel fit counts as a non-rejection, and the success-conditioned rate is reported separately (Table~\ref{tab:v4x-x4-failures}). Tables~\ref{tab:v4x-x4a} and~\ref{tab:v4x-x4b} report the rejection rates, intervals and panel denominators, each with its own procedure's size flag. The bootstrap p-value formula establishes neither exact validity over the composite null nor an analyst's attainable power.

\paragraph{Size at zero strength.}
Each procedure's size is its rejection rate at $a=0$, read against the nominal level \resultnum{v4x.protocol.x4.level}{0.05} (Table~\ref{tab:v4x-x4-calibration}). The rate was compatible with the nominal level at \resultnum{v4x.x4a.count.calibrated}{8} oracle and \resultnum{v4x.x4b.count.calibrated}{1} fitted-panel zero-strength cells (machine label \texttt{calibrated}), above it at \resultnum{v4x.x4a.count.liberal}{1} oracle and \resultnum{v4x.x4b.count.liberal}{1} fitted-panel cells (\texttt{liberal}), and below it at \resultnum{v4x.x4a.count.conservative}{0} oracle and \resultnum{v4x.x4b.count.conservative}{1} fitted-panel cells (\texttt{conservative}). The oracle procedure was above nominal at $r=0.70$; the fitted-panel procedure was above nominal at $r=0.50$, compatible at 0.70 and below at 0.90. Compatibility applies to the specified cell and does not establish calibration over the composite null. Each flag is carried to the positive-strength rates of the same procedure at the same $r$ as a record of that zero-strength cell; it does not certify the size of the procedure at other parameters of B. For the fitted-panel procedure the declaration also names the quantity at each $r$: a rejection rate where the size is above nominal, a power where it is compatible and a lower bound on power where it is below. The tables keep these frozen labels, but every positive-strength entry is reported here as a rejection rate: the size was checked at one null configuration only, and a lower bound on power would need a specified valid procedure to compare with, which the declaration does not name.
The fitted-panel procedure is not a valid single-panel test at $r=0.50$: without any influence it rejected B in \resultnum{j.x4b.zero_rejection.r050}{0.300} of panels, with interval \resultnum{j.x4b.zero_rejection.r050.lower}{0.235}--\resultnum{j.x4b.zero_rejection.r050.upper}{0.365}, against the nominal level of \resultnum{v4x.protocol.x4.level}{0.05}. This zero-strength counterexample is the evidence that the procedure does not control size in this design. At very weak influence the two procedures also gave sharply different rejection rates: at $r=0.70$, where the fitted procedure's zero-strength rate was compatible with the nominal level (\resultnum{v4x.x4.x4b.r070.a0000.x4b.power_pu}{0.035}), it rejected B in \resultnum{v4x.x4.x4b.r070.a0010.x4b.power_pu}{0.33} of panels at $a=0.001$, a cell where the oracle procedure rejected in \resultnum{v4x.x4.x4a.r070.a0010.x4a.power_or}{0.051} and the reference AUC was \resultnum{j.x4a.r070.a0010.auc_ref}{0.504}. Those panels contain influence, so these are rejection rates under an alternative, not sizes. The oracle procedure is one particular rule, not an upper bound on what other procedures reject, and a reference AUC near 0.5 for its score does not mean that the two laws coincide. The source of the discrepancy was not isolated; candidates include fitting B to a single panel, a parametric bootstrap at a fitted trait weight at or near the edge of its parameter box, simulation precision and the treatment of failed fits. Read after the results became available: on the strength of the zero-strength counterexample, fitting a persistence model to one panel and calibrating the oracle score by a parametric bootstrap at the fitted parameter does not give a valid single-panel test in this design, and no valid single-panel procedure was evaluated.

\paragraph{Normal-theory reference.}
The normal-theory power at each cell's own reference AUC lay inside the oracle rejection interval in \resultnum{v4x.x4a.count.describes}{25} cells within the declared reference-AUC range of \resultnum{v4x.x4.summary.normal_theory_in_auc_range.range.i000}{0.55} to \resultnum{v4x.x4.summary.normal_theory_in_auc_range.range.i001}{0.95} (machine label \texttt{describes}) and above it in \resultnum{v4x.x4a.count.overstates}{1} cell (\texttt{overstates}; Table~\ref{tab:v4x-x4-normal}). This records interval compatibility in those cells; it does not validate normal theory.

\paragraph{Paired difference.}
The fitted-panel minus oracle rejection rate, on the same panels, had an interval entirely below zero in \resultnum{v4x.x4b.count.removes-power}{5} cells (``negative''; \texttt{removes-power}), inside the declared margin of $\pm\resultnum{v4x.x4.summary.single_panel_effect.margin}{0.05}$ in \resultnum{v4x.x4b.count.no-material-change}{2} cells (``within margin''; \texttt{no-material-change}), and met neither criterion in \resultnum{v4x.x4b.count.inconclusive}{10} cells (``neither''; frozen label \texttt{inconclusive}; Table~\ref{tab:v4x-x4-paired}). Most of the latter are intervals entirely above zero, where the fitted-panel procedure rejected more often than the oracle one. No reading is made at $r=0.50$, where the fitted-panel size was above nominal. The difference combines fitted-panel estimation, simulation precision and failure treatment and does not isolate the effect of fitting alone.
Read after the results became available: at the cell of the released first boundary against B, the oracle procedure rejected B in \resultnum{j.x4a.boundary_rejection.min}{0.171} (at $r=\resultnum{j.x4a.boundary_rejection.min.r}{0.55}$) to \resultnum{j.x4a.boundary_rejection.max}{0.368} (at $r=\resultnum{j.x4a.boundary_rejection.max.r}{0.90}$) of single panels for $r\ge0.55$, and in \resultnum{j.x4a.boundary_rejection.r050}{0.447} at $r=0.50$, where the boundary lies at a higher strength; its size was above nominal at $r=0.70$ (Table~\ref{tab:v4x-selected-comparisons}).

% Generated by v4_exploratory_export; edited during the journal revision
% (agent S, 2026-10-13): captions, headers and categorical labels translated
% into plain English. Every number keeps its \resultnum binding.
\begin{table*}[htbp]
\setlength{\tabcolsep}{2pt}
\caption{Oracle single-panel test (X4a), a known-parameter procedure: the share of 1,000 new panels with influence whose LDA score exceeds the 0.95 quantile of 4,000 scores of the fitted B (population-matched, known parameters), with bootstrap 95\% interval and the panel denominator. The discriminant direction is trained against the true alternative, which a real analyst does not know. Size at $a=0$: this procedure's zero-strength rejection rate at the same $r$ (Table~\ref{tab:v4x-x4-calibration}): compatible with 0.05 (machine label calibrated) or above 0.05 (liberal); the flag applies to every strength at that $r$. Fit gap: the frozen fit of B failed, so the cell has no LDA and no test.}\label{tab:v4x-x4a}
\centering\footnotesize
\begin{tabular}{@{}p{0.090\textwidth}p{0.090\textwidth}p{0.160\textwidth}p{0.240\textwidth}p{0.120\textwidth}p{0.220\textwidth}@{}}
\toprule
\raggedright $r$ & \raggedright $a$ & \raggedright \mbox{Rejection} \mbox{rate} & \raggedright \mbox{Interval} & \raggedright \mbox{Panels} & \raggedright \mbox{Size} \mbox{at} \mbox{$a=0$} \tabularnewline
\midrule
\raggedright \resultnum{v4x.x4.x4a.r050.a0000.cell.r}{0.5} & \raggedright \resultnum{v4x.x4.x4a.r050.a0000.cell.a}{0} & \raggedright \resultnum{v4x.x4.x4a.r050.a0000.x4a.power_or}{0.056} & \raggedright [\resultnum{v4x.x4.x4a.r050.a0000.x4a.power_or_interval_95.i000}{0.043}, \resultnum{v4x.x4.x4a.r050.a0000.x4a.power_or_interval_95.i001}{0.071}] & \raggedright \resultnum{v4x.x4.x4a.r050.a0000.x4a.test_panels}{1000} & \raggedright compatible \tabularnewline
\raggedright \resultnum{v4x.x4.x4a.r050.a0010.cell.r}{0.5} & \raggedright \resultnum{v4x.x4.x4a.r050.a0010.cell.a}{0.001} & \raggedright \resultnum{v4x.x4.x4a.r050.a0010.x4a.power_or}{0.067} & \raggedright [\resultnum{v4x.x4.x4a.r050.a0010.x4a.power_or_interval_95.i000}{0.051}, \resultnum{v4x.x4.x4a.r050.a0010.x4a.power_or_interval_95.i001}{0.083}] & \raggedright \resultnum{v4x.x4.x4a.r050.a0010.x4a.test_panels}{1000} & \raggedright compatible \tabularnewline
\raggedright \resultnum{v4x.x4.x4a.r050.a0030.cell.r}{0.5} & \raggedright \resultnum{v4x.x4.x4a.r050.a0030.cell.a}{0.003} & \raggedright \resultnum{v4x.x4.x4a.r050.a0030.x4a.power_or}{0.058} & \raggedright [\resultnum{v4x.x4.x4a.r050.a0030.x4a.power_or_interval_95.i000}{0.044}, \resultnum{v4x.x4.x4a.r050.a0030.x4a.power_or_interval_95.i001}{0.073}] & \raggedright \resultnum{v4x.x4.x4a.r050.a0030.x4a.test_panels}{1000} & \raggedright compatible \tabularnewline
\raggedright \resultnum{v4x.x4.x4a.r050.a0100.cell.r}{0.5} & \raggedright \resultnum{v4x.x4.x4a.r050.a0100.cell.a}{0.01} & \raggedright \resultnum{v4x.x4.x4a.r050.a0100.x4a.power_or}{0.059} & \raggedright [\resultnum{v4x.x4.x4a.r050.a0100.x4a.power_or_interval_95.i000}{0.044}, \resultnum{v4x.x4.x4a.r050.a0100.x4a.power_or_interval_95.i001}{0.075}] & \raggedright \resultnum{v4x.x4.x4a.r050.a0100.x4a.test_panels}{1000} & \raggedright compatible \tabularnewline
\raggedright \resultnum{v4x.x4.x4a.r050.a0300.cell.r}{0.5} & \raggedright \resultnum{v4x.x4.x4a.r050.a0300.cell.a}{0.03} & \raggedright \resultnum{v4x.x4.x4a.r050.a0300.x4a.power_or}{0.067} & \raggedright [\resultnum{v4x.x4.x4a.r050.a0300.x4a.power_or_interval_95.i000}{0.051}, \resultnum{v4x.x4.x4a.r050.a0300.x4a.power_or_interval_95.i001}{0.085}] & \raggedright \resultnum{v4x.x4.x4a.r050.a0300.x4a.test_panels}{1000} & \raggedright compatible \tabularnewline
\raggedright \resultnum{v4x.x4.x4a.r050.a1000.cell.r}{0.5} & \raggedright \resultnum{v4x.x4.x4a.r050.a1000.cell.a}{0.1} & \raggedright \resultnum{v4x.x4.x4a.r050.a1000.x4a.power_or}{0.142} & \raggedright [\resultnum{v4x.x4.x4a.r050.a1000.x4a.power_or_interval_95.i000}{0.12}, \resultnum{v4x.x4.x4a.r050.a1000.x4a.power_or_interval_95.i001}{0.165}] & \raggedright \resultnum{v4x.x4.x4a.r050.a1000.x4a.test_panels}{1000} & \raggedright compatible \tabularnewline
\raggedright \resultnum{v4x.x4.x4a.r050.a3000.cell.r}{0.5} & \raggedright \resultnum{v4x.x4.x4a.r050.a3000.cell.a}{0.3} & \raggedright \resultnum{v4x.x4.x4a.r050.a3000.x4a.power_or}{0.447} & \raggedright [\resultnum{v4x.x4.x4a.r050.a3000.x4a.power_or_interval_95.i000}{0.415}, \resultnum{v4x.x4.x4a.r050.a3000.x4a.power_or_interval_95.i001}{0.477025}] & \raggedright \resultnum{v4x.x4.x4a.r050.a3000.x4a.test_panels}{1000} & \raggedright compatible \tabularnewline
\raggedright \resultnum{v4x.x4.x4a.r050.a6000.cell.r}{0.5} & \raggedright \resultnum{v4x.x4.x4a.r050.a6000.cell.a}{0.6} & \raggedright \resultnum{v4x.x4.x4a.r050.a6000.x4a.power_or}{0.87} & \raggedright [\resultnum{v4x.x4.x4a.r050.a6000.x4a.power_or_interval_95.i000}{0.849975}, \resultnum{v4x.x4.x4a.r050.a6000.x4a.power_or_interval_95.i001}{0.891}] & \raggedright \resultnum{v4x.x4.x4a.r050.a6000.x4a.test_panels}{1000} & \raggedright compatible \tabularnewline
\raggedright \resultnum{v4x.x4.x4a.r050.a9000.cell.r}{0.5} & \raggedright \resultnum{v4x.x4.x4a.r050.a9000.cell.a}{0.9} & \raggedright \resultnum{v4x.x4.x4a.r050.a9000.x4a.power_or}{0.972} & \raggedright [\resultnum{v4x.x4.x4a.r050.a9000.x4a.power_or_interval_95.i000}{0.962}, \resultnum{v4x.x4.x4a.r050.a9000.x4a.power_or_interval_95.i001}{0.982}] & \raggedright \resultnum{v4x.x4.x4a.r050.a9000.x4a.test_panels}{1000} & \raggedright compatible \tabularnewline
\raggedright \resultnum{v4x.x4.x4a.r055.a0000.cell.r}{0.55} & \raggedright \resultnum{v4x.x4.x4a.r055.a0000.cell.a}{0} & \raggedright \resultnum{v4x.x4.x4a.r055.a0000.x4a.power_or}{0.047} & \raggedright [\resultnum{v4x.x4.x4a.r055.a0000.x4a.power_or_interval_95.i000}{0.035}, \resultnum{v4x.x4.x4a.r055.a0000.x4a.power_or_interval_95.i001}{0.060025}] & \raggedright \resultnum{v4x.x4.x4a.r055.a0000.x4a.test_panels}{1000} & \raggedright compatible \tabularnewline
\raggedright \resultnum{v4x.x4.x4a.r055.a0010.cell.r}{0.55} & \raggedright \resultnum{v4x.x4.x4a.r055.a0010.cell.a}{0.001} & \raggedright \resultnum{v4x.x4.x4a.r055.a0010.x4a.power_or}{0.056} & \raggedright [\resultnum{v4x.x4.x4a.r055.a0010.x4a.power_or_interval_95.i000}{0.042}, \resultnum{v4x.x4.x4a.r055.a0010.x4a.power_or_interval_95.i001}{0.069025}] & \raggedright \resultnum{v4x.x4.x4a.r055.a0010.x4a.test_panels}{1000} & \raggedright compatible \tabularnewline
\raggedright \resultnum{v4x.x4.x4a.r055.a0030.cell.r}{0.55} & \raggedright \resultnum{v4x.x4.x4a.r055.a0030.cell.a}{0.003} & \raggedright \resultnum{v4x.x4.x4a.r055.a0030.x4a.power_or}{0.044} & \raggedright [\resultnum{v4x.x4.x4a.r055.a0030.x4a.power_or_interval_95.i000}{0.032}, \resultnum{v4x.x4.x4a.r055.a0030.x4a.power_or_interval_95.i001}{0.056}] & \raggedright \resultnum{v4x.x4.x4a.r055.a0030.x4a.test_panels}{1000} & \raggedright compatible \tabularnewline
\raggedright \resultnum{v4x.x4.x4a.r055.a0100.cell.r}{0.55} & \raggedright \resultnum{v4x.x4.x4a.r055.a0100.cell.a}{0.01} & \raggedright \resultnum{v4x.x4.x4a.r055.a0100.x4a.power_or}{0.068} & \raggedright [\resultnum{v4x.x4.x4a.r055.a0100.x4a.power_or_interval_95.i000}{0.054}, \resultnum{v4x.x4.x4a.r055.a0100.x4a.power_or_interval_95.i001}{0.083}] & \raggedright \resultnum{v4x.x4.x4a.r055.a0100.x4a.test_panels}{1000} & \raggedright compatible \tabularnewline
\raggedright \resultnum{v4x.x4.x4a.r055.a0300.cell.r}{0.55} & \raggedright \resultnum{v4x.x4.x4a.r055.a0300.cell.a}{0.03} & \raggedright \resultnum{v4x.x4.x4a.r055.a0300.x4a.power_or}{0.083} & \raggedright [\resultnum{v4x.x4.x4a.r055.a0300.x4a.power_or_interval_95.i000}{0.066}, \resultnum{v4x.x4.x4a.r055.a0300.x4a.power_or_interval_95.i001}{0.1}] & \raggedright \resultnum{v4x.x4.x4a.r055.a0300.x4a.test_panels}{1000} & \raggedright compatible \tabularnewline
\raggedright \resultnum{v4x.x4.x4a.r055.a1000.cell.r}{0.55} & \raggedright \resultnum{v4x.x4.x4a.r055.a1000.cell.a}{0.1} & \raggedright \resultnum{v4x.x4.x4a.r055.a1000.x4a.power_or}{0.171} & \raggedright [\resultnum{v4x.x4.x4a.r055.a1000.x4a.power_or_interval_95.i000}{0.149}, \resultnum{v4x.x4.x4a.r055.a1000.x4a.power_or_interval_95.i001}{0.193}] & \raggedright \resultnum{v4x.x4.x4a.r055.a1000.x4a.test_panels}{1000} & \raggedright compatible \tabularnewline
\raggedright \resultnum{v4x.x4.x4a.r055.a3000.cell.r}{0.55} & \raggedright \resultnum{v4x.x4.x4a.r055.a3000.cell.a}{0.3} & \raggedright \resultnum{v4x.x4.x4a.r055.a3000.x4a.power_or}{0.506} & \raggedright [\resultnum{v4x.x4.x4a.r055.a3000.x4a.power_or_interval_95.i000}{0.476}, \resultnum{v4x.x4.x4a.r055.a3000.x4a.power_or_interval_95.i001}{0.537}] & \raggedright \resultnum{v4x.x4.x4a.r055.a3000.x4a.test_panels}{1000} & \raggedright compatible \tabularnewline
\raggedright \resultnum{v4x.x4.x4a.r055.a6000.cell.r}{0.55} & \raggedright \resultnum{v4x.x4.x4a.r055.a6000.cell.a}{0.6} & \raggedright \resultnum{v4x.x4.x4a.r055.a6000.x4a.power_or}{0.917} & \raggedright [\resultnum{v4x.x4.x4a.r055.a6000.x4a.power_or_interval_95.i000}{0.899}, \resultnum{v4x.x4.x4a.r055.a6000.x4a.power_or_interval_95.i001}{0.934}] & \raggedright \resultnum{v4x.x4.x4a.r055.a6000.x4a.test_panels}{1000} & \raggedright compatible \tabularnewline
\raggedright \resultnum{v4x.x4.x4a.r055.a9000.cell.r}{0.55} & \raggedright \resultnum{v4x.x4.x4a.r055.a9000.cell.a}{0.9} & \raggedright \resultnum{v4x.x4.x4a.r055.a9000.x4a.power_or}{0.988} & \raggedright [\resultnum{v4x.x4.x4a.r055.a9000.x4a.power_or_interval_95.i000}{0.981}, \resultnum{v4x.x4.x4a.r055.a9000.x4a.power_or_interval_95.i001}{0.994}] & \raggedright \resultnum{v4x.x4.x4a.r055.a9000.x4a.test_panels}{1000} & \raggedright compatible \tabularnewline
\raggedright \resultnum{v4x.x4.x4a.r060.a0000.cell.r}{0.6} & \raggedright \resultnum{v4x.x4.x4a.r060.a0000.cell.a}{0} & \raggedright \resultnum{v4x.x4.x4a.r060.a0000.x4a.power_or}{0.054} & \raggedright [\resultnum{v4x.x4.x4a.r060.a0000.x4a.power_or_interval_95.i000}{0.04}, \resultnum{v4x.x4.x4a.r060.a0000.x4a.power_or_interval_95.i001}{0.069}] & \raggedright \resultnum{v4x.x4.x4a.r060.a0000.x4a.test_panels}{1000} & \raggedright compatible \tabularnewline
\raggedright \resultnum{v4x.x4.x4a.r060.a0010.cell.r}{0.6} & \raggedright \resultnum{v4x.x4.x4a.r060.a0010.cell.a}{0.001} & \raggedright \resultnum{v4x.x4.x4a.r060.a0010.x4a.power_or}{0.062} & \raggedright [\resultnum{v4x.x4.x4a.r060.a0010.x4a.power_or_interval_95.i000}{0.048}, \resultnum{v4x.x4.x4a.r060.a0010.x4a.power_or_interval_95.i001}{0.076}] & \raggedright \resultnum{v4x.x4.x4a.r060.a0010.x4a.test_panels}{1000} & \raggedright compatible \tabularnewline
\raggedright \resultnum{v4x.x4.x4a.r060.a0030.cell.r}{0.6} & \raggedright \resultnum{v4x.x4.x4a.r060.a0030.cell.a}{0.003} & \raggedright \resultnum{v4x.x4.x4a.r060.a0030.x4a.power_or}{0.06} & \raggedright [\resultnum{v4x.x4.x4a.r060.a0030.x4a.power_or_interval_95.i000}{0.046}, \resultnum{v4x.x4.x4a.r060.a0030.x4a.power_or_interval_95.i001}{0.074}] & \raggedright \resultnum{v4x.x4.x4a.r060.a0030.x4a.test_panels}{1000} & \raggedright compatible \tabularnewline
\raggedright \resultnum{v4x.x4.x4a.r060.a0100.cell.r}{0.6} & \raggedright \resultnum{v4x.x4.x4a.r060.a0100.cell.a}{0.01} & \raggedright \resultnum{v4x.x4.x4a.r060.a0100.x4a.power_or}{0.069} & \raggedright [\resultnum{v4x.x4.x4a.r060.a0100.x4a.power_or_interval_95.i000}{0.054}, \resultnum{v4x.x4.x4a.r060.a0100.x4a.power_or_interval_95.i001}{0.084}] & \raggedright \resultnum{v4x.x4.x4a.r060.a0100.x4a.test_panels}{1000} & \raggedright compatible \tabularnewline
\raggedright \resultnum{v4x.x4.x4a.r060.a0300.cell.r}{0.6} & \raggedright \resultnum{v4x.x4.x4a.r060.a0300.cell.a}{0.03} & \raggedright \resultnum{v4x.x4.x4a.r060.a0300.x4a.power_or}{0.091} & \raggedright [\resultnum{v4x.x4.x4a.r060.a0300.x4a.power_or_interval_95.i000}{0.074}, \resultnum{v4x.x4.x4a.r060.a0300.x4a.power_or_interval_95.i001}{0.107}] & \raggedright \resultnum{v4x.x4.x4a.r060.a0300.x4a.test_panels}{1000} & \raggedright compatible \tabularnewline
\raggedright \resultnum{v4x.x4.x4a.r060.a1000.cell.r}{0.6} & \raggedright \resultnum{v4x.x4.x4a.r060.a1000.cell.a}{0.1} & \raggedright \resultnum{v4x.x4.x4a.r060.a1000.x4a.power_or}{0.205} & \raggedright [\resultnum{v4x.x4.x4a.r060.a1000.x4a.power_or_interval_95.i000}{0.18}, \resultnum{v4x.x4.x4a.r060.a1000.x4a.power_or_interval_95.i001}{0.229}] & \raggedright \resultnum{v4x.x4.x4a.r060.a1000.x4a.test_panels}{1000} & \raggedright compatible \tabularnewline
\raggedright \resultnum{v4x.x4.x4a.r060.a3000.cell.r}{0.6} & \raggedright \resultnum{v4x.x4.x4a.r060.a3000.cell.a}{0.3} & \raggedright \resultnum{v4x.x4.x4a.r060.a3000.x4a.power_or}{0.645} & \raggedright [\resultnum{v4x.x4.x4a.r060.a3000.x4a.power_or_interval_95.i000}{0.614}, \resultnum{v4x.x4.x4a.r060.a3000.x4a.power_or_interval_95.i001}{0.675}] & \raggedright \resultnum{v4x.x4.x4a.r060.a3000.x4a.test_panels}{1000} & \raggedright compatible \tabularnewline
\raggedright \resultnum{v4x.x4.x4a.r060.a6000.cell.r}{0.6} & \raggedright \resultnum{v4x.x4.x4a.r060.a6000.cell.a}{0.6} & \raggedright fit gap & \raggedright  & \raggedright  & \raggedright compatible \tabularnewline
\raggedright \resultnum{v4x.x4.x4a.r060.a9000.cell.r}{0.6} & \raggedright \resultnum{v4x.x4.x4a.r060.a9000.cell.a}{0.9} & \raggedright fit gap & \raggedright  & \raggedright  & \raggedright compatible \tabularnewline
\raggedright \resultnum{v4x.x4.x4a.r065.a0000.cell.r}{0.65} & \raggedright \resultnum{v4x.x4.x4a.r065.a0000.cell.a}{0} & \raggedright \resultnum{v4x.x4.x4a.r065.a0000.x4a.power_or}{0.05} & \raggedright [\resultnum{v4x.x4.x4a.r065.a0000.x4a.power_or_interval_95.i000}{0.037}, \resultnum{v4x.x4.x4a.r065.a0000.x4a.power_or_interval_95.i001}{0.064}] & \raggedright \resultnum{v4x.x4.x4a.r065.a0000.x4a.test_panels}{1000} & \raggedright compatible \tabularnewline
\raggedright \resultnum{v4x.x4.x4a.r065.a0010.cell.r}{0.65} & \raggedright \resultnum{v4x.x4.x4a.r065.a0010.cell.a}{0.001} & \raggedright \resultnum{v4x.x4.x4a.r065.a0010.x4a.power_or}{0.051} & \raggedright [\resultnum{v4x.x4.x4a.r065.a0010.x4a.power_or_interval_95.i000}{0.038}, \resultnum{v4x.x4.x4a.r065.a0010.x4a.power_or_interval_95.i001}{0.064}] & \raggedright \resultnum{v4x.x4.x4a.r065.a0010.x4a.test_panels}{1000} & \raggedright compatible \tabularnewline
\raggedright \resultnum{v4x.x4.x4a.r065.a0030.cell.r}{0.65} & \raggedright \resultnum{v4x.x4.x4a.r065.a0030.cell.a}{0.003} & \raggedright \resultnum{v4x.x4.x4a.r065.a0030.x4a.power_or}{0.054} & \raggedright [\resultnum{v4x.x4.x4a.r065.a0030.x4a.power_or_interval_95.i000}{0.040975}, \resultnum{v4x.x4.x4a.r065.a0030.x4a.power_or_interval_95.i001}{0.067}] & \raggedright \resultnum{v4x.x4.x4a.r065.a0030.x4a.test_panels}{1000} & \raggedright compatible \tabularnewline
\raggedright \resultnum{v4x.x4.x4a.r065.a0100.cell.r}{0.65} & \raggedright \resultnum{v4x.x4.x4a.r065.a0100.cell.a}{0.01} & \raggedright \resultnum{v4x.x4.x4a.r065.a0100.x4a.power_or}{0.065} & \raggedright [\resultnum{v4x.x4.x4a.r065.a0100.x4a.power_or_interval_95.i000}{0.051}, \resultnum{v4x.x4.x4a.r065.a0100.x4a.power_or_interval_95.i001}{0.08}] & \raggedright \resultnum{v4x.x4.x4a.r065.a0100.x4a.test_panels}{1000} & \raggedright compatible \tabularnewline
\raggedright \resultnum{v4x.x4.x4a.r065.a0300.cell.r}{0.65} & \raggedright \resultnum{v4x.x4.x4a.r065.a0300.cell.a}{0.03} & \raggedright \resultnum{v4x.x4.x4a.r065.a0300.x4a.power_or}{0.099} & \raggedright [\resultnum{v4x.x4.x4a.r065.a0300.x4a.power_or_interval_95.i000}{0.081}, \resultnum{v4x.x4.x4a.r065.a0300.x4a.power_or_interval_95.i001}{0.117}] & \raggedright \resultnum{v4x.x4.x4a.r065.a0300.x4a.test_panels}{1000} & \raggedright compatible \tabularnewline
\raggedright \resultnum{v4x.x4.x4a.r065.a1000.cell.r}{0.65} & \raggedright \resultnum{v4x.x4.x4a.r065.a1000.cell.a}{0.1} & \raggedright \resultnum{v4x.x4.x4a.r065.a1000.x4a.power_or}{0.209} & \raggedright [\resultnum{v4x.x4.x4a.r065.a1000.x4a.power_or_interval_95.i000}{0.184}, \resultnum{v4x.x4.x4a.r065.a1000.x4a.power_or_interval_95.i001}{0.234}] & \raggedright \resultnum{v4x.x4.x4a.r065.a1000.x4a.test_panels}{1000} & \raggedright compatible \tabularnewline
\raggedright \resultnum{v4x.x4.x4a.r065.a3000.cell.r}{0.65} & \raggedright \resultnum{v4x.x4.x4a.r065.a3000.cell.a}{0.3} & \raggedright \resultnum{v4x.x4.x4a.r065.a3000.x4a.power_or}{0.716} & \raggedright [\resultnum{v4x.x4.x4a.r065.a3000.x4a.power_or_interval_95.i000}{0.687}, \resultnum{v4x.x4.x4a.r065.a3000.x4a.power_or_interval_95.i001}{0.743}] & \raggedright \resultnum{v4x.x4.x4a.r065.a3000.x4a.test_panels}{1000} & \raggedright compatible \tabularnewline
\raggedright \resultnum{v4x.x4.x4a.r065.a6000.cell.r}{0.65} & \raggedright \resultnum{v4x.x4.x4a.r065.a6000.cell.a}{0.6} & \raggedright \resultnum{v4x.x4.x4a.r065.a6000.x4a.power_or}{0.99} & \raggedright [\resultnum{v4x.x4.x4a.r065.a6000.x4a.power_or_interval_95.i000}{0.984}, \resultnum{v4x.x4.x4a.r065.a6000.x4a.power_or_interval_95.i001}{0.996}] & \raggedright \resultnum{v4x.x4.x4a.r065.a6000.x4a.test_panels}{1000} & \raggedright compatible \tabularnewline
\raggedright \resultnum{v4x.x4.x4a.r065.a9000.cell.r}{0.65} & \raggedright \resultnum{v4x.x4.x4a.r065.a9000.cell.a}{0.9} & \raggedright \resultnum{v4x.x4.x4a.r065.a9000.x4a.power_or}{0.998} & \raggedright [\resultnum{v4x.x4.x4a.r065.a9000.x4a.power_or_interval_95.i000}{0.995}, \resultnum{v4x.x4.x4a.r065.a9000.x4a.power_or_interval_95.i001}{1}] & \raggedright \resultnum{v4x.x4.x4a.r065.a9000.x4a.test_panels}{1000} & \raggedright compatible \tabularnewline
\raggedright \resultnum{v4x.x4.x4a.r070.a0000.cell.r}{0.7} & \raggedright \resultnum{v4x.x4.x4a.r070.a0000.cell.a}{0} & \raggedright \resultnum{v4x.x4.x4a.r070.a0000.x4a.power_or}{0.069} & \raggedright [\resultnum{v4x.x4.x4a.r070.a0000.x4a.power_or_interval_95.i000}{0.055}, \resultnum{v4x.x4.x4a.r070.a0000.x4a.power_or_interval_95.i001}{0.086}] & \raggedright \resultnum{v4x.x4.x4a.r070.a0000.x4a.test_panels}{1000} & \raggedright above \tabularnewline
\raggedright \resultnum{v4x.x4.x4a.r070.a0010.cell.r}{0.7} & \raggedright \resultnum{v4x.x4.x4a.r070.a0010.cell.a}{0.001} & \raggedright \resultnum{v4x.x4.x4a.r070.a0010.x4a.power_or}{0.051} & \raggedright [\resultnum{v4x.x4.x4a.r070.a0010.x4a.power_or_interval_95.i000}{0.037}, \resultnum{v4x.x4.x4a.r070.a0010.x4a.power_or_interval_95.i001}{0.067}] & \raggedright \resultnum{v4x.x4.x4a.r070.a0010.x4a.test_panels}{1000} & \raggedright above \tabularnewline
\raggedright \resultnum{v4x.x4.x4a.r070.a0030.cell.r}{0.7} & \raggedright \resultnum{v4x.x4.x4a.r070.a0030.cell.a}{0.003} & \raggedright \resultnum{v4x.x4.x4a.r070.a0030.x4a.power_or}{0.05} & \raggedright [\resultnum{v4x.x4.x4a.r070.a0030.x4a.power_or_interval_95.i000}{0.036}, \resultnum{v4x.x4.x4a.r070.a0030.x4a.power_or_interval_95.i001}{0.065}] & \raggedright \resultnum{v4x.x4.x4a.r070.a0030.x4a.test_panels}{1000} & \raggedright above \tabularnewline
\raggedright \resultnum{v4x.x4.x4a.r070.a0100.cell.r}{0.7} & \raggedright \resultnum{v4x.x4.x4a.r070.a0100.cell.a}{0.01} & \raggedright \resultnum{v4x.x4.x4a.r070.a0100.x4a.power_or}{0.058} & \raggedright [\resultnum{v4x.x4.x4a.r070.a0100.x4a.power_or_interval_95.i000}{0.044}, \resultnum{v4x.x4.x4a.r070.a0100.x4a.power_or_interval_95.i001}{0.075}] & \raggedright \resultnum{v4x.x4.x4a.r070.a0100.x4a.test_panels}{1000} & \raggedright above \tabularnewline
\raggedright \resultnum{v4x.x4.x4a.r070.a0300.cell.r}{0.7} & \raggedright \resultnum{v4x.x4.x4a.r070.a0300.cell.a}{0.03} & \raggedright \resultnum{v4x.x4.x4a.r070.a0300.x4a.power_or}{0.089} & \raggedright [\resultnum{v4x.x4.x4a.r070.a0300.x4a.power_or_interval_95.i000}{0.071}, \resultnum{v4x.x4.x4a.r070.a0300.x4a.power_or_interval_95.i001}{0.108}] & \raggedright \resultnum{v4x.x4.x4a.r070.a0300.x4a.test_panels}{1000} & \raggedright above \tabularnewline
\raggedright \resultnum{v4x.x4.x4a.r070.a1000.cell.r}{0.7} & \raggedright \resultnum{v4x.x4.x4a.r070.a1000.cell.a}{0.1} & \raggedright \resultnum{v4x.x4.x4a.r070.a1000.x4a.power_or}{0.199} & \raggedright [\resultnum{v4x.x4.x4a.r070.a1000.x4a.power_or_interval_95.i000}{0.175}, \resultnum{v4x.x4.x4a.r070.a1000.x4a.power_or_interval_95.i001}{0.225}] & \raggedright \resultnum{v4x.x4.x4a.r070.a1000.x4a.test_panels}{1000} & \raggedright above \tabularnewline
\raggedright \resultnum{v4x.x4.x4a.r070.a3000.cell.r}{0.7} & \raggedright \resultnum{v4x.x4.x4a.r070.a3000.cell.a}{0.3} & \raggedright \resultnum{v4x.x4.x4a.r070.a3000.x4a.power_or}{0.759} & \raggedright [\resultnum{v4x.x4.x4a.r070.a3000.x4a.power_or_interval_95.i000}{0.732}, \resultnum{v4x.x4.x4a.r070.a3000.x4a.power_or_interval_95.i001}{0.785}] & \raggedright \resultnum{v4x.x4.x4a.r070.a3000.x4a.test_panels}{1000} & \raggedright above \tabularnewline
\raggedright \resultnum{v4x.x4.x4a.r070.a6000.cell.r}{0.7} & \raggedright \resultnum{v4x.x4.x4a.r070.a6000.cell.a}{0.6} & \raggedright fit gap & \raggedright  & \raggedright  & \raggedright above \tabularnewline
\raggedright \resultnum{v4x.x4.x4a.r070.a9000.cell.r}{0.7} & \raggedright \resultnum{v4x.x4.x4a.r070.a9000.cell.a}{0.9} & \raggedright \resultnum{v4x.x4.x4a.r070.a9000.x4a.power_or}{1} & \raggedright [\resultnum{v4x.x4.x4a.r070.a9000.x4a.power_or_interval_95.i000}{1}, \resultnum{v4x.x4.x4a.r070.a9000.x4a.power_or_interval_95.i001}{1}] & \raggedright \resultnum{v4x.x4.x4a.r070.a9000.x4a.test_panels}{1000} & \raggedright above \tabularnewline
\raggedright \resultnum{v4x.x4.x4a.r075.a0000.cell.r}{0.75} & \raggedright \resultnum{v4x.x4.x4a.r075.a0000.cell.a}{0} & \raggedright \resultnum{v4x.x4.x4a.r075.a0000.x4a.power_or}{0.052} & \raggedright [\resultnum{v4x.x4.x4a.r075.a0000.x4a.power_or_interval_95.i000}{0.039}, \resultnum{v4x.x4.x4a.r075.a0000.x4a.power_or_interval_95.i001}{0.066}] & \raggedright \resultnum{v4x.x4.x4a.r075.a0000.x4a.test_panels}{1000} & \raggedright compatible \tabularnewline
\raggedright \resultnum{v4x.x4.x4a.r075.a0010.cell.r}{0.75} & \raggedright \resultnum{v4x.x4.x4a.r075.a0010.cell.a}{0.001} & \raggedright \resultnum{v4x.x4.x4a.r075.a0010.x4a.power_or}{0.042} & \raggedright [\resultnum{v4x.x4.x4a.r075.a0010.x4a.power_or_interval_95.i000}{0.031}, \resultnum{v4x.x4.x4a.r075.a0010.x4a.power_or_interval_95.i001}{0.054}] & \raggedright \resultnum{v4x.x4.x4a.r075.a0010.x4a.test_panels}{1000} & \raggedright compatible \tabularnewline
\raggedright \resultnum{v4x.x4.x4a.r075.a0030.cell.r}{0.75} & \raggedright \resultnum{v4x.x4.x4a.r075.a0030.cell.a}{0.003} & \raggedright \resultnum{v4x.x4.x4a.r075.a0030.x4a.power_or}{0.066} & \raggedright [\resultnum{v4x.x4.x4a.r075.a0030.x4a.power_or_interval_95.i000}{0.051}, \resultnum{v4x.x4.x4a.r075.a0030.x4a.power_or_interval_95.i001}{0.081}] & \raggedright \resultnum{v4x.x4.x4a.r075.a0030.x4a.test_panels}{1000} & \raggedright compatible \tabularnewline
\raggedright \resultnum{v4x.x4.x4a.r075.a0100.cell.r}{0.75} & \raggedright \resultnum{v4x.x4.x4a.r075.a0100.cell.a}{0.01} & \raggedright \resultnum{v4x.x4.x4a.r075.a0100.x4a.power_or}{0.076} & \raggedright [\resultnum{v4x.x4.x4a.r075.a0100.x4a.power_or_interval_95.i000}{0.06}, \resultnum{v4x.x4.x4a.r075.a0100.x4a.power_or_interval_95.i001}{0.092}] & \raggedright \resultnum{v4x.x4.x4a.r075.a0100.x4a.test_panels}{1000} & \raggedright compatible \tabularnewline
\raggedright \resultnum{v4x.x4.x4a.r075.a0300.cell.r}{0.75} & \raggedright \resultnum{v4x.x4.x4a.r075.a0300.cell.a}{0.03} & \raggedright \resultnum{v4x.x4.x4a.r075.a0300.x4a.power_or}{0.121} & \raggedright [\resultnum{v4x.x4.x4a.r075.a0300.x4a.power_or_interval_95.i000}{0.102}, \resultnum{v4x.x4.x4a.r075.a0300.x4a.power_or_interval_95.i001}{0.141}] & \raggedright \resultnum{v4x.x4.x4a.r075.a0300.x4a.test_panels}{1000} & \raggedright compatible \tabularnewline
\raggedright \resultnum{v4x.x4.x4a.r075.a1000.cell.r}{0.75} & \raggedright \resultnum{v4x.x4.x4a.r075.a1000.cell.a}{0.1} & \raggedright \resultnum{v4x.x4.x4a.r075.a1000.x4a.power_or}{0.282} & \raggedright [\resultnum{v4x.x4.x4a.r075.a1000.x4a.power_or_interval_95.i000}{0.255}, \resultnum{v4x.x4.x4a.r075.a1000.x4a.power_or_interval_95.i001}{0.308025}] & \raggedright \resultnum{v4x.x4.x4a.r075.a1000.x4a.test_panels}{1000} & \raggedright compatible \tabularnewline
\raggedright \resultnum{v4x.x4.x4a.r075.a3000.cell.r}{0.75} & \raggedright \resultnum{v4x.x4.x4a.r075.a3000.cell.a}{0.3} & \raggedright \resultnum{v4x.x4.x4a.r075.a3000.x4a.power_or}{0.826} & \raggedright [\resultnum{v4x.x4.x4a.r075.a3000.x4a.power_or_interval_95.i000}{0.802}, \resultnum{v4x.x4.x4a.r075.a3000.x4a.power_or_interval_95.i001}{0.848}] & \raggedright \resultnum{v4x.x4.x4a.r075.a3000.x4a.test_panels}{1000} & \raggedright compatible \tabularnewline
\botrule
\end{tabular}
\end{table*}
\begin{table*}[htbp]
\setlength{\tabcolsep}{2pt}
\addtocounter{table}{-\csname @ne\endcsname}
\providecommand{\theHtable}{\thetable}
\renewcommand{\theHtable}{v4x-x4a-part1}
\caption{Oracle single-panel test (X4a) (continued, part 2 of 2)}\label{tab:v4x-x4a-part1}
\centering\footnotesize
\begin{tabular}{@{}p{0.090\textwidth}p{0.090\textwidth}p{0.160\textwidth}p{0.240\textwidth}p{0.120\textwidth}p{0.220\textwidth}@{}}
\toprule
\raggedright $r$ & \raggedright $a$ & \raggedright \mbox{Rejection} \mbox{rate} & \raggedright \mbox{Interval} & \raggedright \mbox{Panels} & \raggedright \mbox{Size} \mbox{at} \mbox{$a=0$} \tabularnewline
\midrule
\raggedright \resultnum{v4x.x4.x4a.r075.a6000.cell.r}{0.75} & \raggedright \resultnum{v4x.x4.x4a.r075.a6000.cell.a}{0.6} & \raggedright \resultnum{v4x.x4.x4a.r075.a6000.x4a.power_or}{0.996} & \raggedright [\resultnum{v4x.x4.x4a.r075.a6000.x4a.power_or_interval_95.i000}{0.992}, \resultnum{v4x.x4.x4a.r075.a6000.x4a.power_or_interval_95.i001}{0.999}] & \raggedright \resultnum{v4x.x4.x4a.r075.a6000.x4a.test_panels}{1000} & \raggedright compatible \tabularnewline
\raggedright \resultnum{v4x.x4.x4a.r075.a9000.cell.r}{0.75} & \raggedright \resultnum{v4x.x4.x4a.r075.a9000.cell.a}{0.9} & \raggedright \resultnum{v4x.x4.x4a.r075.a9000.x4a.power_or}{1} & \raggedright [\resultnum{v4x.x4.x4a.r075.a9000.x4a.power_or_interval_95.i000}{1}, \resultnum{v4x.x4.x4a.r075.a9000.x4a.power_or_interval_95.i001}{1}] & \raggedright \resultnum{v4x.x4.x4a.r075.a9000.x4a.test_panels}{1000} & \raggedright compatible \tabularnewline
\raggedright \resultnum{v4x.x4.x4a.r080.a0000.cell.r}{0.8} & \raggedright \resultnum{v4x.x4.x4a.r080.a0000.cell.a}{0} & \raggedright \resultnum{v4x.x4.x4a.r080.a0000.x4a.power_or}{0.054} & \raggedright [\resultnum{v4x.x4.x4a.r080.a0000.x4a.power_or_interval_95.i000}{0.04}, \resultnum{v4x.x4.x4a.r080.a0000.x4a.power_or_interval_95.i001}{0.069}] & \raggedright \resultnum{v4x.x4.x4a.r080.a0000.x4a.test_panels}{1000} & \raggedright compatible \tabularnewline
\raggedright \resultnum{v4x.x4.x4a.r080.a0010.cell.r}{0.8} & \raggedright \resultnum{v4x.x4.x4a.r080.a0010.cell.a}{0.001} & \raggedright \resultnum{v4x.x4.x4a.r080.a0010.x4a.power_or}{0.054} & \raggedright [\resultnum{v4x.x4.x4a.r080.a0010.x4a.power_or_interval_95.i000}{0.04}, \resultnum{v4x.x4.x4a.r080.a0010.x4a.power_or_interval_95.i001}{0.069}] & \raggedright \resultnum{v4x.x4.x4a.r080.a0010.x4a.test_panels}{1000} & \raggedright compatible \tabularnewline
\raggedright \resultnum{v4x.x4.x4a.r080.a0030.cell.r}{0.8} & \raggedright \resultnum{v4x.x4.x4a.r080.a0030.cell.a}{0.003} & \raggedright \resultnum{v4x.x4.x4a.r080.a0030.x4a.power_or}{0.057} & \raggedright [\resultnum{v4x.x4.x4a.r080.a0030.x4a.power_or_interval_95.i000}{0.043}, \resultnum{v4x.x4.x4a.r080.a0030.x4a.power_or_interval_95.i001}{0.071}] & \raggedright \resultnum{v4x.x4.x4a.r080.a0030.x4a.test_panels}{1000} & \raggedright compatible \tabularnewline
\raggedright \resultnum{v4x.x4.x4a.r080.a0100.cell.r}{0.8} & \raggedright \resultnum{v4x.x4.x4a.r080.a0100.cell.a}{0.01} & \raggedright \resultnum{v4x.x4.x4a.r080.a0100.x4a.power_or}{0.069} & \raggedright [\resultnum{v4x.x4.x4a.r080.a0100.x4a.power_or_interval_95.i000}{0.053}, \resultnum{v4x.x4.x4a.r080.a0100.x4a.power_or_interval_95.i001}{0.085}] & \raggedright \resultnum{v4x.x4.x4a.r080.a0100.x4a.test_panels}{1000} & \raggedright compatible \tabularnewline
\raggedright \resultnum{v4x.x4.x4a.r080.a0300.cell.r}{0.8} & \raggedright \resultnum{v4x.x4.x4a.r080.a0300.cell.a}{0.03} & \raggedright \resultnum{v4x.x4.x4a.r080.a0300.x4a.power_or}{0.107} & \raggedright [\resultnum{v4x.x4.x4a.r080.a0300.x4a.power_or_interval_95.i000}{0.089}, \resultnum{v4x.x4.x4a.r080.a0300.x4a.power_or_interval_95.i001}{0.126}] & \raggedright \resultnum{v4x.x4.x4a.r080.a0300.x4a.test_panels}{1000} & \raggedright compatible \tabularnewline
\raggedright \resultnum{v4x.x4.x4a.r080.a1000.cell.r}{0.8} & \raggedright \resultnum{v4x.x4.x4a.r080.a1000.cell.a}{0.1} & \raggedright \resultnum{v4x.x4.x4a.r080.a1000.x4a.power_or}{0.301} & \raggedright [\resultnum{v4x.x4.x4a.r080.a1000.x4a.power_or_interval_95.i000}{0.272975}, \resultnum{v4x.x4.x4a.r080.a1000.x4a.power_or_interval_95.i001}{0.329}] & \raggedright \resultnum{v4x.x4.x4a.r080.a1000.x4a.test_panels}{1000} & \raggedright compatible \tabularnewline
\raggedright \resultnum{v4x.x4.x4a.r080.a3000.cell.r}{0.8} & \raggedright \resultnum{v4x.x4.x4a.r080.a3000.cell.a}{0.3} & \raggedright \resultnum{v4x.x4.x4a.r080.a3000.x4a.power_or}{0.874} & \raggedright [\resultnum{v4x.x4.x4a.r080.a3000.x4a.power_or_interval_95.i000}{0.852}, \resultnum{v4x.x4.x4a.r080.a3000.x4a.power_or_interval_95.i001}{0.894}] & \raggedright \resultnum{v4x.x4.x4a.r080.a3000.x4a.test_panels}{1000} & \raggedright compatible \tabularnewline
\raggedright \resultnum{v4x.x4.x4a.r080.a6000.cell.r}{0.8} & \raggedright \resultnum{v4x.x4.x4a.r080.a6000.cell.a}{0.6} & \raggedright \resultnum{v4x.x4.x4a.r080.a6000.x4a.power_or}{0.999} & \raggedright [\resultnum{v4x.x4.x4a.r080.a6000.x4a.power_or_interval_95.i000}{0.997}, \resultnum{v4x.x4.x4a.r080.a6000.x4a.power_or_interval_95.i001}{1}] & \raggedright \resultnum{v4x.x4.x4a.r080.a6000.x4a.test_panels}{1000} & \raggedright compatible \tabularnewline
\raggedright \resultnum{v4x.x4.x4a.r080.a9000.cell.r}{0.8} & \raggedright \resultnum{v4x.x4.x4a.r080.a9000.cell.a}{0.9} & \raggedright fit gap & \raggedright  & \raggedright  & \raggedright compatible \tabularnewline
\raggedright \resultnum{v4x.x4.x4a.r085.a0000.cell.r}{0.85} & \raggedright \resultnum{v4x.x4.x4a.r085.a0000.cell.a}{0} & \raggedright \resultnum{v4x.x4.x4a.r085.a0000.x4a.power_or}{0.049} & \raggedright [\resultnum{v4x.x4.x4a.r085.a0000.x4a.power_or_interval_95.i000}{0.036}, \resultnum{v4x.x4.x4a.r085.a0000.x4a.power_or_interval_95.i001}{0.062}] & \raggedright \resultnum{v4x.x4.x4a.r085.a0000.x4a.test_panels}{1000} & \raggedright compatible \tabularnewline
\raggedright \resultnum{v4x.x4.x4a.r085.a0010.cell.r}{0.85} & \raggedright \resultnum{v4x.x4.x4a.r085.a0010.cell.a}{0.001} & \raggedright \resultnum{v4x.x4.x4a.r085.a0010.x4a.power_or}{0.051} & \raggedright [\resultnum{v4x.x4.x4a.r085.a0010.x4a.power_or_interval_95.i000}{0.038}, \resultnum{v4x.x4.x4a.r085.a0010.x4a.power_or_interval_95.i001}{0.066}] & \raggedright \resultnum{v4x.x4.x4a.r085.a0010.x4a.test_panels}{1000} & \raggedright compatible \tabularnewline
\raggedright \resultnum{v4x.x4.x4a.r085.a0030.cell.r}{0.85} & \raggedright \resultnum{v4x.x4.x4a.r085.a0030.cell.a}{0.003} & \raggedright \resultnum{v4x.x4.x4a.r085.a0030.x4a.power_or}{0.057} & \raggedright [\resultnum{v4x.x4.x4a.r085.a0030.x4a.power_or_interval_95.i000}{0.043}, \resultnum{v4x.x4.x4a.r085.a0030.x4a.power_or_interval_95.i001}{0.072}] & \raggedright \resultnum{v4x.x4.x4a.r085.a0030.x4a.test_panels}{1000} & \raggedright compatible \tabularnewline
\raggedright \resultnum{v4x.x4.x4a.r085.a0100.cell.r}{0.85} & \raggedright \resultnum{v4x.x4.x4a.r085.a0100.cell.a}{0.01} & \raggedright \resultnum{v4x.x4.x4a.r085.a0100.x4a.power_or}{0.067} & \raggedright [\resultnum{v4x.x4.x4a.r085.a0100.x4a.power_or_interval_95.i000}{0.052}, \resultnum{v4x.x4.x4a.r085.a0100.x4a.power_or_interval_95.i001}{0.083}] & \raggedright \resultnum{v4x.x4.x4a.r085.a0100.x4a.test_panels}{1000} & \raggedright compatible \tabularnewline
\raggedright \resultnum{v4x.x4.x4a.r085.a0300.cell.r}{0.85} & \raggedright \resultnum{v4x.x4.x4a.r085.a0300.cell.a}{0.03} & \raggedright \resultnum{v4x.x4.x4a.r085.a0300.x4a.power_or}{0.118} & \raggedright [\resultnum{v4x.x4.x4a.r085.a0300.x4a.power_or_interval_95.i000}{0.098}, \resultnum{v4x.x4.x4a.r085.a0300.x4a.power_or_interval_95.i001}{0.138}] & \raggedright \resultnum{v4x.x4.x4a.r085.a0300.x4a.test_panels}{1000} & \raggedright compatible \tabularnewline
\raggedright \resultnum{v4x.x4.x4a.r085.a1000.cell.r}{0.85} & \raggedright \resultnum{v4x.x4.x4a.r085.a1000.cell.a}{0.1} & \raggedright \resultnum{v4x.x4.x4a.r085.a1000.x4a.power_or}{0.321} & \raggedright [\resultnum{v4x.x4.x4a.r085.a1000.x4a.power_or_interval_95.i000}{0.293}, \resultnum{v4x.x4.x4a.r085.a1000.x4a.power_or_interval_95.i001}{0.350025}] & \raggedright \resultnum{v4x.x4.x4a.r085.a1000.x4a.test_panels}{1000} & \raggedright compatible \tabularnewline
\raggedright \resultnum{v4x.x4.x4a.r085.a3000.cell.r}{0.85} & \raggedright \resultnum{v4x.x4.x4a.r085.a3000.cell.a}{0.3} & \raggedright \resultnum{v4x.x4.x4a.r085.a3000.x4a.power_or}{0.928} & \raggedright [\resultnum{v4x.x4.x4a.r085.a3000.x4a.power_or_interval_95.i000}{0.912}, \resultnum{v4x.x4.x4a.r085.a3000.x4a.power_or_interval_95.i001}{0.944}] & \raggedright \resultnum{v4x.x4.x4a.r085.a3000.x4a.test_panels}{1000} & \raggedright compatible \tabularnewline
\raggedright \resultnum{v4x.x4.x4a.r085.a6000.cell.r}{0.85} & \raggedright \resultnum{v4x.x4.x4a.r085.a6000.cell.a}{0.6} & \raggedright \resultnum{v4x.x4.x4a.r085.a6000.x4a.power_or}{0.999} & \raggedright [\resultnum{v4x.x4.x4a.r085.a6000.x4a.power_or_interval_95.i000}{0.997}, \resultnum{v4x.x4.x4a.r085.a6000.x4a.power_or_interval_95.i001}{1}] & \raggedright \resultnum{v4x.x4.x4a.r085.a6000.x4a.test_panels}{1000} & \raggedright compatible \tabularnewline
\raggedright \resultnum{v4x.x4.x4a.r085.a9000.cell.r}{0.85} & \raggedright \resultnum{v4x.x4.x4a.r085.a9000.cell.a}{0.9} & \raggedright \resultnum{v4x.x4.x4a.r085.a9000.x4a.power_or}{1} & \raggedright [\resultnum{v4x.x4.x4a.r085.a9000.x4a.power_or_interval_95.i000}{1}, \resultnum{v4x.x4.x4a.r085.a9000.x4a.power_or_interval_95.i001}{1}] & \raggedright \resultnum{v4x.x4.x4a.r085.a9000.x4a.test_panels}{1000} & \raggedright compatible \tabularnewline
\raggedright \resultnum{v4x.x4.x4a.r090.a0000.cell.r}{0.9} & \raggedright \resultnum{v4x.x4.x4a.r090.a0000.cell.a}{0} & \raggedright \resultnum{v4x.x4.x4a.r090.a0000.x4a.power_or}{0.046} & \raggedright [\resultnum{v4x.x4.x4a.r090.a0000.x4a.power_or_interval_95.i000}{0.034}, \resultnum{v4x.x4.x4a.r090.a0000.x4a.power_or_interval_95.i001}{0.06}] & \raggedright \resultnum{v4x.x4.x4a.r090.a0000.x4a.test_panels}{1000} & \raggedright compatible \tabularnewline
\raggedright \resultnum{v4x.x4.x4a.r090.a0010.cell.r}{0.9} & \raggedright \resultnum{v4x.x4.x4a.r090.a0010.cell.a}{0.001} & \raggedright \resultnum{v4x.x4.x4a.r090.a0010.x4a.power_or}{0.048} & \raggedright [\resultnum{v4x.x4.x4a.r090.a0010.x4a.power_or_interval_95.i000}{0.035}, \resultnum{v4x.x4.x4a.r090.a0010.x4a.power_or_interval_95.i001}{0.063}] & \raggedright \resultnum{v4x.x4.x4a.r090.a0010.x4a.test_panels}{1000} & \raggedright compatible \tabularnewline
\raggedright \resultnum{v4x.x4.x4a.r090.a0030.cell.r}{0.9} & \raggedright \resultnum{v4x.x4.x4a.r090.a0030.cell.a}{0.003} & \raggedright \resultnum{v4x.x4.x4a.r090.a0030.x4a.power_or}{0.046} & \raggedright [\resultnum{v4x.x4.x4a.r090.a0030.x4a.power_or_interval_95.i000}{0.033}, \resultnum{v4x.x4.x4a.r090.a0030.x4a.power_or_interval_95.i001}{0.061}] & \raggedright \resultnum{v4x.x4.x4a.r090.a0030.x4a.test_panels}{1000} & \raggedright compatible \tabularnewline
\raggedright \resultnum{v4x.x4.x4a.r090.a0100.cell.r}{0.9} & \raggedright \resultnum{v4x.x4.x4a.r090.a0100.cell.a}{0.01} & \raggedright \resultnum{v4x.x4.x4a.r090.a0100.x4a.power_or}{0.066} & \raggedright [\resultnum{v4x.x4.x4a.r090.a0100.x4a.power_or_interval_95.i000}{0.051}, \resultnum{v4x.x4.x4a.r090.a0100.x4a.power_or_interval_95.i001}{0.083}] & \raggedright \resultnum{v4x.x4.x4a.r090.a0100.x4a.test_panels}{1000} & \raggedright compatible \tabularnewline
\raggedright \resultnum{v4x.x4.x4a.r090.a0300.cell.r}{0.9} & \raggedright \resultnum{v4x.x4.x4a.r090.a0300.cell.a}{0.03} & \raggedright \resultnum{v4x.x4.x4a.r090.a0300.x4a.power_or}{0.11} & \raggedright [\resultnum{v4x.x4.x4a.r090.a0300.x4a.power_or_interval_95.i000}{0.091}, \resultnum{v4x.x4.x4a.r090.a0300.x4a.power_or_interval_95.i001}{0.13}] & \raggedright \resultnum{v4x.x4.x4a.r090.a0300.x4a.test_panels}{1000} & \raggedright compatible \tabularnewline
\raggedright \resultnum{v4x.x4.x4a.r090.a1000.cell.r}{0.9} & \raggedright \resultnum{v4x.x4.x4a.r090.a1000.cell.a}{0.1} & \raggedright \resultnum{v4x.x4.x4a.r090.a1000.x4a.power_or}{0.368} & \raggedright [\resultnum{v4x.x4.x4a.r090.a1000.x4a.power_or_interval_95.i000}{0.338}, \resultnum{v4x.x4.x4a.r090.a1000.x4a.power_or_interval_95.i001}{0.399}] & \raggedright \resultnum{v4x.x4.x4a.r090.a1000.x4a.test_panels}{1000} & \raggedright compatible \tabularnewline
\raggedright \resultnum{v4x.x4.x4a.r090.a3000.cell.r}{0.9} & \raggedright \resultnum{v4x.x4.x4a.r090.a3000.cell.a}{0.3} & \raggedright \resultnum{v4x.x4.x4a.r090.a3000.x4a.power_or}{0.944} & \raggedright [\resultnum{v4x.x4.x4a.r090.a3000.x4a.power_or_interval_95.i000}{0.929}, \resultnum{v4x.x4.x4a.r090.a3000.x4a.power_or_interval_95.i001}{0.957}] & \raggedright \resultnum{v4x.x4.x4a.r090.a3000.x4a.test_panels}{1000} & \raggedright compatible \tabularnewline
\raggedright \resultnum{v4x.x4.x4a.r090.a6000.cell.r}{0.9} & \raggedright \resultnum{v4x.x4.x4a.r090.a6000.cell.a}{0.6} & \raggedright \resultnum{v4x.x4.x4a.r090.a6000.x4a.power_or}{1} & \raggedright [\resultnum{v4x.x4.x4a.r090.a6000.x4a.power_or_interval_95.i000}{1}, \resultnum{v4x.x4.x4a.r090.a6000.x4a.power_or_interval_95.i001}{1}] & \raggedright \resultnum{v4x.x4.x4a.r090.a6000.x4a.test_panels}{1000} & \raggedright compatible \tabularnewline
\raggedright \resultnum{v4x.x4.x4a.r090.a9000.cell.r}{0.9} & \raggedright \resultnum{v4x.x4.x4a.r090.a9000.cell.a}{0.9} & \raggedright \resultnum{v4x.x4.x4a.r090.a9000.x4a.power_or}{1} & \raggedright [\resultnum{v4x.x4.x4a.r090.a9000.x4a.power_or_interval_95.i000}{1}, \resultnum{v4x.x4.x4a.r090.a9000.x4a.power_or_interval_95.i001}{1}] & \raggedright \resultnum{v4x.x4.x4a.r090.a9000.x4a.test_panels}{1000} & \raggedright compatible \tabularnewline
\botrule
\end{tabular}
\end{table*}

% Generated by v4_exploratory_export; edited during the journal revision
% (agent S, 2026-10-13): captions, headers and categorical labels translated
% into plain English. Every number keeps its \resultnum binding.
\begin{table*}[htbp]
\setlength{\tabcolsep}{2pt}
\caption{Fitted single-panel test (X4b), a fitted-panel procedure declared only at $r=0.50$, 0.70 and 0.90: B is fitted to the lag-1 and lag-2 correlations of each single panel with influence (256 simulated panels per evaluation, no validation), 199 new B panels at the fitted parameter give a parametric-bootstrap p-value for the same oracle LDA score, and B is rejected when $p\le0.05$; failed single-panel fits count as non-rejections (Table~\ref{tab:v4x-x4-failures}). Size at $a=0$ and quantity: this procedure's zero-strength rejection rate at the same $r$ (Table~\ref{tab:v4x-x4-calibration}) and the label the declaration then gives to the positive-strength rates: rejection rate where the size is above 0.05 (machine labels liberal, rejection-rate), power where it is compatible with 0.05 (calibrated, power), power as a lower bound where it is below 0.05 (conservative, power-lower-bound). At $r=0.50$ the procedure rejected B in a share of \resultnum{j.x4b.zero_rejection.r050}{0.300} of panels without influence, so it is not a valid test there; its rates at that share are not powers. The frozen labels are kept, but every entry is a rejection rate: the size was checked at one null configuration only, and the label ``power, lower bound'' names no valid procedure whose power it would bound.}\label{tab:v4x-x4b}
\centering\footnotesize
\begin{tabular}{@{}p{0.080\textwidth}p{0.080\textwidth}p{0.149\textwidth}p{0.219\textwidth}p{0.100\textwidth}p{0.309\textwidth}@{}}
\toprule
\raggedright $r$ & \raggedright $a$ & \raggedright \mbox{Rejection} \mbox{rate} & \raggedright \mbox{Interval} & \raggedright \mbox{Panels} & \raggedright \mbox{Size} \mbox{at} \mbox{$a=0$;} \mbox{quantity} \tabularnewline
\midrule
\raggedright \resultnum{v4x.x4.x4b.r050.a0000.cell.r}{0.5} & \raggedright \resultnum{v4x.x4.x4b.r050.a0000.cell.a}{0} & \raggedright \resultnum{v4x.x4.x4b.r050.a0000.x4b.power_pu}{0.3} & \raggedright [\resultnum{v4x.x4.x4b.r050.a0000.x4b.power_pu_interval_95.i000}{0.235}, \resultnum{v4x.x4.x4b.r050.a0000.x4b.power_pu_interval_95.i001}{0.365}] & \raggedright \resultnum{v4x.x4.x4b.r050.a0000.x4b.test_panels}{200} & \raggedright above: rate \tabularnewline
\raggedright \resultnum{v4x.x4.x4b.r050.a0010.cell.r}{0.5} & \raggedright \resultnum{v4x.x4.x4b.r050.a0010.cell.a}{0.001} & \raggedright \resultnum{v4x.x4.x4b.r050.a0010.x4b.power_pu}{0.445} & \raggedright [\resultnum{v4x.x4.x4b.r050.a0010.x4b.power_pu_interval_95.i000}{0.37}, \resultnum{v4x.x4.x4b.r050.a0010.x4b.power_pu_interval_95.i001}{0.52}] & \raggedright \resultnum{v4x.x4.x4b.r050.a0010.x4b.test_panels}{200} & \raggedright above: rate \tabularnewline
\raggedright \resultnum{v4x.x4.x4b.r050.a0030.cell.r}{0.5} & \raggedright \resultnum{v4x.x4.x4b.r050.a0030.cell.a}{0.003} & \raggedright \resultnum{v4x.x4.x4b.r050.a0030.x4b.power_pu}{0.43} & \raggedright [\resultnum{v4x.x4.x4b.r050.a0030.x4b.power_pu_interval_95.i000}{0.365}, \resultnum{v4x.x4.x4b.r050.a0030.x4b.power_pu_interval_95.i001}{0.5}] & \raggedright \resultnum{v4x.x4.x4b.r050.a0030.x4b.test_panels}{200} & \raggedright above: rate \tabularnewline
\raggedright \resultnum{v4x.x4.x4b.r050.a0100.cell.r}{0.5} & \raggedright \resultnum{v4x.x4.x4b.r050.a0100.cell.a}{0.01} & \raggedright \resultnum{v4x.x4.x4b.r050.a0100.x4b.power_pu}{0.42} & \raggedright [\resultnum{v4x.x4.x4b.r050.a0100.x4b.power_pu_interval_95.i000}{0.355}, \resultnum{v4x.x4.x4b.r050.a0100.x4b.power_pu_interval_95.i001}{0.485}] & \raggedright \resultnum{v4x.x4.x4b.r050.a0100.x4b.test_panels}{200} & \raggedright above: rate \tabularnewline
\raggedright \resultnum{v4x.x4.x4b.r050.a0300.cell.r}{0.5} & \raggedright \resultnum{v4x.x4.x4b.r050.a0300.cell.a}{0.03} & \raggedright \resultnum{v4x.x4.x4b.r050.a0300.x4b.power_pu}{0.4} & \raggedright [\resultnum{v4x.x4.x4b.r050.a0300.x4b.power_pu_interval_95.i000}{0.34}, \resultnum{v4x.x4.x4b.r050.a0300.x4b.power_pu_interval_95.i001}{0.46}] & \raggedright \resultnum{v4x.x4.x4b.r050.a0300.x4b.test_panels}{200} & \raggedright above: rate \tabularnewline
\raggedright \resultnum{v4x.x4.x4b.r050.a1000.cell.r}{0.5} & \raggedright \resultnum{v4x.x4.x4b.r050.a1000.cell.a}{0.1} & \raggedright \resultnum{v4x.x4.x4b.r050.a1000.x4b.power_pu}{0.415} & \raggedright [\resultnum{v4x.x4.x4b.r050.a1000.x4b.power_pu_interval_95.i000}{0.345}, \resultnum{v4x.x4.x4b.r050.a1000.x4b.power_pu_interval_95.i001}{0.49}] & \raggedright \resultnum{v4x.x4.x4b.r050.a1000.x4b.test_panels}{200} & \raggedright above: rate \tabularnewline
\raggedright \resultnum{v4x.x4.x4b.r050.a3000.cell.r}{0.5} & \raggedright \resultnum{v4x.x4.x4b.r050.a3000.cell.a}{0.3} & \raggedright \resultnum{v4x.x4.x4b.r050.a3000.x4b.power_pu}{0.5} & \raggedright [\resultnum{v4x.x4.x4b.r050.a3000.x4b.power_pu_interval_95.i000}{0.425}, \resultnum{v4x.x4.x4b.r050.a3000.x4b.power_pu_interval_95.i001}{0.575}] & \raggedright \resultnum{v4x.x4.x4b.r050.a3000.x4b.test_panels}{200} & \raggedright above: rate \tabularnewline
\raggedright \resultnum{v4x.x4.x4b.r050.a6000.cell.r}{0.5} & \raggedright \resultnum{v4x.x4.x4b.r050.a6000.cell.a}{0.6} & \raggedright \resultnum{v4x.x4.x4b.r050.a6000.x4b.power_pu}{0.49} & \raggedright [\resultnum{v4x.x4.x4b.r050.a6000.x4b.power_pu_interval_95.i000}{0.42}, \resultnum{v4x.x4.x4b.r050.a6000.x4b.power_pu_interval_95.i001}{0.565}] & \raggedright \resultnum{v4x.x4.x4b.r050.a6000.x4b.test_panels}{200} & \raggedright above: rate \tabularnewline
\raggedright \resultnum{v4x.x4.x4b.r050.a9000.cell.r}{0.5} & \raggedright \resultnum{v4x.x4.x4b.r050.a9000.cell.a}{0.9} & \raggedright \resultnum{v4x.x4.x4b.r050.a9000.x4b.power_pu}{0.29} & \raggedright [\resultnum{v4x.x4.x4b.r050.a9000.x4b.power_pu_interval_95.i000}{0.225}, \resultnum{v4x.x4.x4b.r050.a9000.x4b.power_pu_interval_95.i001}{0.355}] & \raggedright \resultnum{v4x.x4.x4b.r050.a9000.x4b.test_panels}{200} & \raggedright above: rate \tabularnewline
\raggedright \resultnum{v4x.x4.x4b.r070.a0000.cell.r}{0.7} & \raggedright \resultnum{v4x.x4.x4b.r070.a0000.cell.a}{0} & \raggedright \resultnum{v4x.x4.x4b.r070.a0000.x4b.power_pu}{0.035} & \raggedright [\resultnum{v4x.x4.x4b.r070.a0000.x4b.power_pu_interval_95.i000}{0.01}, \resultnum{v4x.x4.x4b.r070.a0000.x4b.power_pu_interval_95.i001}{0.060125}] & \raggedright \resultnum{v4x.x4.x4b.r070.a0000.x4b.test_panels}{200} & \raggedright compatible: power \tabularnewline
\raggedright \resultnum{v4x.x4.x4b.r070.a0010.cell.r}{0.7} & \raggedright \resultnum{v4x.x4.x4b.r070.a0010.cell.a}{0.001} & \raggedright \resultnum{v4x.x4.x4b.r070.a0010.x4b.power_pu}{0.33} & \raggedright [\resultnum{v4x.x4.x4b.r070.a0010.x4b.power_pu_interval_95.i000}{0.27}, \resultnum{v4x.x4.x4b.r070.a0010.x4b.power_pu_interval_95.i001}{0.395}] & \raggedright \resultnum{v4x.x4.x4b.r070.a0010.x4b.test_panels}{200} & \raggedright compatible: power \tabularnewline
\raggedright \resultnum{v4x.x4.x4b.r070.a0030.cell.r}{0.7} & \raggedright \resultnum{v4x.x4.x4b.r070.a0030.cell.a}{0.003} & \raggedright \resultnum{v4x.x4.x4b.r070.a0030.x4b.power_pu}{0.335} & \raggedright [\resultnum{v4x.x4.x4b.r070.a0030.x4b.power_pu_interval_95.i000}{0.275}, \resultnum{v4x.x4.x4b.r070.a0030.x4b.power_pu_interval_95.i001}{0.4}] & \raggedright \resultnum{v4x.x4.x4b.r070.a0030.x4b.test_panels}{200} & \raggedright compatible: power \tabularnewline
\raggedright \resultnum{v4x.x4.x4b.r070.a0100.cell.r}{0.7} & \raggedright \resultnum{v4x.x4.x4b.r070.a0100.cell.a}{0.01} & \raggedright \resultnum{v4x.x4.x4b.r070.a0100.x4b.power_pu}{0.365} & \raggedright [\resultnum{v4x.x4.x4b.r070.a0100.x4b.power_pu_interval_95.i000}{0.295}, \resultnum{v4x.x4.x4b.r070.a0100.x4b.power_pu_interval_95.i001}{0.435}] & \raggedright \resultnum{v4x.x4.x4b.r070.a0100.x4b.test_panels}{200} & \raggedright compatible: power \tabularnewline
\raggedright \resultnum{v4x.x4.x4b.r070.a0300.cell.r}{0.7} & \raggedright \resultnum{v4x.x4.x4b.r070.a0300.cell.a}{0.03} & \raggedright \resultnum{v4x.x4.x4b.r070.a0300.x4b.power_pu}{0.355} & \raggedright [\resultnum{v4x.x4.x4b.r070.a0300.x4b.power_pu_interval_95.i000}{0.285}, \resultnum{v4x.x4.x4b.r070.a0300.x4b.power_pu_interval_95.i001}{0.420125}] & \raggedright \resultnum{v4x.x4.x4b.r070.a0300.x4b.test_panels}{200} & \raggedright compatible: power \tabularnewline
\raggedright \resultnum{v4x.x4.x4b.r070.a1000.cell.r}{0.7} & \raggedright \resultnum{v4x.x4.x4b.r070.a1000.cell.a}{0.1} & \raggedright \resultnum{v4x.x4.x4b.r070.a1000.x4b.power_pu}{0.385} & \raggedright [\resultnum{v4x.x4.x4b.r070.a1000.x4b.power_pu_interval_95.i000}{0.32}, \resultnum{v4x.x4.x4b.r070.a1000.x4b.power_pu_interval_95.i001}{0.455}] & \raggedright \resultnum{v4x.x4.x4b.r070.a1000.x4b.test_panels}{200} & \raggedright compatible: power \tabularnewline
\raggedright \resultnum{v4x.x4.x4b.r070.a3000.cell.r}{0.7} & \raggedright \resultnum{v4x.x4.x4b.r070.a3000.cell.a}{0.3} & \raggedright \resultnum{v4x.x4.x4b.r070.a3000.x4b.power_pu}{0.515} & \raggedright [\resultnum{v4x.x4.x4b.r070.a3000.x4b.power_pu_interval_95.i000}{0.45}, \resultnum{v4x.x4.x4b.r070.a3000.x4b.power_pu_interval_95.i001}{0.585125}] & \raggedright \resultnum{v4x.x4.x4b.r070.a3000.x4b.test_panels}{200} & \raggedright compatible: power \tabularnewline
\raggedright \resultnum{v4x.x4.x4b.r070.a6000.cell.r}{0.7} & \raggedright \resultnum{v4x.x4.x4b.r070.a6000.cell.a}{0.6} & \raggedright fit gap & \raggedright  & \raggedright  & \raggedright compatible: power \tabularnewline
\raggedright \resultnum{v4x.x4.x4b.r070.a9000.cell.r}{0.7} & \raggedright \resultnum{v4x.x4.x4b.r070.a9000.cell.a}{0.9} & \raggedright \resultnum{v4x.x4.x4b.r070.a9000.x4b.power_pu}{0.415} & \raggedright [\resultnum{v4x.x4.x4b.r070.a9000.x4b.power_pu_interval_95.i000}{0.35}, \resultnum{v4x.x4.x4b.r070.a9000.x4b.power_pu_interval_95.i001}{0.485}] & \raggedright \resultnum{v4x.x4.x4b.r070.a9000.x4b.test_panels}{200} & \raggedright compatible: power \tabularnewline
\raggedright \resultnum{v4x.x4.x4b.r090.a0000.cell.r}{0.9} & \raggedright \resultnum{v4x.x4.x4b.r090.a0000.cell.a}{0} & \raggedright \resultnum{v4x.x4.x4b.r090.a0000.x4b.power_pu}{0.02} & \raggedright [\resultnum{v4x.x4.x4b.r090.a0000.x4b.power_pu_interval_95.i000}{0.005}, \resultnum{v4x.x4.x4b.r090.a0000.x4b.power_pu_interval_95.i001}{0.04}] & \raggedright \resultnum{v4x.x4.x4b.r090.a0000.x4b.test_panels}{200} & \raggedright below: power, lower bound \tabularnewline
\raggedright \resultnum{v4x.x4.x4b.r090.a0010.cell.r}{0.9} & \raggedright \resultnum{v4x.x4.x4b.r090.a0010.cell.a}{0.001} & \raggedright \resultnum{v4x.x4.x4b.r090.a0010.x4b.power_pu}{0.16} & \raggedright [\resultnum{v4x.x4.x4b.r090.a0010.x4b.power_pu_interval_95.i000}{0.11}, \resultnum{v4x.x4.x4b.r090.a0010.x4b.power_pu_interval_95.i001}{0.215}] & \raggedright \resultnum{v4x.x4.x4b.r090.a0010.x4b.test_panels}{200} & \raggedright below: power, lower bound \tabularnewline
\raggedright \resultnum{v4x.x4.x4b.r090.a0030.cell.r}{0.9} & \raggedright \resultnum{v4x.x4.x4b.r090.a0030.cell.a}{0.003} & \raggedright \resultnum{v4x.x4.x4b.r090.a0030.x4b.power_pu}{0.205} & \raggedright [\resultnum{v4x.x4.x4b.r090.a0030.x4b.power_pu_interval_95.i000}{0.15}, \resultnum{v4x.x4.x4b.r090.a0030.x4b.power_pu_interval_95.i001}{0.27}] & \raggedright \resultnum{v4x.x4.x4b.r090.a0030.x4b.test_panels}{200} & \raggedright below: power, lower bound \tabularnewline
\raggedright \resultnum{v4x.x4.x4b.r090.a0100.cell.r}{0.9} & \raggedright \resultnum{v4x.x4.x4b.r090.a0100.cell.a}{0.01} & \raggedright \resultnum{v4x.x4.x4b.r090.a0100.x4b.power_pu}{0.23} & \raggedright [\resultnum{v4x.x4.x4b.r090.a0100.x4b.power_pu_interval_95.i000}{0.17}, \resultnum{v4x.x4.x4b.r090.a0100.x4b.power_pu_interval_95.i001}{0.295}] & \raggedright \resultnum{v4x.x4.x4b.r090.a0100.x4b.test_panels}{200} & \raggedright below: power, lower bound \tabularnewline
\raggedright \resultnum{v4x.x4.x4b.r090.a0300.cell.r}{0.9} & \raggedright \resultnum{v4x.x4.x4b.r090.a0300.cell.a}{0.03} & \raggedright \resultnum{v4x.x4.x4b.r090.a0300.x4b.power_pu}{0.29} & \raggedright [\resultnum{v4x.x4.x4b.r090.a0300.x4b.power_pu_interval_95.i000}{0.23}, \resultnum{v4x.x4.x4b.r090.a0300.x4b.power_pu_interval_95.i001}{0.35}] & \raggedright \resultnum{v4x.x4.x4b.r090.a0300.x4b.test_panels}{200} & \raggedright below: power, lower bound \tabularnewline
\raggedright \resultnum{v4x.x4.x4b.r090.a1000.cell.r}{0.9} & \raggedright \resultnum{v4x.x4.x4b.r090.a1000.cell.a}{0.1} & \raggedright \resultnum{v4x.x4.x4b.r090.a1000.x4b.power_pu}{0.415} & \raggedright [\resultnum{v4x.x4.x4b.r090.a1000.x4b.power_pu_interval_95.i000}{0.35}, \resultnum{v4x.x4.x4b.r090.a1000.x4b.power_pu_interval_95.i001}{0.485}] & \raggedright \resultnum{v4x.x4.x4b.r090.a1000.x4b.test_panels}{200} & \raggedright below: power, lower bound \tabularnewline
\raggedright \resultnum{v4x.x4.x4b.r090.a3000.cell.r}{0.9} & \raggedright \resultnum{v4x.x4.x4b.r090.a3000.cell.a}{0.3} & \raggedright \resultnum{v4x.x4.x4b.r090.a3000.x4b.power_pu}{0.74} & \raggedright [\resultnum{v4x.x4.x4b.r090.a3000.x4b.power_pu_interval_95.i000}{0.675}, \resultnum{v4x.x4.x4b.r090.a3000.x4b.power_pu_interval_95.i001}{0.805}] & \raggedright \resultnum{v4x.x4.x4b.r090.a3000.x4b.test_panels}{200} & \raggedright below: power, lower bound \tabularnewline
\raggedright \resultnum{v4x.x4.x4b.r090.a6000.cell.r}{0.9} & \raggedright \resultnum{v4x.x4.x4b.r090.a6000.cell.a}{0.6} & \raggedright \resultnum{v4x.x4.x4b.r090.a6000.x4b.power_pu}{0.915} & \raggedright [\resultnum{v4x.x4.x4b.r090.a6000.x4b.power_pu_interval_95.i000}{0.875}, \resultnum{v4x.x4.x4b.r090.a6000.x4b.power_pu_interval_95.i001}{0.95}] & \raggedright \resultnum{v4x.x4.x4b.r090.a6000.x4b.test_panels}{200} & \raggedright below: power, lower bound \tabularnewline
\raggedright \resultnum{v4x.x4.x4b.r090.a9000.cell.r}{0.9} & \raggedright \resultnum{v4x.x4.x4b.r090.a9000.cell.a}{0.9} & \raggedright \resultnum{v4x.x4.x4b.r090.a9000.x4b.power_pu}{0.74} & \raggedright [\resultnum{v4x.x4.x4b.r090.a9000.x4b.power_pu_interval_95.i000}{0.675}, \resultnum{v4x.x4.x4b.r090.a9000.x4b.power_pu_interval_95.i001}{0.8}] & \raggedright \resultnum{v4x.x4.x4b.r090.a9000.x4b.test_panels}{200} & \raggedright below: power, lower bound \tabularnewline
\botrule
\end{tabular}
\end{table*}

% Generated by v4_exploratory_export; edited during the journal revision
% (agent S, 2026-10-13): captions, headers and categorical labels translated
% into plain English. Every number keeps its \resultnum binding.
\begin{table*}[htbp]
\setlength{\tabcolsep}{2pt}
\caption{Size of the two single-panel procedures: rejection rate at zero strength with bootstrap 95\% interval and panel denominator. Flag under the declared rule: compatible with 0.05 when the interval contains the nominal level 0.05 (machine label calibrated), above 0.05 when its lower end exceeds 0.05 (liberal), below 0.05 when its upper end is below 0.05 (conservative). Compatibility applies to that zero-strength cell only and does not establish calibration over the composite null; the flag is carried to every positive-strength rate of the same procedure at the same $r$ (Tables~\ref{tab:v4x-x4a}, \ref{tab:v4x-x4b}, \ref{tab:v4x-x4-paired} and~\ref{tab:v4x-selected-comparisons}).}\label{tab:v4x-x4-calibration}
\centering\footnotesize
\begin{tabular}{@{}p{0.190\textwidth}p{0.080\textwidth}p{0.190\textwidth}p{0.100\textwidth}p{0.260\textwidth}p{0.100\textwidth}@{}}
\toprule
\raggedright \mbox{Procedure} & \raggedright $r$ & \raggedright \mbox{Size} \mbox{flag} & \raggedright \mbox{Rate} & \raggedright \mbox{Interval} & \raggedright \mbox{Panels} \tabularnewline
\midrule
\raggedright oracle (X4a) & \raggedright \resultnum{v4x.x4.x4a.r050.a0000.cell.r}{0.5} & \raggedright compatible & \raggedright \resultnum{v4x.x4.x4a.r050.a0000.x4a.power_or}{0.056} & \raggedright [\resultnum{v4x.x4.x4a.r050.a0000.x4a.power_or_interval_95.i000}{0.043}, \resultnum{v4x.x4.x4a.r050.a0000.x4a.power_or_interval_95.i001}{0.071}] & \raggedright \resultnum{v4x.x4.x4a.r050.a0000.x4a.test_panels}{1000} \tabularnewline
\raggedright oracle (X4a) & \raggedright \resultnum{v4x.x4.x4a.r055.a0000.cell.r}{0.55} & \raggedright compatible & \raggedright \resultnum{v4x.x4.x4a.r055.a0000.x4a.power_or}{0.047} & \raggedright [\resultnum{v4x.x4.x4a.r055.a0000.x4a.power_or_interval_95.i000}{0.035}, \resultnum{v4x.x4.x4a.r055.a0000.x4a.power_or_interval_95.i001}{0.060025}] & \raggedright \resultnum{v4x.x4.x4a.r055.a0000.x4a.test_panels}{1000} \tabularnewline
\raggedright oracle (X4a) & \raggedright \resultnum{v4x.x4.x4a.r060.a0000.cell.r}{0.6} & \raggedright compatible & \raggedright \resultnum{v4x.x4.x4a.r060.a0000.x4a.power_or}{0.054} & \raggedright [\resultnum{v4x.x4.x4a.r060.a0000.x4a.power_or_interval_95.i000}{0.04}, \resultnum{v4x.x4.x4a.r060.a0000.x4a.power_or_interval_95.i001}{0.069}] & \raggedright \resultnum{v4x.x4.x4a.r060.a0000.x4a.test_panels}{1000} \tabularnewline
\raggedright oracle (X4a) & \raggedright \resultnum{v4x.x4.x4a.r065.a0000.cell.r}{0.65} & \raggedright compatible & \raggedright \resultnum{v4x.x4.x4a.r065.a0000.x4a.power_or}{0.05} & \raggedright [\resultnum{v4x.x4.x4a.r065.a0000.x4a.power_or_interval_95.i000}{0.037}, \resultnum{v4x.x4.x4a.r065.a0000.x4a.power_or_interval_95.i001}{0.064}] & \raggedright \resultnum{v4x.x4.x4a.r065.a0000.x4a.test_panels}{1000} \tabularnewline
\raggedright oracle (X4a) & \raggedright \resultnum{v4x.x4.x4a.r070.a0000.cell.r}{0.7} & \raggedright above & \raggedright \resultnum{v4x.x4.x4a.r070.a0000.x4a.power_or}{0.069} & \raggedright [\resultnum{v4x.x4.x4a.r070.a0000.x4a.power_or_interval_95.i000}{0.055}, \resultnum{v4x.x4.x4a.r070.a0000.x4a.power_or_interval_95.i001}{0.086}] & \raggedright \resultnum{v4x.x4.x4a.r070.a0000.x4a.test_panels}{1000} \tabularnewline
\raggedright oracle (X4a) & \raggedright \resultnum{v4x.x4.x4a.r075.a0000.cell.r}{0.75} & \raggedright compatible & \raggedright \resultnum{v4x.x4.x4a.r075.a0000.x4a.power_or}{0.052} & \raggedright [\resultnum{v4x.x4.x4a.r075.a0000.x4a.power_or_interval_95.i000}{0.039}, \resultnum{v4x.x4.x4a.r075.a0000.x4a.power_or_interval_95.i001}{0.066}] & \raggedright \resultnum{v4x.x4.x4a.r075.a0000.x4a.test_panels}{1000} \tabularnewline
\raggedright oracle (X4a) & \raggedright \resultnum{v4x.x4.x4a.r080.a0000.cell.r}{0.8} & \raggedright compatible & \raggedright \resultnum{v4x.x4.x4a.r080.a0000.x4a.power_or}{0.054} & \raggedright [\resultnum{v4x.x4.x4a.r080.a0000.x4a.power_or_interval_95.i000}{0.04}, \resultnum{v4x.x4.x4a.r080.a0000.x4a.power_or_interval_95.i001}{0.069}] & \raggedright \resultnum{v4x.x4.x4a.r080.a0000.x4a.test_panels}{1000} \tabularnewline
\raggedright oracle (X4a) & \raggedright \resultnum{v4x.x4.x4a.r085.a0000.cell.r}{0.85} & \raggedright compatible & \raggedright \resultnum{v4x.x4.x4a.r085.a0000.x4a.power_or}{0.049} & \raggedright [\resultnum{v4x.x4.x4a.r085.a0000.x4a.power_or_interval_95.i000}{0.036}, \resultnum{v4x.x4.x4a.r085.a0000.x4a.power_or_interval_95.i001}{0.062}] & \raggedright \resultnum{v4x.x4.x4a.r085.a0000.x4a.test_panels}{1000} \tabularnewline
\raggedright oracle (X4a) & \raggedright \resultnum{v4x.x4.x4a.r090.a0000.cell.r}{0.9} & \raggedright compatible & \raggedright \resultnum{v4x.x4.x4a.r090.a0000.x4a.power_or}{0.046} & \raggedright [\resultnum{v4x.x4.x4a.r090.a0000.x4a.power_or_interval_95.i000}{0.034}, \resultnum{v4x.x4.x4a.r090.a0000.x4a.power_or_interval_95.i001}{0.06}] & \raggedright \resultnum{v4x.x4.x4a.r090.a0000.x4a.test_panels}{1000} \tabularnewline
\raggedright fitted panel (X4b) & \raggedright \resultnum{v4x.x4.x4b.r050.a0000.cell.r}{0.5} & \raggedright above & \raggedright \resultnum{v4x.x4.x4b.r050.a0000.x4b.power_pu}{0.3} & \raggedright [\resultnum{v4x.x4.x4b.r050.a0000.x4b.power_pu_interval_95.i000}{0.235}, \resultnum{v4x.x4.x4b.r050.a0000.x4b.power_pu_interval_95.i001}{0.365}] & \raggedright \resultnum{v4x.x4.x4b.r050.a0000.x4b.test_panels}{200} \tabularnewline
\raggedright fitted panel (X4b) & \raggedright \resultnum{v4x.x4.x4b.r070.a0000.cell.r}{0.7} & \raggedright compatible & \raggedright \resultnum{v4x.x4.x4b.r070.a0000.x4b.power_pu}{0.035} & \raggedright [\resultnum{v4x.x4.x4b.r070.a0000.x4b.power_pu_interval_95.i000}{0.01}, \resultnum{v4x.x4.x4b.r070.a0000.x4b.power_pu_interval_95.i001}{0.060125}] & \raggedright \resultnum{v4x.x4.x4b.r070.a0000.x4b.test_panels}{200} \tabularnewline
\raggedright fitted panel (X4b) & \raggedright \resultnum{v4x.x4.x4b.r090.a0000.cell.r}{0.9} & \raggedright below & \raggedright \resultnum{v4x.x4.x4b.r090.a0000.x4b.power_pu}{0.02} & \raggedright [\resultnum{v4x.x4.x4b.r090.a0000.x4b.power_pu_interval_95.i000}{0.005}, \resultnum{v4x.x4.x4b.r090.a0000.x4b.power_pu_interval_95.i001}{0.04}] & \raggedright \resultnum{v4x.x4.x4b.r090.a0000.x4b.test_panels}{200} \tabularnewline
\botrule
\end{tabular}
\end{table*}

% Generated by v4_exploratory_export; edited during the journal revision
% (agent S, 2026-10-13): captions, headers and categorical labels translated
% into plain English. Every number keeps its \resultnum binding.
\begin{table*}[htbp]
\setlength{\tabcolsep}{2pt}
\caption{Normal-theory reference for the oracle single-panel test (X4a): the power $\Phi(\sqrt2\,\Phi^{-1}(\mathrm{AUC})-1.6449)$ evaluated at each cell's own reference AUC, and its position relative to the 95\% interval of the cell's oracle rejection rate (Table~\ref{tab:v4x-x4a}): inside, above or below (machine labels describes, overstates, understates). Interval compatibility does not validate normal theory. The oracle procedure's zero-strength rate was above 0.05 at $r=0.70$ and compatible with 0.05 at the other shares (Table~\ref{tab:v4x-x4-calibration}). Fit gap: no LDA in that cell.}\label{tab:v4x-x4-normal}
\centering\footnotesize
\begin{tabular}{@{}p{0.160\textwidth}p{0.160\textwidth}p{0.130\textwidth}p{0.160\textwidth}p{0.310\textwidth}@{}}
\toprule
\raggedright $a,r$ & \raggedright \mbox{Normal} \mbox{reference} & \raggedright \mbox{Position} & \raggedright \mbox{Reference} \mbox{AUC} & \raggedright \mbox{AUC} \mbox{interval} \tabularnewline
\midrule
\raggedright \resultnum{v4x.x4.x4a.r050.a0000.cell.a}{0}, \resultnum{v4x.x4.x4a.r050.a0000.cell.r}{0.5} & \raggedright \resultnum{v4x.x4.x4a.r050.a0000.x4a.normal_theory_power.at_auc_ref}{0.051028} & \raggedright inside & \raggedright \resultnum{v4x.x4.x4a.r050.a0000.x4a.auc_ref}{0.502802} & \raggedright [\resultnum{v4x.x4.x4a.r050.a0000.x4a.auc_ref_interval_95.i000}{0.489356}, \resultnum{v4x.x4.x4a.r050.a0000.x4a.auc_ref_interval_95.i001}{0.515179}] \tabularnewline
\raggedright \resultnum{v4x.x4.x4a.r050.a0010.cell.a}{0.001}, \resultnum{v4x.x4.x4a.r050.a0010.cell.r}{0.5} & \raggedright \resultnum{v4x.x4.x4a.r050.a0010.x4a.normal_theory_power.at_auc_ref}{0.0517391} & \raggedright inside & \raggedright \resultnum{v4x.x4.x4a.r050.a0010.x4a.auc_ref}{0.504705} & \raggedright [\resultnum{v4x.x4.x4a.r050.a0010.x4a.auc_ref_interval_95.i000}{0.491839}, \resultnum{v4x.x4.x4a.r050.a0010.x4a.auc_ref_interval_95.i001}{0.516625}] \tabularnewline
\raggedright \resultnum{v4x.x4.x4a.r050.a0030.cell.a}{0.003}, \resultnum{v4x.x4.x4a.r050.a0030.cell.r}{0.5} & \raggedright \resultnum{v4x.x4.x4a.r050.a0030.x4a.normal_theory_power.at_auc_ref}{0.0525067} & \raggedright inside & \raggedright \resultnum{v4x.x4.x4a.r050.a0030.x4a.auc_ref}{0.506736} & \raggedright [\resultnum{v4x.x4.x4a.r050.a0030.x4a.auc_ref_interval_95.i000}{0.494385}, \resultnum{v4x.x4.x4a.r050.a0030.x4a.auc_ref_interval_95.i001}{0.518492}] \tabularnewline
\raggedright \resultnum{v4x.x4.x4a.r050.a0100.cell.a}{0.01}, \resultnum{v4x.x4.x4a.r050.a0100.cell.r}{0.5} & \raggedright \resultnum{v4x.x4.x4a.r050.a0100.x4a.normal_theory_power.at_auc_ref}{0.0578406} & \raggedright inside & \raggedright \resultnum{v4x.x4.x4a.r050.a0100.x4a.auc_ref}{0.520228} & \raggedright [\resultnum{v4x.x4.x4a.r050.a0100.x4a.auc_ref_interval_95.i000}{0.507612}, \resultnum{v4x.x4.x4a.r050.a0100.x4a.auc_ref_interval_95.i001}{0.532693}] \tabularnewline
\raggedright \resultnum{v4x.x4.x4a.r050.a0300.cell.a}{0.03}, \resultnum{v4x.x4.x4a.r050.a0300.cell.r}{0.5} & \raggedright \resultnum{v4x.x4.x4a.r050.a0300.x4a.normal_theory_power.at_auc_ref}{0.0730083} & \raggedright inside & \raggedright \resultnum{v4x.x4.x4a.r050.a0300.x4a.auc_ref}{0.55376} & \raggedright [\resultnum{v4x.x4.x4a.r050.a0300.x4a.auc_ref_interval_95.i000}{0.541611}, \resultnum{v4x.x4.x4a.r050.a0300.x4a.auc_ref_interval_95.i001}{0.566707}] \tabularnewline
\raggedright \resultnum{v4x.x4.x4a.r050.a1000.cell.a}{0.1}, \resultnum{v4x.x4.x4a.r050.a1000.cell.r}{0.5} & \raggedright \resultnum{v4x.x4.x4a.r050.a1000.x4a.normal_theory_power.at_auc_ref}{0.133847} & \raggedright inside & \raggedright \resultnum{v4x.x4.x4a.r050.a1000.x4a.auc_ref}{0.647794} & \raggedright [\resultnum{v4x.x4.x4a.r050.a1000.x4a.auc_ref_interval_95.i000}{0.635456}, \resultnum{v4x.x4.x4a.r050.a1000.x4a.auc_ref_interval_95.i001}{0.660071}] \tabularnewline
\raggedright \resultnum{v4x.x4.x4a.r050.a3000.cell.a}{0.3}, \resultnum{v4x.x4.x4a.r050.a3000.cell.r}{0.5} & \raggedright \resultnum{v4x.x4.x4a.r050.a3000.x4a.normal_theory_power.at_auc_ref}{0.434358} & \raggedright inside & \raggedright \resultnum{v4x.x4.x4a.r050.a3000.x4a.auc_ref}{0.852275} & \raggedright [\resultnum{v4x.x4.x4a.r050.a3000.x4a.auc_ref_interval_95.i000}{0.844115}, \resultnum{v4x.x4.x4a.r050.a3000.x4a.auc_ref_interval_95.i001}{0.86027}] \tabularnewline
\raggedright \resultnum{v4x.x4.x4a.r050.a6000.cell.a}{0.6}, \resultnum{v4x.x4.x4a.r050.a6000.cell.r}{0.5} & \raggedright \resultnum{v4x.x4.x4a.r050.a6000.x4a.normal_theory_power.at_auc_ref}{0.865133} & \raggedright inside & \raggedright \resultnum{v4x.x4.x4a.r050.a6000.x4a.auc_ref}{0.974024} & \raggedright [\resultnum{v4x.x4.x4a.r050.a6000.x4a.auc_ref_interval_95.i000}{0.971179}, \resultnum{v4x.x4.x4a.r050.a6000.x4a.auc_ref_interval_95.i001}{0.976669}] \tabularnewline
\raggedright \resultnum{v4x.x4.x4a.r050.a9000.cell.a}{0.9}, \resultnum{v4x.x4.x4a.r050.a9000.cell.r}{0.5} & \raggedright \resultnum{v4x.x4.x4a.r050.a9000.x4a.normal_theory_power.at_auc_ref}{0.96915} & \raggedright inside & \raggedright \resultnum{v4x.x4.x4a.r050.a9000.x4a.auc_ref}{0.99351} & \raggedright [\resultnum{v4x.x4.x4a.r050.a9000.x4a.auc_ref_interval_95.i000}{0.992334}, \resultnum{v4x.x4.x4a.r050.a9000.x4a.auc_ref_interval_95.i001}{0.994547}] \tabularnewline
\raggedright \resultnum{v4x.x4.x4a.r055.a0000.cell.a}{0}, \resultnum{v4x.x4.x4a.r055.a0000.cell.r}{0.55} & \raggedright \resultnum{v4x.x4.x4a.r055.a0000.x4a.normal_theory_power.at_auc_ref}{0.0500181} & \raggedright inside & \raggedright \resultnum{v4x.x4.x4a.r055.a0000.x4a.auc_ref}{0.500063} & \raggedright [\resultnum{v4x.x4.x4a.r055.a0000.x4a.auc_ref_interval_95.i000}{0.488558}, \resultnum{v4x.x4.x4a.r055.a0000.x4a.auc_ref_interval_95.i001}{0.512989}] \tabularnewline
\raggedright \resultnum{v4x.x4.x4a.r055.a0010.cell.a}{0.001}, \resultnum{v4x.x4.x4a.r055.a0010.cell.r}{0.55} & \raggedright \resultnum{v4x.x4.x4a.r055.a0010.x4a.normal_theory_power.at_auc_ref}{0.0540977} & \raggedright inside & \raggedright \resultnum{v4x.x4.x4a.r055.a0010.x4a.auc_ref}{0.510871} & \raggedright [\resultnum{v4x.x4.x4a.r055.a0010.x4a.auc_ref_interval_95.i000}{0.499328}, \resultnum{v4x.x4.x4a.r055.a0010.x4a.auc_ref_interval_95.i001}{0.523567}] \tabularnewline
\raggedright \resultnum{v4x.x4.x4a.r055.a0030.cell.a}{0.003}, \resultnum{v4x.x4.x4a.r055.a0030.cell.r}{0.55} & \raggedright \resultnum{v4x.x4.x4a.r055.a0030.x4a.normal_theory_power.at_auc_ref}{0.0539507} & \raggedright inside & \raggedright \resultnum{v4x.x4.x4a.r055.a0030.x4a.auc_ref}{0.510493} & \raggedright [\resultnum{v4x.x4.x4a.r055.a0030.x4a.auc_ref_interval_95.i000}{0.498811}, \resultnum{v4x.x4.x4a.r055.a0030.x4a.auc_ref_interval_95.i001}{0.523419}] \tabularnewline
\raggedright \resultnum{v4x.x4.x4a.r055.a0100.cell.a}{0.01}, \resultnum{v4x.x4.x4a.r055.a0100.cell.r}{0.55} & \raggedright \resultnum{v4x.x4.x4a.r055.a0100.x4a.normal_theory_power.at_auc_ref}{0.0641451} & \raggedright inside & \raggedright \resultnum{v4x.x4.x4a.r055.a0100.x4a.auc_ref}{0.534941} & \raggedright [\resultnum{v4x.x4.x4a.r055.a0100.x4a.auc_ref_interval_95.i000}{0.523741}, \resultnum{v4x.x4.x4a.r055.a0100.x4a.auc_ref_interval_95.i001}{0.547768}] \tabularnewline
\raggedright \resultnum{v4x.x4.x4a.r055.a0300.cell.a}{0.03}, \resultnum{v4x.x4.x4a.r055.a0300.cell.r}{0.55} & \raggedright \resultnum{v4x.x4.x4a.r055.a0300.x4a.normal_theory_power.at_auc_ref}{0.0802603} & \raggedright inside & \raggedright \resultnum{v4x.x4.x4a.r055.a0300.x4a.auc_ref}{0.567818} & \raggedright [\resultnum{v4x.x4.x4a.r055.a0300.x4a.auc_ref_interval_95.i000}{0.556245}, \resultnum{v4x.x4.x4a.r055.a0300.x4a.auc_ref_interval_95.i001}{0.579974}] \tabularnewline
\raggedright \resultnum{v4x.x4.x4a.r055.a1000.cell.a}{0.1}, \resultnum{v4x.x4.x4a.r055.a1000.cell.r}{0.55} & \raggedright \resultnum{v4x.x4.x4a.r055.a1000.x4a.normal_theory_power.at_auc_ref}{0.165346} & \raggedright inside & \raggedright \resultnum{v4x.x4.x4a.r055.a1000.x4a.auc_ref}{0.682714} & \raggedright [\resultnum{v4x.x4.x4a.r055.a1000.x4a.auc_ref_interval_95.i000}{0.671661}, \resultnum{v4x.x4.x4a.r055.a1000.x4a.auc_ref_interval_95.i001}{0.693736}] \tabularnewline
\raggedright \resultnum{v4x.x4.x4a.r055.a3000.cell.a}{0.3}, \resultnum{v4x.x4.x4a.r055.a3000.cell.r}{0.55} & \raggedright \resultnum{v4x.x4.x4a.r055.a3000.x4a.normal_theory_power.at_auc_ref}{0.516206} & \raggedright inside & \raggedright \resultnum{v4x.x4.x4a.r055.a3000.x4a.auc_ref}{0.88334} & \raggedright [\resultnum{v4x.x4.x4a.r055.a3000.x4a.auc_ref_interval_95.i000}{0.876039}, \resultnum{v4x.x4.x4a.r055.a3000.x4a.auc_ref_interval_95.i001}{0.890602}] \tabularnewline
\raggedright \resultnum{v4x.x4.x4a.r055.a6000.cell.a}{0.6}, \resultnum{v4x.x4.x4a.r055.a6000.cell.r}{0.55} & \raggedright \resultnum{v4x.x4.x4a.r055.a6000.x4a.normal_theory_power.at_auc_ref}{0.91991} & \raggedright inside & \raggedright \resultnum{v4x.x4.x4a.r055.a6000.x4a.auc_ref}{0.984467} & \raggedright [\resultnum{v4x.x4.x4a.r055.a6000.x4a.auc_ref_interval_95.i000}{0.982268}, \resultnum{v4x.x4.x4a.r055.a6000.x4a.auc_ref_interval_95.i001}{0.98657}] \tabularnewline
\raggedright \resultnum{v4x.x4.x4a.r055.a9000.cell.a}{0.9}, \resultnum{v4x.x4.x4a.r055.a9000.cell.r}{0.55} & \raggedright \resultnum{v4x.x4.x4a.r055.a9000.x4a.normal_theory_power.at_auc_ref}{0.987045} & \raggedright inside & \raggedright \resultnum{v4x.x4.x4a.r055.a9000.x4a.auc_ref}{0.996912} & \raggedright [\resultnum{v4x.x4.x4a.r055.a9000.x4a.auc_ref_interval_95.i000}{0.996131}, \resultnum{v4x.x4.x4a.r055.a9000.x4a.auc_ref_interval_95.i001}{0.997604}] \tabularnewline
\raggedright \resultnum{v4x.x4.x4a.r060.a0000.cell.a}{0}, \resultnum{v4x.x4.x4a.r060.a0000.cell.r}{0.6} & \raggedright \resultnum{v4x.x4.x4a.r060.a0000.x4a.normal_theory_power.at_auc_ref}{0.0531652} & \raggedright inside & \raggedright \resultnum{v4x.x4.x4a.r060.a0000.x4a.auc_ref}{0.50846} & \raggedright [\resultnum{v4x.x4.x4a.r060.a0000.x4a.auc_ref_interval_95.i000}{0.497522}, \resultnum{v4x.x4.x4a.r060.a0000.x4a.auc_ref_interval_95.i001}{0.520657}] \tabularnewline
\raggedright \resultnum{v4x.x4.x4a.r060.a0010.cell.a}{0.001}, \resultnum{v4x.x4.x4a.r060.a0010.cell.r}{0.6} & \raggedright \resultnum{v4x.x4.x4a.r060.a0010.x4a.normal_theory_power.at_auc_ref}{0.0565119} & \raggedright inside & \raggedright \resultnum{v4x.x4.x4a.r060.a0010.x4a.auc_ref}{0.516964} & \raggedright [\resultnum{v4x.x4.x4a.r060.a0010.x4a.auc_ref_interval_95.i000}{0.504318}, \resultnum{v4x.x4.x4a.r060.a0010.x4a.auc_ref_interval_95.i001}{0.529786}] \tabularnewline
\raggedright \resultnum{v4x.x4.x4a.r060.a0030.cell.a}{0.003}, \resultnum{v4x.x4.x4a.r060.a0030.cell.r}{0.6} & \raggedright \resultnum{v4x.x4.x4a.r060.a0030.x4a.normal_theory_power.at_auc_ref}{0.0574738} & \raggedright inside & \raggedright \resultnum{v4x.x4.x4a.r060.a0030.x4a.auc_ref}{0.519333} & \raggedright [\resultnum{v4x.x4.x4a.r060.a0030.x4a.auc_ref_interval_95.i000}{0.506513}, \resultnum{v4x.x4.x4a.r060.a0030.x4a.auc_ref_interval_95.i001}{0.532168}] \tabularnewline
\raggedright \resultnum{v4x.x4.x4a.r060.a0100.cell.a}{0.01}, \resultnum{v4x.x4.x4a.r060.a0100.cell.r}{0.6} & \raggedright \resultnum{v4x.x4.x4a.r060.a0100.x4a.normal_theory_power.at_auc_ref}{0.0638536} & \raggedright inside & \raggedright \resultnum{v4x.x4.x4a.r060.a0100.x4a.auc_ref}{0.534287} & \raggedright [\resultnum{v4x.x4.x4a.r060.a0100.x4a.auc_ref_interval_95.i000}{0.521525}, \resultnum{v4x.x4.x4a.r060.a0100.x4a.auc_ref_interval_95.i001}{0.546911}] \tabularnewline
\raggedright \resultnum{v4x.x4.x4a.r060.a0300.cell.a}{0.03}, \resultnum{v4x.x4.x4a.r060.a0300.cell.r}{0.6} & \raggedright \resultnum{v4x.x4.x4a.r060.a0300.x4a.normal_theory_power.at_auc_ref}{0.0883396} & \raggedright inside & \raggedright \resultnum{v4x.x4.x4a.r060.a0300.x4a.auc_ref}{0.582301} & \raggedright [\resultnum{v4x.x4.x4a.r060.a0300.x4a.auc_ref_interval_95.i000}{0.570073}, \resultnum{v4x.x4.x4a.r060.a0300.x4a.auc_ref_interval_95.i001}{0.594494}] \tabularnewline
\raggedright \resultnum{v4x.x4.x4a.r060.a1000.cell.a}{0.1}, \resultnum{v4x.x4.x4a.r060.a1000.cell.r}{0.6} & \raggedright \resultnum{v4x.x4.x4a.r060.a1000.x4a.normal_theory_power.at_auc_ref}{0.185568} & \raggedright inside & \raggedright \resultnum{v4x.x4.x4a.r060.a1000.x4a.auc_ref}{0.702194} & \raggedright [\resultnum{v4x.x4.x4a.r060.a1000.x4a.auc_ref_interval_95.i000}{0.691113}, \resultnum{v4x.x4.x4a.r060.a1000.x4a.auc_ref_interval_95.i001}{0.713307}] \tabularnewline
\raggedright \resultnum{v4x.x4.x4a.r060.a3000.cell.a}{0.3}, \resultnum{v4x.x4.x4a.r060.a3000.cell.r}{0.6} & \raggedright \resultnum{v4x.x4.x4a.r060.a3000.x4a.normal_theory_power.at_auc_ref}{0.635103} & \raggedright inside & \raggedright \resultnum{v4x.x4.x4a.r060.a3000.x4a.auc_ref}{0.920339} & \raggedright [\resultnum{v4x.x4.x4a.r060.a3000.x4a.auc_ref_interval_95.i000}{0.914751}, \resultnum{v4x.x4.x4a.r060.a3000.x4a.auc_ref_interval_95.i001}{0.926082}] \tabularnewline
\raggedright \resultnum{v4x.x4.x4a.r060.a6000.cell.a}{0.6}, \resultnum{v4x.x4.x4a.r060.a6000.cell.r}{0.6} & \raggedright fit gap & \raggedright  & \raggedright  & \raggedright  \tabularnewline
\raggedright \resultnum{v4x.x4.x4a.r060.a9000.cell.a}{0.9}, \resultnum{v4x.x4.x4a.r060.a9000.cell.r}{0.6} & \raggedright fit gap & \raggedright  & \raggedright  & \raggedright  \tabularnewline
\raggedright \resultnum{v4x.x4.x4a.r065.a0000.cell.a}{0}, \resultnum{v4x.x4.x4a.r065.a0000.cell.r}{0.65} & \raggedright \resultnum{v4x.x4.x4a.r065.a0000.x4a.normal_theory_power.at_auc_ref}{0.0483988} & \raggedright inside & \raggedright \resultnum{v4x.x4.x4a.r065.a0000.x4a.auc_ref}{0.495576} & \raggedright [\resultnum{v4x.x4.x4a.r065.a0000.x4a.auc_ref_interval_95.i000}{0.482302}, \resultnum{v4x.x4.x4a.r065.a0000.x4a.auc_ref_interval_95.i001}{0.507636}] \tabularnewline
\raggedright \resultnum{v4x.x4.x4a.r065.a0010.cell.a}{0.001}, \resultnum{v4x.x4.x4a.r065.a0010.cell.r}{0.65} & \raggedright \resultnum{v4x.x4.x4a.r065.a0010.x4a.normal_theory_power.at_auc_ref}{0.0548187} & \raggedright inside & \raggedright \resultnum{v4x.x4.x4a.r065.a0010.x4a.auc_ref}{0.512713} & \raggedright [\resultnum{v4x.x4.x4a.r065.a0010.x4a.auc_ref_interval_95.i000}{0.499114}, \resultnum{v4x.x4.x4a.r065.a0010.x4a.auc_ref_interval_95.i001}{0.524627}] \tabularnewline
\raggedright \resultnum{v4x.x4.x4a.r065.a0030.cell.a}{0.003}, \resultnum{v4x.x4.x4a.r065.a0030.cell.r}{0.65} & \raggedright \resultnum{v4x.x4.x4a.r065.a0030.x4a.normal_theory_power.at_auc_ref}{0.0613867} & \raggedright inside & \raggedright \resultnum{v4x.x4.x4a.r065.a0030.x4a.auc_ref}{0.528654} & \raggedright [\resultnum{v4x.x4.x4a.r065.a0030.x4a.auc_ref_interval_95.i000}{0.514891}, \resultnum{v4x.x4.x4a.r065.a0030.x4a.auc_ref_interval_95.i001}{0.541106}] \tabularnewline
\raggedright \resultnum{v4x.x4.x4a.r065.a0100.cell.a}{0.01}, \resultnum{v4x.x4.x4a.r065.a0100.cell.r}{0.65} & \raggedright \resultnum{v4x.x4.x4a.r065.a0100.x4a.normal_theory_power.at_auc_ref}{0.0770387} & \raggedright inside & \raggedright \resultnum{v4x.x4.x4a.r065.a0100.x4a.auc_ref}{0.561707} & \raggedright [\resultnum{v4x.x4.x4a.r065.a0100.x4a.auc_ref_interval_95.i000}{0.548665}, \resultnum{v4x.x4.x4a.r065.a0100.x4a.auc_ref_interval_95.i001}{0.574532}] \tabularnewline
\raggedright \resultnum{v4x.x4.x4a.r065.a0300.cell.a}{0.03}, \resultnum{v4x.x4.x4a.r065.a0300.cell.r}{0.65} & \raggedright \resultnum{v4x.x4.x4a.r065.a0300.x4a.normal_theory_power.at_auc_ref}{0.107095} & \raggedright inside & \raggedright \resultnum{v4x.x4.x4a.r065.a0300.x4a.auc_ref}{0.612103} & \raggedright [\resultnum{v4x.x4.x4a.r065.a0300.x4a.auc_ref_interval_95.i000}{0.599064}, \resultnum{v4x.x4.x4a.r065.a0300.x4a.auc_ref_interval_95.i001}{0.624259}] \tabularnewline
\raggedright \resultnum{v4x.x4.x4a.r065.a1000.cell.a}{0.1}, \resultnum{v4x.x4.x4a.r065.a1000.cell.r}{0.65} & \raggedright \resultnum{v4x.x4.x4a.r065.a1000.x4a.normal_theory_power.at_auc_ref}{0.212073} & \raggedright inside & \raggedright \resultnum{v4x.x4.x4a.r065.a1000.x4a.auc_ref}{0.725069} & \raggedright [\resultnum{v4x.x4.x4a.r065.a1000.x4a.auc_ref_interval_95.i000}{0.713585}, \resultnum{v4x.x4.x4a.r065.a1000.x4a.auc_ref_interval_95.i001}{0.735711}] \tabularnewline
\raggedright \resultnum{v4x.x4.x4a.r065.a3000.cell.a}{0.3}, \resultnum{v4x.x4.x4a.r065.a3000.cell.r}{0.65} & \raggedright \resultnum{v4x.x4.x4a.r065.a3000.x4a.normal_theory_power.at_auc_ref}{0.710801} & \raggedright inside & \raggedright \resultnum{v4x.x4.x4a.r065.a3000.x4a.auc_ref}{0.940155} & \raggedright [\resultnum{v4x.x4.x4a.r065.a3000.x4a.auc_ref_interval_95.i000}{0.934844}, \resultnum{v4x.x4.x4a.r065.a3000.x4a.auc_ref_interval_95.i001}{0.944641}] \tabularnewline
\raggedright \resultnum{v4x.x4.x4a.r065.a6000.cell.a}{0.6}, \resultnum{v4x.x4.x4a.r065.a6000.cell.r}{0.65} & \raggedright \resultnum{v4x.x4.x4a.r065.a6000.x4a.normal_theory_power.at_auc_ref}{0.985711} & \raggedright inside & \raggedright \resultnum{v4x.x4.x4a.r065.a6000.x4a.auc_ref}{0.996648} & \raggedright [\resultnum{v4x.x4.x4a.r065.a6000.x4a.auc_ref_interval_95.i000}{0.995892}, \resultnum{v4x.x4.x4a.r065.a6000.x4a.auc_ref_interval_95.i001}{0.997309}] \tabularnewline
\raggedright \resultnum{v4x.x4.x4a.r065.a9000.cell.a}{0.9}, \resultnum{v4x.x4.x4a.r065.a9000.cell.r}{0.65} & \raggedright \resultnum{v4x.x4.x4a.r065.a9000.x4a.normal_theory_power.at_auc_ref}{0.999465} & \raggedright inside & \raggedright \resultnum{v4x.x4.x4a.r065.a9000.x4a.auc_ref}{0.999746} & \raggedright [\resultnum{v4x.x4.x4a.r065.a9000.x4a.auc_ref_interval_95.i000}{0.999625}, \resultnum{v4x.x4.x4a.r065.a9000.x4a.auc_ref_interval_95.i001}{0.999848}] \tabularnewline
\raggedright \resultnum{v4x.x4.x4a.r070.a0000.cell.a}{0}, \resultnum{v4x.x4.x4a.r070.a0000.cell.r}{0.7} & \raggedright \resultnum{v4x.x4.x4a.r070.a0000.x4a.normal_theory_power.at_auc_ref}{0.0496869} & \raggedright below & \raggedright \resultnum{v4x.x4.x4a.r070.a0000.x4a.auc_ref}{0.499154} & \raggedright [\resultnum{v4x.x4.x4a.r070.a0000.x4a.auc_ref_interval_95.i000}{0.48571}, \resultnum{v4x.x4.x4a.r070.a0000.x4a.auc_ref_interval_95.i001}{0.511773}] \tabularnewline
\raggedright \resultnum{v4x.x4.x4a.r070.a0010.cell.a}{0.001}, \resultnum{v4x.x4.x4a.r070.a0010.cell.r}{0.7} & \raggedright \resultnum{v4x.x4.x4a.r070.a0010.x4a.normal_theory_power.at_auc_ref}{0.0515708} & \raggedright inside & \raggedright \resultnum{v4x.x4.x4a.r070.a0010.x4a.auc_ref}{0.504257} & \raggedright [\resultnum{v4x.x4.x4a.r070.a0010.x4a.auc_ref_interval_95.i000}{0.492184}, \resultnum{v4x.x4.x4a.r070.a0010.x4a.auc_ref_interval_95.i001}{0.516688}] \tabularnewline
\raggedright \resultnum{v4x.x4.x4a.r070.a0030.cell.a}{0.003}, \resultnum{v4x.x4.x4a.r070.a0030.cell.r}{0.7} & \raggedright \resultnum{v4x.x4.x4a.r070.a0030.x4a.normal_theory_power.at_auc_ref}{0.0522269} & \raggedright inside & \raggedright \resultnum{v4x.x4.x4a.r070.a0030.x4a.auc_ref}{0.505999} & \raggedright [\resultnum{v4x.x4.x4a.r070.a0030.x4a.auc_ref_interval_95.i000}{0.493904}, \resultnum{v4x.x4.x4a.r070.a0030.x4a.auc_ref_interval_95.i001}{0.51922}] \tabularnewline
\raggedright \resultnum{v4x.x4.x4a.r070.a0100.cell.a}{0.01}, \resultnum{v4x.x4.x4a.r070.a0100.cell.r}{0.7} & \raggedright \resultnum{v4x.x4.x4a.r070.a0100.x4a.normal_theory_power.at_auc_ref}{0.0618119} & \raggedright inside & \raggedright \resultnum{v4x.x4.x4a.r070.a0100.x4a.auc_ref}{0.529638} & \raggedright [\resultnum{v4x.x4.x4a.r070.a0100.x4a.auc_ref_interval_95.i000}{0.517177}, \resultnum{v4x.x4.x4a.r070.a0100.x4a.auc_ref_interval_95.i001}{0.542019}] \tabularnewline
\raggedright \resultnum{v4x.x4.x4a.r070.a0300.cell.a}{0.03}, \resultnum{v4x.x4.x4a.r070.a0300.cell.r}{0.7} & \raggedright \resultnum{v4x.x4.x4a.r070.a0300.x4a.normal_theory_power.at_auc_ref}{0.097031} & \raggedright inside & \raggedright \resultnum{v4x.x4.x4a.r070.a0300.x4a.auc_ref}{0.596707} & \raggedright [\resultnum{v4x.x4.x4a.r070.a0300.x4a.auc_ref_interval_95.i000}{0.584593}, \resultnum{v4x.x4.x4a.r070.a0300.x4a.auc_ref_interval_95.i001}{0.608737}] \tabularnewline
\raggedright \resultnum{v4x.x4.x4a.r070.a1000.cell.a}{0.1}, \resultnum{v4x.x4.x4a.r070.a1000.cell.r}{0.7} & \raggedright \resultnum{v4x.x4.x4a.r070.a1000.x4a.normal_theory_power.at_auc_ref}{0.239062} & \raggedright above & \raggedright \resultnum{v4x.x4.x4a.r070.a1000.x4a.auc_ref}{0.745871} & \raggedright [\resultnum{v4x.x4.x4a.r070.a1000.x4a.auc_ref_interval_95.i000}{0.735521}, \resultnum{v4x.x4.x4a.r070.a1000.x4a.auc_ref_interval_95.i001}{0.756874}] \tabularnewline
\raggedright \resultnum{v4x.x4.x4a.r070.a3000.cell.a}{0.3}, \resultnum{v4x.x4.x4a.r070.a3000.cell.r}{0.7} & \raggedright \resultnum{v4x.x4.x4a.r070.a3000.x4a.normal_theory_power.at_auc_ref}{0.766647} & \raggedright inside & \raggedright \resultnum{v4x.x4.x4a.r070.a3000.x4a.auc_ref}{0.953306} & \raggedright [\resultnum{v4x.x4.x4a.r070.a3000.x4a.auc_ref_interval_95.i000}{0.948872}, \resultnum{v4x.x4.x4a.r070.a3000.x4a.auc_ref_interval_95.i001}{0.957305}] \tabularnewline
\raggedright \resultnum{v4x.x4.x4a.r070.a6000.cell.a}{0.6}, \resultnum{v4x.x4.x4a.r070.a6000.cell.r}{0.7} & \raggedright fit gap & \raggedright  & \raggedright  & \raggedright  \tabularnewline
\raggedright \resultnum{v4x.x4.x4a.r070.a9000.cell.a}{0.9}, \resultnum{v4x.x4.x4a.r070.a9000.cell.r}{0.7} & \raggedright \resultnum{v4x.x4.x4a.r070.a9000.x4a.normal_theory_power.at_auc_ref}{0.999833} & \raggedright below & \raggedright \resultnum{v4x.x4.x4a.r070.a9000.x4a.auc_ref}{0.999892} & \raggedright [\resultnum{v4x.x4.x4a.r070.a9000.x4a.auc_ref_interval_95.i000}{0.999811}, \resultnum{v4x.x4.x4a.r070.a9000.x4a.auc_ref_interval_95.i001}{0.999947}] \tabularnewline
\raggedright \resultnum{v4x.x4.x4a.r075.a0000.cell.a}{0}, \resultnum{v4x.x4.x4a.r075.a0000.cell.r}{0.75} & \raggedright \resultnum{v4x.x4.x4a.r075.a0000.x4a.normal_theory_power.at_auc_ref}{0.0470105} & \raggedright inside & \raggedright \resultnum{v4x.x4.x4a.r075.a0000.x4a.auc_ref}{0.491634} & \raggedright [\resultnum{v4x.x4.x4a.r075.a0000.x4a.auc_ref_interval_95.i000}{0.478386}, \resultnum{v4x.x4.x4a.r075.a0000.x4a.auc_ref_interval_95.i001}{0.504593}] \tabularnewline
\raggedright \resultnum{v4x.x4.x4a.r075.a0010.cell.a}{0.001}, \resultnum{v4x.x4.x4a.r075.a0010.cell.r}{0.75} & \raggedright \resultnum{v4x.x4.x4a.r075.a0010.x4a.normal_theory_power.at_auc_ref}{0.0447757} & \raggedright inside & \raggedright \resultnum{v4x.x4.x4a.r075.a0010.x4a.auc_ref}{0.485089} & \raggedright [\resultnum{v4x.x4.x4a.r075.a0010.x4a.auc_ref_interval_95.i000}{0.471325}, \resultnum{v4x.x4.x4a.r075.a0010.x4a.auc_ref_interval_95.i001}{0.49792}] \tabularnewline
\raggedright \resultnum{v4x.x4.x4a.r075.a0030.cell.a}{0.003}, \resultnum{v4x.x4.x4a.r075.a0030.cell.r}{0.75} & \raggedright \resultnum{v4x.x4.x4a.r075.a0030.x4a.normal_theory_power.at_auc_ref}{0.0489384} & \raggedright below & \raggedright \resultnum{v4x.x4.x4a.r075.a0030.x4a.auc_ref}{0.497085} & \raggedright [\resultnum{v4x.x4.x4a.r075.a0030.x4a.auc_ref_interval_95.i000}{0.484121}, \resultnum{v4x.x4.x4a.r075.a0030.x4a.auc_ref_interval_95.i001}{0.509529}] \tabularnewline
\raggedright \resultnum{v4x.x4.x4a.r075.a0100.cell.a}{0.01}, \resultnum{v4x.x4.x4a.r075.a0100.cell.r}{0.75} & \raggedright \resultnum{v4x.x4.x4a.r075.a0100.x4a.normal_theory_power.at_auc_ref}{0.0689391} & \raggedright inside & \raggedright \resultnum{v4x.x4.x4a.r075.a0100.x4a.auc_ref}{0.545365} & \raggedright [\resultnum{v4x.x4.x4a.r075.a0100.x4a.auc_ref_interval_95.i000}{0.53383}, \resultnum{v4x.x4.x4a.r075.a0100.x4a.auc_ref_interval_95.i001}{0.558139}] \tabularnewline
\raggedright \resultnum{v4x.x4.x4a.r075.a0300.cell.a}{0.03}, \resultnum{v4x.x4.x4a.r075.a0300.cell.r}{0.75} & \raggedright \resultnum{v4x.x4.x4a.r075.a0300.x4a.normal_theory_power.at_auc_ref}{0.112315} & \raggedright inside & \raggedright \resultnum{v4x.x4.x4a.r075.a0300.x4a.auc_ref}{0.619617} & \raggedright [\resultnum{v4x.x4.x4a.r075.a0300.x4a.auc_ref_interval_95.i000}{0.608204}, \resultnum{v4x.x4.x4a.r075.a0300.x4a.auc_ref_interval_95.i001}{0.632408}] \tabularnewline
\raggedright \resultnum{v4x.x4.x4a.r075.a1000.cell.a}{0.1}, \resultnum{v4x.x4.x4a.r075.a1000.cell.r}{0.75} & \raggedright \resultnum{v4x.x4.x4a.r075.a1000.x4a.normal_theory_power.at_auc_ref}{0.296708} & \raggedright inside & \raggedright \resultnum{v4x.x4.x4a.r075.a1000.x4a.auc_ref}{0.783949} & \raggedright [\resultnum{v4x.x4.x4a.r075.a1000.x4a.auc_ref_interval_95.i000}{0.773923}, \resultnum{v4x.x4.x4a.r075.a1000.x4a.auc_ref_interval_95.i001}{0.793465}] \tabularnewline
\raggedright \resultnum{v4x.x4.x4a.r075.a3000.cell.a}{0.3}, \resultnum{v4x.x4.x4a.r075.a3000.cell.r}{0.75} & \raggedright \resultnum{v4x.x4.x4a.r075.a3000.x4a.normal_theory_power.at_auc_ref}{0.835623} & \raggedright inside & \raggedright \resultnum{v4x.x4.x4a.r075.a3000.x4a.auc_ref}{0.968109} & \raggedright [\resultnum{v4x.x4.x4a.r075.a3000.x4a.auc_ref_interval_95.i000}{0.964795}, \resultnum{v4x.x4.x4a.r075.a3000.x4a.auc_ref_interval_95.i001}{0.971232}] \tabularnewline
\botrule
\end{tabular}
\end{table*}
\begin{table*}[htbp]
\setlength{\tabcolsep}{2pt}
\addtocounter{table}{-\csname @ne\endcsname}
\providecommand{\theHtable}{\thetable}
\renewcommand{\theHtable}{v4x-x4-normal-part1}
\caption{Normal-theory reference for the oracle test (X4a) (continued, part 2 of 2)}\label{tab:v4x-x4-normal-part1}
\centering\footnotesize
\begin{tabular}{@{}p{0.160\textwidth}p{0.160\textwidth}p{0.130\textwidth}p{0.160\textwidth}p{0.310\textwidth}@{}}
\toprule
\raggedright $a,r$ & \raggedright \mbox{Normal} \mbox{reference} & \raggedright \mbox{Position} & \raggedright \mbox{Reference} \mbox{AUC} & \raggedright \mbox{AUC} \mbox{interval} \tabularnewline
\midrule
\raggedright \resultnum{v4x.x4.x4a.r075.a6000.cell.a}{0.6}, \resultnum{v4x.x4.x4a.r075.a6000.cell.r}{0.75} & \raggedright \resultnum{v4x.x4.x4a.r075.a6000.x4a.normal_theory_power.at_auc_ref}{0.998941} & \raggedright inside & \raggedright \resultnum{v4x.x4.x4a.r075.a6000.x4a.auc_ref}{0.999575} & \raggedright [\resultnum{v4x.x4.x4a.r075.a6000.x4a.auc_ref_interval_95.i000}{0.999417}, \resultnum{v4x.x4.x4a.r075.a6000.x4a.auc_ref_interval_95.i001}{0.999702}] \tabularnewline
\raggedright \resultnum{v4x.x4.x4a.r075.a9000.cell.a}{0.9}, \resultnum{v4x.x4.x4a.r075.a9000.cell.r}{0.75} & \raggedright \resultnum{v4x.x4.x4a.r075.a9000.x4a.normal_theory_power.at_auc_ref}{0.999996} & \raggedright below & \raggedright \resultnum{v4x.x4.x4a.r075.a9000.x4a.auc_ref}{0.999992} & \raggedright [\resultnum{v4x.x4.x4a.r075.a9000.x4a.auc_ref_interval_95.i000}{0.99998}, \resultnum{v4x.x4.x4a.r075.a9000.x4a.auc_ref_interval_95.i001}{0.999999}] \tabularnewline
\raggedright \resultnum{v4x.x4.x4a.r080.a0000.cell.a}{0}, \resultnum{v4x.x4.x4a.r080.a0000.cell.r}{0.8} & \raggedright \resultnum{v4x.x4.x4a.r080.a0000.x4a.normal_theory_power.at_auc_ref}{0.0525297} & \raggedright inside & \raggedright \resultnum{v4x.x4.x4a.r080.a0000.x4a.auc_ref}{0.506797} & \raggedright [\resultnum{v4x.x4.x4a.r080.a0000.x4a.auc_ref_interval_95.i000}{0.494211}, \resultnum{v4x.x4.x4a.r080.a0000.x4a.auc_ref_interval_95.i001}{0.519438}] \tabularnewline
\raggedright \resultnum{v4x.x4.x4a.r080.a0010.cell.a}{0.001}, \resultnum{v4x.x4.x4a.r080.a0010.cell.r}{0.8} & \raggedright \resultnum{v4x.x4.x4a.r080.a0010.x4a.normal_theory_power.at_auc_ref}{0.0537625} & \raggedright inside & \raggedright \resultnum{v4x.x4.x4a.r080.a0010.x4a.auc_ref}{0.510008} & \raggedright [\resultnum{v4x.x4.x4a.r080.a0010.x4a.auc_ref_interval_95.i000}{0.497935}, \resultnum{v4x.x4.x4a.r080.a0010.x4a.auc_ref_interval_95.i001}{0.521663}] \tabularnewline
\raggedright \resultnum{v4x.x4.x4a.r080.a0030.cell.a}{0.003}, \resultnum{v4x.x4.x4a.r080.a0030.cell.r}{0.8} & \raggedright \resultnum{v4x.x4.x4a.r080.a0030.x4a.normal_theory_power.at_auc_ref}{0.0567936} & \raggedright inside & \raggedright \resultnum{v4x.x4.x4a.r080.a0030.x4a.auc_ref}{0.517661} & \raggedright [\resultnum{v4x.x4.x4a.r080.a0030.x4a.auc_ref_interval_95.i000}{0.505976}, \resultnum{v4x.x4.x4a.r080.a0030.x4a.auc_ref_interval_95.i001}{0.530202}] \tabularnewline
\raggedright \resultnum{v4x.x4.x4a.r080.a0100.cell.a}{0.01}, \resultnum{v4x.x4.x4a.r080.a0100.cell.r}{0.8} & \raggedright \resultnum{v4x.x4.x4a.r080.a0100.x4a.normal_theory_power.at_auc_ref}{0.0729881} & \raggedright inside & \raggedright \resultnum{v4x.x4.x4a.r080.a0100.x4a.auc_ref}{0.553719} & \raggedright [\resultnum{v4x.x4.x4a.r080.a0100.x4a.auc_ref_interval_95.i000}{0.542196}, \resultnum{v4x.x4.x4a.r080.a0100.x4a.auc_ref_interval_95.i001}{0.565979}] \tabularnewline
\raggedright \resultnum{v4x.x4.x4a.r080.a0300.cell.a}{0.03}, \resultnum{v4x.x4.x4a.r080.a0300.cell.r}{0.8} & \raggedright \resultnum{v4x.x4.x4a.r080.a0300.x4a.normal_theory_power.at_auc_ref}{0.105703} & \raggedright inside & \raggedright \resultnum{v4x.x4.x4a.r080.a0300.x4a.auc_ref}{0.610047} & \raggedright [\resultnum{v4x.x4.x4a.r080.a0300.x4a.auc_ref_interval_95.i000}{0.598597}, \resultnum{v4x.x4.x4a.r080.a0300.x4a.auc_ref_interval_95.i001}{0.621413}] \tabularnewline
\raggedright \resultnum{v4x.x4.x4a.r080.a1000.cell.a}{0.1}, \resultnum{v4x.x4.x4a.r080.a1000.cell.r}{0.8} & \raggedright \resultnum{v4x.x4.x4a.r080.a1000.x4a.normal_theory_power.at_auc_ref}{0.290487} & \raggedright inside & \raggedright \resultnum{v4x.x4.x4a.r080.a1000.x4a.auc_ref}{0.780187} & \raggedright [\resultnum{v4x.x4.x4a.r080.a1000.x4a.auc_ref_interval_95.i000}{0.77081}, \resultnum{v4x.x4.x4a.r080.a1000.x4a.auc_ref_interval_95.i001}{0.789852}] \tabularnewline
\raggedright \resultnum{v4x.x4.x4a.r080.a3000.cell.a}{0.3}, \resultnum{v4x.x4.x4a.r080.a3000.cell.r}{0.8} & \raggedright \resultnum{v4x.x4.x4a.r080.a3000.x4a.normal_theory_power.at_auc_ref}{0.882142} & \raggedright inside & \raggedright \resultnum{v4x.x4.x4a.r080.a3000.x4a.auc_ref}{0.977335} & \raggedright [\resultnum{v4x.x4.x4a.r080.a3000.x4a.auc_ref_interval_95.i000}{0.974588}, \resultnum{v4x.x4.x4a.r080.a3000.x4a.auc_ref_interval_95.i001}{0.979927}] \tabularnewline
\raggedright \resultnum{v4x.x4.x4a.r080.a6000.cell.a}{0.6}, \resultnum{v4x.x4.x4a.r080.a6000.cell.r}{0.8} & \raggedright \resultnum{v4x.x4.x4a.r080.a6000.x4a.normal_theory_power.at_auc_ref}{0.99942} & \raggedright inside & \raggedright \resultnum{v4x.x4.x4a.r080.a6000.x4a.auc_ref}{0.99973} & \raggedright [\resultnum{v4x.x4.x4a.r080.a6000.x4a.auc_ref_interval_95.i000}{0.999629}, \resultnum{v4x.x4.x4a.r080.a6000.x4a.auc_ref_interval_95.i001}{0.999816}] \tabularnewline
\raggedright \resultnum{v4x.x4.x4a.r080.a9000.cell.a}{0.9}, \resultnum{v4x.x4.x4a.r080.a9000.cell.r}{0.8} & \raggedright fit gap & \raggedright  & \raggedright  & \raggedright  \tabularnewline
\raggedright \resultnum{v4x.x4.x4a.r085.a0000.cell.a}{0}, \resultnum{v4x.x4.x4a.r085.a0000.cell.r}{0.85} & \raggedright \resultnum{v4x.x4.x4a.r085.a0000.x4a.normal_theory_power.at_auc_ref}{0.0497827} & \raggedright inside & \raggedright \resultnum{v4x.x4.x4a.r085.a0000.x4a.auc_ref}{0.499418} & \raggedright [\resultnum{v4x.x4.x4a.r085.a0000.x4a.auc_ref_interval_95.i000}{0.487753}, \resultnum{v4x.x4.x4a.r085.a0000.x4a.auc_ref_interval_95.i001}{0.512035}] \tabularnewline
\raggedright \resultnum{v4x.x4.x4a.r085.a0010.cell.a}{0.001}, \resultnum{v4x.x4.x4a.r085.a0010.cell.r}{0.85} & \raggedright \resultnum{v4x.x4.x4a.r085.a0010.x4a.normal_theory_power.at_auc_ref}{0.0503297} & \raggedright inside & \raggedright \resultnum{v4x.x4.x4a.r085.a0010.x4a.auc_ref}{0.500912} & \raggedright [\resultnum{v4x.x4.x4a.r085.a0010.x4a.auc_ref_interval_95.i000}{0.489021}, \resultnum{v4x.x4.x4a.r085.a0010.x4a.auc_ref_interval_95.i001}{0.514203}] \tabularnewline
\raggedright \resultnum{v4x.x4.x4a.r085.a0030.cell.a}{0.003}, \resultnum{v4x.x4.x4a.r085.a0030.cell.r}{0.85} & \raggedright \resultnum{v4x.x4.x4a.r085.a0030.x4a.normal_theory_power.at_auc_ref}{0.0576028} & \raggedright inside & \raggedright \resultnum{v4x.x4.x4a.r085.a0030.x4a.auc_ref}{0.519648} & \raggedright [\resultnum{v4x.x4.x4a.r085.a0030.x4a.auc_ref_interval_95.i000}{0.506983}, \resultnum{v4x.x4.x4a.r085.a0030.x4a.auc_ref_interval_95.i001}{0.533327}] \tabularnewline
\raggedright \resultnum{v4x.x4.x4a.r085.a0100.cell.a}{0.01}, \resultnum{v4x.x4.x4a.r085.a0100.cell.r}{0.85} & \raggedright \resultnum{v4x.x4.x4a.r085.a0100.x4a.normal_theory_power.at_auc_ref}{0.0740426} & \raggedright inside & \raggedright \resultnum{v4x.x4.x4a.r085.a0100.x4a.auc_ref}{0.555832} & \raggedright [\resultnum{v4x.x4.x4a.r085.a0100.x4a.auc_ref_interval_95.i000}{0.542701}, \resultnum{v4x.x4.x4a.r085.a0100.x4a.auc_ref_interval_95.i001}{0.568873}] \tabularnewline
\raggedright \resultnum{v4x.x4.x4a.r085.a0300.cell.a}{0.03}, \resultnum{v4x.x4.x4a.r085.a0300.cell.r}{0.85} & \raggedright \resultnum{v4x.x4.x4a.r085.a0300.x4a.normal_theory_power.at_auc_ref}{0.123281} & \raggedright inside & \raggedright \resultnum{v4x.x4.x4a.r085.a0300.x4a.auc_ref}{0.63449} & \raggedright [\resultnum{v4x.x4.x4a.r085.a0300.x4a.auc_ref_interval_95.i000}{0.62137}, \resultnum{v4x.x4.x4a.r085.a0300.x4a.auc_ref_interval_95.i001}{0.646649}] \tabularnewline
\raggedright \resultnum{v4x.x4.x4a.r085.a1000.cell.a}{0.1}, \resultnum{v4x.x4.x4a.r085.a1000.cell.r}{0.85} & \raggedright \resultnum{v4x.x4.x4a.r085.a1000.x4a.normal_theory_power.at_auc_ref}{0.318669} & \raggedright inside & \raggedright \resultnum{v4x.x4.x4a.r085.a1000.x4a.auc_ref}{0.796667} & \raggedright [\resultnum{v4x.x4.x4a.r085.a1000.x4a.auc_ref_interval_95.i000}{0.78664}, \resultnum{v4x.x4.x4a.r085.a1000.x4a.auc_ref_interval_95.i001}{0.806526}] \tabularnewline
\raggedright \resultnum{v4x.x4.x4a.r085.a3000.cell.a}{0.3}, \resultnum{v4x.x4.x4a.r085.a3000.cell.r}{0.85} & \raggedright \resultnum{v4x.x4.x4a.r085.a3000.x4a.normal_theory_power.at_auc_ref}{0.922569} & \raggedright inside & \raggedright \resultnum{v4x.x4.x4a.r085.a3000.x4a.auc_ref}{0.98496} & \raggedright [\resultnum{v4x.x4.x4a.r085.a3000.x4a.auc_ref_interval_95.i000}{0.982817}, \resultnum{v4x.x4.x4a.r085.a3000.x4a.auc_ref_interval_95.i001}{0.986861}] \tabularnewline
\raggedright \resultnum{v4x.x4.x4a.r085.a6000.cell.a}{0.6}, \resultnum{v4x.x4.x4a.r085.a6000.cell.r}{0.85} & \raggedright \resultnum{v4x.x4.x4a.r085.a6000.x4a.normal_theory_power.at_auc_ref}{0.999606} & \raggedright inside & \raggedright \resultnum{v4x.x4.x4a.r085.a6000.x4a.auc_ref}{0.999798} & \raggedright [\resultnum{v4x.x4.x4a.r085.a6000.x4a.auc_ref_interval_95.i000}{0.99955}, \resultnum{v4x.x4.x4a.r085.a6000.x4a.auc_ref_interval_95.i001}{0.999937}] \tabularnewline
\raggedright \resultnum{v4x.x4.x4a.r085.a9000.cell.a}{0.9}, \resultnum{v4x.x4.x4a.r085.a9000.cell.r}{0.85} & \raggedright \resultnum{v4x.x4.x4a.r085.a9000.x4a.normal_theory_power.at_auc_ref}{0.999999} & \raggedright below & \raggedright \resultnum{v4x.x4.x4a.r085.a9000.x4a.auc_ref}{0.999997} & \raggedright [\resultnum{v4x.x4.x4a.r085.a9000.x4a.auc_ref_interval_95.i000}{0.999991}, \resultnum{v4x.x4.x4a.r085.a9000.x4a.auc_ref_interval_95.i001}{1}] \tabularnewline
\raggedright \resultnum{v4x.x4.x4a.r090.a0000.cell.a}{0}, \resultnum{v4x.x4.x4a.r090.a0000.cell.r}{0.9} & \raggedright \resultnum{v4x.x4.x4a.r090.a0000.x4a.normal_theory_power.at_auc_ref}{0.049955} & \raggedright inside & \raggedright \resultnum{v4x.x4.x4a.r090.a0000.x4a.auc_ref}{0.49989} & \raggedright [\resultnum{v4x.x4.x4a.r090.a0000.x4a.auc_ref_interval_95.i000}{0.487319}, \resultnum{v4x.x4.x4a.r090.a0000.x4a.auc_ref_interval_95.i001}{0.512796}] \tabularnewline
\raggedright \resultnum{v4x.x4.x4a.r090.a0010.cell.a}{0.001}, \resultnum{v4x.x4.x4a.r090.a0010.cell.r}{0.9} & \raggedright \resultnum{v4x.x4.x4a.r090.a0010.x4a.normal_theory_power.at_auc_ref}{0.0563145} & \raggedright inside & \raggedright \resultnum{v4x.x4.x4a.r090.a0010.x4a.auc_ref}{0.516473} & \raggedright [\resultnum{v4x.x4.x4a.r090.a0010.x4a.auc_ref_interval_95.i000}{0.503756}, \resultnum{v4x.x4.x4a.r090.a0010.x4a.auc_ref_interval_95.i001}{0.529789}] \tabularnewline
\raggedright \resultnum{v4x.x4.x4a.r090.a0030.cell.a}{0.003}, \resultnum{v4x.x4.x4a.r090.a0030.cell.r}{0.9} & \raggedright \resultnum{v4x.x4.x4a.r090.a0030.x4a.normal_theory_power.at_auc_ref}{0.0601062} & \raggedright inside & \raggedright \resultnum{v4x.x4.x4a.r090.a0030.x4a.auc_ref}{0.525658} & \raggedright [\resultnum{v4x.x4.x4a.r090.a0030.x4a.auc_ref_interval_95.i000}{0.512989}, \resultnum{v4x.x4.x4a.r090.a0030.x4a.auc_ref_interval_95.i001}{0.539112}] \tabularnewline
\raggedright \resultnum{v4x.x4.x4a.r090.a0100.cell.a}{0.01}, \resultnum{v4x.x4.x4a.r090.a0100.cell.r}{0.9} & \raggedright \resultnum{v4x.x4.x4a.r090.a0100.x4a.normal_theory_power.at_auc_ref}{0.0752761} & \raggedright inside & \raggedright \resultnum{v4x.x4.x4a.r090.a0100.x4a.auc_ref}{0.558274} & \raggedright [\resultnum{v4x.x4.x4a.r090.a0100.x4a.auc_ref_interval_95.i000}{0.545681}, \resultnum{v4x.x4.x4a.r090.a0100.x4a.auc_ref_interval_95.i001}{0.571475}] \tabularnewline
\raggedright \resultnum{v4x.x4.x4a.r090.a0300.cell.a}{0.03}, \resultnum{v4x.x4.x4a.r090.a0300.cell.r}{0.9} & \raggedright \resultnum{v4x.x4.x4a.r090.a0300.x4a.normal_theory_power.at_auc_ref}{0.126158} & \raggedright inside & \raggedright \resultnum{v4x.x4.x4a.r090.a0300.x4a.auc_ref}{0.638205} & \raggedright [\resultnum{v4x.x4.x4a.r090.a0300.x4a.auc_ref_interval_95.i000}{0.625796}, \resultnum{v4x.x4.x4a.r090.a0300.x4a.auc_ref_interval_95.i001}{0.651345}] \tabularnewline
\raggedright \resultnum{v4x.x4.x4a.r090.a1000.cell.a}{0.1}, \resultnum{v4x.x4.x4a.r090.a1000.cell.r}{0.9} & \raggedright \resultnum{v4x.x4.x4a.r090.a1000.x4a.normal_theory_power.at_auc_ref}{0.375647} & \raggedright inside & \raggedright \resultnum{v4x.x4.x4a.r090.a1000.x4a.auc_ref}{0.826138} & \raggedright [\resultnum{v4x.x4.x4a.r090.a1000.x4a.auc_ref_interval_95.i000}{0.816889}, \resultnum{v4x.x4.x4a.r090.a1000.x4a.auc_ref_interval_95.i001}{0.83525}] \tabularnewline
\raggedright \resultnum{v4x.x4.x4a.r090.a3000.cell.a}{0.3}, \resultnum{v4x.x4.x4a.r090.a3000.cell.r}{0.9} & \raggedright \resultnum{v4x.x4.x4a.r090.a3000.x4a.normal_theory_power.at_auc_ref}{0.948507} & \raggedright inside & \raggedright \resultnum{v4x.x4.x4a.r090.a3000.x4a.auc_ref}{0.989723} & \raggedright [\resultnum{v4x.x4.x4a.r090.a3000.x4a.auc_ref_interval_95.i000}{0.988171}, \resultnum{v4x.x4.x4a.r090.a3000.x4a.auc_ref_interval_95.i001}{0.991254}] \tabularnewline
\raggedright \resultnum{v4x.x4.x4a.r090.a6000.cell.a}{0.6}, \resultnum{v4x.x4.x4a.r090.a6000.cell.r}{0.9} & \raggedright \resultnum{v4x.x4.x4a.r090.a6000.x4a.normal_theory_power.at_auc_ref}{0.999966} & \raggedright below & \raggedright \resultnum{v4x.x4.x4a.r090.a6000.x4a.auc_ref}{0.999966} & \raggedright [\resultnum{v4x.x4.x4a.r090.a6000.x4a.auc_ref_interval_95.i000}{0.999942}, \resultnum{v4x.x4.x4a.r090.a6000.x4a.auc_ref_interval_95.i001}{0.999985}] \tabularnewline
\raggedright \resultnum{v4x.x4.x4a.r090.a9000.cell.a}{0.9}, \resultnum{v4x.x4.x4a.r090.a9000.cell.r}{0.9} & \raggedright \resultnum{v4x.x4.x4a.r090.a9000.x4a.normal_theory_power.at_auc_ref}{1} & \raggedright inside & \raggedright \resultnum{v4x.x4.x4a.r090.a9000.x4a.auc_ref}{1} & \raggedright [\resultnum{v4x.x4.x4a.r090.a9000.x4a.auc_ref_interval_95.i000}{1}, \resultnum{v4x.x4.x4a.r090.a9000.x4a.auc_ref_interval_95.i001}{1}] \tabularnewline
\botrule
\end{tabular}
\end{table*}

% Generated by v4_exploratory_export; edited during the journal revision
% (agent S, 2026-10-13): captions, headers and categorical labels translated
% into plain English. Every number keeps its \resultnum binding.
\begin{table*}[htbp]
\setlength{\tabcolsep}{2pt}
\caption{Fitted single-panel test (X4b): failed single-panel fits out of 200 panels, the rejection rate among successful fits only, and the label the declaration gives to the rates at that $r$ from the procedure's size at $a=0$: rejection rate (size above 0.05; machine label rejection-rate), power (size compatible with 0.05), power, lower bound (size below 0.05; power-lower-bound). In the primary rates of Table~\ref{tab:v4x-x4b} failed fits count as non-rejections. The labels are frozen; every entry is reported as a rejection rate, because the size was checked at one null configuration only and no valid procedure is named whose power a lower bound would refer to. Fit gap: the frozen fit of B failed, so the cell has no test.}\label{tab:v4x-x4-failures}
\centering\footnotesize
\begin{tabular}{@{}p{0.150\textwidth}p{0.150\textwidth}p{0.240\textwidth}p{0.160\textwidth}p{0.220\textwidth}@{}}
\toprule
\raggedright $a,r$ & \raggedright \mbox{Failed} \mbox{/} \mbox{panels} & \raggedright \mbox{Rejection} \mbox{rate,} \mbox{successful} \mbox{fits} & \raggedright \mbox{Successful} \mbox{panels} & \raggedright \mbox{Quantity} \tabularnewline
\midrule
\raggedright \resultnum{v4x.x4.x4b.r050.a0000.cell.a}{0}, \resultnum{v4x.x4.x4b.r050.a0000.cell.r}{0.5} & \raggedright \resultnum{v4x.x4.x4b.r050.a0000.x4b.fit_failures}{5}/\resultnum{v4x.x4.x4b.r050.a0000.x4b.test_panels}{200} & \raggedright \resultnum{v4x.x4.x4b.r050.a0000.x4b.power_pu_without_failures}{0.307692} & \raggedright \resultnum{v4x.x4.x4b.r050.a0000.successful_panels}{195} & \raggedright rejection rate \tabularnewline
\raggedright \resultnum{v4x.x4.x4b.r050.a0010.cell.a}{0.001}, \resultnum{v4x.x4.x4b.r050.a0010.cell.r}{0.5} & \raggedright \resultnum{v4x.x4.x4b.r050.a0010.x4b.fit_failures}{8}/\resultnum{v4x.x4.x4b.r050.a0010.x4b.test_panels}{200} & \raggedright \resultnum{v4x.x4.x4b.r050.a0010.x4b.power_pu_without_failures}{0.463542} & \raggedright \resultnum{v4x.x4.x4b.r050.a0010.successful_panels}{192} & \raggedright rejection rate \tabularnewline
\raggedright \resultnum{v4x.x4.x4b.r050.a0030.cell.a}{0.003}, \resultnum{v4x.x4.x4b.r050.a0030.cell.r}{0.5} & \raggedright \resultnum{v4x.x4.x4b.r050.a0030.x4b.fit_failures}{11}/\resultnum{v4x.x4.x4b.r050.a0030.x4b.test_panels}{200} & \raggedright \resultnum{v4x.x4.x4b.r050.a0030.x4b.power_pu_without_failures}{0.455026} & \raggedright \resultnum{v4x.x4.x4b.r050.a0030.successful_panels}{189} & \raggedright rejection rate \tabularnewline
\raggedright \resultnum{v4x.x4.x4b.r050.a0100.cell.a}{0.01}, \resultnum{v4x.x4.x4b.r050.a0100.cell.r}{0.5} & \raggedright \resultnum{v4x.x4.x4b.r050.a0100.x4b.fit_failures}{12}/\resultnum{v4x.x4.x4b.r050.a0100.x4b.test_panels}{200} & \raggedright \resultnum{v4x.x4.x4b.r050.a0100.x4b.power_pu_without_failures}{0.446809} & \raggedright \resultnum{v4x.x4.x4b.r050.a0100.successful_panels}{188} & \raggedright rejection rate \tabularnewline
\raggedright \resultnum{v4x.x4.x4b.r050.a0300.cell.a}{0.03}, \resultnum{v4x.x4.x4b.r050.a0300.cell.r}{0.5} & \raggedright \resultnum{v4x.x4.x4b.r050.a0300.x4b.fit_failures}{20}/\resultnum{v4x.x4.x4b.r050.a0300.x4b.test_panels}{200} & \raggedright \resultnum{v4x.x4.x4b.r050.a0300.x4b.power_pu_without_failures}{0.444444} & \raggedright \resultnum{v4x.x4.x4b.r050.a0300.successful_panels}{180} & \raggedright rejection rate \tabularnewline
\raggedright \resultnum{v4x.x4.x4b.r050.a1000.cell.a}{0.1}, \resultnum{v4x.x4.x4b.r050.a1000.cell.r}{0.5} & \raggedright \resultnum{v4x.x4.x4b.r050.a1000.x4b.fit_failures}{23}/\resultnum{v4x.x4.x4b.r050.a1000.x4b.test_panels}{200} & \raggedright \resultnum{v4x.x4.x4b.r050.a1000.x4b.power_pu_without_failures}{0.468927} & \raggedright \resultnum{v4x.x4.x4b.r050.a1000.successful_panels}{177} & \raggedright rejection rate \tabularnewline
\raggedright \resultnum{v4x.x4.x4b.r050.a3000.cell.a}{0.3}, \resultnum{v4x.x4.x4b.r050.a3000.cell.r}{0.5} & \raggedright \resultnum{v4x.x4.x4b.r050.a3000.x4b.fit_failures}{36}/\resultnum{v4x.x4.x4b.r050.a3000.x4b.test_panels}{200} & \raggedright \resultnum{v4x.x4.x4b.r050.a3000.x4b.power_pu_without_failures}{0.609756} & \raggedright \resultnum{v4x.x4.x4b.r050.a3000.successful_panels}{164} & \raggedright rejection rate \tabularnewline
\raggedright \resultnum{v4x.x4.x4b.r050.a6000.cell.a}{0.6}, \resultnum{v4x.x4.x4b.r050.a6000.cell.r}{0.5} & \raggedright \resultnum{v4x.x4.x4b.r050.a6000.x4b.fit_failures}{62}/\resultnum{v4x.x4.x4b.r050.a6000.x4b.test_panels}{200} & \raggedright \resultnum{v4x.x4.x4b.r050.a6000.x4b.power_pu_without_failures}{0.710145} & \raggedright \resultnum{v4x.x4.x4b.r050.a6000.successful_panels}{138} & \raggedright rejection rate \tabularnewline
\raggedright \resultnum{v4x.x4.x4b.r050.a9000.cell.a}{0.9}, \resultnum{v4x.x4.x4b.r050.a9000.cell.r}{0.5} & \raggedright \resultnum{v4x.x4.x4b.r050.a9000.x4b.fit_failures}{127}/\resultnum{v4x.x4.x4b.r050.a9000.x4b.test_panels}{200} & \raggedright \resultnum{v4x.x4.x4b.r050.a9000.x4b.power_pu_without_failures}{0.794521} & \raggedright \resultnum{v4x.x4.x4b.r050.a9000.successful_panels}{73} & \raggedright rejection rate \tabularnewline
\raggedright \resultnum{v4x.x4.x4b.r070.a0000.cell.a}{0}, \resultnum{v4x.x4.x4b.r070.a0000.cell.r}{0.7} & \raggedright \resultnum{v4x.x4.x4b.r070.a0000.x4b.fit_failures}{2}/\resultnum{v4x.x4.x4b.r070.a0000.x4b.test_panels}{200} & \raggedright \resultnum{v4x.x4.x4b.r070.a0000.x4b.power_pu_without_failures}{0.0353535} & \raggedright \resultnum{v4x.x4.x4b.r070.a0000.successful_panels}{198} & \raggedright power \tabularnewline
\raggedright \resultnum{v4x.x4.x4b.r070.a0010.cell.a}{0.001}, \resultnum{v4x.x4.x4b.r070.a0010.cell.r}{0.7} & \raggedright \resultnum{v4x.x4.x4b.r070.a0010.x4b.fit_failures}{1}/\resultnum{v4x.x4.x4b.r070.a0010.x4b.test_panels}{200} & \raggedright \resultnum{v4x.x4.x4b.r070.a0010.x4b.power_pu_without_failures}{0.331658} & \raggedright \resultnum{v4x.x4.x4b.r070.a0010.successful_panels}{199} & \raggedright power \tabularnewline
\raggedright \resultnum{v4x.x4.x4b.r070.a0030.cell.a}{0.003}, \resultnum{v4x.x4.x4b.r070.a0030.cell.r}{0.7} & \raggedright \resultnum{v4x.x4.x4b.r070.a0030.x4b.fit_failures}{2}/\resultnum{v4x.x4.x4b.r070.a0030.x4b.test_panels}{200} & \raggedright \resultnum{v4x.x4.x4b.r070.a0030.x4b.power_pu_without_failures}{0.338384} & \raggedright \resultnum{v4x.x4.x4b.r070.a0030.successful_panels}{198} & \raggedright power \tabularnewline
\raggedright \resultnum{v4x.x4.x4b.r070.a0100.cell.a}{0.01}, \resultnum{v4x.x4.x4b.r070.a0100.cell.r}{0.7} & \raggedright \resultnum{v4x.x4.x4b.r070.a0100.x4b.fit_failures}{5}/\resultnum{v4x.x4.x4b.r070.a0100.x4b.test_panels}{200} & \raggedright \resultnum{v4x.x4.x4b.r070.a0100.x4b.power_pu_without_failures}{0.374359} & \raggedright \resultnum{v4x.x4.x4b.r070.a0100.successful_panels}{195} & \raggedright power \tabularnewline
\raggedright \resultnum{v4x.x4.x4b.r070.a0300.cell.a}{0.03}, \resultnum{v4x.x4.x4b.r070.a0300.cell.r}{0.7} & \raggedright \resultnum{v4x.x4.x4b.r070.a0300.x4b.fit_failures}{5}/\resultnum{v4x.x4.x4b.r070.a0300.x4b.test_panels}{200} & \raggedright \resultnum{v4x.x4.x4b.r070.a0300.x4b.power_pu_without_failures}{0.364103} & \raggedright \resultnum{v4x.x4.x4b.r070.a0300.successful_panels}{195} & \raggedright power \tabularnewline
\raggedright \resultnum{v4x.x4.x4b.r070.a1000.cell.a}{0.1}, \resultnum{v4x.x4.x4b.r070.a1000.cell.r}{0.7} & \raggedright \resultnum{v4x.x4.x4b.r070.a1000.x4b.fit_failures}{7}/\resultnum{v4x.x4.x4b.r070.a1000.x4b.test_panels}{200} & \raggedright \resultnum{v4x.x4.x4b.r070.a1000.x4b.power_pu_without_failures}{0.398964} & \raggedright \resultnum{v4x.x4.x4b.r070.a1000.successful_panels}{193} & \raggedright power \tabularnewline
\raggedright \resultnum{v4x.x4.x4b.r070.a3000.cell.a}{0.3}, \resultnum{v4x.x4.x4b.r070.a3000.cell.r}{0.7} & \raggedright \resultnum{v4x.x4.x4b.r070.a3000.x4b.fit_failures}{17}/\resultnum{v4x.x4.x4b.r070.a3000.x4b.test_panels}{200} & \raggedright \resultnum{v4x.x4.x4b.r070.a3000.x4b.power_pu_without_failures}{0.562842} & \raggedright \resultnum{v4x.x4.x4b.r070.a3000.successful_panels}{183} & \raggedright power \tabularnewline
\raggedright \resultnum{v4x.x4.x4b.r070.a6000.cell.a}{0.6}, \resultnum{v4x.x4.x4b.r070.a6000.cell.r}{0.7} & \raggedright fit gap & \raggedright  & \raggedright  & \raggedright power \tabularnewline
\raggedright \resultnum{v4x.x4.x4b.r070.a9000.cell.a}{0.9}, \resultnum{v4x.x4.x4b.r070.a9000.cell.r}{0.7} & \raggedright \resultnum{v4x.x4.x4b.r070.a9000.x4b.fit_failures}{105}/\resultnum{v4x.x4.x4b.r070.a9000.x4b.test_panels}{200} & \raggedright \resultnum{v4x.x4.x4b.r070.a9000.x4b.power_pu_without_failures}{0.873684} & \raggedright \resultnum{v4x.x4.x4b.r070.a9000.successful_panels}{95} & \raggedright power \tabularnewline
\raggedright \resultnum{v4x.x4.x4b.r090.a0000.cell.a}{0}, \resultnum{v4x.x4.x4b.r090.a0000.cell.r}{0.9} & \raggedright \resultnum{v4x.x4.x4b.r090.a0000.x4b.fit_failures}{0}/\resultnum{v4x.x4.x4b.r090.a0000.x4b.test_panels}{200} & \raggedright \resultnum{v4x.x4.x4b.r090.a0000.x4b.power_pu_without_failures}{0.02} & \raggedright \resultnum{v4x.x4.x4b.r090.a0000.successful_panels}{200} & \raggedright power, lower bound \tabularnewline
\raggedright \resultnum{v4x.x4.x4b.r090.a0010.cell.a}{0.001}, \resultnum{v4x.x4.x4b.r090.a0010.cell.r}{0.9} & \raggedright \resultnum{v4x.x4.x4b.r090.a0010.x4b.fit_failures}{0}/\resultnum{v4x.x4.x4b.r090.a0010.x4b.test_panels}{200} & \raggedright \resultnum{v4x.x4.x4b.r090.a0010.x4b.power_pu_without_failures}{0.16} & \raggedright \resultnum{v4x.x4.x4b.r090.a0010.successful_panels}{200} & \raggedright power, lower bound \tabularnewline
\raggedright \resultnum{v4x.x4.x4b.r090.a0030.cell.a}{0.003}, \resultnum{v4x.x4.x4b.r090.a0030.cell.r}{0.9} & \raggedright \resultnum{v4x.x4.x4b.r090.a0030.x4b.fit_failures}{0}/\resultnum{v4x.x4.x4b.r090.a0030.x4b.test_panels}{200} & \raggedright \resultnum{v4x.x4.x4b.r090.a0030.x4b.power_pu_without_failures}{0.205} & \raggedright \resultnum{v4x.x4.x4b.r090.a0030.successful_panels}{200} & \raggedright power, lower bound \tabularnewline
\raggedright \resultnum{v4x.x4.x4b.r090.a0100.cell.a}{0.01}, \resultnum{v4x.x4.x4b.r090.a0100.cell.r}{0.9} & \raggedright \resultnum{v4x.x4.x4b.r090.a0100.x4b.fit_failures}{0}/\resultnum{v4x.x4.x4b.r090.a0100.x4b.test_panels}{200} & \raggedright \resultnum{v4x.x4.x4b.r090.a0100.x4b.power_pu_without_failures}{0.23} & \raggedright \resultnum{v4x.x4.x4b.r090.a0100.successful_panels}{200} & \raggedright power, lower bound \tabularnewline
\raggedright \resultnum{v4x.x4.x4b.r090.a0300.cell.a}{0.03}, \resultnum{v4x.x4.x4b.r090.a0300.cell.r}{0.9} & \raggedright \resultnum{v4x.x4.x4b.r090.a0300.x4b.fit_failures}{0}/\resultnum{v4x.x4.x4b.r090.a0300.x4b.test_panels}{200} & \raggedright \resultnum{v4x.x4.x4b.r090.a0300.x4b.power_pu_without_failures}{0.29} & \raggedright \resultnum{v4x.x4.x4b.r090.a0300.successful_panels}{200} & \raggedright power, lower bound \tabularnewline
\raggedright \resultnum{v4x.x4.x4b.r090.a1000.cell.a}{0.1}, \resultnum{v4x.x4.x4b.r090.a1000.cell.r}{0.9} & \raggedright \resultnum{v4x.x4.x4b.r090.a1000.x4b.fit_failures}{0}/\resultnum{v4x.x4.x4b.r090.a1000.x4b.test_panels}{200} & \raggedright \resultnum{v4x.x4.x4b.r090.a1000.x4b.power_pu_without_failures}{0.415} & \raggedright \resultnum{v4x.x4.x4b.r090.a1000.successful_panels}{200} & \raggedright power, lower bound \tabularnewline
\raggedright \resultnum{v4x.x4.x4b.r090.a3000.cell.a}{0.3}, \resultnum{v4x.x4.x4b.r090.a3000.cell.r}{0.9} & \raggedright \resultnum{v4x.x4.x4b.r090.a3000.x4b.fit_failures}{1}/\resultnum{v4x.x4.x4b.r090.a3000.x4b.test_panels}{200} & \raggedright \resultnum{v4x.x4.x4b.r090.a3000.x4b.power_pu_without_failures}{0.743719} & \raggedright \resultnum{v4x.x4.x4b.r090.a3000.successful_panels}{199} & \raggedright power, lower bound \tabularnewline
\raggedright \resultnum{v4x.x4.x4b.r090.a6000.cell.a}{0.6}, \resultnum{v4x.x4.x4b.r090.a6000.cell.r}{0.9} & \raggedright \resultnum{v4x.x4.x4b.r090.a6000.x4b.fit_failures}{15}/\resultnum{v4x.x4.x4b.r090.a6000.x4b.test_panels}{200} & \raggedright \resultnum{v4x.x4.x4b.r090.a6000.x4b.power_pu_without_failures}{0.989189} & \raggedright \resultnum{v4x.x4.x4b.r090.a6000.successful_panels}{185} & \raggedright power, lower bound \tabularnewline
\raggedright \resultnum{v4x.x4.x4b.r090.a9000.cell.a}{0.9}, \resultnum{v4x.x4.x4b.r090.a9000.cell.r}{0.9} & \raggedright \resultnum{v4x.x4.x4b.r090.a9000.x4b.fit_failures}{52}/\resultnum{v4x.x4.x4b.r090.a9000.x4b.test_panels}{200} & \raggedright \resultnum{v4x.x4.x4b.r090.a9000.x4b.power_pu_without_failures}{1} & \raggedright \resultnum{v4x.x4.x4b.r090.a9000.successful_panels}{148} & \raggedright power, lower bound \tabularnewline
\botrule
\end{tabular}
\end{table*}

% Generated by v4_exploratory_export; edited during the journal revision
% (agent S, 2026-10-13): captions, headers and categorical labels translated
% into plain English. Every number keeps its \resultnum binding.
\begin{table*}[htbp]
\setlength{\tabcolsep}{2pt}
\caption{Paired difference between the two single-panel procedures on the same 200 panels: fitted-panel (X4b) minus oracle (X4a) rejection rate, with bootstrap 95\% interval. Reading under the declared rule: negative, interval entirely below zero (machine label removes-power); within margin, interval inside $\pm$0.05 (no-material-change); neither, any other interval, including one entirely above zero (frozen label inconclusive). No reading is made where the fitted-panel size at $a=0$ is above 0.05 ($r=0.50$). The difference combines fitting B to one panel, simulation precision and the treatment of failed fits, so it does not isolate the effect of fitting. The last two columns give each procedure's own size flag at that $r$ (Table~\ref{tab:v4x-x4-calibration}).}\label{tab:v4x-x4-paired}
\centering\footnotesize
\begin{tabular}{@{}p{0.137\textwidth}p{0.107\textwidth}p{0.176\textwidth}p{0.185\textwidth}p{0.166\textwidth}p{0.166\textwidth}@{}}
\toprule
\raggedright $a,r$ & \raggedright \mbox{Difference} & \raggedright \mbox{Interval} & \raggedright \mbox{Reading} & \raggedright \mbox{Fitted} \mbox{(X4b)} \mbox{size} & \raggedright \mbox{Oracle} \mbox{(X4a)} \mbox{size} \tabularnewline
\midrule
\raggedright \resultnum{v4x.x4.x4b.r050.a0000.cell.a}{0}, \resultnum{v4x.x4.x4b.r050.a0000.cell.r}{0.5} & \raggedright \resultnum{v4x.x4.x4b.r050.a0000.x4b.difference}{0.24} & \raggedright [\resultnum{v4x.x4.x4b.r050.a0000.x4b.difference_interval_95.i000}{0.18}, \resultnum{v4x.x4.x4b.r050.a0000.x4b.difference_interval_95.i001}{0.31}] & \raggedright not read & \raggedright above & \raggedright compatible \tabularnewline
\raggedright \resultnum{v4x.x4.x4b.r050.a0010.cell.a}{0.001}, \resultnum{v4x.x4.x4b.r050.a0010.cell.r}{0.5} & \raggedright \resultnum{v4x.x4.x4b.r050.a0010.x4b.difference}{0.37} & \raggedright [\resultnum{v4x.x4.x4b.r050.a0010.x4b.difference_interval_95.i000}{0.3}, \resultnum{v4x.x4.x4b.r050.a0010.x4b.difference_interval_95.i001}{0.44}] & \raggedright not read & \raggedright above & \raggedright compatible \tabularnewline
\raggedright \resultnum{v4x.x4.x4b.r050.a0030.cell.a}{0.003}, \resultnum{v4x.x4.x4b.r050.a0030.cell.r}{0.5} & \raggedright \resultnum{v4x.x4.x4b.r050.a0030.x4b.difference}{0.36} & \raggedright [\resultnum{v4x.x4.x4b.r050.a0030.x4b.difference_interval_95.i000}{0.29}, \resultnum{v4x.x4.x4b.r050.a0030.x4b.difference_interval_95.i001}{0.43}] & \raggedright not read & \raggedright above & \raggedright compatible \tabularnewline
\raggedright \resultnum{v4x.x4.x4b.r050.a0100.cell.a}{0.01}, \resultnum{v4x.x4.x4b.r050.a0100.cell.r}{0.5} & \raggedright \resultnum{v4x.x4.x4b.r050.a0100.x4b.difference}{0.37} & \raggedright [\resultnum{v4x.x4.x4b.r050.a0100.x4b.difference_interval_95.i000}{0.3}, \resultnum{v4x.x4.x4b.r050.a0100.x4b.difference_interval_95.i001}{0.44}] & \raggedright not read & \raggedright above & \raggedright compatible \tabularnewline
\raggedright \resultnum{v4x.x4.x4b.r050.a0300.cell.a}{0.03}, \resultnum{v4x.x4.x4b.r050.a0300.cell.r}{0.5} & \raggedright \resultnum{v4x.x4.x4b.r050.a0300.x4b.difference}{0.33} & \raggedright [\resultnum{v4x.x4.x4b.r050.a0300.x4b.difference_interval_95.i000}{0.265}, \resultnum{v4x.x4.x4b.r050.a0300.x4b.difference_interval_95.i001}{0.395}] & \raggedright not read & \raggedright above & \raggedright compatible \tabularnewline
\raggedright \resultnum{v4x.x4.x4b.r050.a1000.cell.a}{0.1}, \resultnum{v4x.x4.x4b.r050.a1000.cell.r}{0.5} & \raggedright \resultnum{v4x.x4.x4b.r050.a1000.x4b.difference}{0.245} & \raggedright [\resultnum{v4x.x4.x4b.r050.a1000.x4b.difference_interval_95.i000}{0.165}, \resultnum{v4x.x4.x4b.r050.a1000.x4b.difference_interval_95.i001}{0.325}] & \raggedright not read & \raggedright above & \raggedright compatible \tabularnewline
\raggedright \resultnum{v4x.x4.x4b.r050.a3000.cell.a}{0.3}, \resultnum{v4x.x4.x4b.r050.a3000.cell.r}{0.5} & \raggedright \resultnum{v4x.x4.x4b.r050.a3000.x4b.difference}{0.045} & \raggedright [$\resultnum{v4x.x4.x4b.r050.a3000.x4b.difference_interval_95.i000}{-0.05}$, \resultnum{v4x.x4.x4b.r050.a3000.x4b.difference_interval_95.i001}{0.13}] & \raggedright not read & \raggedright above & \raggedright compatible \tabularnewline
\raggedright \resultnum{v4x.x4.x4b.r050.a6000.cell.a}{0.6}, \resultnum{v4x.x4.x4b.r050.a6000.cell.r}{0.5} & \raggedright $\resultnum{v4x.x4.x4b.r050.a6000.x4b.difference}{-0.36}$ & \raggedright [$\resultnum{v4x.x4.x4b.r050.a6000.x4b.difference_interval_95.i000}{-0.445}$, $\resultnum{v4x.x4.x4b.r050.a6000.x4b.difference_interval_95.i001}{-0.27}$] & \raggedright not read & \raggedright above & \raggedright compatible \tabularnewline
\raggedright \resultnum{v4x.x4.x4b.r050.a9000.cell.a}{0.9}, \resultnum{v4x.x4.x4b.r050.a9000.cell.r}{0.5} & \raggedright $\resultnum{v4x.x4.x4b.r050.a9000.x4b.difference}{-0.68}$ & \raggedright [$\resultnum{v4x.x4.x4b.r050.a9000.x4b.difference_interval_95.i000}{-0.75}$, $\resultnum{v4x.x4.x4b.r050.a9000.x4b.difference_interval_95.i001}{-0.605}$] & \raggedright not read & \raggedright above & \raggedright compatible \tabularnewline
\raggedright \resultnum{v4x.x4.x4b.r070.a0000.cell.a}{0}, \resultnum{v4x.x4.x4b.r070.a0000.cell.r}{0.7} & \raggedright $\resultnum{v4x.x4.x4b.r070.a0000.x4b.difference}{-0.005}$ & \raggedright [$\resultnum{v4x.x4.x4b.r070.a0000.x4b.difference_interval_95.i000}{-0.03}$, \resultnum{v4x.x4.x4b.r070.a0000.x4b.difference_interval_95.i001}{0.015}] & \raggedright within margin & \raggedright compatible & \raggedright above \tabularnewline
\raggedright \resultnum{v4x.x4.x4b.r070.a0010.cell.a}{0.001}, \resultnum{v4x.x4.x4b.r070.a0010.cell.r}{0.7} & \raggedright \resultnum{v4x.x4.x4b.r070.a0010.x4b.difference}{0.275} & \raggedright [\resultnum{v4x.x4.x4b.r070.a0010.x4b.difference_interval_95.i000}{0.215}, \resultnum{v4x.x4.x4b.r070.a0010.x4b.difference_interval_95.i001}{0.335}] & \raggedright neither & \raggedright compatible & \raggedright above \tabularnewline
\raggedright \resultnum{v4x.x4.x4b.r070.a0030.cell.a}{0.003}, \resultnum{v4x.x4.x4b.r070.a0030.cell.r}{0.7} & \raggedright \resultnum{v4x.x4.x4b.r070.a0030.x4b.difference}{0.275} & \raggedright [\resultnum{v4x.x4.x4b.r070.a0030.x4b.difference_interval_95.i000}{0.215}, \resultnum{v4x.x4.x4b.r070.a0030.x4b.difference_interval_95.i001}{0.335}] & \raggedright neither & \raggedright compatible & \raggedright above \tabularnewline
\raggedright \resultnum{v4x.x4.x4b.r070.a0100.cell.a}{0.01}, \resultnum{v4x.x4.x4b.r070.a0100.cell.r}{0.7} & \raggedright \resultnum{v4x.x4.x4b.r070.a0100.x4b.difference}{0.295} & \raggedright [\resultnum{v4x.x4.x4b.r070.a0100.x4b.difference_interval_95.i000}{0.225}, \resultnum{v4x.x4.x4b.r070.a0100.x4b.difference_interval_95.i001}{0.365}] & \raggedright neither & \raggedright compatible & \raggedright above \tabularnewline
\raggedright \resultnum{v4x.x4.x4b.r070.a0300.cell.a}{0.03}, \resultnum{v4x.x4.x4b.r070.a0300.cell.r}{0.7} & \raggedright \resultnum{v4x.x4.x4b.r070.a0300.x4b.difference}{0.255} & \raggedright [\resultnum{v4x.x4.x4b.r070.a0300.x4b.difference_interval_95.i000}{0.175}, \resultnum{v4x.x4.x4b.r070.a0300.x4b.difference_interval_95.i001}{0.335}] & \raggedright neither & \raggedright compatible & \raggedright above \tabularnewline
\raggedright \resultnum{v4x.x4.x4b.r070.a1000.cell.a}{0.1}, \resultnum{v4x.x4.x4b.r070.a1000.cell.r}{0.7} & \raggedright \resultnum{v4x.x4.x4b.r070.a1000.x4b.difference}{0.18} & \raggedright [\resultnum{v4x.x4.x4b.r070.a1000.x4b.difference_interval_95.i000}{0.095}, \resultnum{v4x.x4.x4b.r070.a1000.x4b.difference_interval_95.i001}{0.265}] & \raggedright neither & \raggedright compatible & \raggedright above \tabularnewline
\raggedright \resultnum{v4x.x4.x4b.r070.a3000.cell.a}{0.3}, \resultnum{v4x.x4.x4b.r070.a3000.cell.r}{0.7} & \raggedright $\resultnum{v4x.x4.x4b.r070.a3000.x4b.difference}{-0.235}$ & \raggedright [$\resultnum{v4x.x4.x4b.r070.a3000.x4b.difference_interval_95.i000}{-0.315}$, $\resultnum{v4x.x4.x4b.r070.a3000.x4b.difference_interval_95.i001}{-0.155}$] & \raggedright negative & \raggedright compatible & \raggedright above \tabularnewline
\raggedright \resultnum{v4x.x4.x4b.r070.a6000.cell.a}{0.6}, \resultnum{v4x.x4.x4b.r070.a6000.cell.r}{0.7} & \raggedright fit gap & \raggedright  & \raggedright  & \raggedright compatible & \raggedright above \tabularnewline
\raggedright \resultnum{v4x.x4.x4b.r070.a9000.cell.a}{0.9}, \resultnum{v4x.x4.x4b.r070.a9000.cell.r}{0.7} & \raggedright $\resultnum{v4x.x4.x4b.r070.a9000.x4b.difference}{-0.585}$ & \raggedright [$\resultnum{v4x.x4.x4b.r070.a9000.x4b.difference_interval_95.i000}{-0.65}$, $\resultnum{v4x.x4.x4b.r070.a9000.x4b.difference_interval_95.i001}{-0.515}$] & \raggedright negative & \raggedright compatible & \raggedright above \tabularnewline
\raggedright \resultnum{v4x.x4.x4b.r090.a0000.cell.a}{0}, \resultnum{v4x.x4.x4b.r090.a0000.cell.r}{0.9} & \raggedright $\resultnum{v4x.x4.x4b.r090.a0000.x4b.difference}{-0.015}$ & \raggedright [$\resultnum{v4x.x4.x4b.r090.a0000.x4b.difference_interval_95.i000}{-0.05}$, \resultnum{v4x.x4.x4b.r090.a0000.x4b.difference_interval_95.i001}{0.015}] & \raggedright within margin & \raggedright below & \raggedright compatible \tabularnewline
\raggedright \resultnum{v4x.x4.x4b.r090.a0010.cell.a}{0.001}, \resultnum{v4x.x4.x4b.r090.a0010.cell.r}{0.9} & \raggedright \resultnum{v4x.x4.x4b.r090.a0010.x4b.difference}{0.12} & \raggedright [\resultnum{v4x.x4.x4b.r090.a0010.x4b.difference_interval_95.i000}{0.065}, \resultnum{v4x.x4.x4b.r090.a0010.x4b.difference_interval_95.i001}{0.18}] & \raggedright neither & \raggedright below & \raggedright compatible \tabularnewline
\raggedright \resultnum{v4x.x4.x4b.r090.a0030.cell.a}{0.003}, \resultnum{v4x.x4.x4b.r090.a0030.cell.r}{0.9} & \raggedright \resultnum{v4x.x4.x4b.r090.a0030.x4b.difference}{0.175} & \raggedright [\resultnum{v4x.x4.x4b.r090.a0030.x4b.difference_interval_95.i000}{0.11}, \resultnum{v4x.x4.x4b.r090.a0030.x4b.difference_interval_95.i001}{0.245}] & \raggedright neither & \raggedright below & \raggedright compatible \tabularnewline
\raggedright \resultnum{v4x.x4.x4b.r090.a0100.cell.a}{0.01}, \resultnum{v4x.x4.x4b.r090.a0100.cell.r}{0.9} & \raggedright \resultnum{v4x.x4.x4b.r090.a0100.x4b.difference}{0.17} & \raggedright [\resultnum{v4x.x4.x4b.r090.a0100.x4b.difference_interval_95.i000}{0.105}, \resultnum{v4x.x4.x4b.r090.a0100.x4b.difference_interval_95.i001}{0.24}] & \raggedright neither & \raggedright below & \raggedright compatible \tabularnewline
\raggedright \resultnum{v4x.x4.x4b.r090.a0300.cell.a}{0.03}, \resultnum{v4x.x4.x4b.r090.a0300.cell.r}{0.9} & \raggedright \resultnum{v4x.x4.x4b.r090.a0300.x4b.difference}{0.195} & \raggedright [\resultnum{v4x.x4.x4b.r090.a0300.x4b.difference_interval_95.i000}{0.13}, \resultnum{v4x.x4.x4b.r090.a0300.x4b.difference_interval_95.i001}{0.26}] & \raggedright neither & \raggedright below & \raggedright compatible \tabularnewline
\raggedright \resultnum{v4x.x4.x4b.r090.a1000.cell.a}{0.1}, \resultnum{v4x.x4.x4b.r090.a1000.cell.r}{0.9} & \raggedright \resultnum{v4x.x4.x4b.r090.a1000.x4b.difference}{0.005} & \raggedright [$\resultnum{v4x.x4.x4b.r090.a1000.x4b.difference_interval_95.i000}{-0.085}$, \resultnum{v4x.x4.x4b.r090.a1000.x4b.difference_interval_95.i001}{0.095}] & \raggedright neither & \raggedright below & \raggedright compatible \tabularnewline
\raggedright \resultnum{v4x.x4.x4b.r090.a3000.cell.a}{0.3}, \resultnum{v4x.x4.x4b.r090.a3000.cell.r}{0.9} & \raggedright $\resultnum{v4x.x4.x4b.r090.a3000.x4b.difference}{-0.22}$ & \raggedright [$\resultnum{v4x.x4.x4b.r090.a3000.x4b.difference_interval_95.i000}{-0.285}$, $\resultnum{v4x.x4.x4b.r090.a3000.x4b.difference_interval_95.i001}{-0.155}$] & \raggedright negative & \raggedright below & \raggedright compatible \tabularnewline
\raggedright \resultnum{v4x.x4.x4b.r090.a6000.cell.a}{0.6}, \resultnum{v4x.x4.x4b.r090.a6000.cell.r}{0.9} & \raggedright $\resultnum{v4x.x4.x4b.r090.a6000.x4b.difference}{-0.085}$ & \raggedright [$\resultnum{v4x.x4.x4b.r090.a6000.x4b.difference_interval_95.i000}{-0.125}$, $\resultnum{v4x.x4.x4b.r090.a6000.x4b.difference_interval_95.i001}{-0.05}$] & \raggedright negative & \raggedright below & \raggedright compatible \tabularnewline
\raggedright \resultnum{v4x.x4.x4b.r090.a9000.cell.a}{0.9}, \resultnum{v4x.x4.x4b.r090.a9000.cell.r}{0.9} & \raggedright $\resultnum{v4x.x4.x4b.r090.a9000.x4b.difference}{-0.26}$ & \raggedright [$\resultnum{v4x.x4.x4b.r090.a9000.x4b.difference_interval_95.i000}{-0.325}$, $\resultnum{v4x.x4.x4b.r090.a9000.x4b.difference_interval_95.i001}{-0.2}$] & \raggedright negative & \raggedright below & \raggedright compatible \tabularnewline
\botrule
\end{tabular}
\end{table*}

\clearpage

\subsection{Channel-misspecification stress test (X5)}\label{sec:x5}\phantomsection\label{supp:x5}
% Moved from sections/supplement-exploratory-x1-x2.tex; numbers from results/x5.json.
X5, declared on 28 September 2026 before any confirmatory result, fits B, correctly specified for the assumed channel, to the influence arm at $a=0$ observed through a misspecified channel, and reads the frozen pipeline's AUC as the separation produced by misspecification alone. % EX14
The declared conditions, fixed before any run, are: a linear trend in the signal share across the eleven waves, $r_t=r+A(t-5)/5$ with amplitude $A=0.03$ or 0.08; a step at wave 5, $r_t=r-A$ before and $r+A$ from wave 5, with the same amplitudes (wave 5 corresponds to 2004, when the SHP added a refreshment sample); panel conditioning in the panels only, with increments 0.10 and 0.20; and household attrition that depends on the latent state, with retention odds ratios of 1.5 and 2 per unit of the household's mean latent state, so that households with lower latent states leave more often, and the average retention kept at 0.89. In the SHP, those who left the panel had reported less political interest and less trust in the federal government than those who stayed \cite{lipps2007,voorpostel2010}, a bias Voorpostel judged mild; these studies motivate the direction of the attrition scenario, not its size, and all magnitudes are stress scenarios, not estimates for the SHP.
All 72 cells finished, with 36 estimates and 36 gaps across eight conditions and nine values of $r$. % EX14 RC13
With no influence at all, a linear trend in the signal share across the eleven waves gave AUCs from \resultnum{j.x5.trend003.min}{0.680} to \resultnum{j.x5.trend003.max}{0.710} for an amplitude of 0.03 (from $r-0.03$ in the first wave to $r+0.03$ in the last) and from \resultnum{j.x5.trend008.min}{0.884} to \resultnum{j.x5.trend008.max}{0.907} for 0.08, and trust-dependent attrition with a retention odds ratio of 2 gave AUCs from \resultnum{j.x5.attrition2.min}{0.578} at $r=0.90$ to \resultnum{j.x5.attrition2.max}{0.820} at $r=0.50$ (\resultnum{j.x5.attrition15.min}{0.540} to \resultnum{j.x5.attrition15.max}{0.640} for an odds ratio of \resultnum{j.x5.attrition.odds_ratio_small}{1.5}). % EX14
The wave-five step and conditioning conditions each produced nine fit-failure gaps at each of their two amplitudes, giving 36 gaps in total. % EX14 RC14
In 35 of these 36 cells both fitting attempts of B ended on the face of its parameter box with no trait and then failed the independent validation, which for such a fit includes lag two; in the remaining cell the second attempt ended ill-conditioned. In these cells, the declared B fitting procedure did not obtain a validated match to the misspecified lag-one and lag-two correlations at that $r$. These are failed matching outcomes rather than AUC estimates, not proof that B cannot reproduce those correlations.
All 36 intervals computed with critical value 2.7729 excluded 0.5. % EX15 RC15
In this known-$a=0$ experiment the corresponding zero-strength check would invalidate every estimated comparison; a real-data analyst has no known-null arm on which to apply that check. % EX15 RC16

Read after the results became available: in the \resultnum{j.x5.signature.cells}{36} estimated cells, the panels without influence had lower mean lag correlations than the fitted B at lags three to five, with a shortfall that grew with the lag (negative at all three lags in \resultnum{j.x5.signature.negative.n}{36} cells and growing with the lag in \resultnum{j.x5.signature.increasing.n}{36}; Table~\ref{tab:x5-lags} gives each cell's differences, B's fitted $\phi$ and trait weight and the AUC), the sign pattern that influence leaves against known-share B at every first-boundary cell (Section~\ref{sec:x3}). In the median, the shortfall rose less steeply with the lag than under influence: the ratio of the lag-five to the lag-three shortfall had medians over the nine shares of \resultnum{j.x5.signature.ratio_l5_l3.median.mnar-ln15}{4.0} and \resultnum{j.x5.signature.ratio_l5_l3.median.mnar-ln2}{4.4} for attrition (odds ratios 1.5 and 2) and \resultnum{j.x5.signature.ratio_l5_l3.median.trend-003}{4.6} and \resultnum{j.x5.signature.ratio_l5_l3.median.trend-008}{4.8} for the trends (amplitudes 0.03 and 0.08), against \resultnum{j.x3.lag_shortfall.ratio_l5_l3.min}{5.8} to \resultnum{j.x3.lag_shortfall.ratio_l5_l3.max}{6.9} for influence; individual cells overlap (Table~\ref{tab:x5-lags}), and the difference was not tested. Both sets of means are over check panels. Each stress-test AUC comes from a discriminant trained on that condition's own panels; the response of a discriminant trained against influence to these departures was not evaluated.

\begin{table}[tp]
\caption{Mean lag differences under channel misspecification on check panels at $a=\resultnum{j.grid.a0000}{0}$ with B fitted at the known share. The influence-trained score was not evaluated on these panels.}\label{tab:x5-lags}
\centering\scriptsize
\setlength{\tabcolsep}{5pt}
\begin{tabular}{>{\raggedright\arraybackslash}p{36mm}lllllll}
\toprule
Condition & $r$ & $\phi$ & $\lambda$ & Lag three & Lag four & Lag five & AUC \\
\midrule
Signal-share trend, $\pm\resultnum{j.x5.trend-003.amount}{0.03}$ & \resultnum{j.b.r050.r}{0.50} & \resultnum{j.x5.trend-003.r050.phi}{0.836} & \resultnum{j.x5.trend-003.r050.lambda}{0.079} & $\resultnum{j.x5.trend-003.r050.l3}{-0.0010}$ & $\resultnum{j.x5.trend-003.r050.l4}{-0.0028}$ & $\resultnum{j.x5.trend-003.r050.l5}{-0.0051}$ & \resultnum{j.x5.trend-003.r050.auc}{0.684} \\
Signal-share trend, $\pm\resultnum{j.x5.trend-003.amount}{0.03}$ & \resultnum{j.b.r055.r}{0.55} & \resultnum{j.x5.trend-003.r055.phi}{0.835} & \resultnum{j.x5.trend-003.r055.lambda}{0.081} & $\resultnum{j.x5.trend-003.r055.l3}{-0.0013}$ & $\resultnum{j.x5.trend-003.r055.l4}{-0.0032}$ & $\resultnum{j.x5.trend-003.r055.l5}{-0.0059}$ & \resultnum{j.x5.trend-003.r055.auc}{0.710} \\
Signal-share trend, $\pm\resultnum{j.x5.trend-003.amount}{0.03}$ & \resultnum{j.b.r060.r}{0.60} & \resultnum{j.x5.trend-003.r060.phi}{0.837} & \resultnum{j.x5.trend-003.r060.lambda}{0.071} & $\resultnum{j.x5.trend-003.r060.l3}{-0.0013}$ & $\resultnum{j.x5.trend-003.r060.l4}{-0.0031}$ & $\resultnum{j.x5.trend-003.r060.l5}{-0.0054}$ & \resultnum{j.x5.trend-003.r060.auc}{0.690} \\
Signal-share trend, $\pm\resultnum{j.x5.trend-003.amount}{0.03}$ & \resultnum{j.b.r065.r}{0.65} & \resultnum{j.x5.trend-003.r065.phi}{0.838} & \resultnum{j.x5.trend-003.r065.lambda}{0.065} & $\resultnum{j.x5.trend-003.r065.l3}{-0.0009}$ & $\resultnum{j.x5.trend-003.r065.l4}{-0.0029}$ & $\resultnum{j.x5.trend-003.r065.l5}{-0.0053}$ & \resultnum{j.x5.trend-003.r065.auc}{0.693} \\
Signal-share trend, $\pm\resultnum{j.x5.trend-003.amount}{0.03}$ & \resultnum{j.b.r070.r}{0.70} & \resultnum{j.x5.trend-003.r070.phi}{0.839} & \resultnum{j.x5.trend-003.r070.lambda}{0.061} & $\resultnum{j.x5.trend-003.r070.l3}{-0.0012}$ & $\resultnum{j.x5.trend-003.r070.l4}{-0.0031}$ & $\resultnum{j.x5.trend-003.r070.l5}{-0.0053}$ & \resultnum{j.x5.trend-003.r070.auc}{0.688} \\
Signal-share trend, $\pm\resultnum{j.x5.trend-003.amount}{0.03}$ & \resultnum{j.b.r075.r}{0.75} & \resultnum{j.x5.trend-003.r075.phi}{0.839} & \resultnum{j.x5.trend-003.r075.lambda}{0.059} & $\resultnum{j.x5.trend-003.r075.l3}{-0.0008}$ & $\resultnum{j.x5.trend-003.r075.l4}{-0.0025}$ & $\resultnum{j.x5.trend-003.r075.l5}{-0.0050}$ & \resultnum{j.x5.trend-003.r075.auc}{0.680} \\
Signal-share trend, $\pm\resultnum{j.x5.trend-003.amount}{0.03}$ & \resultnum{j.b.r080.r}{0.80} & \resultnum{j.x5.trend-003.r080.phi}{0.840} & \resultnum{j.x5.trend-003.r080.lambda}{0.055} & $\resultnum{j.x5.trend-003.r080.l3}{-0.0011}$ & $\resultnum{j.x5.trend-003.r080.l4}{-0.0030}$ & $\resultnum{j.x5.trend-003.r080.l5}{-0.0053}$ & \resultnum{j.x5.trend-003.r080.auc}{0.691} \\
Signal-share trend, $\pm\resultnum{j.x5.trend-003.amount}{0.03}$ & \resultnum{j.b.r085.r}{0.85} & \resultnum{j.x5.trend-003.r085.phi}{0.841} & \resultnum{j.x5.trend-003.r085.lambda}{0.052} & $\resultnum{j.x5.trend-003.r085.l3}{-0.0012}$ & $\resultnum{j.x5.trend-003.r085.l4}{-0.0031}$ & $\resultnum{j.x5.trend-003.r085.l5}{-0.0055}$ & \resultnum{j.x5.trend-003.r085.auc}{0.696} \\
Signal-share trend, $\pm\resultnum{j.x5.trend-003.amount}{0.03}$ & \resultnum{j.b.r090.r}{0.90} & \resultnum{j.x5.trend-003.r090.phi}{0.841} & \resultnum{j.x5.trend-003.r090.lambda}{0.049} & $\resultnum{j.x5.trend-003.r090.l3}{-0.0014}$ & $\resultnum{j.x5.trend-003.r090.l4}{-0.0034}$ & $\resultnum{j.x5.trend-003.r090.l5}{-0.0059}$ & \resultnum{j.x5.trend-003.r090.auc}{0.701} \\
Signal-share trend, $\pm\resultnum{j.x5.trend-008.amount}{0.08}$ & \resultnum{j.b.r050.r}{0.50} & \resultnum{j.x5.trend-008.r050.phi}{0.812} & \resultnum{j.x5.trend-008.r050.lambda}{0.181} & $\resultnum{j.x5.trend-008.r050.l3}{-0.0027}$ & $\resultnum{j.x5.trend-008.r050.l4}{-0.0074}$ & $\resultnum{j.x5.trend-008.r050.l5}{-0.0134}$ & \resultnum{j.x5.trend-008.r050.auc}{0.891} \\
Signal-share trend, $\pm\resultnum{j.x5.trend-008.amount}{0.08}$ & \resultnum{j.b.r055.r}{0.55} & \resultnum{j.x5.trend-008.r055.phi}{0.814} & \resultnum{j.x5.trend-008.r055.lambda}{0.175} & $\resultnum{j.x5.trend-008.r055.l3}{-0.0031}$ & $\resultnum{j.x5.trend-008.r055.l4}{-0.0078}$ & $\resultnum{j.x5.trend-008.r055.l5}{-0.0142}$ & \resultnum{j.x5.trend-008.r055.auc}{0.907} \\
Signal-share trend, $\pm\resultnum{j.x5.trend-008.amount}{0.08}$ & \resultnum{j.b.r060.r}{0.60} & \resultnum{j.x5.trend-008.r060.phi}{0.819} & \resultnum{j.x5.trend-008.r060.lambda}{0.154} & $\resultnum{j.x5.trend-008.r060.l3}{-0.0028}$ & $\resultnum{j.x5.trend-008.r060.l4}{-0.0073}$ & $\resultnum{j.x5.trend-008.r060.l5}{-0.0130}$ & \resultnum{j.x5.trend-008.r060.auc}{0.884} \\
Signal-share trend, $\pm\resultnum{j.x5.trend-008.amount}{0.08}$ & \resultnum{j.b.r065.r}{0.65} & \resultnum{j.x5.trend-008.r065.phi}{0.819} & \resultnum{j.x5.trend-008.r065.lambda}{0.153} & $\resultnum{j.x5.trend-008.r065.l3}{-0.0027}$ & $\resultnum{j.x5.trend-008.r065.l4}{-0.0077}$ & $\resultnum{j.x5.trend-008.r065.l5}{-0.0140}$ & \resultnum{j.x5.trend-008.r065.auc}{0.903} \\
Signal-share trend, $\pm\resultnum{j.x5.trend-008.amount}{0.08}$ & \resultnum{j.b.r070.r}{0.70} & \resultnum{j.x5.trend-008.r070.phi}{0.823} & \resultnum{j.x5.trend-008.r070.lambda}{0.138} & $\resultnum{j.x5.trend-008.r070.l3}{-0.0029}$ & $\resultnum{j.x5.trend-008.r070.l4}{-0.0076}$ & $\resultnum{j.x5.trend-008.r070.l5}{-0.0134}$ & \resultnum{j.x5.trend-008.r070.auc}{0.887} \\
Signal-share trend, $\pm\resultnum{j.x5.trend-008.amount}{0.08}$ & \resultnum{j.b.r075.r}{0.75} & \resultnum{j.x5.trend-008.r075.phi}{0.824} & \resultnum{j.x5.trend-008.r075.lambda}{0.133} & $\resultnum{j.x5.trend-008.r075.l3}{-0.0025}$ & $\resultnum{j.x5.trend-008.r075.l4}{-0.0071}$ & $\resultnum{j.x5.trend-008.r075.l5}{-0.0132}$ & \resultnum{j.x5.trend-008.r075.auc}{0.886} \\
Signal-share trend, $\pm\resultnum{j.x5.trend-008.amount}{0.08}$ & \resultnum{j.b.r080.r}{0.80} & \resultnum{j.x5.trend-008.r080.phi}{0.825} & \resultnum{j.x5.trend-008.r080.lambda}{0.126} & $\resultnum{j.x5.trend-008.r080.l3}{-0.0029}$ & $\resultnum{j.x5.trend-008.r080.l4}{-0.0077}$ & $\resultnum{j.x5.trend-008.r080.l5}{-0.0138}$ & \resultnum{j.x5.trend-008.r080.auc}{0.897} \\
Signal-share trend, $\pm\resultnum{j.x5.trend-008.amount}{0.08}$ & \resultnum{j.b.r085.r}{0.85} & \resultnum{j.x5.trend-008.r085.phi}{0.827} & \resultnum{j.x5.trend-008.r085.lambda}{0.119} & $\resultnum{j.x5.trend-008.r085.l3}{-0.0029}$ & $\resultnum{j.x5.trend-008.r085.l4}{-0.0078}$ & $\resultnum{j.x5.trend-008.r085.l5}{-0.0138}$ & \resultnum{j.x5.trend-008.r085.auc}{0.896} \\
Signal-share trend, $\pm\resultnum{j.x5.trend-008.amount}{0.08}$ & \resultnum{j.b.r090.r}{0.90} & \resultnum{j.x5.trend-008.r090.phi}{0.828} & \resultnum{j.x5.trend-008.r090.lambda}{0.113} & $\resultnum{j.x5.trend-008.r090.l3}{-0.0031}$ & $\resultnum{j.x5.trend-008.r090.l4}{-0.0081}$ & $\resultnum{j.x5.trend-008.r090.l5}{-0.0143}$ & \resultnum{j.x5.trend-008.r090.auc}{0.902} \\
Attrition, odds ratio \resultnum{j.x5.mnar-ln15.amount}{1.5} & \resultnum{j.b.r050.r}{0.50} & \resultnum{j.x5.mnar-ln15.r050.phi}{0.836} & \resultnum{j.x5.mnar-ln15.r050.lambda}{0.061} & $\resultnum{j.x5.mnar-ln15.r050.l3}{-0.0008}$ & $\resultnum{j.x5.mnar-ln15.r050.l4}{-0.0022}$ & $\resultnum{j.x5.mnar-ln15.r050.l5}{-0.0038}$ & \resultnum{j.x5.mnar-ln15.r050.auc}{0.640} \\
Attrition, odds ratio \resultnum{j.x5.mnar-ln15.amount}{1.5} & \resultnum{j.b.r055.r}{0.55} & \resultnum{j.x5.mnar-ln15.r055.phi}{0.838} & \resultnum{j.x5.mnar-ln15.r055.lambda}{0.051} & $\resultnum{j.x5.mnar-ln15.r055.l3}{-0.0009}$ & $\resultnum{j.x5.mnar-ln15.r055.l4}{-0.0019}$ & $\resultnum{j.x5.mnar-ln15.r055.l5}{-0.0036}$ & \resultnum{j.x5.mnar-ln15.r055.auc}{0.628} \\
Attrition, odds ratio \resultnum{j.x5.mnar-ln15.amount}{1.5} & \resultnum{j.b.r060.r}{0.60} & \resultnum{j.x5.mnar-ln15.r060.phi}{0.840} & \resultnum{j.x5.mnar-ln15.r060.lambda}{0.043} & $\resultnum{j.x5.mnar-ln15.r060.l3}{-0.0008}$ & $\resultnum{j.x5.mnar-ln15.r060.l4}{-0.0018}$ & $\resultnum{j.x5.mnar-ln15.r060.l5}{-0.0031}$ & \resultnum{j.x5.mnar-ln15.r060.auc}{0.610} \\
Attrition, odds ratio \resultnum{j.x5.mnar-ln15.amount}{1.5} & \resultnum{j.b.r065.r}{0.65} & \resultnum{j.x5.mnar-ln15.r065.phi}{0.840} & \resultnum{j.x5.mnar-ln15.r065.lambda}{0.039} & $\resultnum{j.x5.mnar-ln15.r065.l3}{-0.0005}$ & $\resultnum{j.x5.mnar-ln15.r065.l4}{-0.0017}$ & $\resultnum{j.x5.mnar-ln15.r065.l5}{-0.0032}$ & \resultnum{j.x5.mnar-ln15.r065.auc}{0.621} \\
Attrition, odds ratio \resultnum{j.x5.mnar-ln15.amount}{1.5} & \resultnum{j.b.r070.r}{0.70} & \resultnum{j.x5.mnar-ln15.r070.phi}{0.843} & \resultnum{j.x5.mnar-ln15.r070.lambda}{0.027} & $\resultnum{j.x5.mnar-ln15.r070.l3}{-0.0006}$ & $\resultnum{j.x5.mnar-ln15.r070.l4}{-0.0015}$ & $\resultnum{j.x5.mnar-ln15.r070.l5}{-0.0024}$ & \resultnum{j.x5.mnar-ln15.r070.auc}{0.588} \\
Attrition, odds ratio \resultnum{j.x5.mnar-ln15.amount}{1.5} & \resultnum{j.b.r075.r}{0.75} & \resultnum{j.x5.mnar-ln15.r075.phi}{0.844} & \resultnum{j.x5.mnar-ln15.r075.lambda}{0.020} & $\resultnum{j.x5.mnar-ln15.r075.l3}{-0.0001}$ & $\resultnum{j.x5.mnar-ln15.r075.l4}{-0.0007}$ & $\resultnum{j.x5.mnar-ln15.r075.l5}{-0.0016}$ & \resultnum{j.x5.mnar-ln15.r075.auc}{0.559} \\
Attrition, odds ratio \resultnum{j.x5.mnar-ln15.amount}{1.5} & \resultnum{j.b.r080.r}{0.80} & \resultnum{j.x5.mnar-ln15.r080.phi}{0.846} & \resultnum{j.x5.mnar-ln15.r080.lambda}{0.014} & $\resultnum{j.x5.mnar-ln15.r080.l3}{-0.0002}$ & $\resultnum{j.x5.mnar-ln15.r080.l4}{-0.0008}$ & $\resultnum{j.x5.mnar-ln15.r080.l5}{-0.0014}$ & \resultnum{j.x5.mnar-ln15.r080.auc}{0.558} \\
Attrition, odds ratio \resultnum{j.x5.mnar-ln15.amount}{1.5} & \resultnum{j.b.r085.r}{0.85} & \resultnum{j.x5.mnar-ln15.r085.phi}{0.847} & \resultnum{j.x5.mnar-ln15.r085.lambda}{0.006} & $\resultnum{j.x5.mnar-ln15.r085.l3}{-0.0005}$ & $\resultnum{j.x5.mnar-ln15.r085.l4}{-0.0009}$ & $\resultnum{j.x5.mnar-ln15.r085.l5}{-0.0013}$ & \resultnum{j.x5.mnar-ln15.r085.auc}{0.545} \\
Attrition, odds ratio \resultnum{j.x5.mnar-ln15.amount}{1.5} & \resultnum{j.b.r090.r}{0.90} & \resultnum{j.x5.mnar-ln15.r090.phi}{0.848} & \resultnum{j.x5.mnar-ln15.r090.lambda}{0.000} & $\resultnum{j.x5.mnar-ln15.r090.l3}{-0.0007}$ & $\resultnum{j.x5.mnar-ln15.r090.l4}{-0.0011}$ & $\resultnum{j.x5.mnar-ln15.r090.l5}{-0.0013}$ & \resultnum{j.x5.mnar-ln15.r090.auc}{0.540} \\
Attrition, odds ratio \resultnum{j.x5.mnar-ln2.amount}{2} & \resultnum{j.b.r050.r}{0.50} & \resultnum{j.x5.mnar-ln2.r050.phi}{0.812} & \resultnum{j.x5.mnar-ln2.r050.lambda}{0.142} & $\resultnum{j.x5.mnar-ln2.r050.l3}{-0.0022}$ & $\resultnum{j.x5.mnar-ln2.r050.l4}{-0.0056}$ & $\resultnum{j.x5.mnar-ln2.r050.l5}{-0.0099}$ & \resultnum{j.x5.mnar-ln2.r050.auc}{0.820} \\
Attrition, odds ratio \resultnum{j.x5.mnar-ln2.amount}{2} & \resultnum{j.b.r055.r}{0.55} & \resultnum{j.x5.mnar-ln2.r055.phi}{0.818} & \resultnum{j.x5.mnar-ln2.r055.lambda}{0.119} & $\resultnum{j.x5.mnar-ln2.r055.l3}{-0.0021}$ & $\resultnum{j.x5.mnar-ln2.r055.l4}{-0.0051}$ & $\resultnum{j.x5.mnar-ln2.r055.l5}{-0.0090}$ & \resultnum{j.x5.mnar-ln2.r055.auc}{0.799} \\
Attrition, odds ratio \resultnum{j.x5.mnar-ln2.amount}{2} & \resultnum{j.b.r060.r}{0.60} & \resultnum{j.x5.mnar-ln2.r060.phi}{0.820} & \resultnum{j.x5.mnar-ln2.r060.lambda}{0.112} & $\resultnum{j.x5.mnar-ln2.r060.l3}{-0.0021}$ & $\resultnum{j.x5.mnar-ln2.r060.l4}{-0.0053}$ & $\resultnum{j.x5.mnar-ln2.r060.l5}{-0.0090}$ & \resultnum{j.x5.mnar-ln2.r060.auc}{0.797} \\
Attrition, odds ratio \resultnum{j.x5.mnar-ln2.amount}{2} & \resultnum{j.b.r065.r}{0.65} & \resultnum{j.x5.mnar-ln2.r065.phi}{0.825} & \resultnum{j.x5.mnar-ln2.r065.lambda}{0.094} & $\resultnum{j.x5.mnar-ln2.r065.l3}{-0.0015}$ & $\resultnum{j.x5.mnar-ln2.r065.l4}{-0.0046}$ & $\resultnum{j.x5.mnar-ln2.r065.l5}{-0.0082}$ & \resultnum{j.x5.mnar-ln2.r065.auc}{0.778} \\
Attrition, odds ratio \resultnum{j.x5.mnar-ln2.amount}{2} & \resultnum{j.b.r070.r}{0.70} & \resultnum{j.x5.mnar-ln2.r070.phi}{0.830} & \resultnum{j.x5.mnar-ln2.r070.lambda}{0.074} & $\resultnum{j.x5.mnar-ln2.r070.l3}{-0.0016}$ & $\resultnum{j.x5.mnar-ln2.r070.l4}{-0.0041}$ & $\resultnum{j.x5.mnar-ln2.r070.l5}{-0.0070}$ & \resultnum{j.x5.mnar-ln2.r070.auc}{0.741} \\
Attrition, odds ratio \resultnum{j.x5.mnar-ln2.amount}{2} & \resultnum{j.b.r075.r}{0.75} & \resultnum{j.x5.mnar-ln2.r075.phi}{0.833} & \resultnum{j.x5.mnar-ln2.r075.lambda}{0.057} & $\resultnum{j.x5.mnar-ln2.r075.l3}{-0.0010}$ & $\resultnum{j.x5.mnar-ln2.r075.l4}{-0.0030}$ & $\resultnum{j.x5.mnar-ln2.r075.l5}{-0.0055}$ & \resultnum{j.x5.mnar-ln2.r075.auc}{0.693} \\
Attrition, odds ratio \resultnum{j.x5.mnar-ln2.amount}{2} & \resultnum{j.b.r080.r}{0.80} & \resultnum{j.x5.mnar-ln2.r080.phi}{0.838} & \resultnum{j.x5.mnar-ln2.r080.lambda}{0.035} & $\resultnum{j.x5.mnar-ln2.r080.l3}{-0.0007}$ & $\resultnum{j.x5.mnar-ln2.r080.l4}{-0.0023}$ & $\resultnum{j.x5.mnar-ln2.r080.l5}{-0.0040}$ & \resultnum{j.x5.mnar-ln2.r080.auc}{0.649} \\
Attrition, odds ratio \resultnum{j.x5.mnar-ln2.amount}{2} & \resultnum{j.b.r085.r}{0.85} & \resultnum{j.x5.mnar-ln2.r085.phi}{0.843} & \resultnum{j.x5.mnar-ln2.r085.lambda}{0.013} & $\resultnum{j.x5.mnar-ln2.r085.l3}{-0.0010}$ & $\resultnum{j.x5.mnar-ln2.r085.l4}{-0.0020}$ & $\resultnum{j.x5.mnar-ln2.r085.l5}{-0.0030}$ & \resultnum{j.x5.mnar-ln2.r085.auc}{0.602} \\
Attrition, odds ratio \resultnum{j.x5.mnar-ln2.amount}{2} & \resultnum{j.b.r090.r}{0.90} & \resultnum{j.x5.mnar-ln2.r090.phi}{0.845} & \resultnum{j.x5.mnar-ln2.r090.lambda}{0.000} & $\resultnum{j.x5.mnar-ln2.r090.l3}{-0.0011}$ & $\resultnum{j.x5.mnar-ln2.r090.l4}{-0.0019}$ & $\resultnum{j.x5.mnar-ln2.r090.l5}{-0.0025}$ & \resultnum{j.x5.mnar-ln2.r090.auc}{0.578} \\
\botrule
\end{tabular}
\par\smallskip\parbox{\linewidth}{Lag differences are the influence-free arm minus fitted B, averaged over check panels. AUCs are from the recorded stress classifier. Only estimated cells are shown; the wave-step and conditioning fits were unavailable.}
\end{table}


The boundary is therefore conditional on a correctly specified measurement channel (for a realistic channel shared by both arms, see Section~\ref{sec:shp-reading}). The attrition result is partly a property of the summary statistic: pooled correlations over the available pairs make no adjustment for selective exit, whereas a likelihood-based estimate of the lag structure under missingness at random would adjust for the part of the selection that runs through the observed answers. How attrition in the SHP depends on trust was not estimated. The comparison of each condition with the released curve against B is in Table~\ref{tab:v4x-stress-comparisons} (Section~\ref{sec:x-readings}). % EX15 EX12

\subsection{Band coverage check (X6)}\label{sec:x6}

\paragraph{Design and denominators.}
X6 asks whether the frozen intervals and the per-share band cover the frozen estimand at positive strength, which no gate tests. It ran the whole unit at $r=0.70$ (X6a: 200 independent repetitions of all nine strengths, with the frozen pipeline and common random numbers across strengths within a repetition) and the two other cells of the dispersion check, $(a,r)=(0.03,0.50)$ and $(0.03,0.90)$ (X6b: 200 independent pipelines each, pointwise readings only). At the operational cut, X6a had recorded \resultnum{v4x.x6.x6a.recorded_prose}{1,008} of \resultnum{v4x.x6.x6a.units_prose}{1,008} jobs, with \resultnum{v4x.x6.x6a.incomplete}{0} incomplete, and X6b \resultnum{v4x.x6.x6b.recorded}{602} of \resultnum{v4x.x6.x6b.units}{602}, with \resultnum{v4x.x6.x6b.incomplete}{0} incomplete (Table~\ref{tab:v4x-x6-inventory}). There were \resultnum{v4x.x6.x6a.without_curve}{0} repetitions without a curve and \resultnum{v4x.x6.x6a.repetitions.errors}{0} recorded repetition errors; these counts are reported outside the coverage denominators. Table~\ref{tab:v4x-x6-denominators} states the conditioning of each reference: complete gap-free units for the X6a primary readings and successful cell pipelines for X6b. In \resultnum{v4x.x6.x6a.strengths.r070.a0000.gaps}{34} of the 200 X6a repetitions the band had a gap somewhere on the grid, mostly at $a=0.60$ and 0.90, so the primary denominator is the remaining complete gap-free units.

\paragraph{Truths.}
The primary truth is the frozen estimand read literally (decision item D1): the expected AUC over the declared matching and training draws, conditional on successful fits, estimated as the mean of 100 independent frozen pipelines per cell, which include matching with \resultnum{v4x.x3.sizes.matching_prose}{1,024} pairs and the declared \resultnum{v4x.x3.sizes.retry_matching_prose}{4,096}-pair retries. The secondary truth uses approximate population matching with \resultnum{v4x.x6.sizes.population_pairs_prose}{16,384} pairs; its Monte Carlo error omits the error of the population fit. Table~\ref{tab:v4x-x6-truth} reports retry fractions and truth Monte Carlo errors against the declared precision target of \resultnum{v4x.x6.truth_se_target}{0.002}. That $r=0.70$ was examined during calibration remains relevant to adaptation, even though X6 reads no boundary.

\paragraph{Both readings of joint coverage.}
Joint coverage requires the band to cover the truth at all \resultnum{v4x.protocol.grid_strengths}{9} strengths, including $a=\resultnum{v4x.protocol.zero}{0}$. Against the procedural truth, the band covered jointly in \resultnum{v4x.x6.x6a.band.vs_d1.primary.covered}{162} of \resultnum{v4x.x6.x6a.band.vs_d1.primary.total}{166} complete gap-free bands, with Wilson 95\% interval \resultnum{v4x.x6.x6a.band.vs_d1.primary.wilson95.prose_lower}{0.940}--\resultnum{v4x.x6.x6a.band.vs_d1.primary.wilson95.prose_upper}{0.991}: under the declared criterion, which requires the Wilson upper end to reach \resultnum{v4x.protocol.x6.nominal}{0.95}, we did not detect undercoverage (protocol label \texttt{adequate}). The criterion does not demonstrate nominal coverage, and the critical-value floor alone does not remove bias, variance error or the conditioning on success. With incomplete bands counted as misses, jointly and at every strength, the band covered in \resultnum{v4x.x6.x6a.band.vs_d1.gaps_as_misses.covered}{162} of \resultnum{v4x.x6.x6a.band.vs_d1.gaps_as_misses.total}{200} repetitions with a curve, with Wilson interval \resultnum{j.x6.gaps_as_misses.wilson.lower}{0.750}--\resultnum{j.x6.gaps_as_misses.wilson.upper}{0.858}, which the declared rule classifies as \texttt{failure} (upper end below \resultnum{v4x.protocol.x6.failure}{0.9}). Both readings are prespecified; neither is unconditional coverage over all planned repetitions. They condition on different events. The primary reading keeps only the bands that are complete and gap-free over the whole grid, so it selects repetitions by fit success at every strength, whereas the procedural truth conditions on success cell by cell. The reading with gaps as misses keeps every recorded repetition, none of which ended in an error or without a curve, so it is the rate at which the procedure delivered a complete band that covered, not the coverage of the bands it delivered. Among the bands it delivered, no undercoverage was detected in the settings evaluated; this does not show nominal coverage, and neither reading speaks to the shares that were not checked.
Across all coverage references against the procedural truth, the reading was ``no undercoverage detected'' in \resultnum{v4x.x6.x6a.count.adequate}{19} X6a and \resultnum{v4x.x6.x6b.count.adequate}{4} X6b references and ``failure'' in \resultnum{v4x.x6.x6a.count.failure}{10} X6a and \resultnum{v4x.x6.x6b.count.failure}{0} X6b references (Table~\ref{tab:v4x-x6}). The failure count includes the separately tabulated gaps-as-misses rows; it counts references, not distinct cells.

\paragraph{The two cells away from $r=0.70$.}
X6b checks the two cells of the dispersion check outside $r=0.70$. Both are at $a=0.03$, the strength at which the region label of a first boundary is decided, and at $r=0.90$ the upper band at $a=0.03$ came closest to 0.70 (Section~\ref{sec:additional-b}). Against the procedural truth, the pointwise interval covered in \resultnum{v4x.x6.x6b.cells.r050.a0300.vs_d1.pointwise.covered}{193} of \resultnum{v4x.x6.x6b.cells.r050.a0300.vs_d1.pointwise.total}{200} pipelines at $(0.03,0.50)$ and in \resultnum{v4x.x6.x6b.cells.r090.a0300.vs_d1.pointwise.covered}{189} of \resultnum{v4x.x6.x6b.cells.r090.a0300.vs_d1.pointwise.total}{200} at $(0.03,0.90)$, both read as no undercoverage detected; the procedural truths of both cells lie well below 0.70 (Table~\ref{tab:v4x-x6-truth}). Together with the whole unit at $r=0.70$, these two cells are the evidence on coverage at positive strength; no other share was checked.

\paragraph{Precision sensitivity.}
The declaration repeats each coverage reading at the procedural truth plus and minus two Monte Carlo errors, and marks a reading that changes there as inconclusive by truth precision, only where a truth's Monte Carlo error exceeds the target of \resultnum{v4x.x6.truth_se_target}{0.002}. Every truth met the target (Table~\ref{tab:v4x-x6-truth}), so the sensitivity was not computed, and the released fields for it are empty (Table~\ref{tab:v4x-x6-precision}). The robustness of the coverage readings to the estimation error of the truth is therefore not assessed beyond that precision target. No conservative flag (Wilson lower end above 0.99) was raised.

\paragraph{Matching diagnostics.}
Table~\ref{tab:v4x-x6-diagnostics} reports the standardized errors and their uncovered-bias flags, the mean shift of the local nuisance draws, the empirical matching bias and the difference between the procedural and population truths, together with fit residuals, amplification and the ratio of mean squared error to mean squared standard error. These diagnostics keep their own success conditioning and reference-precision limits.

% Generated by v4_exploratory_export; edited during the journal revision
% (agent S, 2026-10-13): captions, headers and categorical labels translated
% into plain English. Every number keeps its \resultnum binding.
\begin{table*}[htbp]
\setlength{\tabcolsep}{2pt}
\caption{Band coverage check (X6): job inventory at the declared cut, including the repetitions and the independent truth and population-match jobs. Incomplete jobs would be reported separately from errors and fit gaps; none occurred.}\label{tab:v4x-x6-inventory}
\centering\footnotesize
\begin{tabular}{@{}p{0.460\textwidth}p{0.150\textwidth}p{0.150\textwidth}p{0.150\textwidth}@{}}
\toprule
\raggedright \mbox{Part} & \raggedright \mbox{Declared} \mbox{jobs} & \raggedright \mbox{Recorded} \mbox{jobs} & \raggedright \mbox{Incomplete} \mbox{jobs} \tabularnewline
\midrule
\raggedright Whole unit at $r=0.70$, 200 repetitions, with truths (X6a) & \raggedright \resultnum{v4x.x6.x6a.units}{1008} & \raggedright \resultnum{v4x.x6.x6a.recorded}{1008} & \raggedright \resultnum{v4x.x6.x6a.incomplete}{0} \tabularnewline
\raggedright Cells $(0.03, 0.50)$ and $(0.03, 0.90)$, with truths (X6b) & \raggedright \resultnum{v4x.x6.x6b.units}{602} & \raggedright \resultnum{v4x.x6.x6b.recorded}{602} & \raggedright \resultnum{v4x.x6.x6b.incomplete}{0} \tabularnewline
\botrule
\end{tabular}
\end{table*}

% Generated by v4_exploratory_export; edited during the journal revision
% (agent S, 2026-10-13): captions, headers and categorical labels translated
% into plain English. Every number keeps its \resultnum binding.
\begin{table*}[htbp]
\setlength{\tabcolsep}{2pt}
\caption{Band coverage check (X6): coverage of the success-conditioned procedural truth (the frozen estimand; machine label D1) or of the approximate population-match truth (population; pop), with Wilson 95\% intervals. Interval: the pointwise interval (PW), the band at that strength, or the joint band, which must cover the truth at all nine strengths, $a=0$ included. The primary denominator of X6a is the complete gap-free bands (\resultnum{v4x.x6.x6a.band.vs_d1.primary.total}{166} of the \resultnum{v4x.x6.x6a.band.vs_d1.gaps_as_misses.total}{200} repetitions at $r=0.70$); in the rows with gaps as misses (joint+gap, PW+gap) every repetition with a curve counts, and an incomplete band is a miss jointly and at every strength. For the two X6b cells the denominator is the successful cell pipelines. Reading under the declared rule: no undercoverage detected, Wilson upper end at least 0.95 (machine label adequate; this does not demonstrate nominal coverage); failure, upper end below 0.90 (failure). Both readings of the joint band against the procedural truth are given: no undercoverage detected among complete gap-free bands, failure with gaps as misses.}\label{tab:v4x-x6}
\centering\footnotesize
\begin{tabular}{@{}p{0.108\textwidth}p{0.118\textwidth}p{0.128\textwidth}p{0.069\textwidth}p{0.098\textwidth}p{0.197\textwidth}p{0.207\textwidth}@{}}
\toprule
\raggedright $a,r$ & \raggedright \mbox{Interval} & \raggedright \mbox{Truth} & \raggedright \mbox{Gaps} \mbox{as} \mbox{misses} & \raggedright \mbox{Covered} \mbox{/} \mbox{total} & \raggedright \mbox{Wilson} \mbox{interval} & \raggedright \mbox{Reading} \tabularnewline
\midrule
\raggedright all nine strengths & \raggedright joint band & \raggedright procedural & \raggedright no & \raggedright \resultnum{v4x.x6.x6a.band.vs_d1.primary.covered}{162}/\resultnum{v4x.x6.x6a.band.vs_d1.primary.total}{166} & \raggedright [\resultnum{v4x.x6.x6a.band.vs_d1.primary.wilson95.i000}{0.939689}, \resultnum{v4x.x6.x6a.band.vs_d1.primary.wilson95.i001}{0.99059}] & \raggedright no undercoverage detected \tabularnewline
\raggedright all nine strengths & \raggedright joint band & \raggedright procedural & \raggedright yes & \raggedright \resultnum{v4x.x6.x6a.band.vs_d1.gaps_as_misses.covered}{162}/\resultnum{v4x.x6.x6a.band.vs_d1.gaps_as_misses.total}{200} & \raggedright [\resultnum{v4x.x6.x6a.band.vs_d1.gaps_as_misses.wilson95.i000}{0.749988}, \resultnum{v4x.x6.x6a.band.vs_d1.gaps_as_misses.wilson95.i001}{0.858328}] & \raggedright failure \tabularnewline
\raggedright all nine strengths & \raggedright joint band & \raggedright population & \raggedright no & \raggedright \resultnum{v4x.x6.x6a.band.vs_pop.primary.covered}{160}/\resultnum{v4x.x6.x6a.band.vs_pop.primary.total}{166} & \raggedright [\resultnum{v4x.x6.x6a.band.vs_pop.primary.wilson95.i000}{0.923397}, \resultnum{v4x.x6.x6a.band.vs_pop.primary.wilson95.i001}{0.983331}] & \raggedright no undercoverage detected \tabularnewline
\raggedright all nine strengths & \raggedright joint band & \raggedright population & \raggedright yes & \raggedright \resultnum{v4x.x6.x6a.band.vs_pop.gaps_as_misses.covered}{160}/\resultnum{v4x.x6.x6a.band.vs_pop.gaps_as_misses.total}{200} & \raggedright [\resultnum{v4x.x6.x6a.band.vs_pop.gaps_as_misses.wilson95.i000}{0.739145}, \resultnum{v4x.x6.x6a.band.vs_pop.gaps_as_misses.wilson95.i001}{0.849548}] & \raggedright failure \tabularnewline
\raggedright \resultnum{v4x.x6.x6a.strengths.r070.a0000.a}{0}, \resultnum{v4x.x6.x6a.strengths.r070.a0000.r}{0.7} & \raggedright pointwise & \raggedright procedural & \raggedright no & \raggedright \resultnum{v4x.x6.x6a.strengths.r070.a0000.vs_d1.pointwise.covered}{157}/\resultnum{v4x.x6.x6a.strengths.r070.a0000.vs_d1.pointwise.total}{166} & \raggedright [\resultnum{v4x.x6.x6a.strengths.r070.a0000.vs_d1.pointwise.wilson95.i000}{0.900184}, \resultnum{v4x.x6.x6a.strengths.r070.a0000.vs_d1.pointwise.wilson95.i001}{0.971217}] & \raggedright no undercoverage detected \tabularnewline
\raggedright \resultnum{v4x.x6.x6a.strengths.r070.a0000.a}{0}, \resultnum{v4x.x6.x6a.strengths.r070.a0000.r}{0.7} & \raggedright pointwise & \raggedright procedural & \raggedright yes & \raggedright \resultnum{v4x.x6.x6a.strengths.r070.a0000.vs_d1.pointwise_gaps_as_misses.covered}{157}/\resultnum{v4x.x6.x6a.strengths.r070.a0000.vs_d1.pointwise_gaps_as_misses.total}{200} & \raggedright [\resultnum{v4x.x6.x6a.strengths.r070.a0000.vs_d1.pointwise_gaps_as_misses.wilson95.i000}{0.722977}, \resultnum{v4x.x6.x6a.strengths.r070.a0000.vs_d1.pointwise_gaps_as_misses.wilson95.i001}{0.836281}] & \raggedright failure \tabularnewline
\raggedright \resultnum{v4x.x6.x6a.strengths.r070.a0000.a}{0}, \resultnum{v4x.x6.x6a.strengths.r070.a0000.r}{0.7} & \raggedright band & \raggedright procedural & \raggedright no & \raggedright \resultnum{v4x.x6.x6a.strengths.r070.a0000.vs_d1.band_interval.covered}{165}/\resultnum{v4x.x6.x6a.strengths.r070.a0000.vs_d1.band_interval.total}{166} & \raggedright [\resultnum{v4x.x6.x6a.strengths.r070.a0000.vs_d1.band_interval.wilson95.i000}{0.966671}, \resultnum{v4x.x6.x6a.strengths.r070.a0000.vs_d1.band_interval.wilson95.i001}{0.998936}] & \raggedright no undercoverage detected \tabularnewline
\raggedright \resultnum{v4x.x6.x6a.strengths.r070.a0000.a}{0}, \resultnum{v4x.x6.x6a.strengths.r070.a0000.r}{0.7} & \raggedright pointwise & \raggedright population & \raggedright no & \raggedright \resultnum{v4x.x6.x6a.strengths.r070.a0000.vs_pop.pointwise.covered}{157}/\resultnum{v4x.x6.x6a.strengths.r070.a0000.vs_pop.pointwise.total}{166} & \raggedright [\resultnum{v4x.x6.x6a.strengths.r070.a0000.vs_pop.pointwise.wilson95.i000}{0.900184}, \resultnum{v4x.x6.x6a.strengths.r070.a0000.vs_pop.pointwise.wilson95.i001}{0.971217}] & \raggedright no undercoverage detected \tabularnewline
\raggedright \resultnum{v4x.x6.x6a.strengths.r070.a0000.a}{0}, \resultnum{v4x.x6.x6a.strengths.r070.a0000.r}{0.7} & \raggedright pointwise & \raggedright population & \raggedright yes & \raggedright \resultnum{v4x.x6.x6a.strengths.r070.a0000.vs_pop.pointwise_gaps_as_misses.covered}{157}/\resultnum{v4x.x6.x6a.strengths.r070.a0000.vs_pop.pointwise_gaps_as_misses.total}{200} & \raggedright [\resultnum{v4x.x6.x6a.strengths.r070.a0000.vs_pop.pointwise_gaps_as_misses.wilson95.i000}{0.722977}, \resultnum{v4x.x6.x6a.strengths.r070.a0000.vs_pop.pointwise_gaps_as_misses.wilson95.i001}{0.836281}] & \raggedright failure \tabularnewline
\raggedright \resultnum{v4x.x6.x6a.strengths.r070.a0000.a}{0}, \resultnum{v4x.x6.x6a.strengths.r070.a0000.r}{0.7} & \raggedright band & \raggedright population & \raggedright no & \raggedright \resultnum{v4x.x6.x6a.strengths.r070.a0000.vs_pop.band_interval.covered}{165}/\resultnum{v4x.x6.x6a.strengths.r070.a0000.vs_pop.band_interval.total}{166} & \raggedright [\resultnum{v4x.x6.x6a.strengths.r070.a0000.vs_pop.band_interval.wilson95.i000}{0.966671}, \resultnum{v4x.x6.x6a.strengths.r070.a0000.vs_pop.band_interval.wilson95.i001}{0.998936}] & \raggedright no undercoverage detected \tabularnewline
\raggedright \resultnum{v4x.x6.x6a.strengths.r070.a0010.a}{0.001}, \resultnum{v4x.x6.x6a.strengths.r070.a0010.r}{0.7} & \raggedright pointwise & \raggedright procedural & \raggedright no & \raggedright \resultnum{v4x.x6.x6a.strengths.r070.a0010.vs_d1.pointwise.covered}{158}/\resultnum{v4x.x6.x6a.strengths.r070.a0010.vs_d1.pointwise.total}{166} & \raggedright [\resultnum{v4x.x6.x6a.strengths.r070.a0010.vs_d1.pointwise.wilson95.i000}{0.907796}, \resultnum{v4x.x6.x6a.strengths.r070.a0010.vs_d1.pointwise.wilson95.i001}{0.975381}] & \raggedright no undercoverage detected \tabularnewline
\raggedright \resultnum{v4x.x6.x6a.strengths.r070.a0010.a}{0.001}, \resultnum{v4x.x6.x6a.strengths.r070.a0010.r}{0.7} & \raggedright pointwise & \raggedright procedural & \raggedright yes & \raggedright \resultnum{v4x.x6.x6a.strengths.r070.a0010.vs_d1.pointwise_gaps_as_misses.covered}{158}/\resultnum{v4x.x6.x6a.strengths.r070.a0010.vs_d1.pointwise_gaps_as_misses.total}{200} & \raggedright [\resultnum{v4x.x6.x6a.strengths.r070.a0010.vs_d1.pointwise_gaps_as_misses.wilson95.i000}{0.728354}, \resultnum{v4x.x6.x6a.strengths.r070.a0010.vs_d1.pointwise_gaps_as_misses.wilson95.i001}{0.840716}] & \raggedright failure \tabularnewline
\raggedright \resultnum{v4x.x6.x6a.strengths.r070.a0010.a}{0.001}, \resultnum{v4x.x6.x6a.strengths.r070.a0010.r}{0.7} & \raggedright band & \raggedright procedural & \raggedright no & \raggedright \resultnum{v4x.x6.x6a.strengths.r070.a0010.vs_d1.band_interval.covered}{166}/\resultnum{v4x.x6.x6a.strengths.r070.a0010.vs_d1.band_interval.total}{166} & \raggedright [\resultnum{v4x.x6.x6a.strengths.r070.a0010.vs_d1.band_interval.wilson95.i000}{0.977382}, \resultnum{v4x.x6.x6a.strengths.r070.a0010.vs_d1.band_interval.wilson95.i001}{1}] & \raggedright no undercoverage detected \tabularnewline
\raggedright \resultnum{v4x.x6.x6a.strengths.r070.a0010.a}{0.001}, \resultnum{v4x.x6.x6a.strengths.r070.a0010.r}{0.7} & \raggedright pointwise & \raggedright population & \raggedright no & \raggedright \resultnum{v4x.x6.x6a.strengths.r070.a0010.vs_pop.pointwise.covered}{157}/\resultnum{v4x.x6.x6a.strengths.r070.a0010.vs_pop.pointwise.total}{166} & \raggedright [\resultnum{v4x.x6.x6a.strengths.r070.a0010.vs_pop.pointwise.wilson95.i000}{0.900184}, \resultnum{v4x.x6.x6a.strengths.r070.a0010.vs_pop.pointwise.wilson95.i001}{0.971217}] & \raggedright no undercoverage detected \tabularnewline
\raggedright \resultnum{v4x.x6.x6a.strengths.r070.a0010.a}{0.001}, \resultnum{v4x.x6.x6a.strengths.r070.a0010.r}{0.7} & \raggedright pointwise & \raggedright population & \raggedright yes & \raggedright \resultnum{v4x.x6.x6a.strengths.r070.a0010.vs_pop.pointwise_gaps_as_misses.covered}{157}/\resultnum{v4x.x6.x6a.strengths.r070.a0010.vs_pop.pointwise_gaps_as_misses.total}{200} & \raggedright [\resultnum{v4x.x6.x6a.strengths.r070.a0010.vs_pop.pointwise_gaps_as_misses.wilson95.i000}{0.722977}, \resultnum{v4x.x6.x6a.strengths.r070.a0010.vs_pop.pointwise_gaps_as_misses.wilson95.i001}{0.836281}] & \raggedright failure \tabularnewline
\raggedright \resultnum{v4x.x6.x6a.strengths.r070.a0010.a}{0.001}, \resultnum{v4x.x6.x6a.strengths.r070.a0010.r}{0.7} & \raggedright band & \raggedright population & \raggedright no & \raggedright \resultnum{v4x.x6.x6a.strengths.r070.a0010.vs_pop.band_interval.covered}{166}/\resultnum{v4x.x6.x6a.strengths.r070.a0010.vs_pop.band_interval.total}{166} & \raggedright [\resultnum{v4x.x6.x6a.strengths.r070.a0010.vs_pop.band_interval.wilson95.i000}{0.977382}, \resultnum{v4x.x6.x6a.strengths.r070.a0010.vs_pop.band_interval.wilson95.i001}{1}] & \raggedright no undercoverage detected \tabularnewline
\raggedright \resultnum{v4x.x6.x6a.strengths.r070.a0030.a}{0.003}, \resultnum{v4x.x6.x6a.strengths.r070.a0030.r}{0.7} & \raggedright pointwise & \raggedright procedural & \raggedright no & \raggedright \resultnum{v4x.x6.x6a.strengths.r070.a0030.vs_d1.pointwise.covered}{160}/\resultnum{v4x.x6.x6a.strengths.r070.a0030.vs_d1.pointwise.total}{166} & \raggedright [\resultnum{v4x.x6.x6a.strengths.r070.a0030.vs_d1.pointwise.wilson95.i000}{0.923397}, \resultnum{v4x.x6.x6a.strengths.r070.a0030.vs_d1.pointwise.wilson95.i001}{0.983331}] & \raggedright no undercoverage detected \tabularnewline
\raggedright \resultnum{v4x.x6.x6a.strengths.r070.a0030.a}{0.003}, \resultnum{v4x.x6.x6a.strengths.r070.a0030.r}{0.7} & \raggedright pointwise & \raggedright procedural & \raggedright yes & \raggedright \resultnum{v4x.x6.x6a.strengths.r070.a0030.vs_d1.pointwise_gaps_as_misses.covered}{160}/\resultnum{v4x.x6.x6a.strengths.r070.a0030.vs_d1.pointwise_gaps_as_misses.total}{200} & \raggedright [\resultnum{v4x.x6.x6a.strengths.r070.a0030.vs_d1.pointwise_gaps_as_misses.wilson95.i000}{0.739145}, \resultnum{v4x.x6.x6a.strengths.r070.a0030.vs_d1.pointwise_gaps_as_misses.wilson95.i001}{0.849548}] & \raggedright failure \tabularnewline
\raggedright \resultnum{v4x.x6.x6a.strengths.r070.a0030.a}{0.003}, \resultnum{v4x.x6.x6a.strengths.r070.a0030.r}{0.7} & \raggedright band & \raggedright procedural & \raggedright no & \raggedright \resultnum{v4x.x6.x6a.strengths.r070.a0030.vs_d1.band_interval.covered}{165}/\resultnum{v4x.x6.x6a.strengths.r070.a0030.vs_d1.band_interval.total}{166} & \raggedright [\resultnum{v4x.x6.x6a.strengths.r070.a0030.vs_d1.band_interval.wilson95.i000}{0.966671}, \resultnum{v4x.x6.x6a.strengths.r070.a0030.vs_d1.band_interval.wilson95.i001}{0.998936}] & \raggedright no undercoverage detected \tabularnewline
\raggedright \resultnum{v4x.x6.x6a.strengths.r070.a0030.a}{0.003}, \resultnum{v4x.x6.x6a.strengths.r070.a0030.r}{0.7} & \raggedright pointwise & \raggedright population & \raggedright no & \raggedright \resultnum{v4x.x6.x6a.strengths.r070.a0030.vs_pop.pointwise.covered}{161}/\resultnum{v4x.x6.x6a.strengths.r070.a0030.vs_pop.pointwise.total}{166} & \raggedright [\resultnum{v4x.x6.x6a.strengths.r070.a0030.vs_pop.pointwise.wilson95.i000}{0.931437}, \resultnum{v4x.x6.x6a.strengths.r070.a0030.vs_pop.pointwise.wilson95.i001}{0.987067}] & \raggedright no undercoverage detected \tabularnewline
\raggedright \resultnum{v4x.x6.x6a.strengths.r070.a0030.a}{0.003}, \resultnum{v4x.x6.x6a.strengths.r070.a0030.r}{0.7} & \raggedright pointwise & \raggedright population & \raggedright yes & \raggedright \resultnum{v4x.x6.x6a.strengths.r070.a0030.vs_pop.pointwise_gaps_as_misses.covered}{161}/\resultnum{v4x.x6.x6a.strengths.r070.a0030.vs_pop.pointwise_gaps_as_misses.total}{200} & \raggedright [\resultnum{v4x.x6.x6a.strengths.r070.a0030.vs_pop.pointwise_gaps_as_misses.wilson95.i000}{0.74456}, \resultnum{v4x.x6.x6a.strengths.r070.a0030.vs_pop.pointwise_gaps_as_misses.wilson95.i001}{0.853945}] & \raggedright failure \tabularnewline
\raggedright \resultnum{v4x.x6.x6a.strengths.r070.a0030.a}{0.003}, \resultnum{v4x.x6.x6a.strengths.r070.a0030.r}{0.7} & \raggedright band & \raggedright population & \raggedright no & \raggedright \resultnum{v4x.x6.x6a.strengths.r070.a0030.vs_pop.band_interval.covered}{165}/\resultnum{v4x.x6.x6a.strengths.r070.a0030.vs_pop.band_interval.total}{166} & \raggedright [\resultnum{v4x.x6.x6a.strengths.r070.a0030.vs_pop.band_interval.wilson95.i000}{0.966671}, \resultnum{v4x.x6.x6a.strengths.r070.a0030.vs_pop.band_interval.wilson95.i001}{0.998936}] & \raggedright no undercoverage detected \tabularnewline
\raggedright \resultnum{v4x.x6.x6a.strengths.r070.a0100.a}{0.01}, \resultnum{v4x.x6.x6a.strengths.r070.a0100.r}{0.7} & \raggedright pointwise & \raggedright procedural & \raggedright no & \raggedright \resultnum{v4x.x6.x6a.strengths.r070.a0100.vs_d1.pointwise.covered}{162}/\resultnum{v4x.x6.x6a.strengths.r070.a0100.vs_d1.pointwise.total}{166} & \raggedright [\resultnum{v4x.x6.x6a.strengths.r070.a0100.vs_d1.pointwise.wilson95.i000}{0.939689}, \resultnum{v4x.x6.x6a.strengths.r070.a0100.vs_d1.pointwise.wilson95.i001}{0.99059}] & \raggedright no undercoverage detected \tabularnewline
\raggedright \resultnum{v4x.x6.x6a.strengths.r070.a0100.a}{0.01}, \resultnum{v4x.x6.x6a.strengths.r070.a0100.r}{0.7} & \raggedright pointwise & \raggedright procedural & \raggedright yes & \raggedright \resultnum{v4x.x6.x6a.strengths.r070.a0100.vs_d1.pointwise_gaps_as_misses.covered}{162}/\resultnum{v4x.x6.x6a.strengths.r070.a0100.vs_d1.pointwise_gaps_as_misses.total}{200} & \raggedright [\resultnum{v4x.x6.x6a.strengths.r070.a0100.vs_d1.pointwise_gaps_as_misses.wilson95.i000}{0.749988}, \resultnum{v4x.x6.x6a.strengths.r070.a0100.vs_d1.pointwise_gaps_as_misses.wilson95.i001}{0.858328}] & \raggedright failure \tabularnewline
\raggedright \resultnum{v4x.x6.x6a.strengths.r070.a0100.a}{0.01}, \resultnum{v4x.x6.x6a.strengths.r070.a0100.r}{0.7} & \raggedright band & \raggedright procedural & \raggedright no & \raggedright \resultnum{v4x.x6.x6a.strengths.r070.a0100.vs_d1.band_interval.covered}{165}/\resultnum{v4x.x6.x6a.strengths.r070.a0100.vs_d1.band_interval.total}{166} & \raggedright [\resultnum{v4x.x6.x6a.strengths.r070.a0100.vs_d1.band_interval.wilson95.i000}{0.966671}, \resultnum{v4x.x6.x6a.strengths.r070.a0100.vs_d1.band_interval.wilson95.i001}{0.998936}] & \raggedright no undercoverage detected \tabularnewline
\raggedright \resultnum{v4x.x6.x6a.strengths.r070.a0100.a}{0.01}, \resultnum{v4x.x6.x6a.strengths.r070.a0100.r}{0.7} & \raggedright pointwise & \raggedright population & \raggedright no & \raggedright \resultnum{v4x.x6.x6a.strengths.r070.a0100.vs_pop.pointwise.covered}{154}/\resultnum{v4x.x6.x6a.strengths.r070.a0100.vs_pop.pointwise.total}{166} & \raggedright [\resultnum{v4x.x6.x6a.strengths.r070.a0100.vs_pop.pointwise.wilson95.i000}{0.877907}, \resultnum{v4x.x6.x6a.strengths.r070.a0100.vs_pop.pointwise.wilson95.i001}{0.958167}] & \raggedright no undercoverage detected \tabularnewline
\raggedright \resultnum{v4x.x6.x6a.strengths.r070.a0100.a}{0.01}, \resultnum{v4x.x6.x6a.strengths.r070.a0100.r}{0.7} & \raggedright pointwise & \raggedright population & \raggedright yes & \raggedright \resultnum{v4x.x6.x6a.strengths.r070.a0100.vs_pop.pointwise_gaps_as_misses.covered}{154}/\resultnum{v4x.x6.x6a.strengths.r070.a0100.vs_pop.pointwise_gaps_as_misses.total}{200} & \raggedright [\resultnum{v4x.x6.x6a.strengths.r070.a0100.vs_pop.pointwise_gaps_as_misses.wilson95.i000}{0.706917}, \resultnum{v4x.x6.x6a.strengths.r070.a0100.vs_pop.pointwise_gaps_as_misses.wilson95.i001}{0.822907}] & \raggedright failure \tabularnewline
\raggedright \resultnum{v4x.x6.x6a.strengths.r070.a0100.a}{0.01}, \resultnum{v4x.x6.x6a.strengths.r070.a0100.r}{0.7} & \raggedright band & \raggedright population & \raggedright no & \raggedright \resultnum{v4x.x6.x6a.strengths.r070.a0100.vs_pop.band_interval.covered}{164}/\resultnum{v4x.x6.x6a.strengths.r070.a0100.vs_pop.band_interval.total}{166} & \raggedright [\resultnum{v4x.x6.x6a.strengths.r070.a0100.vs_pop.band_interval.wilson95.i000}{0.957141}, \resultnum{v4x.x6.x6a.strengths.r070.a0100.vs_pop.band_interval.wilson95.i001}{0.99669}] & \raggedright no undercoverage detected \tabularnewline
\botrule
\end{tabular}
\end{table*}
\begin{table*}[htbp]
\setlength{\tabcolsep}{2pt}
\addtocounter{table}{-\csname @ne\endcsname}
\providecommand{\theHtable}{\thetable}
\renewcommand{\theHtable}{v4x-x6-part1}
\caption{Band coverage check (X6) (continued, part 2 of 3)}\label{tab:v4x-x6-part1}
\centering\footnotesize
\begin{tabular}{@{}p{0.108\textwidth}p{0.118\textwidth}p{0.128\textwidth}p{0.069\textwidth}p{0.098\textwidth}p{0.197\textwidth}p{0.207\textwidth}@{}}
\toprule
\raggedright $a,r$ & \raggedright \mbox{Interval} & \raggedright \mbox{Truth} & \raggedright \mbox{Gaps} \mbox{as} \mbox{misses} & \raggedright \mbox{Covered} \mbox{/} \mbox{total} & \raggedright \mbox{Wilson} \mbox{interval} & \raggedright \mbox{Reading} \tabularnewline
\midrule
\raggedright \resultnum{v4x.x6.x6a.strengths.r070.a0300.a}{0.03}, \resultnum{v4x.x6.x6a.strengths.r070.a0300.r}{0.7} & \raggedright pointwise & \raggedright procedural & \raggedright no & \raggedright \resultnum{v4x.x6.x6a.strengths.r070.a0300.vs_d1.pointwise.covered}{159}/\resultnum{v4x.x6.x6a.strengths.r070.a0300.vs_d1.pointwise.total}{166} & \raggedright [\resultnum{v4x.x6.x6a.strengths.r070.a0300.vs_d1.pointwise.wilson95.i000}{0.915527}, \resultnum{v4x.x6.x6a.strengths.r070.a0300.vs_d1.pointwise.wilson95.i001}{0.979426}] & \raggedright no undercoverage detected \tabularnewline
\raggedright \resultnum{v4x.x6.x6a.strengths.r070.a0300.a}{0.03}, \resultnum{v4x.x6.x6a.strengths.r070.a0300.r}{0.7} & \raggedright pointwise & \raggedright procedural & \raggedright yes & \raggedright \resultnum{v4x.x6.x6a.strengths.r070.a0300.vs_d1.pointwise_gaps_as_misses.covered}{159}/\resultnum{v4x.x6.x6a.strengths.r070.a0300.vs_d1.pointwise_gaps_as_misses.total}{200} & \raggedright [\resultnum{v4x.x6.x6a.strengths.r070.a0300.vs_d1.pointwise_gaps_as_misses.wilson95.i000}{0.733743}, \resultnum{v4x.x6.x6a.strengths.r070.a0300.vs_d1.pointwise_gaps_as_misses.wilson95.i001}{0.845138}] & \raggedright failure \tabularnewline
\raggedright \resultnum{v4x.x6.x6a.strengths.r070.a0300.a}{0.03}, \resultnum{v4x.x6.x6a.strengths.r070.a0300.r}{0.7} & \raggedright band & \raggedright procedural & \raggedright no & \raggedright \resultnum{v4x.x6.x6a.strengths.r070.a0300.vs_d1.band_interval.covered}{166}/\resultnum{v4x.x6.x6a.strengths.r070.a0300.vs_d1.band_interval.total}{166} & \raggedright [\resultnum{v4x.x6.x6a.strengths.r070.a0300.vs_d1.band_interval.wilson95.i000}{0.977382}, \resultnum{v4x.x6.x6a.strengths.r070.a0300.vs_d1.band_interval.wilson95.i001}{1}] & \raggedright no undercoverage detected \tabularnewline
\raggedright \resultnum{v4x.x6.x6a.strengths.r070.a0300.a}{0.03}, \resultnum{v4x.x6.x6a.strengths.r070.a0300.r}{0.7} & \raggedright pointwise & \raggedright population & \raggedright no & \raggedright \resultnum{v4x.x6.x6a.strengths.r070.a0300.vs_pop.pointwise.covered}{156}/\resultnum{v4x.x6.x6a.strengths.r070.a0300.vs_pop.pointwise.total}{166} & \raggedright [\resultnum{v4x.x6.x6a.strengths.r070.a0300.vs_pop.pointwise.wilson95.i000}{0.892673}, \resultnum{v4x.x6.x6a.strengths.r070.a0300.vs_pop.pointwise.wilson95.i001}{0.966953}] & \raggedright no undercoverage detected \tabularnewline
\raggedright \resultnum{v4x.x6.x6a.strengths.r070.a0300.a}{0.03}, \resultnum{v4x.x6.x6a.strengths.r070.a0300.r}{0.7} & \raggedright pointwise & \raggedright population & \raggedright yes & \raggedright \resultnum{v4x.x6.x6a.strengths.r070.a0300.vs_pop.pointwise_gaps_as_misses.covered}{156}/\resultnum{v4x.x6.x6a.strengths.r070.a0300.vs_pop.pointwise_gaps_as_misses.total}{200} & \raggedright [\resultnum{v4x.x6.x6a.strengths.r070.a0300.vs_pop.pointwise_gaps_as_misses.wilson95.i000}{0.717612}, \resultnum{v4x.x6.x6a.strengths.r070.a0300.vs_pop.pointwise_gaps_as_misses.wilson95.i001}{0.831835}] & \raggedright failure \tabularnewline
\raggedright \resultnum{v4x.x6.x6a.strengths.r070.a0300.a}{0.03}, \resultnum{v4x.x6.x6a.strengths.r070.a0300.r}{0.7} & \raggedright band & \raggedright population & \raggedright no & \raggedright \resultnum{v4x.x6.x6a.strengths.r070.a0300.vs_pop.band_interval.covered}{165}/\resultnum{v4x.x6.x6a.strengths.r070.a0300.vs_pop.band_interval.total}{166} & \raggedright [\resultnum{v4x.x6.x6a.strengths.r070.a0300.vs_pop.band_interval.wilson95.i000}{0.966671}, \resultnum{v4x.x6.x6a.strengths.r070.a0300.vs_pop.band_interval.wilson95.i001}{0.998936}] & \raggedright no undercoverage detected \tabularnewline
\raggedright \resultnum{v4x.x6.x6a.strengths.r070.a1000.a}{0.1}, \resultnum{v4x.x6.x6a.strengths.r070.a1000.r}{0.7} & \raggedright pointwise & \raggedright procedural & \raggedright no & \raggedright \resultnum{v4x.x6.x6a.strengths.r070.a1000.vs_d1.pointwise.covered}{161}/\resultnum{v4x.x6.x6a.strengths.r070.a1000.vs_d1.pointwise.total}{166} & \raggedright [\resultnum{v4x.x6.x6a.strengths.r070.a1000.vs_d1.pointwise.wilson95.i000}{0.931437}, \resultnum{v4x.x6.x6a.strengths.r070.a1000.vs_d1.pointwise.wilson95.i001}{0.987067}] & \raggedright no undercoverage detected \tabularnewline
\raggedright \resultnum{v4x.x6.x6a.strengths.r070.a1000.a}{0.1}, \resultnum{v4x.x6.x6a.strengths.r070.a1000.r}{0.7} & \raggedright pointwise & \raggedright procedural & \raggedright yes & \raggedright \resultnum{v4x.x6.x6a.strengths.r070.a1000.vs_d1.pointwise_gaps_as_misses.covered}{161}/\resultnum{v4x.x6.x6a.strengths.r070.a1000.vs_d1.pointwise_gaps_as_misses.total}{200} & \raggedright [\resultnum{v4x.x6.x6a.strengths.r070.a1000.vs_d1.pointwise_gaps_as_misses.wilson95.i000}{0.74456}, \resultnum{v4x.x6.x6a.strengths.r070.a1000.vs_d1.pointwise_gaps_as_misses.wilson95.i001}{0.853945}] & \raggedright failure \tabularnewline
\raggedright \resultnum{v4x.x6.x6a.strengths.r070.a1000.a}{0.1}, \resultnum{v4x.x6.x6a.strengths.r070.a1000.r}{0.7} & \raggedright band & \raggedright procedural & \raggedright no & \raggedright \resultnum{v4x.x6.x6a.strengths.r070.a1000.vs_d1.band_interval.covered}{165}/\resultnum{v4x.x6.x6a.strengths.r070.a1000.vs_d1.band_interval.total}{166} & \raggedright [\resultnum{v4x.x6.x6a.strengths.r070.a1000.vs_d1.band_interval.wilson95.i000}{0.966671}, \resultnum{v4x.x6.x6a.strengths.r070.a1000.vs_d1.band_interval.wilson95.i001}{0.998936}] & \raggedright no undercoverage detected \tabularnewline
\raggedright \resultnum{v4x.x6.x6a.strengths.r070.a1000.a}{0.1}, \resultnum{v4x.x6.x6a.strengths.r070.a1000.r}{0.7} & \raggedright pointwise & \raggedright population & \raggedright no & \raggedright \resultnum{v4x.x6.x6a.strengths.r070.a1000.vs_pop.pointwise.covered}{158}/\resultnum{v4x.x6.x6a.strengths.r070.a1000.vs_pop.pointwise.total}{166} & \raggedright [\resultnum{v4x.x6.x6a.strengths.r070.a1000.vs_pop.pointwise.wilson95.i000}{0.907796}, \resultnum{v4x.x6.x6a.strengths.r070.a1000.vs_pop.pointwise.wilson95.i001}{0.975381}] & \raggedright no undercoverage detected \tabularnewline
\raggedright \resultnum{v4x.x6.x6a.strengths.r070.a1000.a}{0.1}, \resultnum{v4x.x6.x6a.strengths.r070.a1000.r}{0.7} & \raggedright pointwise & \raggedright population & \raggedright yes & \raggedright \resultnum{v4x.x6.x6a.strengths.r070.a1000.vs_pop.pointwise_gaps_as_misses.covered}{158}/\resultnum{v4x.x6.x6a.strengths.r070.a1000.vs_pop.pointwise_gaps_as_misses.total}{200} & \raggedright [\resultnum{v4x.x6.x6a.strengths.r070.a1000.vs_pop.pointwise_gaps_as_misses.wilson95.i000}{0.728354}, \resultnum{v4x.x6.x6a.strengths.r070.a1000.vs_pop.pointwise_gaps_as_misses.wilson95.i001}{0.840716}] & \raggedright failure \tabularnewline
\raggedright \resultnum{v4x.x6.x6a.strengths.r070.a1000.a}{0.1}, \resultnum{v4x.x6.x6a.strengths.r070.a1000.r}{0.7} & \raggedright band & \raggedright population & \raggedright no & \raggedright \resultnum{v4x.x6.x6a.strengths.r070.a1000.vs_pop.band_interval.covered}{165}/\resultnum{v4x.x6.x6a.strengths.r070.a1000.vs_pop.band_interval.total}{166} & \raggedright [\resultnum{v4x.x6.x6a.strengths.r070.a1000.vs_pop.band_interval.wilson95.i000}{0.966671}, \resultnum{v4x.x6.x6a.strengths.r070.a1000.vs_pop.band_interval.wilson95.i001}{0.998936}] & \raggedright no undercoverage detected \tabularnewline
\raggedright \resultnum{v4x.x6.x6a.strengths.r070.a3000.a}{0.3}, \resultnum{v4x.x6.x6a.strengths.r070.a3000.r}{0.7} & \raggedright pointwise & \raggedright procedural & \raggedright no & \raggedright \resultnum{v4x.x6.x6a.strengths.r070.a3000.vs_d1.pointwise.covered}{159}/\resultnum{v4x.x6.x6a.strengths.r070.a3000.vs_d1.pointwise.total}{166} & \raggedright [\resultnum{v4x.x6.x6a.strengths.r070.a3000.vs_d1.pointwise.wilson95.i000}{0.915527}, \resultnum{v4x.x6.x6a.strengths.r070.a3000.vs_d1.pointwise.wilson95.i001}{0.979426}] & \raggedright no undercoverage detected \tabularnewline
\raggedright \resultnum{v4x.x6.x6a.strengths.r070.a3000.a}{0.3}, \resultnum{v4x.x6.x6a.strengths.r070.a3000.r}{0.7} & \raggedright pointwise & \raggedright procedural & \raggedright yes & \raggedright \resultnum{v4x.x6.x6a.strengths.r070.a3000.vs_d1.pointwise_gaps_as_misses.covered}{159}/\resultnum{v4x.x6.x6a.strengths.r070.a3000.vs_d1.pointwise_gaps_as_misses.total}{200} & \raggedright [\resultnum{v4x.x6.x6a.strengths.r070.a3000.vs_d1.pointwise_gaps_as_misses.wilson95.i000}{0.733743}, \resultnum{v4x.x6.x6a.strengths.r070.a3000.vs_d1.pointwise_gaps_as_misses.wilson95.i001}{0.845138}] & \raggedright failure \tabularnewline
\raggedright \resultnum{v4x.x6.x6a.strengths.r070.a3000.a}{0.3}, \resultnum{v4x.x6.x6a.strengths.r070.a3000.r}{0.7} & \raggedright band & \raggedright procedural & \raggedright no & \raggedright \resultnum{v4x.x6.x6a.strengths.r070.a3000.vs_d1.band_interval.covered}{166}/\resultnum{v4x.x6.x6a.strengths.r070.a3000.vs_d1.band_interval.total}{166} & \raggedright [\resultnum{v4x.x6.x6a.strengths.r070.a3000.vs_d1.band_interval.wilson95.i000}{0.977382}, \resultnum{v4x.x6.x6a.strengths.r070.a3000.vs_d1.band_interval.wilson95.i001}{1}] & \raggedright no undercoverage detected \tabularnewline
\raggedright \resultnum{v4x.x6.x6a.strengths.r070.a3000.a}{0.3}, \resultnum{v4x.x6.x6a.strengths.r070.a3000.r}{0.7} & \raggedright pointwise & \raggedright population & \raggedright no & \raggedright \resultnum{v4x.x6.x6a.strengths.r070.a3000.vs_pop.pointwise.covered}{158}/\resultnum{v4x.x6.x6a.strengths.r070.a3000.vs_pop.pointwise.total}{166} & \raggedright [\resultnum{v4x.x6.x6a.strengths.r070.a3000.vs_pop.pointwise.wilson95.i000}{0.907796}, \resultnum{v4x.x6.x6a.strengths.r070.a3000.vs_pop.pointwise.wilson95.i001}{0.975381}] & \raggedright no undercoverage detected \tabularnewline
\raggedright \resultnum{v4x.x6.x6a.strengths.r070.a3000.a}{0.3}, \resultnum{v4x.x6.x6a.strengths.r070.a3000.r}{0.7} & \raggedright pointwise & \raggedright population & \raggedright yes & \raggedright \resultnum{v4x.x6.x6a.strengths.r070.a3000.vs_pop.pointwise_gaps_as_misses.covered}{158}/\resultnum{v4x.x6.x6a.strengths.r070.a3000.vs_pop.pointwise_gaps_as_misses.total}{200} & \raggedright [\resultnum{v4x.x6.x6a.strengths.r070.a3000.vs_pop.pointwise_gaps_as_misses.wilson95.i000}{0.728354}, \resultnum{v4x.x6.x6a.strengths.r070.a3000.vs_pop.pointwise_gaps_as_misses.wilson95.i001}{0.840716}] & \raggedright failure \tabularnewline
\raggedright \resultnum{v4x.x6.x6a.strengths.r070.a3000.a}{0.3}, \resultnum{v4x.x6.x6a.strengths.r070.a3000.r}{0.7} & \raggedright band & \raggedright population & \raggedright no & \raggedright \resultnum{v4x.x6.x6a.strengths.r070.a3000.vs_pop.band_interval.covered}{166}/\resultnum{v4x.x6.x6a.strengths.r070.a3000.vs_pop.band_interval.total}{166} & \raggedright [\resultnum{v4x.x6.x6a.strengths.r070.a3000.vs_pop.band_interval.wilson95.i000}{0.977382}, \resultnum{v4x.x6.x6a.strengths.r070.a3000.vs_pop.band_interval.wilson95.i001}{1}] & \raggedright no undercoverage detected \tabularnewline
\raggedright \resultnum{v4x.x6.x6a.strengths.r070.a6000.a}{0.6}, \resultnum{v4x.x6.x6a.strengths.r070.a6000.r}{0.7} & \raggedright pointwise & \raggedright procedural & \raggedright no & \raggedright \resultnum{v4x.x6.x6a.strengths.r070.a6000.vs_d1.pointwise.covered}{158}/\resultnum{v4x.x6.x6a.strengths.r070.a6000.vs_d1.pointwise.total}{166} & \raggedright [\resultnum{v4x.x6.x6a.strengths.r070.a6000.vs_d1.pointwise.wilson95.i000}{0.907796}, \resultnum{v4x.x6.x6a.strengths.r070.a6000.vs_d1.pointwise.wilson95.i001}{0.975381}] & \raggedright no undercoverage detected \tabularnewline
\raggedright \resultnum{v4x.x6.x6a.strengths.r070.a6000.a}{0.6}, \resultnum{v4x.x6.x6a.strengths.r070.a6000.r}{0.7} & \raggedright pointwise & \raggedright procedural & \raggedright yes & \raggedright \resultnum{v4x.x6.x6a.strengths.r070.a6000.vs_d1.pointwise_gaps_as_misses.covered}{158}/\resultnum{v4x.x6.x6a.strengths.r070.a6000.vs_d1.pointwise_gaps_as_misses.total}{200} & \raggedright [\resultnum{v4x.x6.x6a.strengths.r070.a6000.vs_d1.pointwise_gaps_as_misses.wilson95.i000}{0.728354}, \resultnum{v4x.x6.x6a.strengths.r070.a6000.vs_d1.pointwise_gaps_as_misses.wilson95.i001}{0.840716}] & \raggedright failure \tabularnewline
\raggedright \resultnum{v4x.x6.x6a.strengths.r070.a6000.a}{0.6}, \resultnum{v4x.x6.x6a.strengths.r070.a6000.r}{0.7} & \raggedright band & \raggedright procedural & \raggedright no & \raggedright \resultnum{v4x.x6.x6a.strengths.r070.a6000.vs_d1.band_interval.covered}{165}/\resultnum{v4x.x6.x6a.strengths.r070.a6000.vs_d1.band_interval.total}{166} & \raggedright [\resultnum{v4x.x6.x6a.strengths.r070.a6000.vs_d1.band_interval.wilson95.i000}{0.966671}, \resultnum{v4x.x6.x6a.strengths.r070.a6000.vs_d1.band_interval.wilson95.i001}{0.998936}] & \raggedright no undercoverage detected \tabularnewline
\raggedright \resultnum{v4x.x6.x6a.strengths.r070.a6000.a}{0.6}, \resultnum{v4x.x6.x6a.strengths.r070.a6000.r}{0.7} & \raggedright pointwise & \raggedright population & \raggedright no & \raggedright \resultnum{v4x.x6.x6a.strengths.r070.a6000.vs_pop.pointwise.covered}{158}/\resultnum{v4x.x6.x6a.strengths.r070.a6000.vs_pop.pointwise.total}{166} & \raggedright [\resultnum{v4x.x6.x6a.strengths.r070.a6000.vs_pop.pointwise.wilson95.i000}{0.907796}, \resultnum{v4x.x6.x6a.strengths.r070.a6000.vs_pop.pointwise.wilson95.i001}{0.975381}] & \raggedright no undercoverage detected \tabularnewline
\raggedright \resultnum{v4x.x6.x6a.strengths.r070.a6000.a}{0.6}, \resultnum{v4x.x6.x6a.strengths.r070.a6000.r}{0.7} & \raggedright pointwise & \raggedright population & \raggedright yes & \raggedright \resultnum{v4x.x6.x6a.strengths.r070.a6000.vs_pop.pointwise_gaps_as_misses.covered}{158}/\resultnum{v4x.x6.x6a.strengths.r070.a6000.vs_pop.pointwise_gaps_as_misses.total}{200} & \raggedright [\resultnum{v4x.x6.x6a.strengths.r070.a6000.vs_pop.pointwise_gaps_as_misses.wilson95.i000}{0.728354}, \resultnum{v4x.x6.x6a.strengths.r070.a6000.vs_pop.pointwise_gaps_as_misses.wilson95.i001}{0.840716}] & \raggedright failure \tabularnewline
\raggedright \resultnum{v4x.x6.x6a.strengths.r070.a6000.a}{0.6}, \resultnum{v4x.x6.x6a.strengths.r070.a6000.r}{0.7} & \raggedright band & \raggedright population & \raggedright no & \raggedright \resultnum{v4x.x6.x6a.strengths.r070.a6000.vs_pop.band_interval.covered}{165}/\resultnum{v4x.x6.x6a.strengths.r070.a6000.vs_pop.band_interval.total}{166} & \raggedright [\resultnum{v4x.x6.x6a.strengths.r070.a6000.vs_pop.band_interval.wilson95.i000}{0.966671}, \resultnum{v4x.x6.x6a.strengths.r070.a6000.vs_pop.band_interval.wilson95.i001}{0.998936}] & \raggedright no undercoverage detected \tabularnewline
\raggedright \resultnum{v4x.x6.x6a.strengths.r070.a9000.a}{0.9}, \resultnum{v4x.x6.x6a.strengths.r070.a9000.r}{0.7} & \raggedright pointwise & \raggedright procedural & \raggedright no & \raggedright \resultnum{v4x.x6.x6a.strengths.r070.a9000.vs_d1.pointwise.covered}{166}/\resultnum{v4x.x6.x6a.strengths.r070.a9000.vs_d1.pointwise.total}{166} & \raggedright [\resultnum{v4x.x6.x6a.strengths.r070.a9000.vs_d1.pointwise.wilson95.i000}{0.977382}, \resultnum{v4x.x6.x6a.strengths.r070.a9000.vs_d1.pointwise.wilson95.i001}{1}] & \raggedright no undercoverage detected \tabularnewline
\raggedright \resultnum{v4x.x6.x6a.strengths.r070.a9000.a}{0.9}, \resultnum{v4x.x6.x6a.strengths.r070.a9000.r}{0.7} & \raggedright pointwise & \raggedright procedural & \raggedright yes & \raggedright \resultnum{v4x.x6.x6a.strengths.r070.a9000.vs_d1.pointwise_gaps_as_misses.covered}{166}/\resultnum{v4x.x6.x6a.strengths.r070.a9000.vs_d1.pointwise_gaps_as_misses.total}{200} & \raggedright [\resultnum{v4x.x6.x6a.strengths.r070.a9000.vs_d1.pointwise_gaps_as_misses.wilson95.i000}{0.771841}, \resultnum{v4x.x6.x6a.strengths.r070.a9000.vs_d1.pointwise_gaps_as_misses.wilson95.i001}{0.875721}] & \raggedright failure \tabularnewline
\raggedright \resultnum{v4x.x6.x6a.strengths.r070.a9000.a}{0.9}, \resultnum{v4x.x6.x6a.strengths.r070.a9000.r}{0.7} & \raggedright band & \raggedright procedural & \raggedright no & \raggedright \resultnum{v4x.x6.x6a.strengths.r070.a9000.vs_d1.band_interval.covered}{166}/\resultnum{v4x.x6.x6a.strengths.r070.a9000.vs_d1.band_interval.total}{166} & \raggedright [\resultnum{v4x.x6.x6a.strengths.r070.a9000.vs_d1.band_interval.wilson95.i000}{0.977382}, \resultnum{v4x.x6.x6a.strengths.r070.a9000.vs_d1.band_interval.wilson95.i001}{1}] & \raggedright no undercoverage detected \tabularnewline
\raggedright \resultnum{v4x.x6.x6a.strengths.r070.a9000.a}{0.9}, \resultnum{v4x.x6.x6a.strengths.r070.a9000.r}{0.7} & \raggedright pointwise & \raggedright population & \raggedright no & \raggedright \resultnum{v4x.x6.x6a.strengths.r070.a9000.vs_pop.pointwise.covered}{166}/\resultnum{v4x.x6.x6a.strengths.r070.a9000.vs_pop.pointwise.total}{166} & \raggedright [\resultnum{v4x.x6.x6a.strengths.r070.a9000.vs_pop.pointwise.wilson95.i000}{0.977382}, \resultnum{v4x.x6.x6a.strengths.r070.a9000.vs_pop.pointwise.wilson95.i001}{1}] & \raggedright no undercoverage detected \tabularnewline
\raggedright \resultnum{v4x.x6.x6a.strengths.r070.a9000.a}{0.9}, \resultnum{v4x.x6.x6a.strengths.r070.a9000.r}{0.7} & \raggedright pointwise & \raggedright population & \raggedright yes & \raggedright \resultnum{v4x.x6.x6a.strengths.r070.a9000.vs_pop.pointwise_gaps_as_misses.covered}{166}/\resultnum{v4x.x6.x6a.strengths.r070.a9000.vs_pop.pointwise_gaps_as_misses.total}{200} & \raggedright [\resultnum{v4x.x6.x6a.strengths.r070.a9000.vs_pop.pointwise_gaps_as_misses.wilson95.i000}{0.771841}, \resultnum{v4x.x6.x6a.strengths.r070.a9000.vs_pop.pointwise_gaps_as_misses.wilson95.i001}{0.875721}] & \raggedright failure \tabularnewline
\raggedright \resultnum{v4x.x6.x6a.strengths.r070.a9000.a}{0.9}, \resultnum{v4x.x6.x6a.strengths.r070.a9000.r}{0.7} & \raggedright band & \raggedright population & \raggedright no & \raggedright \resultnum{v4x.x6.x6a.strengths.r070.a9000.vs_pop.band_interval.covered}{166}/\resultnum{v4x.x6.x6a.strengths.r070.a9000.vs_pop.band_interval.total}{166} & \raggedright [\resultnum{v4x.x6.x6a.strengths.r070.a9000.vs_pop.band_interval.wilson95.i000}{0.977382}, \resultnum{v4x.x6.x6a.strengths.r070.a9000.vs_pop.band_interval.wilson95.i001}{1}] & \raggedright no undercoverage detected \tabularnewline
\raggedright \resultnum{v4x.x6.x6b.cells.r050.a0300.a}{0.03}, \resultnum{v4x.x6.x6b.cells.r050.a0300.r}{0.5} & \raggedright pointwise & \raggedright procedural & \raggedright no & \raggedright \resultnum{v4x.x6.x6b.cells.r050.a0300.vs_d1.pointwise.covered}{193}/\resultnum{v4x.x6.x6b.cells.r050.a0300.vs_d1.pointwise.total}{200} & \raggedright [\resultnum{v4x.x6.x6b.cells.r050.a0300.vs_d1.pointwise.wilson95.i000}{0.929529}, \resultnum{v4x.x6.x6b.cells.r050.a0300.vs_d1.pointwise.wilson95.i001}{0.982944}] & \raggedright no undercoverage detected \tabularnewline
\raggedright \resultnum{v4x.x6.x6b.cells.r050.a0300.a}{0.03}, \resultnum{v4x.x6.x6b.cells.r050.a0300.r}{0.5} & \raggedright pointwise & \raggedright procedural & \raggedright yes & \raggedright \resultnum{v4x.x6.x6b.cells.r050.a0300.vs_d1.pointwise_gaps_as_misses.covered}{193}/\resultnum{v4x.x6.x6b.cells.r050.a0300.vs_d1.pointwise_gaps_as_misses.total}{200} & \raggedright [\resultnum{v4x.x6.x6b.cells.r050.a0300.vs_d1.pointwise_gaps_as_misses.wilson95.i000}{0.929529}, \resultnum{v4x.x6.x6b.cells.r050.a0300.vs_d1.pointwise_gaps_as_misses.wilson95.i001}{0.982944}] & \raggedright no undercoverage detected \tabularnewline
\botrule
\end{tabular}
\end{table*}
\begin{table*}[htbp]
\setlength{\tabcolsep}{2pt}
\addtocounter{table}{-\csname @ne\endcsname}
\providecommand{\theHtable}{\thetable}
\renewcommand{\theHtable}{v4x-x6-part2}
\caption{Band coverage check (X6) (continued, part 3 of 3)}\label{tab:v4x-x6-part2}
\centering\footnotesize
\begin{tabular}{@{}p{0.108\textwidth}p{0.118\textwidth}p{0.128\textwidth}p{0.069\textwidth}p{0.098\textwidth}p{0.197\textwidth}p{0.207\textwidth}@{}}
\toprule
\raggedright $a,r$ & \raggedright \mbox{Interval} & \raggedright \mbox{Truth} & \raggedright \mbox{Gaps} \mbox{as} \mbox{misses} & \raggedright \mbox{Covered} \mbox{/} \mbox{total} & \raggedright \mbox{Wilson} \mbox{interval} & \raggedright \mbox{Reading} \tabularnewline
\midrule
\raggedright \resultnum{v4x.x6.x6b.cells.r050.a0300.a}{0.03}, \resultnum{v4x.x6.x6b.cells.r050.a0300.r}{0.5} & \raggedright pointwise & \raggedright population & \raggedright no & \raggedright \resultnum{v4x.x6.x6b.cells.r050.a0300.vs_pop.pointwise.covered}{193}/\resultnum{v4x.x6.x6b.cells.r050.a0300.vs_pop.pointwise.total}{200} & \raggedright [\resultnum{v4x.x6.x6b.cells.r050.a0300.vs_pop.pointwise.wilson95.i000}{0.929529}, \resultnum{v4x.x6.x6b.cells.r050.a0300.vs_pop.pointwise.wilson95.i001}{0.982944}] & \raggedright no undercoverage detected \tabularnewline
\raggedright \resultnum{v4x.x6.x6b.cells.r050.a0300.a}{0.03}, \resultnum{v4x.x6.x6b.cells.r050.a0300.r}{0.5} & \raggedright pointwise & \raggedright population & \raggedright yes & \raggedright \resultnum{v4x.x6.x6b.cells.r050.a0300.vs_pop.pointwise_gaps_as_misses.covered}{193}/\resultnum{v4x.x6.x6b.cells.r050.a0300.vs_pop.pointwise_gaps_as_misses.total}{200} & \raggedright [\resultnum{v4x.x6.x6b.cells.r050.a0300.vs_pop.pointwise_gaps_as_misses.wilson95.i000}{0.929529}, \resultnum{v4x.x6.x6b.cells.r050.a0300.vs_pop.pointwise_gaps_as_misses.wilson95.i001}{0.982944}] & \raggedright no undercoverage detected \tabularnewline
\raggedright \resultnum{v4x.x6.x6b.cells.r090.a0300.a}{0.03}, \resultnum{v4x.x6.x6b.cells.r090.a0300.r}{0.9} & \raggedright pointwise & \raggedright procedural & \raggedright no & \raggedright \resultnum{v4x.x6.x6b.cells.r090.a0300.vs_d1.pointwise.covered}{189}/\resultnum{v4x.x6.x6b.cells.r090.a0300.vs_d1.pointwise.total}{200} & \raggedright [\resultnum{v4x.x6.x6b.cells.r090.a0300.vs_d1.pointwise.wilson95.i000}{0.904213}, \resultnum{v4x.x6.x6b.cells.r090.a0300.vs_d1.pointwise.wilson95.i001}{0.969015}] & \raggedright no undercoverage detected \tabularnewline
\raggedright \resultnum{v4x.x6.x6b.cells.r090.a0300.a}{0.03}, \resultnum{v4x.x6.x6b.cells.r090.a0300.r}{0.9} & \raggedright pointwise & \raggedright procedural & \raggedright yes & \raggedright \resultnum{v4x.x6.x6b.cells.r090.a0300.vs_d1.pointwise_gaps_as_misses.covered}{189}/\resultnum{v4x.x6.x6b.cells.r090.a0300.vs_d1.pointwise_gaps_as_misses.total}{200} & \raggedright [\resultnum{v4x.x6.x6b.cells.r090.a0300.vs_d1.pointwise_gaps_as_misses.wilson95.i000}{0.904213}, \resultnum{v4x.x6.x6b.cells.r090.a0300.vs_d1.pointwise_gaps_as_misses.wilson95.i001}{0.969015}] & \raggedright no undercoverage detected \tabularnewline
\raggedright \resultnum{v4x.x6.x6b.cells.r090.a0300.a}{0.03}, \resultnum{v4x.x6.x6b.cells.r090.a0300.r}{0.9} & \raggedright pointwise & \raggedright population & \raggedright no & \raggedright \resultnum{v4x.x6.x6b.cells.r090.a0300.vs_pop.pointwise.covered}{189}/\resultnum{v4x.x6.x6b.cells.r090.a0300.vs_pop.pointwise.total}{200} & \raggedright [\resultnum{v4x.x6.x6b.cells.r090.a0300.vs_pop.pointwise.wilson95.i000}{0.904213}, \resultnum{v4x.x6.x6b.cells.r090.a0300.vs_pop.pointwise.wilson95.i001}{0.969015}] & \raggedright no undercoverage detected \tabularnewline
\raggedright \resultnum{v4x.x6.x6b.cells.r090.a0300.a}{0.03}, \resultnum{v4x.x6.x6b.cells.r090.a0300.r}{0.9} & \raggedright pointwise & \raggedright population & \raggedright yes & \raggedright \resultnum{v4x.x6.x6b.cells.r090.a0300.vs_pop.pointwise_gaps_as_misses.covered}{189}/\resultnum{v4x.x6.x6b.cells.r090.a0300.vs_pop.pointwise_gaps_as_misses.total}{200} & \raggedright [\resultnum{v4x.x6.x6b.cells.r090.a0300.vs_pop.pointwise_gaps_as_misses.wilson95.i000}{0.904213}, \resultnum{v4x.x6.x6b.cells.r090.a0300.vs_pop.pointwise_gaps_as_misses.wilson95.i001}{0.969015}] & \raggedright no undercoverage detected \tabularnewline
\botrule
\end{tabular}
\end{table*}

% Generated by v4_exploratory_export; edited during the journal revision
% (agent S, 2026-10-13): captions, headers and categorical labels translated
% into plain English. Every number keeps its \resultnum binding.
\begin{table*}[htbp]
\setlength{\tabcolsep}{2pt}
\caption{Band coverage check (X6): denominators. For the whole unit at $r=0.70$ (X6a), complete units are the repetitions whose band is complete and gap-free over the whole grid, successful cells are the repetitions whose pipeline succeeded at that strength, and units with any grid gap count the repetitions whose band had a gap at some strength of the grid (\resultnum{v4x.x6.x6a.strengths.r070.a0000.gaps}{34} of \resultnum{v4x.x6.x6a.band.vs_d1.gaps_as_misses.total}{200}). That count is a property of the whole repetition and is therefore repeated at every strength; it is not the number of failed cells at that strength. A successful cell of an incomplete repetition is counted among successful cells but does not enter the primary coverage, which is conditioned on complete units; this is why successful cells can exceed complete units. For the two X6b cells, the unit is the cell pipeline and the gap column counts failed cell pipelines. Errors would be reported outside both denominators.}\label{tab:v4x-x6-denominators}
\centering\footnotesize
\begin{tabular}{@{}p{0.150\textwidth}p{0.210\textwidth}p{0.150\textwidth}p{0.170\textwidth}p{0.100\textwidth}p{0.120\textwidth}@{}}
\toprule
\raggedright $a,r$ & \raggedright \mbox{Complete} \mbox{units} \mbox{(X6a)} \mbox{or} \mbox{cells} \mbox{(X6b)} & \raggedright \mbox{Successful} \mbox{cells} & \raggedright \mbox{Units} \mbox{with} \mbox{any} \mbox{grid} \mbox{gap} & \raggedright \mbox{Errors} & \raggedright \mbox{Retry} \mbox{fraction} \tabularnewline
\midrule
\raggedright \resultnum{v4x.x6.x6a.strengths.r070.a0000.a}{0}, \resultnum{v4x.x6.x6a.strengths.r070.a0000.r}{0.7} & \raggedright \resultnum{v4x.x6.x6a.strengths.r070.a0000.complete}{166} & \raggedright \resultnum{v4x.x6.x6a.strengths.r070.a0000.cell_complete}{200} & \raggedright \resultnum{v4x.x6.x6a.strengths.r070.a0000.gaps}{34} & \raggedright \resultnum{v4x.x6.x6a.strengths.r070.a0000.errors}{0} & \raggedright \resultnum{v4x.x6.x6a.strengths.r070.a0000.retry_fraction}{0} \tabularnewline
\raggedright \resultnum{v4x.x6.x6a.strengths.r070.a0010.a}{0.001}, \resultnum{v4x.x6.x6a.strengths.r070.a0010.r}{0.7} & \raggedright \resultnum{v4x.x6.x6a.strengths.r070.a0010.complete}{166} & \raggedright \resultnum{v4x.x6.x6a.strengths.r070.a0010.cell_complete}{199} & \raggedright \resultnum{v4x.x6.x6a.strengths.r070.a0010.gaps}{34} & \raggedright \resultnum{v4x.x6.x6a.strengths.r070.a0010.errors}{0} & \raggedright \resultnum{v4x.x6.x6a.strengths.r070.a0010.retry_fraction}{0} \tabularnewline
\raggedright \resultnum{v4x.x6.x6a.strengths.r070.a0030.a}{0.003}, \resultnum{v4x.x6.x6a.strengths.r070.a0030.r}{0.7} & \raggedright \resultnum{v4x.x6.x6a.strengths.r070.a0030.complete}{166} & \raggedright \resultnum{v4x.x6.x6a.strengths.r070.a0030.cell_complete}{200} & \raggedright \resultnum{v4x.x6.x6a.strengths.r070.a0030.gaps}{34} & \raggedright \resultnum{v4x.x6.x6a.strengths.r070.a0030.errors}{0} & \raggedright \resultnum{v4x.x6.x6a.strengths.r070.a0030.retry_fraction}{0} \tabularnewline
\raggedright \resultnum{v4x.x6.x6a.strengths.r070.a0100.a}{0.01}, \resultnum{v4x.x6.x6a.strengths.r070.a0100.r}{0.7} & \raggedright \resultnum{v4x.x6.x6a.strengths.r070.a0100.complete}{166} & \raggedright \resultnum{v4x.x6.x6a.strengths.r070.a0100.cell_complete}{200} & \raggedright \resultnum{v4x.x6.x6a.strengths.r070.a0100.gaps}{34} & \raggedright \resultnum{v4x.x6.x6a.strengths.r070.a0100.errors}{0} & \raggedright \resultnum{v4x.x6.x6a.strengths.r070.a0100.retry_fraction}{0} \tabularnewline
\raggedright \resultnum{v4x.x6.x6a.strengths.r070.a0300.a}{0.03}, \resultnum{v4x.x6.x6a.strengths.r070.a0300.r}{0.7} & \raggedright \resultnum{v4x.x6.x6a.strengths.r070.a0300.complete}{166} & \raggedright \resultnum{v4x.x6.x6a.strengths.r070.a0300.cell_complete}{200} & \raggedright \resultnum{v4x.x6.x6a.strengths.r070.a0300.gaps}{34} & \raggedright \resultnum{v4x.x6.x6a.strengths.r070.a0300.errors}{0} & \raggedright \resultnum{v4x.x6.x6a.strengths.r070.a0300.retry_fraction}{0} \tabularnewline
\raggedright \resultnum{v4x.x6.x6a.strengths.r070.a1000.a}{0.1}, \resultnum{v4x.x6.x6a.strengths.r070.a1000.r}{0.7} & \raggedright \resultnum{v4x.x6.x6a.strengths.r070.a1000.complete}{166} & \raggedright \resultnum{v4x.x6.x6a.strengths.r070.a1000.cell_complete}{200} & \raggedright \resultnum{v4x.x6.x6a.strengths.r070.a1000.gaps}{34} & \raggedright \resultnum{v4x.x6.x6a.strengths.r070.a1000.errors}{0} & \raggedright \resultnum{v4x.x6.x6a.strengths.r070.a1000.retry_fraction}{0} \tabularnewline
\raggedright \resultnum{v4x.x6.x6a.strengths.r070.a3000.a}{0.3}, \resultnum{v4x.x6.x6a.strengths.r070.a3000.r}{0.7} & \raggedright \resultnum{v4x.x6.x6a.strengths.r070.a3000.complete}{166} & \raggedright \resultnum{v4x.x6.x6a.strengths.r070.a3000.cell_complete}{199} & \raggedright \resultnum{v4x.x6.x6a.strengths.r070.a3000.gaps}{34} & \raggedright \resultnum{v4x.x6.x6a.strengths.r070.a3000.errors}{0} & \raggedright \resultnum{v4x.x6.x6a.strengths.r070.a3000.retry_fraction}{0} \tabularnewline
\raggedright \resultnum{v4x.x6.x6a.strengths.r070.a6000.a}{0.6}, \resultnum{v4x.x6.x6a.strengths.r070.a6000.r}{0.7} & \raggedright \resultnum{v4x.x6.x6a.strengths.r070.a6000.complete}{166} & \raggedright \resultnum{v4x.x6.x6a.strengths.r070.a6000.cell_complete}{191} & \raggedright \resultnum{v4x.x6.x6a.strengths.r070.a6000.gaps}{34} & \raggedright \resultnum{v4x.x6.x6a.strengths.r070.a6000.errors}{0} & \raggedright \resultnum{v4x.x6.x6a.strengths.r070.a6000.retry_fraction}{0} \tabularnewline
\raggedright \resultnum{v4x.x6.x6a.strengths.r070.a9000.a}{0.9}, \resultnum{v4x.x6.x6a.strengths.r070.a9000.r}{0.7} & \raggedright \resultnum{v4x.x6.x6a.strengths.r070.a9000.complete}{166} & \raggedright \resultnum{v4x.x6.x6a.strengths.r070.a9000.cell_complete}{176} & \raggedright \resultnum{v4x.x6.x6a.strengths.r070.a9000.gaps}{34} & \raggedright \resultnum{v4x.x6.x6a.strengths.r070.a9000.errors}{0} & \raggedright \resultnum{v4x.x6.x6a.strengths.r070.a9000.retry_fraction}{0} \tabularnewline
\raggedright \resultnum{v4x.x6.x6b.cells.r050.a0300.a}{0.03}, \resultnum{v4x.x6.x6b.cells.r050.a0300.r}{0.5} & \raggedright \resultnum{v4x.x6.x6b.cells.r050.a0300.complete}{200} & \raggedright \resultnum{v4x.x6.x6b.cells.r050.a0300.cell_complete}{200} & \raggedright \resultnum{v4x.x6.x6b.cells.r050.a0300.gaps}{0} & \raggedright \resultnum{v4x.x6.x6b.cells.r050.a0300.errors}{0} & \raggedright \resultnum{v4x.x6.x6b.cells.r050.a0300.retry_fraction}{0} \tabularnewline
\raggedright \resultnum{v4x.x6.x6b.cells.r090.a0300.a}{0.03}, \resultnum{v4x.x6.x6b.cells.r090.a0300.r}{0.9} & \raggedright \resultnum{v4x.x6.x6b.cells.r090.a0300.complete}{200} & \raggedright \resultnum{v4x.x6.x6b.cells.r090.a0300.cell_complete}{200} & \raggedright \resultnum{v4x.x6.x6b.cells.r090.a0300.gaps}{0} & \raggedright \resultnum{v4x.x6.x6b.cells.r090.a0300.errors}{0} & \raggedright \resultnum{v4x.x6.x6b.cells.r090.a0300.retry_fraction}{0} \tabularnewline
\botrule
\end{tabular}
\end{table*}

% Generated by v4_exploratory_export; edited during the journal revision
% (agent S, 2026-10-13): captions, headers and categorical labels translated
% into plain English. Every number keeps its \resultnum binding.
\begin{table*}[htbp]
\setlength{\tabcolsep}{2pt}
\caption{Band coverage check (X6): the two truths. The procedural truth (machine label D1) is the mean AUC of 100 independent frozen pipelines per cell, conditional on success and including the declared retries; the population-match truth is the mean of 50 trainings at a fit to 16,384 coupled matching pairs, and its Monte Carlo error omits the error of that population fit. At $a=0$ both truths are 0.5 by construction and were not simulated. The retry fraction does not apply to the population-match truth. Every Monte Carlo error is below the declared target of 0.002.}\label{tab:v4x-x6-truth}
\centering\footnotesize
\begin{tabular}{@{}p{0.150\textwidth}p{0.170\textwidth}p{0.130\textwidth}p{0.150\textwidth}p{0.140\textwidth}p{0.160\textwidth}@{}}
\toprule
\raggedright $a,r$ & \raggedright \mbox{Truth} & \raggedright \mbox{Estimate} & \raggedright \mbox{MC} \mbox{error} & \raggedright \mbox{Retry} \mbox{fraction} & \raggedright \mbox{Successes} \mbox{/} \mbox{trainings} \tabularnewline
\midrule
\raggedright \resultnum{v4x.x6.x6a.strengths.r070.a0000.a}{0}, \resultnum{v4x.x6.x6a.strengths.r070.a0000.r}{0.7} & \raggedright procedural (D1) & \raggedright \resultnum{v4x.x6.x6a.strengths.r070.a0000.truth_d1.value}{0.5} & \raggedright \resultnum{v4x.x6.x6a.strengths.r070.a0000.truth_d1.mcse}{0} & \raggedright not simulated & \raggedright not simulated \tabularnewline
\raggedright \resultnum{v4x.x6.x6a.strengths.r070.a0000.a}{0}, \resultnum{v4x.x6.x6a.strengths.r070.a0000.r}{0.7} & \raggedright population match & \raggedright \resultnum{v4x.x6.x6a.strengths.r070.a0000.truth_pop.value}{0.5} & \raggedright \resultnum{v4x.x6.x6a.strengths.r070.a0000.truth_pop.mcse}{0} & \raggedright not simulated & \raggedright not simulated \tabularnewline
\raggedright \resultnum{v4x.x6.x6a.strengths.r070.a0010.a}{0.001}, \resultnum{v4x.x6.x6a.strengths.r070.a0010.r}{0.7} & \raggedright procedural (D1) & \raggedright \resultnum{v4x.x6.x6a.strengths.r070.a0010.truth_d1.value}{0.507497} & \raggedright \resultnum{v4x.x6.x6a.strengths.r070.a0010.truth_d1.mcse}{0.00103226} & \raggedright \resultnum{v4x.x6.x6a.strengths.r070.a0010.truth_d1.retry_fraction}{0} & \raggedright \resultnum{v4x.x6.x6a.strengths.r070.a0010.truth_d1.successful}{100} \tabularnewline
\raggedright \resultnum{v4x.x6.x6a.strengths.r070.a0010.a}{0.001}, \resultnum{v4x.x6.x6a.strengths.r070.a0010.r}{0.7} & \raggedright population match & \raggedright \resultnum{v4x.x6.x6a.strengths.r070.a0010.truth_pop.value}{0.508423} & \raggedright \resultnum{v4x.x6.x6a.strengths.r070.a0010.truth_pop.mcse}{0.00110905} & \raggedright not applicable & \raggedright \resultnum{v4x.x6.x6a.strengths.r070.a0010.truth_pop.trainings}{50} \tabularnewline
\raggedright \resultnum{v4x.x6.x6a.strengths.r070.a0030.a}{0.003}, \resultnum{v4x.x6.x6a.strengths.r070.a0030.r}{0.7} & \raggedright procedural (D1) & \raggedright \resultnum{v4x.x6.x6a.strengths.r070.a0030.truth_d1.value}{0.518494} & \raggedright \resultnum{v4x.x6.x6a.strengths.r070.a0030.truth_d1.mcse}{0.0010681} & \raggedright \resultnum{v4x.x6.x6a.strengths.r070.a0030.truth_d1.retry_fraction}{0} & \raggedright \resultnum{v4x.x6.x6a.strengths.r070.a0030.truth_d1.successful}{100} \tabularnewline
\raggedright \resultnum{v4x.x6.x6a.strengths.r070.a0030.a}{0.003}, \resultnum{v4x.x6.x6a.strengths.r070.a0030.r}{0.7} & \raggedright population match & \raggedright \resultnum{v4x.x6.x6a.strengths.r070.a0030.truth_pop.value}{0.52054} & \raggedright \resultnum{v4x.x6.x6a.strengths.r070.a0030.truth_pop.mcse}{0.00101405} & \raggedright not applicable & \raggedright \resultnum{v4x.x6.x6a.strengths.r070.a0030.truth_pop.trainings}{50} \tabularnewline
\raggedright \resultnum{v4x.x6.x6a.strengths.r070.a0100.a}{0.01}, \resultnum{v4x.x6.x6a.strengths.r070.a0100.r}{0.7} & \raggedright procedural (D1) & \raggedright \resultnum{v4x.x6.x6a.strengths.r070.a0100.truth_d1.value}{0.547701} & \raggedright \resultnum{v4x.x6.x6a.strengths.r070.a0100.truth_d1.mcse}{0.001012} & \raggedright \resultnum{v4x.x6.x6a.strengths.r070.a0100.truth_d1.retry_fraction}{0} & \raggedright \resultnum{v4x.x6.x6a.strengths.r070.a0100.truth_d1.successful}{100} \tabularnewline
\raggedright \resultnum{v4x.x6.x6a.strengths.r070.a0100.a}{0.01}, \resultnum{v4x.x6.x6a.strengths.r070.a0100.r}{0.7} & \raggedright population match & \raggedright \resultnum{v4x.x6.x6a.strengths.r070.a0100.truth_pop.value}{0.540234} & \raggedright \resultnum{v4x.x6.x6a.strengths.r070.a0100.truth_pop.mcse}{0.000926784} & \raggedright not applicable & \raggedright \resultnum{v4x.x6.x6a.strengths.r070.a0100.truth_pop.trainings}{50} \tabularnewline
\raggedright \resultnum{v4x.x6.x6a.strengths.r070.a0300.a}{0.03}, \resultnum{v4x.x6.x6a.strengths.r070.a0300.r}{0.7} & \raggedright procedural (D1) & \raggedright \resultnum{v4x.x6.x6a.strengths.r070.a0300.truth_d1.value}{0.605416} & \raggedright \resultnum{v4x.x6.x6a.strengths.r070.a0300.truth_d1.mcse}{0.00104722} & \raggedright \resultnum{v4x.x6.x6a.strengths.r070.a0300.truth_d1.retry_fraction}{0} & \raggedright \resultnum{v4x.x6.x6a.strengths.r070.a0300.truth_d1.successful}{100} \tabularnewline
\raggedright \resultnum{v4x.x6.x6a.strengths.r070.a0300.a}{0.03}, \resultnum{v4x.x6.x6a.strengths.r070.a0300.r}{0.7} & \raggedright population match & \raggedright \resultnum{v4x.x6.x6a.strengths.r070.a0300.truth_pop.value}{0.607754} & \raggedright \resultnum{v4x.x6.x6a.strengths.r070.a0300.truth_pop.mcse}{0.000916423} & \raggedright not applicable & \raggedright \resultnum{v4x.x6.x6a.strengths.r070.a0300.truth_pop.trainings}{50} \tabularnewline
\raggedright \resultnum{v4x.x6.x6a.strengths.r070.a1000.a}{0.1}, \resultnum{v4x.x6.x6a.strengths.r070.a1000.r}{0.7} & \raggedright procedural (D1) & \raggedright \resultnum{v4x.x6.x6a.strengths.r070.a1000.truth_d1.value}{0.749949} & \raggedright \resultnum{v4x.x6.x6a.strengths.r070.a1000.truth_d1.mcse}{0.00119131} & \raggedright \resultnum{v4x.x6.x6a.strengths.r070.a1000.truth_d1.retry_fraction}{0} & \raggedright \resultnum{v4x.x6.x6a.strengths.r070.a1000.truth_d1.successful}{100} \tabularnewline
\raggedright \resultnum{v4x.x6.x6a.strengths.r070.a1000.a}{0.1}, \resultnum{v4x.x6.x6a.strengths.r070.a1000.r}{0.7} & \raggedright population match & \raggedright \resultnum{v4x.x6.x6a.strengths.r070.a1000.truth_pop.value}{0.746519} & \raggedright \resultnum{v4x.x6.x6a.strengths.r070.a1000.truth_pop.mcse}{0.000832892} & \raggedright not applicable & \raggedright \resultnum{v4x.x6.x6a.strengths.r070.a1000.truth_pop.trainings}{50} \tabularnewline
\raggedright \resultnum{v4x.x6.x6a.strengths.r070.a3000.a}{0.3}, \resultnum{v4x.x6.x6a.strengths.r070.a3000.r}{0.7} & \raggedright procedural (D1) & \raggedright \resultnum{v4x.x6.x6a.strengths.r070.a3000.truth_d1.value}{0.953599} & \raggedright \resultnum{v4x.x6.x6a.strengths.r070.a3000.truth_d1.mcse}{0.000507309} & \raggedright \resultnum{v4x.x6.x6a.strengths.r070.a3000.truth_d1.retry_fraction}{0} & \raggedright \resultnum{v4x.x6.x6a.strengths.r070.a3000.truth_d1.successful}{98} \tabularnewline
\raggedright \resultnum{v4x.x6.x6a.strengths.r070.a3000.a}{0.3}, \resultnum{v4x.x6.x6a.strengths.r070.a3000.r}{0.7} & \raggedright population match & \raggedright \resultnum{v4x.x6.x6a.strengths.r070.a3000.truth_pop.value}{0.954624} & \raggedright \resultnum{v4x.x6.x6a.strengths.r070.a3000.truth_pop.mcse}{0.000309819} & \raggedright not applicable & \raggedright \resultnum{v4x.x6.x6a.strengths.r070.a3000.truth_pop.trainings}{50} \tabularnewline
\raggedright \resultnum{v4x.x6.x6a.strengths.r070.a6000.a}{0.6}, \resultnum{v4x.x6.x6a.strengths.r070.a6000.r}{0.7} & \raggedright procedural (D1) & \raggedright \resultnum{v4x.x6.x6a.strengths.r070.a6000.truth_d1.value}{0.998271} & \raggedright \resultnum{v4x.x6.x6a.strengths.r070.a6000.truth_d1.mcse}{3.63176e-05} & \raggedright \resultnum{v4x.x6.x6a.strengths.r070.a6000.truth_d1.retry_fraction}{0} & \raggedright \resultnum{v4x.x6.x6a.strengths.r070.a6000.truth_d1.successful}{92} \tabularnewline
\raggedright \resultnum{v4x.x6.x6a.strengths.r070.a6000.a}{0.6}, \resultnum{v4x.x6.x6a.strengths.r070.a6000.r}{0.7} & \raggedright population match & \raggedright \resultnum{v4x.x6.x6a.strengths.r070.a6000.truth_pop.value}{0.998264} & \raggedright \resultnum{v4x.x6.x6a.strengths.r070.a6000.truth_pop.mcse}{3.70024e-05} & \raggedright not applicable & \raggedright \resultnum{v4x.x6.x6a.strengths.r070.a6000.truth_pop.trainings}{50} \tabularnewline
\raggedright \resultnum{v4x.x6.x6a.strengths.r070.a9000.a}{0.9}, \resultnum{v4x.x6.x6a.strengths.r070.a9000.r}{0.7} & \raggedright procedural (D1) & \raggedright \resultnum{v4x.x6.x6a.strengths.r070.a9000.truth_d1.value}{0.999907} & \raggedright \resultnum{v4x.x6.x6a.strengths.r070.a9000.truth_d1.mcse}{3.84673e-06} & \raggedright \resultnum{v4x.x6.x6a.strengths.r070.a9000.truth_d1.retry_fraction}{0} & \raggedright \resultnum{v4x.x6.x6a.strengths.r070.a9000.truth_d1.successful}{81} \tabularnewline
\raggedright \resultnum{v4x.x6.x6a.strengths.r070.a9000.a}{0.9}, \resultnum{v4x.x6.x6a.strengths.r070.a9000.r}{0.7} & \raggedright population match & \raggedright \resultnum{v4x.x6.x6a.strengths.r070.a9000.truth_pop.value}{0.999907} & \raggedright \resultnum{v4x.x6.x6a.strengths.r070.a9000.truth_pop.mcse}{5.35683e-06} & \raggedright not applicable & \raggedright \resultnum{v4x.x6.x6a.strengths.r070.a9000.truth_pop.trainings}{50} \tabularnewline
\raggedright \resultnum{v4x.x6.x6b.cells.r050.a0300.a}{0.03}, \resultnum{v4x.x6.x6b.cells.r050.a0300.r}{0.5} & \raggedright procedural (D1) & \raggedright \resultnum{v4x.x6.x6b.cells.r050.a0300.truth_d1.value}{0.566297} & \raggedright \resultnum{v4x.x6.x6b.cells.r050.a0300.truth_d1.mcse}{0.00142674} & \raggedright \resultnum{v4x.x6.x6b.cells.r050.a0300.truth_d1.retry_fraction}{0} & \raggedright \resultnum{v4x.x6.x6b.cells.r050.a0300.truth_d1.successful}{100} \tabularnewline
\raggedright \resultnum{v4x.x6.x6b.cells.r050.a0300.a}{0.03}, \resultnum{v4x.x6.x6b.cells.r050.a0300.r}{0.5} & \raggedright population match & \raggedright \resultnum{v4x.x6.x6b.cells.r050.a0300.truth_pop.value}{0.568237} & \raggedright \resultnum{v4x.x6.x6b.cells.r050.a0300.truth_pop.mcse}{0.000800405} & \raggedright not applicable & \raggedright \resultnum{v4x.x6.x6b.cells.r050.a0300.truth_pop.trainings}{50} \tabularnewline
\raggedright \resultnum{v4x.x6.x6b.cells.r090.a0300.a}{0.03}, \resultnum{v4x.x6.x6b.cells.r090.a0300.r}{0.9} & \raggedright procedural (D1) & \raggedright \resultnum{v4x.x6.x6b.cells.r090.a0300.truth_d1.value}{0.64128} & \raggedright \resultnum{v4x.x6.x6b.cells.r090.a0300.truth_d1.mcse}{0.000942834} & \raggedright \resultnum{v4x.x6.x6b.cells.r090.a0300.truth_d1.retry_fraction}{0} & \raggedright \resultnum{v4x.x6.x6b.cells.r090.a0300.truth_d1.successful}{100} \tabularnewline
\raggedright \resultnum{v4x.x6.x6b.cells.r090.a0300.a}{0.03}, \resultnum{v4x.x6.x6b.cells.r090.a0300.r}{0.9} & \raggedright population match & \raggedright \resultnum{v4x.x6.x6b.cells.r090.a0300.truth_pop.value}{0.642349} & \raggedright \resultnum{v4x.x6.x6b.cells.r090.a0300.truth_pop.mcse}{0.000970984} & \raggedright not applicable & \raggedright \resultnum{v4x.x6.x6b.cells.r090.a0300.truth_pop.trainings}{50} \tabularnewline
\botrule
\end{tabular}
\end{table*}

% Generated by v4_exploratory_export; edited during the journal revision
% (agent S, 2026-10-13): captions, headers and categorical labels translated
% into plain English. Every number keeps its \resultnum binding.
\begin{table*}[htbp]
\setlength{\tabcolsep}{2pt}
\caption{Band coverage check (X6): conservative flags and truth-precision sensitivity, against the procedural truth. Conservative: Wilson lower end above 0.99 (no flag was raised). The declared sensitivity repeats each reading at the truth $\pm$2 Monte Carlo errors and marks a reading that changes as inconclusive by truth precision, keeping the original as the nominal reading; it applies only where the truth's Monte Carlo error exceeds the target of 0.002. No truth exceeded it (Table~\ref{tab:v4x-x6-truth}), so the sensitivity was not computed and its fields are empty in the released record; the robustness of the readings to the truth's estimation error was not assessed beyond that target.}\label{tab:v4x-x6-precision}
\centering\footnotesize
\begin{tabular}{@{}p{0.150\textwidth}p{0.250\textwidth}p{0.140\textwidth}p{0.220\textwidth}p{0.170\textwidth}@{}}
\toprule
\raggedright $a,r$ & \raggedright \mbox{Reference} & \raggedright \mbox{Conservative} & \raggedright \mbox{Inconclusive} \mbox{by} \mbox{truth} \mbox{precision} & \raggedright \mbox{Nominal} \mbox{reading} \tabularnewline
\midrule
\raggedright \resultnum{v4x.x6.x6a.strengths.r070.a0000.a}{0}, \resultnum{v4x.x6.x6a.strengths.r070.a0000.r}{0.7} & \raggedright pointwise & \raggedright no & \raggedright not computed & \raggedright not computed \tabularnewline
\raggedright \resultnum{v4x.x6.x6a.strengths.r070.a0000.a}{0}, \resultnum{v4x.x6.x6a.strengths.r070.a0000.r}{0.7} & \raggedright pointwise, gaps as misses & \raggedright no & \raggedright not computed & \raggedright not computed \tabularnewline
\raggedright \resultnum{v4x.x6.x6a.strengths.r070.a0000.a}{0}, \resultnum{v4x.x6.x6a.strengths.r070.a0000.r}{0.7} & \raggedright band at this strength & \raggedright no & \raggedright not computed & \raggedright not computed \tabularnewline
\raggedright \resultnum{v4x.x6.x6a.strengths.r070.a0010.a}{0.001}, \resultnum{v4x.x6.x6a.strengths.r070.a0010.r}{0.7} & \raggedright pointwise & \raggedright no & \raggedright not computed & \raggedright not computed \tabularnewline
\raggedright \resultnum{v4x.x6.x6a.strengths.r070.a0010.a}{0.001}, \resultnum{v4x.x6.x6a.strengths.r070.a0010.r}{0.7} & \raggedright pointwise, gaps as misses & \raggedright no & \raggedright not computed & \raggedright not computed \tabularnewline
\raggedright \resultnum{v4x.x6.x6a.strengths.r070.a0010.a}{0.001}, \resultnum{v4x.x6.x6a.strengths.r070.a0010.r}{0.7} & \raggedright band at this strength & \raggedright no & \raggedright not computed & \raggedright not computed \tabularnewline
\raggedright \resultnum{v4x.x6.x6a.strengths.r070.a0030.a}{0.003}, \resultnum{v4x.x6.x6a.strengths.r070.a0030.r}{0.7} & \raggedright pointwise & \raggedright no & \raggedright not computed & \raggedright not computed \tabularnewline
\raggedright \resultnum{v4x.x6.x6a.strengths.r070.a0030.a}{0.003}, \resultnum{v4x.x6.x6a.strengths.r070.a0030.r}{0.7} & \raggedright pointwise, gaps as misses & \raggedright no & \raggedright not computed & \raggedright not computed \tabularnewline
\raggedright \resultnum{v4x.x6.x6a.strengths.r070.a0030.a}{0.003}, \resultnum{v4x.x6.x6a.strengths.r070.a0030.r}{0.7} & \raggedright band at this strength & \raggedright no & \raggedright not computed & \raggedright not computed \tabularnewline
\raggedright \resultnum{v4x.x6.x6a.strengths.r070.a0100.a}{0.01}, \resultnum{v4x.x6.x6a.strengths.r070.a0100.r}{0.7} & \raggedright pointwise & \raggedright no & \raggedright not computed & \raggedright not computed \tabularnewline
\raggedright \resultnum{v4x.x6.x6a.strengths.r070.a0100.a}{0.01}, \resultnum{v4x.x6.x6a.strengths.r070.a0100.r}{0.7} & \raggedright pointwise, gaps as misses & \raggedright no & \raggedright not computed & \raggedright not computed \tabularnewline
\raggedright \resultnum{v4x.x6.x6a.strengths.r070.a0100.a}{0.01}, \resultnum{v4x.x6.x6a.strengths.r070.a0100.r}{0.7} & \raggedright band at this strength & \raggedright no & \raggedright not computed & \raggedright not computed \tabularnewline
\raggedright \resultnum{v4x.x6.x6a.strengths.r070.a0300.a}{0.03}, \resultnum{v4x.x6.x6a.strengths.r070.a0300.r}{0.7} & \raggedright pointwise & \raggedright no & \raggedright not computed & \raggedright not computed \tabularnewline
\raggedright \resultnum{v4x.x6.x6a.strengths.r070.a0300.a}{0.03}, \resultnum{v4x.x6.x6a.strengths.r070.a0300.r}{0.7} & \raggedright pointwise, gaps as misses & \raggedright no & \raggedright not computed & \raggedright not computed \tabularnewline
\raggedright \resultnum{v4x.x6.x6a.strengths.r070.a0300.a}{0.03}, \resultnum{v4x.x6.x6a.strengths.r070.a0300.r}{0.7} & \raggedright band at this strength & \raggedright no & \raggedright not computed & \raggedright not computed \tabularnewline
\raggedright \resultnum{v4x.x6.x6a.strengths.r070.a1000.a}{0.1}, \resultnum{v4x.x6.x6a.strengths.r070.a1000.r}{0.7} & \raggedright pointwise & \raggedright no & \raggedright not computed & \raggedright not computed \tabularnewline
\raggedright \resultnum{v4x.x6.x6a.strengths.r070.a1000.a}{0.1}, \resultnum{v4x.x6.x6a.strengths.r070.a1000.r}{0.7} & \raggedright pointwise, gaps as misses & \raggedright no & \raggedright not computed & \raggedright not computed \tabularnewline
\raggedright \resultnum{v4x.x6.x6a.strengths.r070.a1000.a}{0.1}, \resultnum{v4x.x6.x6a.strengths.r070.a1000.r}{0.7} & \raggedright band at this strength & \raggedright no & \raggedright not computed & \raggedright not computed \tabularnewline
\raggedright \resultnum{v4x.x6.x6a.strengths.r070.a3000.a}{0.3}, \resultnum{v4x.x6.x6a.strengths.r070.a3000.r}{0.7} & \raggedright pointwise & \raggedright no & \raggedright not computed & \raggedright not computed \tabularnewline
\raggedright \resultnum{v4x.x6.x6a.strengths.r070.a3000.a}{0.3}, \resultnum{v4x.x6.x6a.strengths.r070.a3000.r}{0.7} & \raggedright pointwise, gaps as misses & \raggedright no & \raggedright not computed & \raggedright not computed \tabularnewline
\raggedright \resultnum{v4x.x6.x6a.strengths.r070.a3000.a}{0.3}, \resultnum{v4x.x6.x6a.strengths.r070.a3000.r}{0.7} & \raggedright band at this strength & \raggedright no & \raggedright not computed & \raggedright not computed \tabularnewline
\raggedright \resultnum{v4x.x6.x6a.strengths.r070.a6000.a}{0.6}, \resultnum{v4x.x6.x6a.strengths.r070.a6000.r}{0.7} & \raggedright pointwise & \raggedright no & \raggedright not computed & \raggedright not computed \tabularnewline
\raggedright \resultnum{v4x.x6.x6a.strengths.r070.a6000.a}{0.6}, \resultnum{v4x.x6.x6a.strengths.r070.a6000.r}{0.7} & \raggedright pointwise, gaps as misses & \raggedright no & \raggedright not computed & \raggedright not computed \tabularnewline
\raggedright \resultnum{v4x.x6.x6a.strengths.r070.a6000.a}{0.6}, \resultnum{v4x.x6.x6a.strengths.r070.a6000.r}{0.7} & \raggedright band at this strength & \raggedright no & \raggedright not computed & \raggedright not computed \tabularnewline
\raggedright \resultnum{v4x.x6.x6a.strengths.r070.a9000.a}{0.9}, \resultnum{v4x.x6.x6a.strengths.r070.a9000.r}{0.7} & \raggedright pointwise & \raggedright no & \raggedright not computed & \raggedright not computed \tabularnewline
\raggedright \resultnum{v4x.x6.x6a.strengths.r070.a9000.a}{0.9}, \resultnum{v4x.x6.x6a.strengths.r070.a9000.r}{0.7} & \raggedright pointwise, gaps as misses & \raggedright no & \raggedright not computed & \raggedright not computed \tabularnewline
\raggedright \resultnum{v4x.x6.x6a.strengths.r070.a9000.a}{0.9}, \resultnum{v4x.x6.x6a.strengths.r070.a9000.r}{0.7} & \raggedright band at this strength & \raggedright no & \raggedright not computed & \raggedright not computed \tabularnewline
\raggedright \resultnum{v4x.x6.x6b.cells.r050.a0300.a}{0.03}, \resultnum{v4x.x6.x6b.cells.r050.a0300.r}{0.5} & \raggedright pointwise & \raggedright no & \raggedright not computed & \raggedright not computed \tabularnewline
\raggedright \resultnum{v4x.x6.x6b.cells.r050.a0300.a}{0.03}, \resultnum{v4x.x6.x6b.cells.r050.a0300.r}{0.5} & \raggedright pointwise, gaps as misses & \raggedright no & \raggedright not computed & \raggedright not computed \tabularnewline
\raggedright \resultnum{v4x.x6.x6b.cells.r090.a0300.a}{0.03}, \resultnum{v4x.x6.x6b.cells.r090.a0300.r}{0.9} & \raggedright pointwise & \raggedright no & \raggedright not computed & \raggedright not computed \tabularnewline
\raggedright \resultnum{v4x.x6.x6b.cells.r090.a0300.a}{0.03}, \resultnum{v4x.x6.x6b.cells.r090.a0300.r}{0.9} & \raggedright pointwise, gaps as misses & \raggedright no & \raggedright not computed & \raggedright not computed \tabularnewline
\botrule
\end{tabular}
\end{table*}

% Generated by v4_exploratory_export; edited during the journal revision
% (agent S, 2026-10-13): captions, headers and categorical labels translated
% into plain English. Every number keeps its \resultnum binding.
\begin{table*}[htbp]
\setlength{\tabcolsep}{2pt}
\caption{Band coverage check (X6): standardized errors $z=(\mathrm{logit}\,\hat p-\mathrm{logit}\,T)/\mathrm{SE}_{\mathrm{logit}}$ against the procedural truth $T$ and matching-bias diagnostics, by cell. The last column gives the Monte Carlo interval of the mean, the Monte Carlo error, the residual SD or the uncovered-bias flag (raised when the interval of the mean excludes zero and its absolute value exceeds 0.25; none means no flag). The mean shift of the local nuisance draws is the shift of the locally propagated draws (decision item D5); the fit residuals are in standard-error units. These diagnostics condition on success and do not establish coverage.}\label{tab:v4x-x6-diagnostics}
\centering\footnotesize
\begin{tabular}{@{}p{0.150\textwidth}p{0.400\textwidth}p{0.170\textwidth}p{0.220\textwidth}@{}}
\toprule
\raggedright $a,r$ & \raggedright \mbox{Diagnostic} & \raggedright \mbox{Estimate} & \raggedright \mbox{MC} \mbox{interval} \mbox{/} \mbox{error} \mbox{/} \mbox{SD} \mbox{/} \mbox{flag} \tabularnewline
\midrule
\raggedright \resultnum{v4x.x6.x6a.strengths.r070.a0000.a}{0}, \resultnum{v4x.x6.x6a.strengths.r070.a0000.r}{0.7} & \raggedright Mean standardized error & \raggedright $\resultnum{v4x.x6.x6a.strengths.r070.a0000.vs_d1.standardized_error.mean}{-0.0179277}$ & \raggedright [$\resultnum{v4x.x6.x6a.strengths.r070.a0000.vs_d1.standardized_error.mean_interval95.i000}{-0.167402}$, \resultnum{v4x.x6.x6a.strengths.r070.a0000.vs_d1.standardized_error.mean_interval95.i001}{0.131547}] \tabularnewline
\raggedright \resultnum{v4x.x6.x6a.strengths.r070.a0000.a}{0}, \resultnum{v4x.x6.x6a.strengths.r070.a0000.r}{0.7} & \raggedright SD of standardized errors / uncovered-bias flag & \raggedright \resultnum{v4x.x6.x6a.strengths.r070.a0000.vs_d1.standardized_error.sd}{0.982592} & \raggedright none \tabularnewline
\raggedright \resultnum{v4x.x6.x6a.strengths.r070.a0000.a}{0}, \resultnum{v4x.x6.x6a.strengths.r070.a0000.r}{0.7} & \raggedright Share of standardized errors below $-1.96$ & \raggedright \resultnum{v4x.x6.x6a.strengths.r070.a0000.vs_d1.standardized_error.below_minus_1_96}{0.0240964} & \raggedright not applicable \tabularnewline
\raggedright \resultnum{v4x.x6.x6a.strengths.r070.a0000.a}{0}, \resultnum{v4x.x6.x6a.strengths.r070.a0000.r}{0.7} & \raggedright Share of standardized errors above 1.96 & \raggedright \resultnum{v4x.x6.x6a.strengths.r070.a0000.vs_d1.standardized_error.above_1_96}{0.0301205} & \raggedright not applicable \tabularnewline
\raggedright \resultnum{v4x.x6.x6a.strengths.r070.a0000.a}{0}, \resultnum{v4x.x6.x6a.strengths.r070.a0000.r}{0.7} & \raggedright Mean shift of the local nuisance draws (D5) & \raggedright \resultnum{v4x.x6.x6a.strengths.r070.a0000.matching_bias.d5_mean_shift.mean}{0} & \raggedright not applicable \tabularnewline
\raggedright \resultnum{v4x.x6.x6a.strengths.r070.a0000.a}{0}, \resultnum{v4x.x6.x6a.strengths.r070.a0000.r}{0.7} & \raggedright Empirical matching bias & \raggedright $\resultnum{v4x.x6.x6a.strengths.r070.a0000.matching_bias.empirical_bias.mean}{-0.000108141}$ & \raggedright \resultnum{v4x.x6.x6a.strengths.r070.a0000.matching_bias.empirical_bias.mcse}{0.00049064} \tabularnewline
\raggedright \resultnum{v4x.x6.x6a.strengths.r070.a0000.a}{0}, \resultnum{v4x.x6.x6a.strengths.r070.a0000.r}{0.7} & \raggedright Procedural minus population-match truth & \raggedright \resultnum{v4x.x6.x6a.strengths.r070.a0000.matching_bias.truth_d1_minus_pop.value}{0} & \raggedright \resultnum{v4x.x6.x6a.strengths.r070.a0000.matching_bias.truth_d1_minus_pop.se}{0} \tabularnewline
\raggedright \resultnum{v4x.x6.x6a.strengths.r070.a0000.a}{0}, \resultnum{v4x.x6.x6a.strengths.r070.a0000.r}{0.7} & \raggedright Mean amplification & \raggedright \resultnum{v4x.x6.x6a.strengths.r070.a0000.matching_bias.amplification.mean}{0} & \raggedright not applicable \tabularnewline
\raggedright \resultnum{v4x.x6.x6a.strengths.r070.a0000.a}{0}, \resultnum{v4x.x6.x6a.strengths.r070.a0000.r}{0.7} & \raggedright MSE / mean squared SE & \raggedright \resultnum{v4x.x6.x6a.strengths.r070.a0000.matching_bias.mse_to_mean_se2}{0.951575} & \raggedright not applicable \tabularnewline
\raggedright \resultnum{v4x.x6.x6a.strengths.r070.a0000.a}{0}, \resultnum{v4x.x6.x6a.strengths.r070.a0000.r}{0.7} & \raggedright Lag-1 fit residual, mean / SD & \raggedright \resultnum{v4x.x6.x6a.strengths.r070.a0000.matching_bias.fit_residual_se_units.mean.i000}{0} & \raggedright \resultnum{v4x.x6.x6a.strengths.r070.a0000.matching_bias.fit_residual_se_units.sd.i000}{0} \tabularnewline
\raggedright \resultnum{v4x.x6.x6a.strengths.r070.a0000.a}{0}, \resultnum{v4x.x6.x6a.strengths.r070.a0000.r}{0.7} & \raggedright Lag-2 fit residual, mean / SD & \raggedright \resultnum{v4x.x6.x6a.strengths.r070.a0000.matching_bias.fit_residual_se_units.mean.i001}{0} & \raggedright \resultnum{v4x.x6.x6a.strengths.r070.a0000.matching_bias.fit_residual_se_units.sd.i001}{0} \tabularnewline
\raggedright \resultnum{v4x.x6.x6a.strengths.r070.a0010.a}{0.001}, \resultnum{v4x.x6.x6a.strengths.r070.a0010.r}{0.7} & \raggedright Mean standardized error & \raggedright \resultnum{v4x.x6.x6a.strengths.r070.a0010.vs_d1.standardized_error.mean}{0.0922797} & \raggedright [$\resultnum{v4x.x6.x6a.strengths.r070.a0010.vs_d1.standardized_error.mean_interval95.i000}{-0.0488752}$, \resultnum{v4x.x6.x6a.strengths.r070.a0010.vs_d1.standardized_error.mean_interval95.i001}{0.233435}] \tabularnewline
\raggedright \resultnum{v4x.x6.x6a.strengths.r070.a0010.a}{0.001}, \resultnum{v4x.x6.x6a.strengths.r070.a0010.r}{0.7} & \raggedright SD of standardized errors / uncovered-bias flag & \raggedright \resultnum{v4x.x6.x6a.strengths.r070.a0010.vs_d1.standardized_error.sd}{0.927902} & \raggedright none \tabularnewline
\raggedright \resultnum{v4x.x6.x6a.strengths.r070.a0010.a}{0.001}, \resultnum{v4x.x6.x6a.strengths.r070.a0010.r}{0.7} & \raggedright Share of standardized errors below $-1.96$ & \raggedright \resultnum{v4x.x6.x6a.strengths.r070.a0010.vs_d1.standardized_error.below_minus_1_96}{0.0120482} & \raggedright not applicable \tabularnewline
\raggedright \resultnum{v4x.x6.x6a.strengths.r070.a0010.a}{0.001}, \resultnum{v4x.x6.x6a.strengths.r070.a0010.r}{0.7} & \raggedright Share of standardized errors above 1.96 & \raggedright \resultnum{v4x.x6.x6a.strengths.r070.a0010.vs_d1.standardized_error.above_1_96}{0.0361446} & \raggedright not applicable \tabularnewline
\raggedright \resultnum{v4x.x6.x6a.strengths.r070.a0010.a}{0.001}, \resultnum{v4x.x6.x6a.strengths.r070.a0010.r}{0.7} & \raggedright Mean shift of the local nuisance draws (D5) & \raggedright $\resultnum{v4x.x6.x6a.strengths.r070.a0010.matching_bias.d5_mean_shift.mean}{-0.000219701}$ & \raggedright not applicable \tabularnewline
\raggedright \resultnum{v4x.x6.x6a.strengths.r070.a0010.a}{0.001}, \resultnum{v4x.x6.x6a.strengths.r070.a0010.r}{0.7} & \raggedright Empirical matching bias & \raggedright \resultnum{v4x.x6.x6a.strengths.r070.a0010.matching_bias.empirical_bias.mean}{0.00080174} & \raggedright \resultnum{v4x.x6.x6a.strengths.r070.a0010.matching_bias.empirical_bias.mcse}{0.0012593} \tabularnewline
\raggedright \resultnum{v4x.x6.x6a.strengths.r070.a0010.a}{0.001}, \resultnum{v4x.x6.x6a.strengths.r070.a0010.r}{0.7} & \raggedright Procedural minus population-match truth & \raggedright $\resultnum{v4x.x6.x6a.strengths.r070.a0010.matching_bias.truth_d1_minus_pop.value}{-0.000926452}$ & \raggedright \resultnum{v4x.x6.x6a.strengths.r070.a0010.matching_bias.truth_d1_minus_pop.se}{0.0015151} \tabularnewline
\raggedright \resultnum{v4x.x6.x6a.strengths.r070.a0010.a}{0.001}, \resultnum{v4x.x6.x6a.strengths.r070.a0010.r}{0.7} & \raggedright Mean amplification & \raggedright \resultnum{v4x.x6.x6a.strengths.r070.a0010.matching_bias.amplification.mean}{0.0227978} & \raggedright not applicable \tabularnewline
\raggedright \resultnum{v4x.x6.x6a.strengths.r070.a0010.a}{0.001}, \resultnum{v4x.x6.x6a.strengths.r070.a0010.r}{0.7} & \raggedright MSE / mean squared SE & \raggedright \resultnum{v4x.x6.x6a.strengths.r070.a0010.matching_bias.mse_to_mean_se2}{0.778922} & \raggedright not applicable \tabularnewline
\raggedright \resultnum{v4x.x6.x6a.strengths.r070.a0010.a}{0.001}, \resultnum{v4x.x6.x6a.strengths.r070.a0010.r}{0.7} & \raggedright Lag-1 fit residual, mean / SD & \raggedright $\resultnum{v4x.x6.x6a.strengths.r070.a0010.matching_bias.fit_residual_se_units.mean.i000}{-0.00983144}$ & \raggedright \resultnum{v4x.x6.x6a.strengths.r070.a0010.matching_bias.fit_residual_se_units.sd.i000}{0.0767452} \tabularnewline
\raggedright \resultnum{v4x.x6.x6a.strengths.r070.a0010.a}{0.001}, \resultnum{v4x.x6.x6a.strengths.r070.a0010.r}{0.7} & \raggedright Lag-2 fit residual, mean / SD & \raggedright \resultnum{v4x.x6.x6a.strengths.r070.a0010.matching_bias.fit_residual_se_units.mean.i001}{0.000990088} & \raggedright \resultnum{v4x.x6.x6a.strengths.r070.a0010.matching_bias.fit_residual_se_units.sd.i001}{0.145399} \tabularnewline
\raggedright \resultnum{v4x.x6.x6a.strengths.r070.a0030.a}{0.003}, \resultnum{v4x.x6.x6a.strengths.r070.a0030.r}{0.7} & \raggedright Mean standardized error & \raggedright $\resultnum{v4x.x6.x6a.strengths.r070.a0030.vs_d1.standardized_error.mean}{-0.0133531}$ & \raggedright [$\resultnum{v4x.x6.x6a.strengths.r070.a0030.vs_d1.standardized_error.mean_interval95.i000}{-0.144175}$, \resultnum{v4x.x6.x6a.strengths.r070.a0030.vs_d1.standardized_error.mean_interval95.i001}{0.117469}] \tabularnewline
\raggedright \resultnum{v4x.x6.x6a.strengths.r070.a0030.a}{0.003}, \resultnum{v4x.x6.x6a.strengths.r070.a0030.r}{0.7} & \raggedright SD of standardized errors / uncovered-bias flag & \raggedright \resultnum{v4x.x6.x6a.strengths.r070.a0030.vs_d1.standardized_error.sd}{0.859976} & \raggedright none \tabularnewline
\raggedright \resultnum{v4x.x6.x6a.strengths.r070.a0030.a}{0.003}, \resultnum{v4x.x6.x6a.strengths.r070.a0030.r}{0.7} & \raggedright Share of standardized errors below $-1.96$ & \raggedright \resultnum{v4x.x6.x6a.strengths.r070.a0030.vs_d1.standardized_error.below_minus_1_96}{0.0240964} & \raggedright not applicable \tabularnewline
\raggedright \resultnum{v4x.x6.x6a.strengths.r070.a0030.a}{0.003}, \resultnum{v4x.x6.x6a.strengths.r070.a0030.r}{0.7} & \raggedright Share of standardized errors above 1.96 & \raggedright \resultnum{v4x.x6.x6a.strengths.r070.a0030.vs_d1.standardized_error.above_1_96}{0.0120482} & \raggedright not applicable \tabularnewline
\raggedright \resultnum{v4x.x6.x6a.strengths.r070.a0030.a}{0.003}, \resultnum{v4x.x6.x6a.strengths.r070.a0030.r}{0.7} & \raggedright Mean shift of the local nuisance draws (D5) & \raggedright $\resultnum{v4x.x6.x6a.strengths.r070.a0030.matching_bias.d5_mean_shift.mean}{-0.00186318}$ & \raggedright not applicable \tabularnewline
\raggedright \resultnum{v4x.x6.x6a.strengths.r070.a0030.a}{0.003}, \resultnum{v4x.x6.x6a.strengths.r070.a0030.r}{0.7} & \raggedright Empirical matching bias & \raggedright $\resultnum{v4x.x6.x6a.strengths.r070.a0030.matching_bias.empirical_bias.mean}{-0.000105507}$ & \raggedright \resultnum{v4x.x6.x6a.strengths.r070.a0030.matching_bias.empirical_bias.mcse}{0.00129004} \tabularnewline
\raggedright \resultnum{v4x.x6.x6a.strengths.r070.a0030.a}{0.003}, \resultnum{v4x.x6.x6a.strengths.r070.a0030.r}{0.7} & \raggedright Procedural minus population-match truth & \raggedright $\resultnum{v4x.x6.x6a.strengths.r070.a0030.matching_bias.truth_d1_minus_pop.value}{-0.0020458}$ & \raggedright \resultnum{v4x.x6.x6a.strengths.r070.a0030.matching_bias.truth_d1_minus_pop.se}{0.00147279} \tabularnewline
\raggedright \resultnum{v4x.x6.x6a.strengths.r070.a0030.a}{0.003}, \resultnum{v4x.x6.x6a.strengths.r070.a0030.r}{0.7} & \raggedright Mean amplification & \raggedright \resultnum{v4x.x6.x6a.strengths.r070.a0030.matching_bias.amplification.mean}{0.0256389} & \raggedright not applicable \tabularnewline
\raggedright \resultnum{v4x.x6.x6a.strengths.r070.a0030.a}{0.003}, \resultnum{v4x.x6.x6a.strengths.r070.a0030.r}{0.7} & \raggedright MSE / mean squared SE & \raggedright \resultnum{v4x.x6.x6a.strengths.r070.a0030.matching_bias.mse_to_mean_se2}{0.608646} & \raggedright not applicable \tabularnewline
\raggedright \resultnum{v4x.x6.x6a.strengths.r070.a0030.a}{0.003}, \resultnum{v4x.x6.x6a.strengths.r070.a0030.r}{0.7} & \raggedright Lag-1 fit residual, mean / SD & \raggedright $\resultnum{v4x.x6.x6a.strengths.r070.a0030.matching_bias.fit_residual_se_units.mean.i000}{-0.00539802}$ & \raggedright \resultnum{v4x.x6.x6a.strengths.r070.a0030.matching_bias.fit_residual_se_units.sd.i000}{0.0721043} \tabularnewline
\raggedright \resultnum{v4x.x6.x6a.strengths.r070.a0030.a}{0.003}, \resultnum{v4x.x6.x6a.strengths.r070.a0030.r}{0.7} & \raggedright Lag-2 fit residual, mean / SD & \raggedright \resultnum{v4x.x6.x6a.strengths.r070.a0030.matching_bias.fit_residual_se_units.mean.i001}{0.0145717} & \raggedright \resultnum{v4x.x6.x6a.strengths.r070.a0030.matching_bias.fit_residual_se_units.sd.i001}{0.116888} \tabularnewline
\raggedright \resultnum{v4x.x6.x6a.strengths.r070.a0100.a}{0.01}, \resultnum{v4x.x6.x6a.strengths.r070.a0100.r}{0.7} & \raggedright Mean standardized error & \raggedright $\resultnum{v4x.x6.x6a.strengths.r070.a0100.vs_d1.standardized_error.mean}{-0.159141}$ & \raggedright [$\resultnum{v4x.x6.x6a.strengths.r070.a0100.vs_d1.standardized_error.mean_interval95.i000}{-0.301163}$, $\resultnum{v4x.x6.x6a.strengths.r070.a0100.vs_d1.standardized_error.mean_interval95.i001}{-0.0171188}$] \tabularnewline
\raggedright \resultnum{v4x.x6.x6a.strengths.r070.a0100.a}{0.01}, \resultnum{v4x.x6.x6a.strengths.r070.a0100.r}{0.7} & \raggedright SD of standardized errors / uncovered-bias flag & \raggedright \resultnum{v4x.x6.x6a.strengths.r070.a0100.vs_d1.standardized_error.sd}{0.933603} & \raggedright none \tabularnewline
\raggedright \resultnum{v4x.x6.x6a.strengths.r070.a0100.a}{0.01}, \resultnum{v4x.x6.x6a.strengths.r070.a0100.r}{0.7} & \raggedright Share of standardized errors below $-1.96$ & \raggedright \resultnum{v4x.x6.x6a.strengths.r070.a0100.vs_d1.standardized_error.below_minus_1_96}{0.0180723} & \raggedright not applicable \tabularnewline
\raggedright \resultnum{v4x.x6.x6a.strengths.r070.a0100.a}{0.01}, \resultnum{v4x.x6.x6a.strengths.r070.a0100.r}{0.7} & \raggedright Share of standardized errors above 1.96 & \raggedright \resultnum{v4x.x6.x6a.strengths.r070.a0100.vs_d1.standardized_error.above_1_96}{0.0060241} & \raggedright not applicable \tabularnewline
\raggedright \resultnum{v4x.x6.x6a.strengths.r070.a0100.a}{0.01}, \resultnum{v4x.x6.x6a.strengths.r070.a0100.r}{0.7} & \raggedright Mean shift of the local nuisance draws (D5) & \raggedright $\resultnum{v4x.x6.x6a.strengths.r070.a0100.matching_bias.d5_mean_shift.mean}{-0.00104476}$ & \raggedright not applicable \tabularnewline
\raggedright \resultnum{v4x.x6.x6a.strengths.r070.a0100.a}{0.01}, \resultnum{v4x.x6.x6a.strengths.r070.a0100.r}{0.7} & \raggedright Empirical matching bias & \raggedright $\resultnum{v4x.x6.x6a.strengths.r070.a0100.matching_bias.empirical_bias.mean}{-0.00206635}$ & \raggedright \resultnum{v4x.x6.x6a.strengths.r070.a0100.matching_bias.empirical_bias.mcse}{0.00129037} \tabularnewline
\raggedright \resultnum{v4x.x6.x6a.strengths.r070.a0100.a}{0.01}, \resultnum{v4x.x6.x6a.strengths.r070.a0100.r}{0.7} & \raggedright Procedural minus population-match truth & \raggedright \resultnum{v4x.x6.x6a.strengths.r070.a0100.matching_bias.truth_d1_minus_pop.value}{0.0074669} & \raggedright \resultnum{v4x.x6.x6a.strengths.r070.a0100.matching_bias.truth_d1_minus_pop.se}{0.00137225} \tabularnewline
\raggedright \resultnum{v4x.x6.x6a.strengths.r070.a0100.a}{0.01}, \resultnum{v4x.x6.x6a.strengths.r070.a0100.r}{0.7} & \raggedright Mean amplification & \raggedright \resultnum{v4x.x6.x6a.strengths.r070.a0100.matching_bias.amplification.mean}{0.0316512} & \raggedright not applicable \tabularnewline
\raggedright \resultnum{v4x.x6.x6a.strengths.r070.a0100.a}{0.01}, \resultnum{v4x.x6.x6a.strengths.r070.a0100.r}{0.7} & \raggedright MSE / mean squared SE & \raggedright \resultnum{v4x.x6.x6a.strengths.r070.a0100.matching_bias.mse_to_mean_se2}{0.80575} & \raggedright not applicable \tabularnewline
\raggedright \resultnum{v4x.x6.x6a.strengths.r070.a0100.a}{0.01}, \resultnum{v4x.x6.x6a.strengths.r070.a0100.r}{0.7} & \raggedright Lag-1 fit residual, mean / SD & \raggedright $\resultnum{v4x.x6.x6a.strengths.r070.a0100.matching_bias.fit_residual_se_units.mean.i000}{-0.000976358}$ & \raggedright \resultnum{v4x.x6.x6a.strengths.r070.a0100.matching_bias.fit_residual_se_units.sd.i000}{0.0573138} \tabularnewline
\raggedright \resultnum{v4x.x6.x6a.strengths.r070.a0100.a}{0.01}, \resultnum{v4x.x6.x6a.strengths.r070.a0100.r}{0.7} & \raggedright Lag-2 fit residual, mean / SD & \raggedright $\resultnum{v4x.x6.x6a.strengths.r070.a0100.matching_bias.fit_residual_se_units.mean.i001}{-0.00911525}$ & \raggedright \resultnum{v4x.x6.x6a.strengths.r070.a0100.matching_bias.fit_residual_se_units.sd.i001}{0.135314} \tabularnewline
\raggedright \resultnum{v4x.x6.x6a.strengths.r070.a0300.a}{0.03}, \resultnum{v4x.x6.x6a.strengths.r070.a0300.r}{0.7} & \raggedright Mean standardized error & \raggedright $\resultnum{v4x.x6.x6a.strengths.r070.a0300.vs_d1.standardized_error.mean}{-0.149444}$ & \raggedright [$\resultnum{v4x.x6.x6a.strengths.r070.a0300.vs_d1.standardized_error.mean_interval95.i000}{-0.291794}$, $\resultnum{v4x.x6.x6a.strengths.r070.a0300.vs_d1.standardized_error.mean_interval95.i001}{-0.00709367}$] \tabularnewline
\raggedright \resultnum{v4x.x6.x6a.strengths.r070.a0300.a}{0.03}, \resultnum{v4x.x6.x6a.strengths.r070.a0300.r}{0.7} & \raggedright SD of standardized errors / uncovered-bias flag & \raggedright \resultnum{v4x.x6.x6a.strengths.r070.a0300.vs_d1.standardized_error.sd}{0.935758} & \raggedright none \tabularnewline
\raggedright \resultnum{v4x.x6.x6a.strengths.r070.a0300.a}{0.03}, \resultnum{v4x.x6.x6a.strengths.r070.a0300.r}{0.7} & \raggedright Share of standardized errors below $-1.96$ & \raggedright \resultnum{v4x.x6.x6a.strengths.r070.a0300.vs_d1.standardized_error.below_minus_1_96}{0.0240964} & \raggedright not applicable \tabularnewline
\raggedright \resultnum{v4x.x6.x6a.strengths.r070.a0300.a}{0.03}, \resultnum{v4x.x6.x6a.strengths.r070.a0300.r}{0.7} & \raggedright Share of standardized errors above 1.96 & \raggedright \resultnum{v4x.x6.x6a.strengths.r070.a0300.vs_d1.standardized_error.above_1_96}{0.0180723} & \raggedright not applicable \tabularnewline
\botrule
\end{tabular}
\end{table*}
\begin{table*}[htbp]
\setlength{\tabcolsep}{2pt}
\addtocounter{table}{-\csname @ne\endcsname}
\providecommand{\theHtable}{\thetable}
\renewcommand{\theHtable}{v4x-x6-diagnostics-part1}
\caption{Band coverage check (X6): diagnostics (continued, part 2 of 3)}\label{tab:v4x-x6-diagnostics-part1}
\centering\footnotesize
\begin{tabular}{@{}p{0.150\textwidth}p{0.400\textwidth}p{0.170\textwidth}p{0.220\textwidth}@{}}
\toprule
\raggedright $a,r$ & \raggedright \mbox{Diagnostic} & \raggedright \mbox{Estimate} & \raggedright \mbox{MC} \mbox{interval} \mbox{/} \mbox{error} \mbox{/} \mbox{SD} \mbox{/} \mbox{flag} \tabularnewline
\midrule
\raggedright \resultnum{v4x.x6.x6a.strengths.r070.a0300.a}{0.03}, \resultnum{v4x.x6.x6a.strengths.r070.a0300.r}{0.7} & \raggedright Mean shift of the local nuisance draws (D5) & \raggedright $\resultnum{v4x.x6.x6a.strengths.r070.a0300.matching_bias.d5_mean_shift.mean}{-0.000756845}$ & \raggedright not applicable \tabularnewline
\raggedright \resultnum{v4x.x6.x6a.strengths.r070.a0300.a}{0.03}, \resultnum{v4x.x6.x6a.strengths.r070.a0300.r}{0.7} & \raggedright Empirical matching bias & \raggedright $\resultnum{v4x.x6.x6a.strengths.r070.a0300.matching_bias.empirical_bias.mean}{-0.00168515}$ & \raggedright \resultnum{v4x.x6.x6a.strengths.r070.a0300.matching_bias.empirical_bias.mcse}{0.00134042} \tabularnewline
\raggedright \resultnum{v4x.x6.x6a.strengths.r070.a0300.a}{0.03}, \resultnum{v4x.x6.x6a.strengths.r070.a0300.r}{0.7} & \raggedright Procedural minus population-match truth & \raggedright $\resultnum{v4x.x6.x6a.strengths.r070.a0300.matching_bias.truth_d1_minus_pop.value}{-0.00233842}$ & \raggedright \resultnum{v4x.x6.x6a.strengths.r070.a0300.matching_bias.truth_d1_minus_pop.se}{0.00139158} \tabularnewline
\raggedright \resultnum{v4x.x6.x6a.strengths.r070.a0300.a}{0.03}, \resultnum{v4x.x6.x6a.strengths.r070.a0300.r}{0.7} & \raggedright Mean amplification & \raggedright \resultnum{v4x.x6.x6a.strengths.r070.a0300.matching_bias.amplification.mean}{0.0375584} & \raggedright not applicable \tabularnewline
\raggedright \resultnum{v4x.x6.x6a.strengths.r070.a0300.a}{0.03}, \resultnum{v4x.x6.x6a.strengths.r070.a0300.r}{0.7} & \raggedright MSE / mean squared SE & \raggedright \resultnum{v4x.x6.x6a.strengths.r070.a0300.matching_bias.mse_to_mean_se2}{0.870295} & \raggedright not applicable \tabularnewline
\raggedright \resultnum{v4x.x6.x6a.strengths.r070.a0300.a}{0.03}, \resultnum{v4x.x6.x6a.strengths.r070.a0300.r}{0.7} & \raggedright Lag-1 fit residual, mean / SD & \raggedright $\resultnum{v4x.x6.x6a.strengths.r070.a0300.matching_bias.fit_residual_se_units.mean.i000}{-0.00596864}$ & \raggedright \resultnum{v4x.x6.x6a.strengths.r070.a0300.matching_bias.fit_residual_se_units.sd.i000}{0.0712458} \tabularnewline
\raggedright \resultnum{v4x.x6.x6a.strengths.r070.a0300.a}{0.03}, \resultnum{v4x.x6.x6a.strengths.r070.a0300.r}{0.7} & \raggedright Lag-2 fit residual, mean / SD & \raggedright \resultnum{v4x.x6.x6a.strengths.r070.a0300.matching_bias.fit_residual_se_units.mean.i001}{0.016178} & \raggedright \resultnum{v4x.x6.x6a.strengths.r070.a0300.matching_bias.fit_residual_se_units.sd.i001}{0.135724} \tabularnewline
\raggedright \resultnum{v4x.x6.x6a.strengths.r070.a1000.a}{0.1}, \resultnum{v4x.x6.x6a.strengths.r070.a1000.r}{0.7} & \raggedright Mean standardized error & \raggedright \resultnum{v4x.x6.x6a.strengths.r070.a1000.vs_d1.standardized_error.mean}{0.0154353} & \raggedright [$\resultnum{v4x.x6.x6a.strengths.r070.a1000.vs_d1.standardized_error.mean_interval95.i000}{-0.13706}$, \resultnum{v4x.x6.x6a.strengths.r070.a1000.vs_d1.standardized_error.mean_interval95.i001}{0.167931}] \tabularnewline
\raggedright \resultnum{v4x.x6.x6a.strengths.r070.a1000.a}{0.1}, \resultnum{v4x.x6.x6a.strengths.r070.a1000.r}{0.7} & \raggedright SD of standardized errors / uncovered-bias flag & \raggedright \resultnum{v4x.x6.x6a.strengths.r070.a1000.vs_d1.standardized_error.sd}{1.00245} & \raggedright none \tabularnewline
\raggedright \resultnum{v4x.x6.x6a.strengths.r070.a1000.a}{0.1}, \resultnum{v4x.x6.x6a.strengths.r070.a1000.r}{0.7} & \raggedright Share of standardized errors below $-1.96$ & \raggedright \resultnum{v4x.x6.x6a.strengths.r070.a1000.vs_d1.standardized_error.below_minus_1_96}{0.0120482} & \raggedright not applicable \tabularnewline
\raggedright \resultnum{v4x.x6.x6a.strengths.r070.a1000.a}{0.1}, \resultnum{v4x.x6.x6a.strengths.r070.a1000.r}{0.7} & \raggedright Share of standardized errors above 1.96 & \raggedright \resultnum{v4x.x6.x6a.strengths.r070.a1000.vs_d1.standardized_error.above_1_96}{0.0180723} & \raggedright not applicable \tabularnewline
\raggedright \resultnum{v4x.x6.x6a.strengths.r070.a1000.a}{0.1}, \resultnum{v4x.x6.x6a.strengths.r070.a1000.r}{0.7} & \raggedright Mean shift of the local nuisance draws (D5) & \raggedright \resultnum{v4x.x6.x6a.strengths.r070.a1000.matching_bias.d5_mean_shift.mean}{0.000532221} & \raggedright not applicable \tabularnewline
\raggedright \resultnum{v4x.x6.x6a.strengths.r070.a1000.a}{0.1}, \resultnum{v4x.x6.x6a.strengths.r070.a1000.r}{0.7} & \raggedright Empirical matching bias & \raggedright \resultnum{v4x.x6.x6a.strengths.r070.a1000.matching_bias.empirical_bias.mean}{0.00010501} & \raggedright \resultnum{v4x.x6.x6a.strengths.r070.a1000.matching_bias.empirical_bias.mcse}{0.00147809} \tabularnewline
\raggedright \resultnum{v4x.x6.x6a.strengths.r070.a1000.a}{0.1}, \resultnum{v4x.x6.x6a.strengths.r070.a1000.r}{0.7} & \raggedright Procedural minus population-match truth & \raggedright \resultnum{v4x.x6.x6a.strengths.r070.a1000.matching_bias.truth_d1_minus_pop.value}{0.00342957} & \raggedright \resultnum{v4x.x6.x6a.strengths.r070.a1000.matching_bias.truth_d1_minus_pop.se}{0.00145359} \tabularnewline
\raggedright \resultnum{v4x.x6.x6a.strengths.r070.a1000.a}{0.1}, \resultnum{v4x.x6.x6a.strengths.r070.a1000.r}{0.7} & \raggedright Mean amplification & \raggedright \resultnum{v4x.x6.x6a.strengths.r070.a1000.matching_bias.amplification.mean}{0.0562927} & \raggedright not applicable \tabularnewline
\raggedright \resultnum{v4x.x6.x6a.strengths.r070.a1000.a}{0.1}, \resultnum{v4x.x6.x6a.strengths.r070.a1000.r}{0.7} & \raggedright MSE / mean squared SE & \raggedright \resultnum{v4x.x6.x6a.strengths.r070.a1000.matching_bias.mse_to_mean_se2}{0.944881} & \raggedright not applicable \tabularnewline
\raggedright \resultnum{v4x.x6.x6a.strengths.r070.a1000.a}{0.1}, \resultnum{v4x.x6.x6a.strengths.r070.a1000.r}{0.7} & \raggedright Lag-1 fit residual, mean / SD & \raggedright \resultnum{v4x.x6.x6a.strengths.r070.a1000.matching_bias.fit_residual_se_units.mean.i000}{0.0299889} & \raggedright \resultnum{v4x.x6.x6a.strengths.r070.a1000.matching_bias.fit_residual_se_units.sd.i000}{0.092809} \tabularnewline
\raggedright \resultnum{v4x.x6.x6a.strengths.r070.a1000.a}{0.1}, \resultnum{v4x.x6.x6a.strengths.r070.a1000.r}{0.7} & \raggedright Lag-2 fit residual, mean / SD & \raggedright $\resultnum{v4x.x6.x6a.strengths.r070.a1000.matching_bias.fit_residual_se_units.mean.i001}{-0.0102286}$ & \raggedright \resultnum{v4x.x6.x6a.strengths.r070.a1000.matching_bias.fit_residual_se_units.sd.i001}{0.133393} \tabularnewline
\raggedright \resultnum{v4x.x6.x6a.strengths.r070.a3000.a}{0.3}, \resultnum{v4x.x6.x6a.strengths.r070.a3000.r}{0.7} & \raggedright Mean standardized error & \raggedright $\resultnum{v4x.x6.x6a.strengths.r070.a3000.vs_d1.standardized_error.mean}{-0.130574}$ & \raggedright [$\resultnum{v4x.x6.x6a.strengths.r070.a3000.vs_d1.standardized_error.mean_interval95.i000}{-0.282385}$, \resultnum{v4x.x6.x6a.strengths.r070.a3000.vs_d1.standardized_error.mean_interval95.i001}{0.0212364}] \tabularnewline
\raggedright \resultnum{v4x.x6.x6a.strengths.r070.a3000.a}{0.3}, \resultnum{v4x.x6.x6a.strengths.r070.a3000.r}{0.7} & \raggedright SD of standardized errors / uncovered-bias flag & \raggedright \resultnum{v4x.x6.x6a.strengths.r070.a3000.vs_d1.standardized_error.sd}{0.997948} & \raggedright none \tabularnewline
\raggedright \resultnum{v4x.x6.x6a.strengths.r070.a3000.a}{0.3}, \resultnum{v4x.x6.x6a.strengths.r070.a3000.r}{0.7} & \raggedright Share of standardized errors below $-1.96$ & \raggedright \resultnum{v4x.x6.x6a.strengths.r070.a3000.vs_d1.standardized_error.below_minus_1_96}{0.0361446} & \raggedright not applicable \tabularnewline
\raggedright \resultnum{v4x.x6.x6a.strengths.r070.a3000.a}{0.3}, \resultnum{v4x.x6.x6a.strengths.r070.a3000.r}{0.7} & \raggedright Share of standardized errors above 1.96 & \raggedright \resultnum{v4x.x6.x6a.strengths.r070.a3000.vs_d1.standardized_error.above_1_96}{0.0060241} & \raggedright not applicable \tabularnewline
\raggedright \resultnum{v4x.x6.x6a.strengths.r070.a3000.a}{0.3}, \resultnum{v4x.x6.x6a.strengths.r070.a3000.r}{0.7} & \raggedright Mean shift of the local nuisance draws (D5) & \raggedright \resultnum{v4x.x6.x6a.strengths.r070.a3000.matching_bias.d5_mean_shift.mean}{-7.55195e-05} & \raggedright not applicable \tabularnewline
\raggedright \resultnum{v4x.x6.x6a.strengths.r070.a3000.a}{0.3}, \resultnum{v4x.x6.x6a.strengths.r070.a3000.r}{0.7} & \raggedright Empirical matching bias & \raggedright $\resultnum{v4x.x6.x6a.strengths.r070.a3000.matching_bias.empirical_bias.mean}{-0.000632768}$ & \raggedright \resultnum{v4x.x6.x6a.strengths.r070.a3000.matching_bias.empirical_bias.mcse}{0.000612858} \tabularnewline
\raggedright \resultnum{v4x.x6.x6a.strengths.r070.a3000.a}{0.3}, \resultnum{v4x.x6.x6a.strengths.r070.a3000.r}{0.7} & \raggedright Procedural minus population-match truth & \raggedright $\resultnum{v4x.x6.x6a.strengths.r070.a3000.matching_bias.truth_d1_minus_pop.value}{-0.00102495}$ & \raggedright \resultnum{v4x.x6.x6a.strengths.r070.a3000.matching_bias.truth_d1_minus_pop.se}{0.000594433} \tabularnewline
\raggedright \resultnum{v4x.x6.x6a.strengths.r070.a3000.a}{0.3}, \resultnum{v4x.x6.x6a.strengths.r070.a3000.r}{0.7} & \raggedright Mean amplification & \raggedright \resultnum{v4x.x6.x6a.strengths.r070.a3000.matching_bias.amplification.mean}{0.071804} & \raggedright not applicable \tabularnewline
\raggedright \resultnum{v4x.x6.x6a.strengths.r070.a3000.a}{0.3}, \resultnum{v4x.x6.x6a.strengths.r070.a3000.r}{0.7} & \raggedright MSE / mean squared SE & \raggedright \resultnum{v4x.x6.x6a.strengths.r070.a3000.matching_bias.mse_to_mean_se2}{0.998125} & \raggedright not applicable \tabularnewline
\raggedright \resultnum{v4x.x6.x6a.strengths.r070.a3000.a}{0.3}, \resultnum{v4x.x6.x6a.strengths.r070.a3000.r}{0.7} & \raggedright Lag-1 fit residual, mean / SD & \raggedright \resultnum{v4x.x6.x6a.strengths.r070.a3000.matching_bias.fit_residual_se_units.mean.i000}{0.000452055} & \raggedright \resultnum{v4x.x6.x6a.strengths.r070.a3000.matching_bias.fit_residual_se_units.sd.i000}{0.0337667} \tabularnewline
\raggedright \resultnum{v4x.x6.x6a.strengths.r070.a3000.a}{0.3}, \resultnum{v4x.x6.x6a.strengths.r070.a3000.r}{0.7} & \raggedright Lag-2 fit residual, mean / SD & \raggedright $\resultnum{v4x.x6.x6a.strengths.r070.a3000.matching_bias.fit_residual_se_units.mean.i001}{-0.0136411}$ & \raggedright \resultnum{v4x.x6.x6a.strengths.r070.a3000.matching_bias.fit_residual_se_units.sd.i001}{0.104316} \tabularnewline
\raggedright \resultnum{v4x.x6.x6a.strengths.r070.a6000.a}{0.6}, \resultnum{v4x.x6.x6a.strengths.r070.a6000.r}{0.7} & \raggedright Mean standardized error & \raggedright \resultnum{v4x.x6.x6a.strengths.r070.a6000.vs_d1.standardized_error.mean}{0.0882024} & \raggedright [$\resultnum{v4x.x6.x6a.strengths.r070.a6000.vs_d1.standardized_error.mean_interval95.i000}{-0.0607774}$, \resultnum{v4x.x6.x6a.strengths.r070.a6000.vs_d1.standardized_error.mean_interval95.i001}{0.237182}] \tabularnewline
\raggedright \resultnum{v4x.x6.x6a.strengths.r070.a6000.a}{0.6}, \resultnum{v4x.x6.x6a.strengths.r070.a6000.r}{0.7} & \raggedright SD of standardized errors / uncovered-bias flag & \raggedright \resultnum{v4x.x6.x6a.strengths.r070.a6000.vs_d1.standardized_error.sd}{0.97934} & \raggedright none \tabularnewline
\raggedright \resultnum{v4x.x6.x6a.strengths.r070.a6000.a}{0.6}, \resultnum{v4x.x6.x6a.strengths.r070.a6000.r}{0.7} & \raggedright Share of standardized errors below $-1.96$ & \raggedright \resultnum{v4x.x6.x6a.strengths.r070.a6000.vs_d1.standardized_error.below_minus_1_96}{0.0361446} & \raggedright not applicable \tabularnewline
\raggedright \resultnum{v4x.x6.x6a.strengths.r070.a6000.a}{0.6}, \resultnum{v4x.x6.x6a.strengths.r070.a6000.r}{0.7} & \raggedright Share of standardized errors above 1.96 & \raggedright \resultnum{v4x.x6.x6a.strengths.r070.a6000.vs_d1.standardized_error.above_1_96}{0.0120482} & \raggedright not applicable \tabularnewline
\raggedright \resultnum{v4x.x6.x6a.strengths.r070.a6000.a}{0.6}, \resultnum{v4x.x6.x6a.strengths.r070.a6000.r}{0.7} & \raggedright Mean shift of the local nuisance draws (D5) & \raggedright \resultnum{v4x.x6.x6a.strengths.r070.a6000.matching_bias.d5_mean_shift.mean}{-1.36922e-06} & \raggedright not applicable \tabularnewline
\raggedright \resultnum{v4x.x6.x6a.strengths.r070.a6000.a}{0.6}, \resultnum{v4x.x6.x6a.strengths.r070.a6000.r}{0.7} & \raggedright Empirical matching bias & \raggedright \resultnum{v4x.x6.x6a.strengths.r070.a6000.matching_bias.empirical_bias.mean}{2.276e-05} & \raggedright \resultnum{v4x.x6.x6a.strengths.r070.a6000.matching_bias.empirical_bias.mcse}{4.4259e-05} \tabularnewline
\raggedright \resultnum{v4x.x6.x6a.strengths.r070.a6000.a}{0.6}, \resultnum{v4x.x6.x6a.strengths.r070.a6000.r}{0.7} & \raggedright Procedural minus population-match truth & \raggedright \resultnum{v4x.x6.x6a.strengths.r070.a6000.matching_bias.truth_d1_minus_pop.value}{6.30617e-06} & \raggedright \resultnum{v4x.x6.x6a.strengths.r070.a6000.matching_bias.truth_d1_minus_pop.se}{5.18473e-05} \tabularnewline
\raggedright \resultnum{v4x.x6.x6a.strengths.r070.a6000.a}{0.6}, \resultnum{v4x.x6.x6a.strengths.r070.a6000.r}{0.7} & \raggedright Mean amplification & \raggedright \resultnum{v4x.x6.x6a.strengths.r070.a6000.matching_bias.amplification.mean}{0.0821286} & \raggedright not applicable \tabularnewline
\raggedright \resultnum{v4x.x6.x6a.strengths.r070.a6000.a}{0.6}, \resultnum{v4x.x6.x6a.strengths.r070.a6000.r}{0.7} & \raggedright MSE / mean squared SE & \raggedright \resultnum{v4x.x6.x6a.strengths.r070.a6000.matching_bias.mse_to_mean_se2}{0.990375} & \raggedright not applicable \tabularnewline
\raggedright \resultnum{v4x.x6.x6a.strengths.r070.a6000.a}{0.6}, \resultnum{v4x.x6.x6a.strengths.r070.a6000.r}{0.7} & \raggedright Lag-1 fit residual, mean / SD & \raggedright \resultnum{v4x.x6.x6a.strengths.r070.a6000.matching_bias.fit_residual_se_units.mean.i000}{0.00159932} & \raggedright \resultnum{v4x.x6.x6a.strengths.r070.a6000.matching_bias.fit_residual_se_units.sd.i000}{0.0395943} \tabularnewline
\raggedright \resultnum{v4x.x6.x6a.strengths.r070.a6000.a}{0.6}, \resultnum{v4x.x6.x6a.strengths.r070.a6000.r}{0.7} & \raggedright Lag-2 fit residual, mean / SD & \raggedright $\resultnum{v4x.x6.x6a.strengths.r070.a6000.matching_bias.fit_residual_se_units.mean.i001}{-0.040435}$ & \raggedright \resultnum{v4x.x6.x6a.strengths.r070.a6000.matching_bias.fit_residual_se_units.sd.i001}{0.117166} \tabularnewline
\raggedright \resultnum{v4x.x6.x6a.strengths.r070.a9000.a}{0.9}, \resultnum{v4x.x6.x6a.strengths.r070.a9000.r}{0.7} & \raggedright Mean standardized error & \raggedright \resultnum{v4x.x6.x6a.strengths.r070.a9000.vs_d1.standardized_error.mean}{0} & \raggedright [\resultnum{v4x.x6.x6a.strengths.r070.a9000.vs_d1.standardized_error.mean_interval95.i000}{0}, \resultnum{v4x.x6.x6a.strengths.r070.a9000.vs_d1.standardized_error.mean_interval95.i001}{0}] \tabularnewline
\raggedright \resultnum{v4x.x6.x6a.strengths.r070.a9000.a}{0.9}, \resultnum{v4x.x6.x6a.strengths.r070.a9000.r}{0.7} & \raggedright SD of standardized errors / uncovered-bias flag & \raggedright \resultnum{v4x.x6.x6a.strengths.r070.a9000.vs_d1.standardized_error.sd}{0} & \raggedright none \tabularnewline
\raggedright \resultnum{v4x.x6.x6a.strengths.r070.a9000.a}{0.9}, \resultnum{v4x.x6.x6a.strengths.r070.a9000.r}{0.7} & \raggedright Share of standardized errors below $-1.96$ & \raggedright \resultnum{v4x.x6.x6a.strengths.r070.a9000.vs_d1.standardized_error.below_minus_1_96}{0} & \raggedright not applicable \tabularnewline
\raggedright \resultnum{v4x.x6.x6a.strengths.r070.a9000.a}{0.9}, \resultnum{v4x.x6.x6a.strengths.r070.a9000.r}{0.7} & \raggedright Share of standardized errors above 1.96 & \raggedright \resultnum{v4x.x6.x6a.strengths.r070.a9000.vs_d1.standardized_error.above_1_96}{0} & \raggedright not applicable \tabularnewline
\raggedright \resultnum{v4x.x6.x6a.strengths.r070.a9000.a}{0.9}, \resultnum{v4x.x6.x6a.strengths.r070.a9000.r}{0.7} & \raggedright Mean shift of the local nuisance draws (D5) & \raggedright \resultnum{v4x.x6.x6a.strengths.r070.a9000.matching_bias.d5_mean_shift.mean}{-8.39641e-07} & \raggedright not applicable \tabularnewline
\raggedright \resultnum{v4x.x6.x6a.strengths.r070.a9000.a}{0.9}, \resultnum{v4x.x6.x6a.strengths.r070.a9000.r}{0.7} & \raggedright Empirical matching bias & \raggedright \resultnum{v4x.x6.x6a.strengths.r070.a9000.matching_bias.empirical_bias.mean}{4.5314e-06} & \raggedright \resultnum{v4x.x6.x6a.strengths.r070.a9000.matching_bias.empirical_bias.mcse}{4.73525e-06} \tabularnewline
\raggedright \resultnum{v4x.x6.x6a.strengths.r070.a9000.a}{0.9}, \resultnum{v4x.x6.x6a.strengths.r070.a9000.r}{0.7} & \raggedright Procedural minus population-match truth & \raggedright \resultnum{v4x.x6.x6a.strengths.r070.a9000.matching_bias.truth_d1_minus_pop.value}{-1.87423e-07} & \raggedright \resultnum{v4x.x6.x6a.strengths.r070.a9000.matching_bias.truth_d1_minus_pop.se}{6.59492e-06} \tabularnewline
\raggedright \resultnum{v4x.x6.x6a.strengths.r070.a9000.a}{0.9}, \resultnum{v4x.x6.x6a.strengths.r070.a9000.r}{0.7} & \raggedright Mean amplification & \raggedright \resultnum{v4x.x6.x6a.strengths.r070.a9000.matching_bias.amplification.mean}{0.0941873} & \raggedright not applicable \tabularnewline
\raggedright \resultnum{v4x.x6.x6a.strengths.r070.a9000.a}{0.9}, \resultnum{v4x.x6.x6a.strengths.r070.a9000.r}{0.7} & \raggedright MSE / mean squared SE & \raggedright \resultnum{v4x.x6.x6a.strengths.r070.a9000.matching_bias.mse_to_mean_se2}{1.0931} & \raggedright not applicable \tabularnewline
\raggedright \resultnum{v4x.x6.x6a.strengths.r070.a9000.a}{0.9}, \resultnum{v4x.x6.x6a.strengths.r070.a9000.r}{0.7} & \raggedright Lag-1 fit residual, mean / SD & \raggedright \resultnum{v4x.x6.x6a.strengths.r070.a9000.matching_bias.fit_residual_se_units.mean.i000}{0.00137473} & \raggedright \resultnum{v4x.x6.x6a.strengths.r070.a9000.matching_bias.fit_residual_se_units.sd.i000}{0.0519058} \tabularnewline
\raggedright \resultnum{v4x.x6.x6a.strengths.r070.a9000.a}{0.9}, \resultnum{v4x.x6.x6a.strengths.r070.a9000.r}{0.7} & \raggedright Lag-2 fit residual, mean / SD & \raggedright $\resultnum{v4x.x6.x6a.strengths.r070.a9000.matching_bias.fit_residual_se_units.mean.i001}{-0.0146933}$ & \raggedright \resultnum{v4x.x6.x6a.strengths.r070.a9000.matching_bias.fit_residual_se_units.sd.i001}{0.113276} \tabularnewline
\raggedright \resultnum{v4x.x6.x6b.cells.r050.a0300.a}{0.03}, \resultnum{v4x.x6.x6b.cells.r050.a0300.r}{0.5} & \raggedright Mean standardized error & \raggedright \resultnum{v4x.x6.x6b.cells.r050.a0300.vs_d1.standardized_error.mean}{0.0669317} & \raggedright [$\resultnum{v4x.x6.x6b.cells.r050.a0300.vs_d1.standardized_error.mean_interval95.i000}{-0.0646519}$, \resultnum{v4x.x6.x6b.cells.r050.a0300.vs_d1.standardized_error.mean_interval95.i001}{0.198515}] \tabularnewline
\botrule
\end{tabular}
\end{table*}
\begin{table*}[htbp]
\setlength{\tabcolsep}{2pt}
\addtocounter{table}{-\csname @ne\endcsname}
\providecommand{\theHtable}{\thetable}
\renewcommand{\theHtable}{v4x-x6-diagnostics-part2}
\caption{Band coverage check (X6): diagnostics (continued, part 3 of 3)}\label{tab:v4x-x6-diagnostics-part2}
\centering\footnotesize
\begin{tabular}{@{}p{0.150\textwidth}p{0.400\textwidth}p{0.170\textwidth}p{0.220\textwidth}@{}}
\toprule
\raggedright $a,r$ & \raggedright \mbox{Diagnostic} & \raggedright \mbox{Estimate} & \raggedright \mbox{MC} \mbox{interval} \mbox{/} \mbox{error} \mbox{/} \mbox{SD} \mbox{/} \mbox{flag} \tabularnewline
\midrule
\raggedright \resultnum{v4x.x6.x6b.cells.r050.a0300.a}{0.03}, \resultnum{v4x.x6.x6b.cells.r050.a0300.r}{0.5} & \raggedright SD of standardized errors / uncovered-bias flag & \raggedright \resultnum{v4x.x6.x6b.cells.r050.a0300.vs_d1.standardized_error.sd}{0.949442} & \raggedright none \tabularnewline
\raggedright \resultnum{v4x.x6.x6b.cells.r050.a0300.a}{0.03}, \resultnum{v4x.x6.x6b.cells.r050.a0300.r}{0.5} & \raggedright Share of standardized errors below $-1.96$ & \raggedright \resultnum{v4x.x6.x6b.cells.r050.a0300.vs_d1.standardized_error.below_minus_1_96}{0.005} & \raggedright not applicable \tabularnewline
\raggedright \resultnum{v4x.x6.x6b.cells.r050.a0300.a}{0.03}, \resultnum{v4x.x6.x6b.cells.r050.a0300.r}{0.5} & \raggedright Share of standardized errors above 1.96 & \raggedright \resultnum{v4x.x6.x6b.cells.r050.a0300.vs_d1.standardized_error.above_1_96}{0.03} & \raggedright not applicable \tabularnewline
\raggedright \resultnum{v4x.x6.x6b.cells.r050.a0300.a}{0.03}, \resultnum{v4x.x6.x6b.cells.r050.a0300.r}{0.5} & \raggedright Mean shift of the local nuisance draws (D5) & \raggedright $\resultnum{v4x.x6.x6b.cells.r050.a0300.matching_bias.d5_mean_shift.mean}{-0.00124302}$ & \raggedright not applicable \tabularnewline
\raggedright \resultnum{v4x.x6.x6b.cells.r050.a0300.a}{0.03}, \resultnum{v4x.x6.x6b.cells.r050.a0300.r}{0.5} & \raggedright Empirical matching bias & \raggedright \resultnum{v4x.x6.x6b.cells.r050.a0300.matching_bias.empirical_bias.mean}{0.000899131} & \raggedright \resultnum{v4x.x6.x6b.cells.r050.a0300.matching_bias.empirical_bias.mcse}{0.0016859} \tabularnewline
\raggedright \resultnum{v4x.x6.x6b.cells.r050.a0300.a}{0.03}, \resultnum{v4x.x6.x6b.cells.r050.a0300.r}{0.5} & \raggedright Procedural minus population-match truth & \raggedright $\resultnum{v4x.x6.x6b.cells.r050.a0300.matching_bias.truth_d1_minus_pop.value}{-0.00193966}$ & \raggedright \resultnum{v4x.x6.x6b.cells.r050.a0300.matching_bias.truth_d1_minus_pop.se}{0.00163592} \tabularnewline
\raggedright \resultnum{v4x.x6.x6b.cells.r050.a0300.a}{0.03}, \resultnum{v4x.x6.x6b.cells.r050.a0300.r}{0.5} & \raggedright Mean amplification & \raggedright \resultnum{v4x.x6.x6b.cells.r050.a0300.matching_bias.amplification.mean}{0.0450147} & \raggedright not applicable \tabularnewline
\raggedright \resultnum{v4x.x6.x6b.cells.r050.a0300.a}{0.03}, \resultnum{v4x.x6.x6b.cells.r050.a0300.r}{0.5} & \raggedright MSE / mean squared SE & \raggedright \resultnum{v4x.x6.x6b.cells.r050.a0300.matching_bias.mse_to_mean_se2}{0.791098} & \raggedright not applicable \tabularnewline
\raggedright \resultnum{v4x.x6.x6b.cells.r050.a0300.a}{0.03}, \resultnum{v4x.x6.x6b.cells.r050.a0300.r}{0.5} & \raggedright Lag-1 fit residual, mean / SD & \raggedright $\resultnum{v4x.x6.x6b.cells.r050.a0300.matching_bias.fit_residual_se_units.mean.i000}{-0.00593687}$ & \raggedright \resultnum{v4x.x6.x6b.cells.r050.a0300.matching_bias.fit_residual_se_units.sd.i000}{0.051117} \tabularnewline
\raggedright \resultnum{v4x.x6.x6b.cells.r050.a0300.a}{0.03}, \resultnum{v4x.x6.x6b.cells.r050.a0300.r}{0.5} & \raggedright Lag-2 fit residual, mean / SD & \raggedright \resultnum{v4x.x6.x6b.cells.r050.a0300.matching_bias.fit_residual_se_units.mean.i001}{0.00819784} & \raggedright \resultnum{v4x.x6.x6b.cells.r050.a0300.matching_bias.fit_residual_se_units.sd.i001}{0.127718} \tabularnewline
\raggedright \resultnum{v4x.x6.x6b.cells.r090.a0300.a}{0.03}, \resultnum{v4x.x6.x6b.cells.r090.a0300.r}{0.9} & \raggedright Mean standardized error & \raggedright \resultnum{v4x.x6.x6b.cells.r090.a0300.vs_d1.standardized_error.mean}{0.0313658} & \raggedright [$\resultnum{v4x.x6.x6b.cells.r090.a0300.vs_d1.standardized_error.mean_interval95.i000}{-0.105767}$, \resultnum{v4x.x6.x6b.cells.r090.a0300.vs_d1.standardized_error.mean_interval95.i001}{0.168498}] \tabularnewline
\raggedright \resultnum{v4x.x6.x6b.cells.r090.a0300.a}{0.03}, \resultnum{v4x.x6.x6b.cells.r090.a0300.r}{0.9} & \raggedright SD of standardized errors / uncovered-bias flag & \raggedright \resultnum{v4x.x6.x6b.cells.r090.a0300.vs_d1.standardized_error.sd}{0.989482} & \raggedright none \tabularnewline
\raggedright \resultnum{v4x.x6.x6b.cells.r090.a0300.a}{0.03}, \resultnum{v4x.x6.x6b.cells.r090.a0300.r}{0.9} & \raggedright Share of standardized errors below $-1.96$ & \raggedright \resultnum{v4x.x6.x6b.cells.r090.a0300.vs_d1.standardized_error.below_minus_1_96}{0.025} & \raggedright not applicable \tabularnewline
\raggedright \resultnum{v4x.x6.x6b.cells.r090.a0300.a}{0.03}, \resultnum{v4x.x6.x6b.cells.r090.a0300.r}{0.9} & \raggedright Share of standardized errors above 1.96 & \raggedright \resultnum{v4x.x6.x6b.cells.r090.a0300.vs_d1.standardized_error.above_1_96}{0.03} & \raggedright not applicable \tabularnewline
\raggedright \resultnum{v4x.x6.x6b.cells.r090.a0300.a}{0.03}, \resultnum{v4x.x6.x6b.cells.r090.a0300.r}{0.9} & \raggedright Mean shift of the local nuisance draws (D5) & \raggedright \resultnum{v4x.x6.x6b.cells.r090.a0300.matching_bias.d5_mean_shift.mean}{0.000101384} & \raggedright not applicable \tabularnewline
\raggedright \resultnum{v4x.x6.x6b.cells.r090.a0300.a}{0.03}, \resultnum{v4x.x6.x6b.cells.r090.a0300.r}{0.9} & \raggedright Empirical matching bias & \raggedright \resultnum{v4x.x6.x6b.cells.r090.a0300.matching_bias.empirical_bias.mean}{0.000214193} & \raggedright \resultnum{v4x.x6.x6b.cells.r090.a0300.matching_bias.empirical_bias.mcse}{0.00114435} \tabularnewline
\raggedright \resultnum{v4x.x6.x6b.cells.r090.a0300.a}{0.03}, \resultnum{v4x.x6.x6b.cells.r090.a0300.r}{0.9} & \raggedright Procedural minus population-match truth & \raggedright $\resultnum{v4x.x6.x6b.cells.r090.a0300.matching_bias.truth_d1_minus_pop.value}{-0.00106907}$ & \raggedright \resultnum{v4x.x6.x6b.cells.r090.a0300.matching_bias.truth_d1_minus_pop.se}{0.00135342} \tabularnewline
\raggedright \resultnum{v4x.x6.x6b.cells.r090.a0300.a}{0.03}, \resultnum{v4x.x6.x6b.cells.r090.a0300.r}{0.9} & \raggedright Mean amplification & \raggedright \resultnum{v4x.x6.x6b.cells.r090.a0300.matching_bias.amplification.mean}{0.0303112} & \raggedright not applicable \tabularnewline
\raggedright \resultnum{v4x.x6.x6b.cells.r090.a0300.a}{0.03}, \resultnum{v4x.x6.x6b.cells.r090.a0300.r}{0.9} & \raggedright MSE / mean squared SE & \raggedright \resultnum{v4x.x6.x6b.cells.r090.a0300.matching_bias.mse_to_mean_se2}{0.983214} & \raggedright not applicable \tabularnewline
\raggedright \resultnum{v4x.x6.x6b.cells.r090.a0300.a}{0.03}, \resultnum{v4x.x6.x6b.cells.r090.a0300.r}{0.9} & \raggedright Lag-1 fit residual, mean / SD & \raggedright \resultnum{v4x.x6.x6b.cells.r090.a0300.matching_bias.fit_residual_se_units.mean.i000}{0.0254949} & \raggedright \resultnum{v4x.x6.x6b.cells.r090.a0300.matching_bias.fit_residual_se_units.sd.i000}{0.0868791} \tabularnewline
\raggedright \resultnum{v4x.x6.x6b.cells.r090.a0300.a}{0.03}, \resultnum{v4x.x6.x6b.cells.r090.a0300.r}{0.9} & \raggedright Lag-2 fit residual, mean / SD & \raggedright \resultnum{v4x.x6.x6b.cells.r090.a0300.matching_bias.fit_residual_se_units.mean.i001}{0.00253732} & \raggedright \resultnum{v4x.x6.x6b.cells.r090.a0300.matching_bias.fit_residual_se_units.sd.i001}{0.122083} \tabularnewline
\botrule
\end{tabular}
\end{table*}

\clearpage

\subsection{Expanded-feature classifier (X7)}\label{sec:x7}

\paragraph{Features.}\phantomsection\label{sec:x7-features}
X7 keeps the arms, the frozen fit of B on lags one and two and the frozen LDA code, and replaces the three frozen features by a fixed vector of 36 features of each panel, declared before any run and identical at every strength and signal share:
\begin{itemize}
\item the pooled Pearson lag correlations $L_3,\ldots,L_{10}$ of the observed categories, by the frozen lag-profile formula extended to lags 6--10 (8 features);
\item the frozen mobility summary: the mean absolute change of the observed category between consecutive waves within each tercile of the wave's mid-rank position (3 features);
\item the lag-1 transition frequencies between five groups of categories, $\{0,1\}$, $\{2,3\}$, $\{4,5,6\}$, $\{7,8\}$ and $\{9,10\}$, conditional on the row, over all consecutive observed pairs, with the last column of each row dropped (20 features);
\item among persons with at least six observed waves, the 0.1, 0.25, 0.5, 0.75 and 0.9 quantiles (linear interpolation) of the within-person standard deviation of the observed categories (5 features).
\end{itemize}
$L_1$ and $L_2$ stay out, because they are the matched moments. A feature that is not finite in some panel makes the cell a recorded numerical gap, without imputation. The features are standardized with the pooled training mean and standard deviation, and the LDA uses the frozen shrinkage $0.8S+0.2I$, equal priors and no feature search, hyperparameter or other classifier. The sizes are the frozen ones: 2,000 training panels per arm and 4,000 test and 4,000 reference panels per arm. Each panel is generated once and gives both feature vectors, so the frozen-feature and expanded-feature LDAs are evaluated on the same panels. The nuisance uncertainty uses the frozen local propagation with the expanded features, the Jacobian of their means coming from common-random-number finite differences with step 0.01 on regenerated matching panels (3,072 extra persistence-arm panels per cell). The boundaries, bands, censoring and zero-strength check come from the frozen curve code applied to each feature set; the AUC difference per cell has a paired bootstrap interval with 1,000 resamples of the networks.

\paragraph{Readings.}
The paired run gives a frozen-feature first boundary on the exploratory seeds and an expanded-feature first boundary in every unit (Tables~\ref{tab:v4x-x7}, \ref{tab:v4x-x7-boundaries} and~\ref{tab:v4x-x7-deltas}); errors and gaps of the expanded-feature arm are recorded separately from those of the frozen-feature arm, and the paired run had \resultnum{v4x.x7.inventory.gaps}{2} fit gaps, both at $a=0.90$. In \resultnum{v4x.x7.count.far}{9} units the expanded-feature LDA gained relative to the frozen-feature LDA under the declared paired criterion (reading ``gain''; machine label \texttt{far}): either the expanded boundary interval's upper end was below the frozen-feature point boundary, or the paired AUC-difference intervals were positive at every grid strength from the expanded first boundary up to, but excluding, the frozen-feature first boundary. The exact applicable predicate and any censoring are retained by unit. For this predicate, a frozen-feature first boundary not attained without point censoring is treated as beyond the grid.
Read after the results became available: the expanded-feature first boundary was $a=\resultnum{j.x7.first.min}{0.003}$ at $r=\resultnum{j.x7.first.min.r}{0.90}$ to $a=\resultnum{j.x7.first.max}{0.03}$ at $r=\resultnum{j.x7.first.max.r}{0.50, 0.55, 0.60}$, and at $a=0.03$ the expanded features raised the AUC by \resultnum{j.x7.auc_gain.a0300.min}{0.196} (at $r=\resultnum{j.x7.auc_gain.a0300.min.r}{0.50}$) to \resultnum{j.x7.auc_gain.a0300.max}{0.339} (at $r=\resultnum{j.x7.auc_gain.a0300.max.r}{0.85}$).
Changing the feature dimension changes estimation and regularization as well as information, so a richer fitted LDA establishes neither an information ceiling nor Bayes optimality; in population terms the expanded-feature boundary, like the frozen one, bounds the optimal first crossing on the grid from above, and equality of the two boundaries would not prove optimality.
Read after the results became available: richer summaries distinguish these specified laws at lower strengths, and the departure they detect is from persistence matched on only two moments. B is matched only on $L_1$ and $L_2$, and several of the expanded features, such as the lag-1 transition frequencies, the mobility by tercile and the quantiles of within-person dispersion, are moments that B was never matched to. The gain can therefore be described as sensitivity to features that B was not matched to; that description does not show that the gain is unrelated to influence, nor that it is specific to influence. No persistence model matched on those features was evaluated, and the channel-misspecification stress test (Section~\ref{sec:x5}) used only the three frozen features, so how far a misspecified channel alone separates panels under the expanded summary is unknown. The paired AUC difference was positive at $a=0.001$ at every share, with every interval above zero (Table~\ref{tab:v4x-x7-deltas}), a strength at which influence explains a share of about 0.0005 of the residual latent variance (Table~\ref{tab:supp-x2}); a small share of each person's variance can still change the law of a panel of thousands of people detectably, so a small latent share does not mean little observable information. These paired intervals describe evaluation uncertainty given the fitted comparison: at $a=0$, where the expected difference is zero, the intervals excluded zero at two of the nine shares, in opposite directions ($r=0.65$ and 0.80), and at $r=0.50$ and 0.55 the differences at $a=0.001$ are not much larger than those zero-strength excursions.

% Generated by v4_exploratory_export; edited during the journal revision
% (agent S, 2026-10-13): captions, headers and categorical labels translated
% into plain English. Every number keeps its \resultnum binding.
\begin{table*}[htbp]
\setlength{\tabcolsep}{2pt}
\caption{Expanded-feature classifier (X7): first boundaries of the frozen-feature LDA (lags 3--5) and of the expanded-feature LDA (36 features, Section~\ref{sec:x7-features}) on the same exploratory panels, with the interval of the expanded boundary. Reading under the declared paired criterion: gain, the expanded interval's upper end lies below the frozen-feature point boundary or the paired AUC differences are positive at every strength from the expanded boundary up to the frozen one (machine label far, which does not mean distance from an information ceiling). Changing the feature dimension also changes estimation and regularization.}\label{tab:v4x-x7}
\centering\footnotesize
\begin{tabular}{@{}p{0.100\textwidth}p{0.120\textwidth}p{0.220\textwidth}p{0.220\textwidth}p{0.250\textwidth}@{}}
\toprule
\raggedright $r$ & \raggedright \mbox{Reading} & \raggedright \mbox{Frozen-feature} \mbox{first} \mbox{boundary} & \raggedright \mbox{Expanded-feature} \mbox{first} \mbox{boundary} & \raggedright \mbox{Expanded} \mbox{interval} \tabularnewline
\midrule
\raggedright \resultnum{v4x.x7.units.r050.r}{0.5} & \raggedright gain & \raggedright \resultnum{v4x.x7.units.r050.x7.first}{0.3} & \raggedright \resultnum{v4x.x7.units.r050.rich.first}{0.03} & \raggedright [\resultnum{v4x.x7.units.r050.rich.first_interval.i000}{0.03}, \resultnum{v4x.x7.units.r050.rich.first_interval.i001}{0.03}] \tabularnewline
\raggedright \resultnum{v4x.x7.units.r055.r}{0.55} & \raggedright gain & \raggedright \resultnum{v4x.x7.units.r055.x7.first}{0.3} & \raggedright \resultnum{v4x.x7.units.r055.rich.first}{0.03} & \raggedright [\resultnum{v4x.x7.units.r055.rich.first_interval.i000}{0.03}, \resultnum{v4x.x7.units.r055.rich.first_interval.i001}{0.03}] \tabularnewline
\raggedright \resultnum{v4x.x7.units.r060.r}{0.6} & \raggedright gain & \raggedright \resultnum{v4x.x7.units.r060.x7.first}{0.1} & \raggedright \resultnum{v4x.x7.units.r060.rich.first}{0.03} & \raggedright [\resultnum{v4x.x7.units.r060.rich.first_interval.i000}{0.01}, \resultnum{v4x.x7.units.r060.rich.first_interval.i001}{0.03}] \tabularnewline
\raggedright \resultnum{v4x.x7.units.r065.r}{0.65} & \raggedright gain & \raggedright \resultnum{v4x.x7.units.r065.x7.first}{0.1} & \raggedright \resultnum{v4x.x7.units.r065.rich.first}{0.01} & \raggedright [\resultnum{v4x.x7.units.r065.rich.first_interval.i000}{0.01}, \resultnum{v4x.x7.units.r065.rich.first_interval.i001}{0.03}] \tabularnewline
\raggedright \resultnum{v4x.x7.units.r070.r}{0.7} & \raggedright gain & \raggedright \resultnum{v4x.x7.units.r070.x7.first}{0.1} & \raggedright \resultnum{v4x.x7.units.r070.rich.first}{0.01} & \raggedright [\resultnum{v4x.x7.units.r070.rich.first_interval.i000}{0.01}, \resultnum{v4x.x7.units.r070.rich.first_interval.i001}{0.01}] \tabularnewline
\raggedright \resultnum{v4x.x7.units.r075.r}{0.75} & \raggedright gain & \raggedright \resultnum{v4x.x7.units.r075.x7.first}{0.1} & \raggedright \resultnum{v4x.x7.units.r075.rich.first}{0.01} & \raggedright [\resultnum{v4x.x7.units.r075.rich.first_interval.i000}{0.01}, \resultnum{v4x.x7.units.r075.rich.first_interval.i001}{0.01}] \tabularnewline
\raggedright \resultnum{v4x.x7.units.r080.r}{0.8} & \raggedright gain & \raggedright \resultnum{v4x.x7.units.r080.x7.first}{0.1} & \raggedright \resultnum{v4x.x7.units.r080.rich.first}{0.01} & \raggedright [\resultnum{v4x.x7.units.r080.rich.first_interval.i000}{0.01}, \resultnum{v4x.x7.units.r080.rich.first_interval.i001}{0.01}] \tabularnewline
\raggedright \resultnum{v4x.x7.units.r085.r}{0.85} & \raggedright gain & \raggedright \resultnum{v4x.x7.units.r085.x7.first}{0.1} & \raggedright \resultnum{v4x.x7.units.r085.rich.first}{0.01} & \raggedright [\resultnum{v4x.x7.units.r085.rich.first_interval.i000}{0.003}, \resultnum{v4x.x7.units.r085.rich.first_interval.i001}{0.01}] \tabularnewline
\raggedright \resultnum{v4x.x7.units.r090.r}{0.9} & \raggedright gain & \raggedright \resultnum{v4x.x7.units.r090.x7.first}{0.1} & \raggedright \resultnum{v4x.x7.units.r090.rich.first}{0.003} & \raggedright [\resultnum{v4x.x7.units.r090.rich.first_interval.i000}{0.003}, \resultnum{v4x.x7.units.r090.rich.first_interval.i001}{0.003}] \tabularnewline
\botrule
\end{tabular}
\end{table*}

% Generated by v4_exploratory_export; edited during the journal revision
% (agent S, 2026-10-13): captions, headers and categorical labels translated
% into plain English. Every number keeps its \resultnum binding.
\begin{table*}[htbp]
\setlength{\tabcolsep}{2pt}
\caption{Expanded-feature classifier (X7): first and sustained boundaries of both paired LDAs with their boundary intervals and their point and interval censoring under the frozen gap rule. None: no censoring; gap-censored: the boundary cannot be declared across the gap named in the censoring columns (here the fit gaps at $a=0.90$ for $r=0.50$ and 0.55).}\label{tab:v4x-x7-boundaries}
\centering\footnotesize
\begin{tabular}{@{}p{0.079\textwidth}p{0.266\textwidth}p{0.128\textwidth}p{0.168\textwidth}p{0.148\textwidth}p{0.148\textwidth}@{}}
\toprule
\raggedright $r$ & \raggedright \mbox{LDA,} \mbox{boundary} & \raggedright \mbox{Point} & \raggedright \mbox{Interval} & \raggedright \mbox{Point} \mbox{censoring} & \raggedright \mbox{Interval} \mbox{censoring} \tabularnewline
\midrule
\raggedright \resultnum{v4x.x7.units.r050.r}{0.5} & \raggedright frozen features, first & \raggedright \resultnum{v4x.x7.units.r050.x7.first}{0.3} & \raggedright [\resultnum{v4x.x7.units.r050.x7.first_interval.i000}{0.3}, \resultnum{v4x.x7.units.r050.x7.first_interval.i001}{0.3}] & \raggedright none & \raggedright none \tabularnewline
\raggedright \resultnum{v4x.x7.units.r050.r}{0.5} & \raggedright frozen features, sustained & \raggedright gap-censored & \raggedright [censored, censored] & \raggedright gap at $a=0.9$ & \raggedright gap at $a=0.9$ \tabularnewline
\raggedright \resultnum{v4x.x7.units.r050.r}{0.5} & \raggedright expanded features, first & \raggedright \resultnum{v4x.x7.units.r050.rich.first}{0.03} & \raggedright [\resultnum{v4x.x7.units.r050.rich.first_interval.i000}{0.03}, \resultnum{v4x.x7.units.r050.rich.first_interval.i001}{0.03}] & \raggedright none & \raggedright none \tabularnewline
\raggedright \resultnum{v4x.x7.units.r050.r}{0.5} & \raggedright expanded features, sustained & \raggedright gap-censored & \raggedright [censored, censored] & \raggedright gap at $a=0.9$ & \raggedright gap at $a=0.9$ \tabularnewline
\raggedright \resultnum{v4x.x7.units.r055.r}{0.55} & \raggedright frozen features, first & \raggedright \resultnum{v4x.x7.units.r055.x7.first}{0.3} & \raggedright [\resultnum{v4x.x7.units.r055.x7.first_interval.i000}{0.1}, \resultnum{v4x.x7.units.r055.x7.first_interval.i001}{0.3}] & \raggedright none & \raggedright none \tabularnewline
\raggedright \resultnum{v4x.x7.units.r055.r}{0.55} & \raggedright frozen features, sustained & \raggedright gap-censored & \raggedright [censored, censored] & \raggedright gap at $a=0.9$ & \raggedright gap at $a=0.9$ \tabularnewline
\raggedright \resultnum{v4x.x7.units.r055.r}{0.55} & \raggedright expanded features, first & \raggedright \resultnum{v4x.x7.units.r055.rich.first}{0.03} & \raggedright [\resultnum{v4x.x7.units.r055.rich.first_interval.i000}{0.03}, \resultnum{v4x.x7.units.r055.rich.first_interval.i001}{0.03}] & \raggedright none & \raggedright none \tabularnewline
\raggedright \resultnum{v4x.x7.units.r055.r}{0.55} & \raggedright expanded features, sustained & \raggedright gap-censored & \raggedright [censored, censored] & \raggedright gap at $a=0.9$ & \raggedright gap at $a=0.9$ \tabularnewline
\raggedright \resultnum{v4x.x7.units.r060.r}{0.6} & \raggedright frozen features, first & \raggedright \resultnum{v4x.x7.units.r060.x7.first}{0.1} & \raggedright [\resultnum{v4x.x7.units.r060.x7.first_interval.i000}{0.1}, \resultnum{v4x.x7.units.r060.x7.first_interval.i001}{0.3}] & \raggedright none & \raggedright none \tabularnewline
\raggedright \resultnum{v4x.x7.units.r060.r}{0.6} & \raggedright frozen features, sustained & \raggedright \resultnum{v4x.x7.units.r060.x7.sustained}{0.1} & \raggedright [\resultnum{v4x.x7.units.r060.x7.sustained_interval.i000}{0.1}, \resultnum{v4x.x7.units.r060.x7.sustained_interval.i001}{0.3}] & \raggedright none & \raggedright none \tabularnewline
\raggedright \resultnum{v4x.x7.units.r060.r}{0.6} & \raggedright expanded features, first & \raggedright \resultnum{v4x.x7.units.r060.rich.first}{0.03} & \raggedright [\resultnum{v4x.x7.units.r060.rich.first_interval.i000}{0.01}, \resultnum{v4x.x7.units.r060.rich.first_interval.i001}{0.03}] & \raggedright none & \raggedright none \tabularnewline
\raggedright \resultnum{v4x.x7.units.r060.r}{0.6} & \raggedright expanded features, sustained & \raggedright \resultnum{v4x.x7.units.r060.rich.sustained}{0.03} & \raggedright [\resultnum{v4x.x7.units.r060.rich.sustained_interval.i000}{0.01}, \resultnum{v4x.x7.units.r060.rich.sustained_interval.i001}{0.03}] & \raggedright none & \raggedright none \tabularnewline
\raggedright \resultnum{v4x.x7.units.r065.r}{0.65} & \raggedright frozen features, first & \raggedright \resultnum{v4x.x7.units.r065.x7.first}{0.1} & \raggedright [\resultnum{v4x.x7.units.r065.x7.first_interval.i000}{0.1}, \resultnum{v4x.x7.units.r065.x7.first_interval.i001}{0.3}] & \raggedright none & \raggedright none \tabularnewline
\raggedright \resultnum{v4x.x7.units.r065.r}{0.65} & \raggedright frozen features, sustained & \raggedright \resultnum{v4x.x7.units.r065.x7.sustained}{0.1} & \raggedright [\resultnum{v4x.x7.units.r065.x7.sustained_interval.i000}{0.1}, \resultnum{v4x.x7.units.r065.x7.sustained_interval.i001}{0.3}] & \raggedright none & \raggedright none \tabularnewline
\raggedright \resultnum{v4x.x7.units.r065.r}{0.65} & \raggedright expanded features, first & \raggedright \resultnum{v4x.x7.units.r065.rich.first}{0.01} & \raggedright [\resultnum{v4x.x7.units.r065.rich.first_interval.i000}{0.01}, \resultnum{v4x.x7.units.r065.rich.first_interval.i001}{0.03}] & \raggedright none & \raggedright none \tabularnewline
\raggedright \resultnum{v4x.x7.units.r065.r}{0.65} & \raggedright expanded features, sustained & \raggedright \resultnum{v4x.x7.units.r065.rich.sustained}{0.01} & \raggedright [\resultnum{v4x.x7.units.r065.rich.sustained_interval.i000}{0.01}, \resultnum{v4x.x7.units.r065.rich.sustained_interval.i001}{0.03}] & \raggedright none & \raggedright none \tabularnewline
\raggedright \resultnum{v4x.x7.units.r070.r}{0.7} & \raggedright frozen features, first & \raggedright \resultnum{v4x.x7.units.r070.x7.first}{0.1} & \raggedright [\resultnum{v4x.x7.units.r070.x7.first_interval.i000}{0.1}, \resultnum{v4x.x7.units.r070.x7.first_interval.i001}{0.1}] & \raggedright none & \raggedright none \tabularnewline
\raggedright \resultnum{v4x.x7.units.r070.r}{0.7} & \raggedright frozen features, sustained & \raggedright \resultnum{v4x.x7.units.r070.x7.sustained}{0.1} & \raggedright [\resultnum{v4x.x7.units.r070.x7.sustained_interval.i000}{0.1}, \resultnum{v4x.x7.units.r070.x7.sustained_interval.i001}{0.1}] & \raggedright none & \raggedright none \tabularnewline
\raggedright \resultnum{v4x.x7.units.r070.r}{0.7} & \raggedright expanded features, first & \raggedright \resultnum{v4x.x7.units.r070.rich.first}{0.01} & \raggedright [\resultnum{v4x.x7.units.r070.rich.first_interval.i000}{0.01}, \resultnum{v4x.x7.units.r070.rich.first_interval.i001}{0.01}] & \raggedright none & \raggedright none \tabularnewline
\raggedright \resultnum{v4x.x7.units.r070.r}{0.7} & \raggedright expanded features, sustained & \raggedright \resultnum{v4x.x7.units.r070.rich.sustained}{0.01} & \raggedright [\resultnum{v4x.x7.units.r070.rich.sustained_interval.i000}{0.01}, \resultnum{v4x.x7.units.r070.rich.sustained_interval.i001}{0.01}] & \raggedright none & \raggedright none \tabularnewline
\raggedright \resultnum{v4x.x7.units.r075.r}{0.75} & \raggedright frozen features, first & \raggedright \resultnum{v4x.x7.units.r075.x7.first}{0.1} & \raggedright [\resultnum{v4x.x7.units.r075.x7.first_interval.i000}{0.1}, \resultnum{v4x.x7.units.r075.x7.first_interval.i001}{0.1}] & \raggedright none & \raggedright none \tabularnewline
\raggedright \resultnum{v4x.x7.units.r075.r}{0.75} & \raggedright frozen features, sustained & \raggedright \resultnum{v4x.x7.units.r075.x7.sustained}{0.1} & \raggedright [\resultnum{v4x.x7.units.r075.x7.sustained_interval.i000}{0.1}, \resultnum{v4x.x7.units.r075.x7.sustained_interval.i001}{0.1}] & \raggedright none & \raggedright none \tabularnewline
\raggedright \resultnum{v4x.x7.units.r075.r}{0.75} & \raggedright expanded features, first & \raggedright \resultnum{v4x.x7.units.r075.rich.first}{0.01} & \raggedright [\resultnum{v4x.x7.units.r075.rich.first_interval.i000}{0.01}, \resultnum{v4x.x7.units.r075.rich.first_interval.i001}{0.01}] & \raggedright none & \raggedright none \tabularnewline
\raggedright \resultnum{v4x.x7.units.r075.r}{0.75} & \raggedright expanded features, sustained & \raggedright \resultnum{v4x.x7.units.r075.rich.sustained}{0.01} & \raggedright [\resultnum{v4x.x7.units.r075.rich.sustained_interval.i000}{0.01}, \resultnum{v4x.x7.units.r075.rich.sustained_interval.i001}{0.01}] & \raggedright none & \raggedright none \tabularnewline
\raggedright \resultnum{v4x.x7.units.r080.r}{0.8} & \raggedright frozen features, first & \raggedright \resultnum{v4x.x7.units.r080.x7.first}{0.1} & \raggedright [\resultnum{v4x.x7.units.r080.x7.first_interval.i000}{0.1}, \resultnum{v4x.x7.units.r080.x7.first_interval.i001}{0.1}] & \raggedright none & \raggedright none \tabularnewline
\raggedright \resultnum{v4x.x7.units.r080.r}{0.8} & \raggedright frozen features, sustained & \raggedright \resultnum{v4x.x7.units.r080.x7.sustained}{0.1} & \raggedright [\resultnum{v4x.x7.units.r080.x7.sustained_interval.i000}{0.1}, \resultnum{v4x.x7.units.r080.x7.sustained_interval.i001}{0.1}] & \raggedright none & \raggedright none \tabularnewline
\raggedright \resultnum{v4x.x7.units.r080.r}{0.8} & \raggedright expanded features, first & \raggedright \resultnum{v4x.x7.units.r080.rich.first}{0.01} & \raggedright [\resultnum{v4x.x7.units.r080.rich.first_interval.i000}{0.01}, \resultnum{v4x.x7.units.r080.rich.first_interval.i001}{0.01}] & \raggedright none & \raggedright none \tabularnewline
\raggedright \resultnum{v4x.x7.units.r080.r}{0.8} & \raggedright expanded features, sustained & \raggedright \resultnum{v4x.x7.units.r080.rich.sustained}{0.01} & \raggedright [\resultnum{v4x.x7.units.r080.rich.sustained_interval.i000}{0.01}, \resultnum{v4x.x7.units.r080.rich.sustained_interval.i001}{0.01}] & \raggedright none & \raggedright none \tabularnewline
\raggedright \resultnum{v4x.x7.units.r085.r}{0.85} & \raggedright frozen features, first & \raggedright \resultnum{v4x.x7.units.r085.x7.first}{0.1} & \raggedright [\resultnum{v4x.x7.units.r085.x7.first_interval.i000}{0.1}, \resultnum{v4x.x7.units.r085.x7.first_interval.i001}{0.1}] & \raggedright none & \raggedright none \tabularnewline
\raggedright \resultnum{v4x.x7.units.r085.r}{0.85} & \raggedright frozen features, sustained & \raggedright \resultnum{v4x.x7.units.r085.x7.sustained}{0.1} & \raggedright [\resultnum{v4x.x7.units.r085.x7.sustained_interval.i000}{0.1}, \resultnum{v4x.x7.units.r085.x7.sustained_interval.i001}{0.1}] & \raggedright none & \raggedright none \tabularnewline
\raggedright \resultnum{v4x.x7.units.r085.r}{0.85} & \raggedright expanded features, first & \raggedright \resultnum{v4x.x7.units.r085.rich.first}{0.01} & \raggedright [\resultnum{v4x.x7.units.r085.rich.first_interval.i000}{0.003}, \resultnum{v4x.x7.units.r085.rich.first_interval.i001}{0.01}] & \raggedright none & \raggedright none \tabularnewline
\raggedright \resultnum{v4x.x7.units.r085.r}{0.85} & \raggedright expanded features, sustained & \raggedright \resultnum{v4x.x7.units.r085.rich.sustained}{0.01} & \raggedright [\resultnum{v4x.x7.units.r085.rich.sustained_interval.i000}{0.003}, \resultnum{v4x.x7.units.r085.rich.sustained_interval.i001}{0.01}] & \raggedright none & \raggedright none \tabularnewline
\raggedright \resultnum{v4x.x7.units.r090.r}{0.9} & \raggedright frozen features, first & \raggedright \resultnum{v4x.x7.units.r090.x7.first}{0.1} & \raggedright [\resultnum{v4x.x7.units.r090.x7.first_interval.i000}{0.1}, \resultnum{v4x.x7.units.r090.x7.first_interval.i001}{0.1}] & \raggedright none & \raggedright none \tabularnewline
\raggedright \resultnum{v4x.x7.units.r090.r}{0.9} & \raggedright frozen features, sustained & \raggedright \resultnum{v4x.x7.units.r090.x7.sustained}{0.1} & \raggedright [\resultnum{v4x.x7.units.r090.x7.sustained_interval.i000}{0.1}, \resultnum{v4x.x7.units.r090.x7.sustained_interval.i001}{0.1}] & \raggedright none & \raggedright none \tabularnewline
\raggedright \resultnum{v4x.x7.units.r090.r}{0.9} & \raggedright expanded features, first & \raggedright \resultnum{v4x.x7.units.r090.rich.first}{0.003} & \raggedright [\resultnum{v4x.x7.units.r090.rich.first_interval.i000}{0.003}, \resultnum{v4x.x7.units.r090.rich.first_interval.i001}{0.003}] & \raggedright none & \raggedright none \tabularnewline
\raggedright \resultnum{v4x.x7.units.r090.r}{0.9} & \raggedright expanded features, sustained & \raggedright \resultnum{v4x.x7.units.r090.rich.sustained}{0.003} & \raggedright [\resultnum{v4x.x7.units.r090.rich.sustained_interval.i000}{0.003}, \resultnum{v4x.x7.units.r090.rich.sustained_interval.i001}{0.003}] & \raggedright none & \raggedright none \tabularnewline
\botrule
\end{tabular}
\end{table*}

% Generated by v4_exploratory_export; edited during the journal revision
% (agent S, 2026-10-13): captions, headers and categorical labels translated
% into plain English. Every number keeps its \resultnum binding.
\begin{table*}[htbp]
\setlength{\tabcolsep}{2pt}
\caption{Expanded-feature classifier (X7): expanded-feature minus frozen-feature AUC on the same panels, with paired network-bootstrap 95\% intervals (1,000 resamples), by cell. The two fit gaps of the paired run, $(a,r)=(0.90,0.50)$ and $(0.90,0.55)$, have no contrast. The contrasts concern the two declared fitted LDAs, not an information ceiling.}\label{tab:v4x-x7-deltas}
\centering\footnotesize
\begin{tabular}{@{}p{0.140\textwidth}p{0.140\textwidth}p{0.220\textwidth}p{0.400\textwidth}@{}}
\toprule
\raggedright $r$ & \raggedright $a$ & \raggedright $\Delta$AUC & \raggedright \mbox{Interval} \tabularnewline
\midrule
\raggedright \resultnum{v4x.x7.units.r050.r}{0.5} & \raggedright \resultnum{v4x.x7.units.r050.delta_auc.i000.a}{0} & \raggedright \resultnum{v4x.x7.units.r050.delta_auc.i000.delta}{0.00115494} & \raggedright [$\resultnum{v4x.x7.units.r050.delta_auc.i000.interval.i000}{-0.013158}$, \resultnum{v4x.x7.units.r050.delta_auc.i000.interval.i001}{0.0148821}] \tabularnewline
\raggedright \resultnum{v4x.x7.units.r050.r}{0.5} & \raggedright \resultnum{v4x.x7.units.r050.delta_auc.i001.a}{0.001} & \raggedright \resultnum{v4x.x7.units.r050.delta_auc.i001.delta}{0.0353583} & \raggedright [\resultnum{v4x.x7.units.r050.delta_auc.i001.interval.i000}{0.0201266}, \resultnum{v4x.x7.units.r050.delta_auc.i001.interval.i001}{0.0491028}] \tabularnewline
\raggedright \resultnum{v4x.x7.units.r050.r}{0.5} & \raggedright \resultnum{v4x.x7.units.r050.delta_auc.i002.a}{0.003} & \raggedright \resultnum{v4x.x7.units.r050.delta_auc.i002.delta}{0.0657667} & \raggedright [\resultnum{v4x.x7.units.r050.delta_auc.i002.interval.i000}{0.0500842}, \resultnum{v4x.x7.units.r050.delta_auc.i002.interval.i001}{0.0806704}] \tabularnewline
\raggedright \resultnum{v4x.x7.units.r050.r}{0.5} & \raggedright \resultnum{v4x.x7.units.r050.delta_auc.i003.a}{0.01} & \raggedright \resultnum{v4x.x7.units.r050.delta_auc.i003.delta}{0.122845} & \raggedright [\resultnum{v4x.x7.units.r050.delta_auc.i003.interval.i000}{0.107062}, \resultnum{v4x.x7.units.r050.delta_auc.i003.interval.i001}{0.137263}] \tabularnewline
\raggedright \resultnum{v4x.x7.units.r050.r}{0.5} & \raggedright \resultnum{v4x.x7.units.r050.delta_auc.i004.a}{0.03} & \raggedright \resultnum{v4x.x7.units.r050.delta_auc.i004.delta}{0.195917} & \raggedright [\resultnum{v4x.x7.units.r050.delta_auc.i004.interval.i000}{0.181863}, \resultnum{v4x.x7.units.r050.delta_auc.i004.interval.i001}{0.210104}] \tabularnewline
\raggedright \resultnum{v4x.x7.units.r050.r}{0.5} & \raggedright \resultnum{v4x.x7.units.r050.delta_auc.i005.a}{0.1} & \raggedright \resultnum{v4x.x7.units.r050.delta_auc.i005.delta}{0.263348} & \raggedright [\resultnum{v4x.x7.units.r050.delta_auc.i005.interval.i000}{0.251523}, \resultnum{v4x.x7.units.r050.delta_auc.i005.interval.i001}{0.274276}] \tabularnewline
\raggedright \resultnum{v4x.x7.units.r050.r}{0.5} & \raggedright \resultnum{v4x.x7.units.r050.delta_auc.i006.a}{0.3} & \raggedright \resultnum{v4x.x7.units.r050.delta_auc.i006.delta}{0.127883} & \raggedright [\resultnum{v4x.x7.units.r050.delta_auc.i006.interval.i000}{0.119693}, \resultnum{v4x.x7.units.r050.delta_auc.i006.interval.i001}{0.13517}] \tabularnewline
\raggedright \resultnum{v4x.x7.units.r050.r}{0.5} & \raggedright \resultnum{v4x.x7.units.r050.delta_auc.i007.a}{0.6} & \raggedright \resultnum{v4x.x7.units.r050.delta_auc.i007.delta}{0.0262341} & \raggedright [\resultnum{v4x.x7.units.r050.delta_auc.i007.interval.i000}{0.023121}, \resultnum{v4x.x7.units.r050.delta_auc.i007.interval.i001}{0.0291499}] \tabularnewline
\raggedright \resultnum{v4x.x7.units.r055.r}{0.55} & \raggedright \resultnum{v4x.x7.units.r055.delta_auc.i000.a}{0} & \raggedright \resultnum{v4x.x7.units.r055.delta_auc.i000.delta}{0.00758806} & \raggedright [$\resultnum{v4x.x7.units.r055.delta_auc.i000.interval.i000}{-0.00760962}$, \resultnum{v4x.x7.units.r055.delta_auc.i000.interval.i001}{0.0251054}] \tabularnewline
\raggedright \resultnum{v4x.x7.units.r055.r}{0.55} & \raggedright \resultnum{v4x.x7.units.r055.delta_auc.i001.a}{0.001} & \raggedright \resultnum{v4x.x7.units.r055.delta_auc.i001.delta}{0.0321082} & \raggedright [\resultnum{v4x.x7.units.r055.delta_auc.i001.interval.i000}{0.0167819}, \resultnum{v4x.x7.units.r055.delta_auc.i001.interval.i001}{0.049348}] \tabularnewline
\raggedright \resultnum{v4x.x7.units.r055.r}{0.55} & \raggedright \resultnum{v4x.x7.units.r055.delta_auc.i002.a}{0.003} & \raggedright \resultnum{v4x.x7.units.r055.delta_auc.i002.delta}{0.0608152} & \raggedright [\resultnum{v4x.x7.units.r055.delta_auc.i002.interval.i000}{0.0450501}, \resultnum{v4x.x7.units.r055.delta_auc.i002.interval.i001}{0.0777038}] \tabularnewline
\raggedright \resultnum{v4x.x7.units.r055.r}{0.55} & \raggedright \resultnum{v4x.x7.units.r055.delta_auc.i003.a}{0.01} & \raggedright \resultnum{v4x.x7.units.r055.delta_auc.i003.delta}{0.117959} & \raggedright [\resultnum{v4x.x7.units.r055.delta_auc.i003.interval.i000}{0.102271}, \resultnum{v4x.x7.units.r055.delta_auc.i003.interval.i001}{0.133239}] \tabularnewline
\raggedright \resultnum{v4x.x7.units.r055.r}{0.55} & \raggedright \resultnum{v4x.x7.units.r055.delta_auc.i004.a}{0.03} & \raggedright \resultnum{v4x.x7.units.r055.delta_auc.i004.delta}{0.199158} & \raggedright [\resultnum{v4x.x7.units.r055.delta_auc.i004.interval.i000}{0.185136}, \resultnum{v4x.x7.units.r055.delta_auc.i004.interval.i001}{0.212597}] \tabularnewline
\raggedright \resultnum{v4x.x7.units.r055.r}{0.55} & \raggedright \resultnum{v4x.x7.units.r055.delta_auc.i005.a}{0.1} & \raggedright \resultnum{v4x.x7.units.r055.delta_auc.i005.delta}{0.24993} & \raggedright [\resultnum{v4x.x7.units.r055.delta_auc.i005.interval.i000}{0.238658}, \resultnum{v4x.x7.units.r055.delta_auc.i005.interval.i001}{0.260552}] \tabularnewline
\raggedright \resultnum{v4x.x7.units.r055.r}{0.55} & \raggedright \resultnum{v4x.x7.units.r055.delta_auc.i006.a}{0.3} & \raggedright \resultnum{v4x.x7.units.r055.delta_auc.i006.delta}{0.106727} & \raggedright [\resultnum{v4x.x7.units.r055.delta_auc.i006.interval.i000}{0.100005}, \resultnum{v4x.x7.units.r055.delta_auc.i006.interval.i001}{0.113984}] \tabularnewline
\raggedright \resultnum{v4x.x7.units.r055.r}{0.55} & \raggedright \resultnum{v4x.x7.units.r055.delta_auc.i007.a}{0.6} & \raggedright \resultnum{v4x.x7.units.r055.delta_auc.i007.delta}{0.0156547} & \raggedright [\resultnum{v4x.x7.units.r055.delta_auc.i007.interval.i000}{0.0134675}, \resultnum{v4x.x7.units.r055.delta_auc.i007.interval.i001}{0.0178096}] \tabularnewline
\raggedright \resultnum{v4x.x7.units.r060.r}{0.6} & \raggedright \resultnum{v4x.x7.units.r060.delta_auc.i000.a}{0} & \raggedright $\resultnum{v4x.x7.units.r060.delta_auc.i000.delta}{-0.000774}$ & \raggedright [$\resultnum{v4x.x7.units.r060.delta_auc.i000.interval.i000}{-0.0158726}$, \resultnum{v4x.x7.units.r060.delta_auc.i000.interval.i001}{0.014291}] \tabularnewline
\raggedright \resultnum{v4x.x7.units.r060.r}{0.6} & \raggedright \resultnum{v4x.x7.units.r060.delta_auc.i001.a}{0.001} & \raggedright \resultnum{v4x.x7.units.r060.delta_auc.i001.delta}{0.0593226} & \raggedright [\resultnum{v4x.x7.units.r060.delta_auc.i001.interval.i000}{0.0428861}, \resultnum{v4x.x7.units.r060.delta_auc.i001.interval.i001}{0.0755278}] \tabularnewline
\raggedright \resultnum{v4x.x7.units.r060.r}{0.6} & \raggedright \resultnum{v4x.x7.units.r060.delta_auc.i002.a}{0.003} & \raggedright \resultnum{v4x.x7.units.r060.delta_auc.i002.delta}{0.100116} & \raggedright [\resultnum{v4x.x7.units.r060.delta_auc.i002.interval.i000}{0.084005}, \resultnum{v4x.x7.units.r060.delta_auc.i002.interval.i001}{0.115913}] \tabularnewline
\raggedright \resultnum{v4x.x7.units.r060.r}{0.6} & \raggedright \resultnum{v4x.x7.units.r060.delta_auc.i003.a}{0.01} & \raggedright \resultnum{v4x.x7.units.r060.delta_auc.i003.delta}{0.164403} & \raggedright [\resultnum{v4x.x7.units.r060.delta_auc.i003.interval.i000}{0.149312}, \resultnum{v4x.x7.units.r060.delta_auc.i003.interval.i001}{0.179613}] \tabularnewline
\raggedright \resultnum{v4x.x7.units.r060.r}{0.6} & \raggedright \resultnum{v4x.x7.units.r060.delta_auc.i004.a}{0.03} & \raggedright \resultnum{v4x.x7.units.r060.delta_auc.i004.delta}{0.234533} & \raggedright [\resultnum{v4x.x7.units.r060.delta_auc.i004.interval.i000}{0.22152}, \resultnum{v4x.x7.units.r060.delta_auc.i004.interval.i001}{0.24819}] \tabularnewline
\raggedright \resultnum{v4x.x7.units.r060.r}{0.6} & \raggedright \resultnum{v4x.x7.units.r060.delta_auc.i005.a}{0.1} & \raggedright \resultnum{v4x.x7.units.r060.delta_auc.i005.delta}{0.243493} & \raggedright [\resultnum{v4x.x7.units.r060.delta_auc.i005.interval.i000}{0.233428}, \resultnum{v4x.x7.units.r060.delta_auc.i005.interval.i001}{0.253997}] \tabularnewline
\raggedright \resultnum{v4x.x7.units.r060.r}{0.6} & \raggedright \resultnum{v4x.x7.units.r060.delta_auc.i006.a}{0.3} & \raggedright \resultnum{v4x.x7.units.r060.delta_auc.i006.delta}{0.0777485} & \raggedright [\resultnum{v4x.x7.units.r060.delta_auc.i006.interval.i000}{0.0722064}, \resultnum{v4x.x7.units.r060.delta_auc.i006.interval.i001}{0.0835989}] \tabularnewline
\raggedright \resultnum{v4x.x7.units.r060.r}{0.6} & \raggedright \resultnum{v4x.x7.units.r060.delta_auc.i007.a}{0.6} & \raggedright \resultnum{v4x.x7.units.r060.delta_auc.i007.delta}{0.00729975} & \raggedright [\resultnum{v4x.x7.units.r060.delta_auc.i007.interval.i000}{0.00608993}, \resultnum{v4x.x7.units.r060.delta_auc.i007.interval.i001}{0.00857353}] \tabularnewline
\raggedright \resultnum{v4x.x7.units.r060.r}{0.6} & \raggedright \resultnum{v4x.x7.units.r060.delta_auc.i008.a}{0.9} & \raggedright \resultnum{v4x.x7.units.r060.delta_auc.i008.delta}{0.000994563} & \raggedright [\resultnum{v4x.x7.units.r060.delta_auc.i008.interval.i000}{0.000692359}, \resultnum{v4x.x7.units.r060.delta_auc.i008.interval.i001}{0.00135486}] \tabularnewline
\raggedright \resultnum{v4x.x7.units.r065.r}{0.65} & \raggedright \resultnum{v4x.x7.units.r065.delta_auc.i000.a}{0} & \raggedright \resultnum{v4x.x7.units.r065.delta_auc.i000.delta}{0.013916} & \raggedright [\resultnum{v4x.x7.units.r065.delta_auc.i000.interval.i000}{0.000184975}, \resultnum{v4x.x7.units.r065.delta_auc.i000.interval.i001}{0.0274991}] \tabularnewline
\raggedright \resultnum{v4x.x7.units.r065.r}{0.65} & \raggedright \resultnum{v4x.x7.units.r065.delta_auc.i001.a}{0.001} & \raggedright \resultnum{v4x.x7.units.r065.delta_auc.i001.delta}{0.0720102} & \raggedright [\resultnum{v4x.x7.units.r065.delta_auc.i001.interval.i000}{0.0581912}, \resultnum{v4x.x7.units.r065.delta_auc.i001.interval.i001}{0.0855384}] \tabularnewline
\raggedright \resultnum{v4x.x7.units.r065.r}{0.65} & \raggedright \resultnum{v4x.x7.units.r065.delta_auc.i002.a}{0.003} & \raggedright \resultnum{v4x.x7.units.r065.delta_auc.i002.delta}{0.108291} & \raggedright [\resultnum{v4x.x7.units.r065.delta_auc.i002.interval.i000}{0.0944298}, \resultnum{v4x.x7.units.r065.delta_auc.i002.interval.i001}{0.123209}] \tabularnewline
\raggedright \resultnum{v4x.x7.units.r065.r}{0.65} & \raggedright \resultnum{v4x.x7.units.r065.delta_auc.i003.a}{0.01} & \raggedright \resultnum{v4x.x7.units.r065.delta_auc.i003.delta}{0.183008} & \raggedright [\resultnum{v4x.x7.units.r065.delta_auc.i003.interval.i000}{0.169619}, \resultnum{v4x.x7.units.r065.delta_auc.i003.interval.i001}{0.197547}] \tabularnewline
\raggedright \resultnum{v4x.x7.units.r065.r}{0.65} & \raggedright \resultnum{v4x.x7.units.r065.delta_auc.i004.a}{0.03} & \raggedright \resultnum{v4x.x7.units.r065.delta_auc.i004.delta}{0.273072} & \raggedright [\resultnum{v4x.x7.units.r065.delta_auc.i004.interval.i000}{0.260776}, \resultnum{v4x.x7.units.r065.delta_auc.i004.interval.i001}{0.286279}] \tabularnewline
\raggedright \resultnum{v4x.x7.units.r065.r}{0.65} & \raggedright \resultnum{v4x.x7.units.r065.delta_auc.i005.a}{0.1} & \raggedright \resultnum{v4x.x7.units.r065.delta_auc.i005.delta}{0.266939} & \raggedright [\resultnum{v4x.x7.units.r065.delta_auc.i005.interval.i000}{0.256875}, \resultnum{v4x.x7.units.r065.delta_auc.i005.interval.i001}{0.277086}] \tabularnewline
\raggedright \resultnum{v4x.x7.units.r065.r}{0.65} & \raggedright \resultnum{v4x.x7.units.r065.delta_auc.i006.a}{0.3} & \raggedright \resultnum{v4x.x7.units.r065.delta_auc.i006.delta}{0.0608268} & \raggedright [\resultnum{v4x.x7.units.r065.delta_auc.i006.interval.i000}{0.0562613}, \resultnum{v4x.x7.units.r065.delta_auc.i006.interval.i001}{0.0655785}] \tabularnewline
\raggedright \resultnum{v4x.x7.units.r065.r}{0.65} & \raggedright \resultnum{v4x.x7.units.r065.delta_auc.i007.a}{0.6} & \raggedright \resultnum{v4x.x7.units.r065.delta_auc.i007.delta}{0.00373294} & \raggedright [\resultnum{v4x.x7.units.r065.delta_auc.i007.interval.i000}{0.00305521}, \resultnum{v4x.x7.units.r065.delta_auc.i007.interval.i001}{0.00449177}] \tabularnewline
\raggedright \resultnum{v4x.x7.units.r065.r}{0.65} & \raggedright \resultnum{v4x.x7.units.r065.delta_auc.i008.a}{0.9} & \raggedright \resultnum{v4x.x7.units.r065.delta_auc.i008.delta}{0.000267375} & \raggedright [\resultnum{v4x.x7.units.r065.delta_auc.i008.interval.i000}{0.000152392}, \resultnum{v4x.x7.units.r065.delta_auc.i008.interval.i001}{0.000401273}] \tabularnewline
\raggedright \resultnum{v4x.x7.units.r070.r}{0.7} & \raggedright \resultnum{v4x.x7.units.r070.delta_auc.i000.a}{0} & \raggedright \resultnum{v4x.x7.units.r070.delta_auc.i000.delta}{0.00685556} & \raggedright [$\resultnum{v4x.x7.units.r070.delta_auc.i000.interval.i000}{-0.00931754}$, \resultnum{v4x.x7.units.r070.delta_auc.i000.interval.i001}{0.0222272}] \tabularnewline
\raggedright \resultnum{v4x.x7.units.r070.r}{0.7} & \raggedright \resultnum{v4x.x7.units.r070.delta_auc.i001.a}{0.001} & \raggedright \resultnum{v4x.x7.units.r070.delta_auc.i001.delta}{0.0660752} & \raggedright [\resultnum{v4x.x7.units.r070.delta_auc.i001.interval.i000}{0.0501335}, \resultnum{v4x.x7.units.r070.delta_auc.i001.interval.i001}{0.083557}] \tabularnewline
\raggedright \resultnum{v4x.x7.units.r070.r}{0.7} & \raggedright \resultnum{v4x.x7.units.r070.delta_auc.i002.a}{0.003} & \raggedright \resultnum{v4x.x7.units.r070.delta_auc.i002.delta}{0.115053} & \raggedright [\resultnum{v4x.x7.units.r070.delta_auc.i002.interval.i000}{0.0991056}, \resultnum{v4x.x7.units.r070.delta_auc.i002.interval.i001}{0.131179}] \tabularnewline
\raggedright \resultnum{v4x.x7.units.r070.r}{0.7} & \raggedright \resultnum{v4x.x7.units.r070.delta_auc.i003.a}{0.01} & \raggedright \resultnum{v4x.x7.units.r070.delta_auc.i003.delta}{0.190072} & \raggedright [\resultnum{v4x.x7.units.r070.delta_auc.i003.interval.i000}{0.175386}, \resultnum{v4x.x7.units.r070.delta_auc.i003.interval.i001}{0.204146}] \tabularnewline
\raggedright \resultnum{v4x.x7.units.r070.r}{0.7} & \raggedright \resultnum{v4x.x7.units.r070.delta_auc.i004.a}{0.03} & \raggedright \resultnum{v4x.x7.units.r070.delta_auc.i004.delta}{0.271992} & \raggedright [\resultnum{v4x.x7.units.r070.delta_auc.i004.interval.i000}{0.259088}, \resultnum{v4x.x7.units.r070.delta_auc.i004.interval.i001}{0.284457}] \tabularnewline
\raggedright \resultnum{v4x.x7.units.r070.r}{0.7} & \raggedright \resultnum{v4x.x7.units.r070.delta_auc.i005.a}{0.1} & \raggedright \resultnum{v4x.x7.units.r070.delta_auc.i005.delta}{0.238399} & \raggedright [\resultnum{v4x.x7.units.r070.delta_auc.i005.interval.i000}{0.228165}, \resultnum{v4x.x7.units.r070.delta_auc.i005.interval.i001}{0.248944}] \tabularnewline
\raggedright \resultnum{v4x.x7.units.r070.r}{0.7} & \raggedright \resultnum{v4x.x7.units.r070.delta_auc.i006.a}{0.3} & \raggedright \resultnum{v4x.x7.units.r070.delta_auc.i006.delta}{0.0448589} & \raggedright [\resultnum{v4x.x7.units.r070.delta_auc.i006.interval.i000}{0.0406897}, \resultnum{v4x.x7.units.r070.delta_auc.i006.interval.i001}{0.0491794}] \tabularnewline
\raggedright \resultnum{v4x.x7.units.r070.r}{0.7} & \raggedright \resultnum{v4x.x7.units.r070.delta_auc.i007.a}{0.6} & \raggedright \resultnum{v4x.x7.units.r070.delta_auc.i007.delta}{0.00157913} & \raggedright [\resultnum{v4x.x7.units.r070.delta_auc.i007.interval.i000}{0.00117372}, \resultnum{v4x.x7.units.r070.delta_auc.i007.interval.i001}{0.00199335}] \tabularnewline
\raggedright \resultnum{v4x.x7.units.r070.r}{0.7} & \raggedright \resultnum{v4x.x7.units.r070.delta_auc.i008.a}{0.9} & \raggedright \resultnum{v4x.x7.units.r070.delta_auc.i008.delta}{5.98125e-05} & \raggedright [\resultnum{v4x.x7.units.r070.delta_auc.i008.interval.i000}{2.64266e-05}, \resultnum{v4x.x7.units.r070.delta_auc.i008.interval.i001}{0.000100125}] \tabularnewline
\raggedright \resultnum{v4x.x7.units.r075.r}{0.75} & \raggedright \resultnum{v4x.x7.units.r075.delta_auc.i000.a}{0} & \raggedright \resultnum{v4x.x7.units.r075.delta_auc.i000.delta}{0.00455031} & \raggedright [$\resultnum{v4x.x7.units.r075.delta_auc.i000.interval.i000}{-0.0126181}$, \resultnum{v4x.x7.units.r075.delta_auc.i000.interval.i001}{0.0212804}] \tabularnewline
\raggedright \resultnum{v4x.x7.units.r075.r}{0.75} & \raggedright \resultnum{v4x.x7.units.r075.delta_auc.i001.a}{0.001} & \raggedright \resultnum{v4x.x7.units.r075.delta_auc.i001.delta}{0.0790354} & \raggedright [\resultnum{v4x.x7.units.r075.delta_auc.i001.interval.i000}{0.0624958}, \resultnum{v4x.x7.units.r075.delta_auc.i001.interval.i001}{0.0945164}] \tabularnewline
\raggedright \resultnum{v4x.x7.units.r075.r}{0.75} & \raggedright \resultnum{v4x.x7.units.r075.delta_auc.i002.a}{0.003} & \raggedright \resultnum{v4x.x7.units.r075.delta_auc.i002.delta}{0.132473} & \raggedright [\resultnum{v4x.x7.units.r075.delta_auc.i002.interval.i000}{0.116755}, \resultnum{v4x.x7.units.r075.delta_auc.i002.interval.i001}{0.147707}] \tabularnewline
\raggedright \resultnum{v4x.x7.units.r075.r}{0.75} & \raggedright \resultnum{v4x.x7.units.r075.delta_auc.i003.a}{0.01} & \raggedright \resultnum{v4x.x7.units.r075.delta_auc.i003.delta}{0.216569} & \raggedright [\resultnum{v4x.x7.units.r075.delta_auc.i003.interval.i000}{0.201888}, \resultnum{v4x.x7.units.r075.delta_auc.i003.interval.i001}{0.230765}] \tabularnewline
\raggedright \resultnum{v4x.x7.units.r075.r}{0.75} & \raggedright \resultnum{v4x.x7.units.r075.delta_auc.i004.a}{0.03} & \raggedright \resultnum{v4x.x7.units.r075.delta_auc.i004.delta}{0.293496} & \raggedright [\resultnum{v4x.x7.units.r075.delta_auc.i004.interval.i000}{0.281177}, \resultnum{v4x.x7.units.r075.delta_auc.i004.interval.i001}{0.305564}] \tabularnewline
\raggedright \resultnum{v4x.x7.units.r075.r}{0.75} & \raggedright \resultnum{v4x.x7.units.r075.delta_auc.i005.a}{0.1} & \raggedright \resultnum{v4x.x7.units.r075.delta_auc.i005.delta}{0.22565} & \raggedright [\resultnum{v4x.x7.units.r075.delta_auc.i005.interval.i000}{0.215622}, \resultnum{v4x.x7.units.r075.delta_auc.i005.interval.i001}{0.235425}] \tabularnewline
\raggedright \resultnum{v4x.x7.units.r075.r}{0.75} & \raggedright \resultnum{v4x.x7.units.r075.delta_auc.i006.a}{0.3} & \raggedright \resultnum{v4x.x7.units.r075.delta_auc.i006.delta}{0.030858} & \raggedright [\resultnum{v4x.x7.units.r075.delta_auc.i006.interval.i000}{0.0276905}, \resultnum{v4x.x7.units.r075.delta_auc.i006.interval.i001}{0.0339923}] \tabularnewline
\raggedright \resultnum{v4x.x7.units.r075.r}{0.75} & \raggedright \resultnum{v4x.x7.units.r075.delta_auc.i007.a}{0.6} & \raggedright \resultnum{v4x.x7.units.r075.delta_auc.i007.delta}{0.000548313} & \raggedright [\resultnum{v4x.x7.units.r075.delta_auc.i007.interval.i000}{0.000386998}, \resultnum{v4x.x7.units.r075.delta_auc.i007.interval.i001}{0.000740159}] \tabularnewline
\raggedright \resultnum{v4x.x7.units.r075.r}{0.75} & \raggedright \resultnum{v4x.x7.units.r075.delta_auc.i008.a}{0.9} & \raggedright \resultnum{v4x.x7.units.r075.delta_auc.i008.delta}{2.49375e-05} & \raggedright [\resultnum{v4x.x7.units.r075.delta_auc.i008.interval.i000}{8.49687e-06}, \resultnum{v4x.x7.units.r075.delta_auc.i008.interval.i001}{4.77531e-05}] \tabularnewline
\botrule
\end{tabular}
\end{table*}
\begin{table*}[htbp]
\setlength{\tabcolsep}{2pt}
\addtocounter{table}{-\csname @ne\endcsname}
\providecommand{\theHtable}{\thetable}
\renewcommand{\theHtable}{v4x-x7-deltas-part1}
\caption{Expanded-feature classifier (X7): AUC differences (continued, part 2 of 2)}\label{tab:v4x-x7-deltas-part1}
\centering\footnotesize
\begin{tabular}{@{}p{0.140\textwidth}p{0.140\textwidth}p{0.220\textwidth}p{0.400\textwidth}@{}}
\toprule
\raggedright $r$ & \raggedright $a$ & \raggedright $\Delta$AUC & \raggedright \mbox{Interval} \tabularnewline
\midrule
\raggedright \resultnum{v4x.x7.units.r080.r}{0.8} & \raggedright \resultnum{v4x.x7.units.r080.delta_auc.i000.a}{0} & \raggedright $\resultnum{v4x.x7.units.r080.delta_auc.i000.delta}{-0.0188126}$ & \raggedright [$\resultnum{v4x.x7.units.r080.delta_auc.i000.interval.i000}{-0.0341549}$, $\resultnum{v4x.x7.units.r080.delta_auc.i000.interval.i001}{-0.00476359}$] \tabularnewline
\raggedright \resultnum{v4x.x7.units.r080.r}{0.8} & \raggedright \resultnum{v4x.x7.units.r080.delta_auc.i001.a}{0.001} & \raggedright \resultnum{v4x.x7.units.r080.delta_auc.i001.delta}{0.080784} & \raggedright [\resultnum{v4x.x7.units.r080.delta_auc.i001.interval.i000}{0.0627973}, \resultnum{v4x.x7.units.r080.delta_auc.i001.interval.i001}{0.0980934}] \tabularnewline
\raggedright \resultnum{v4x.x7.units.r080.r}{0.8} & \raggedright \resultnum{v4x.x7.units.r080.delta_auc.i002.a}{0.003} & \raggedright \resultnum{v4x.x7.units.r080.delta_auc.i002.delta}{0.150712} & \raggedright [\resultnum{v4x.x7.units.r080.delta_auc.i002.interval.i000}{0.13469}, \resultnum{v4x.x7.units.r080.delta_auc.i002.interval.i001}{0.167658}] \tabularnewline
\raggedright \resultnum{v4x.x7.units.r080.r}{0.8} & \raggedright \resultnum{v4x.x7.units.r080.delta_auc.i003.a}{0.01} & \raggedright \resultnum{v4x.x7.units.r080.delta_auc.i003.delta}{0.254148} & \raggedright [\resultnum{v4x.x7.units.r080.delta_auc.i003.interval.i000}{0.239998}, \resultnum{v4x.x7.units.r080.delta_auc.i003.interval.i001}{0.270179}] \tabularnewline
\raggedright \resultnum{v4x.x7.units.r080.r}{0.8} & \raggedright \resultnum{v4x.x7.units.r080.delta_auc.i004.a}{0.03} & \raggedright \resultnum{v4x.x7.units.r080.delta_auc.i004.delta}{0.327689} & \raggedright [\resultnum{v4x.x7.units.r080.delta_auc.i004.interval.i000}{0.31552}, \resultnum{v4x.x7.units.r080.delta_auc.i004.interval.i001}{0.340956}] \tabularnewline
\raggedright \resultnum{v4x.x7.units.r080.r}{0.8} & \raggedright \resultnum{v4x.x7.units.r080.delta_auc.i005.a}{0.1} & \raggedright \resultnum{v4x.x7.units.r080.delta_auc.i005.delta}{0.228639} & \raggedright [\resultnum{v4x.x7.units.r080.delta_auc.i005.interval.i000}{0.219358}, \resultnum{v4x.x7.units.r080.delta_auc.i005.interval.i001}{0.239007}] \tabularnewline
\raggedright \resultnum{v4x.x7.units.r080.r}{0.8} & \raggedright \resultnum{v4x.x7.units.r080.delta_auc.i006.a}{0.3} & \raggedright \resultnum{v4x.x7.units.r080.delta_auc.i006.delta}{0.0283355} & \raggedright [\resultnum{v4x.x7.units.r080.delta_auc.i006.interval.i000}{0.0253965}, \resultnum{v4x.x7.units.r080.delta_auc.i006.interval.i001}{0.0312821}] \tabularnewline
\raggedright \resultnum{v4x.x7.units.r080.r}{0.8} & \raggedright \resultnum{v4x.x7.units.r080.delta_auc.i007.a}{0.6} & \raggedright \resultnum{v4x.x7.units.r080.delta_auc.i007.delta}{0.0004005} & \raggedright [\resultnum{v4x.x7.units.r080.delta_auc.i007.interval.i000}{0.000254466}, \resultnum{v4x.x7.units.r080.delta_auc.i007.interval.i001}{0.000560202}] \tabularnewline
\raggedright \resultnum{v4x.x7.units.r080.r}{0.8} & \raggedright \resultnum{v4x.x7.units.r080.delta_auc.i008.a}{0.9} & \raggedright \resultnum{v4x.x7.units.r080.delta_auc.i008.delta}{7.5e-06} & \raggedright [\resultnum{v4x.x7.units.r080.delta_auc.i008.interval.i000}{1.25e-06}, \resultnum{v4x.x7.units.r080.delta_auc.i008.interval.i001}{1.62531e-05}] \tabularnewline
\raggedright \resultnum{v4x.x7.units.r085.r}{0.85} & \raggedright \resultnum{v4x.x7.units.r085.delta_auc.i000.a}{0} & \raggedright \resultnum{v4x.x7.units.r085.delta_auc.i000.delta}{0.0126361} & \raggedright [$\resultnum{v4x.x7.units.r085.delta_auc.i000.interval.i000}{-0.00410872}$, \resultnum{v4x.x7.units.r085.delta_auc.i000.interval.i001}{0.0274255}] \tabularnewline
\raggedright \resultnum{v4x.x7.units.r085.r}{0.85} & \raggedright \resultnum{v4x.x7.units.r085.delta_auc.i001.a}{0.001} & \raggedright \resultnum{v4x.x7.units.r085.delta_auc.i001.delta}{0.105782} & \raggedright [\resultnum{v4x.x7.units.r085.delta_auc.i001.interval.i000}{0.0903723}, \resultnum{v4x.x7.units.r085.delta_auc.i001.interval.i001}{0.121331}] \tabularnewline
\raggedright \resultnum{v4x.x7.units.r085.r}{0.85} & \raggedright \resultnum{v4x.x7.units.r085.delta_auc.i002.a}{0.003} & \raggedright \resultnum{v4x.x7.units.r085.delta_auc.i002.delta}{0.181696} & \raggedright [\resultnum{v4x.x7.units.r085.delta_auc.i002.interval.i000}{0.16616}, \resultnum{v4x.x7.units.r085.delta_auc.i002.interval.i001}{0.196099}] \tabularnewline
\raggedright \resultnum{v4x.x7.units.r085.r}{0.85} & \raggedright \resultnum{v4x.x7.units.r085.delta_auc.i003.a}{0.01} & \raggedright \resultnum{v4x.x7.units.r085.delta_auc.i003.delta}{0.286376} & \raggedright [\resultnum{v4x.x7.units.r085.delta_auc.i003.interval.i000}{0.272406}, \resultnum{v4x.x7.units.r085.delta_auc.i003.interval.i001}{0.300246}] \tabularnewline
\raggedright \resultnum{v4x.x7.units.r085.r}{0.85} & \raggedright \resultnum{v4x.x7.units.r085.delta_auc.i004.a}{0.03} & \raggedright \resultnum{v4x.x7.units.r085.delta_auc.i004.delta}{0.338956} & \raggedright [\resultnum{v4x.x7.units.r085.delta_auc.i004.interval.i000}{0.326453}, \resultnum{v4x.x7.units.r085.delta_auc.i004.interval.i001}{0.351157}] \tabularnewline
\raggedright \resultnum{v4x.x7.units.r085.r}{0.85} & \raggedright \resultnum{v4x.x7.units.r085.delta_auc.i005.a}{0.1} & \raggedright \resultnum{v4x.x7.units.r085.delta_auc.i005.delta}{0.205532} & \raggedright [\resultnum{v4x.x7.units.r085.delta_auc.i005.interval.i000}{0.196426}, \resultnum{v4x.x7.units.r085.delta_auc.i005.interval.i001}{0.215337}] \tabularnewline
\raggedright \resultnum{v4x.x7.units.r085.r}{0.85} & \raggedright \resultnum{v4x.x7.units.r085.delta_auc.i006.a}{0.3} & \raggedright \resultnum{v4x.x7.units.r085.delta_auc.i006.delta}{0.0143213} & \raggedright [\resultnum{v4x.x7.units.r085.delta_auc.i006.interval.i000}{0.0125952}, \resultnum{v4x.x7.units.r085.delta_auc.i006.interval.i001}{0.0162084}] \tabularnewline
\raggedright \resultnum{v4x.x7.units.r085.r}{0.85} & \raggedright \resultnum{v4x.x7.units.r085.delta_auc.i007.a}{0.6} & \raggedright \resultnum{v4x.x7.units.r085.delta_auc.i007.delta}{9.28125e-05} & \raggedright [\resultnum{v4x.x7.units.r085.delta_auc.i007.interval.i000}{5.01234e-05}, \resultnum{v4x.x7.units.r085.delta_auc.i007.interval.i001}{0.000143069}] \tabularnewline
\raggedright \resultnum{v4x.x7.units.r085.r}{0.85} & \raggedright \resultnum{v4x.x7.units.r085.delta_auc.i008.a}{0.9} & \raggedright \resultnum{v4x.x7.units.r085.delta_auc.i008.delta}{2.875e-06} & \raggedright [\resultnum{v4x.x7.units.r085.delta_auc.i008.interval.i000}{0}, \resultnum{v4x.x7.units.r085.delta_auc.i008.interval.i001}{9.50312e-06}] \tabularnewline
\raggedright \resultnum{v4x.x7.units.r090.r}{0.9} & \raggedright \resultnum{v4x.x7.units.r090.delta_auc.i000.a}{0} & \raggedright $\resultnum{v4x.x7.units.r090.delta_auc.i000.delta}{-0.008023}$ & \raggedright [$\resultnum{v4x.x7.units.r090.delta_auc.i000.interval.i000}{-0.024507}$, \resultnum{v4x.x7.units.r090.delta_auc.i000.interval.i001}{0.00757633}] \tabularnewline
\raggedright \resultnum{v4x.x7.units.r090.r}{0.9} & \raggedright \resultnum{v4x.x7.units.r090.delta_auc.i001.a}{0.001} & \raggedright \resultnum{v4x.x7.units.r090.delta_auc.i001.delta}{0.119499} & \raggedright [\resultnum{v4x.x7.units.r090.delta_auc.i001.interval.i000}{0.10298}, \resultnum{v4x.x7.units.r090.delta_auc.i001.interval.i001}{0.136394}] \tabularnewline
\raggedright \resultnum{v4x.x7.units.r090.r}{0.9} & \raggedright \resultnum{v4x.x7.units.r090.delta_auc.i002.a}{0.003} & \raggedright \resultnum{v4x.x7.units.r090.delta_auc.i002.delta}{0.20338} & \raggedright [\resultnum{v4x.x7.units.r090.delta_auc.i002.interval.i000}{0.187203}, \resultnum{v4x.x7.units.r090.delta_auc.i002.interval.i001}{0.220323}] \tabularnewline
\raggedright \resultnum{v4x.x7.units.r090.r}{0.9} & \raggedright \resultnum{v4x.x7.units.r090.delta_auc.i003.a}{0.01} & \raggedright \resultnum{v4x.x7.units.r090.delta_auc.i003.delta}{0.304275} & \raggedright [\resultnum{v4x.x7.units.r090.delta_auc.i003.interval.i000}{0.290494}, \resultnum{v4x.x7.units.r090.delta_auc.i003.interval.i001}{0.319104}] \tabularnewline
\raggedright \resultnum{v4x.x7.units.r090.r}{0.9} & \raggedright \resultnum{v4x.x7.units.r090.delta_auc.i004.a}{0.03} & \raggedright \resultnum{v4x.x7.units.r090.delta_auc.i004.delta}{0.335997} & \raggedright [\resultnum{v4x.x7.units.r090.delta_auc.i004.interval.i000}{0.3244}, \resultnum{v4x.x7.units.r090.delta_auc.i004.interval.i001}{0.347855}] \tabularnewline
\raggedright \resultnum{v4x.x7.units.r090.r}{0.9} & \raggedright \resultnum{v4x.x7.units.r090.delta_auc.i005.a}{0.1} & \raggedright \resultnum{v4x.x7.units.r090.delta_auc.i005.delta}{0.174195} & \raggedright [\resultnum{v4x.x7.units.r090.delta_auc.i005.interval.i000}{0.165482}, \resultnum{v4x.x7.units.r090.delta_auc.i005.interval.i001}{0.182926}] \tabularnewline
\raggedright \resultnum{v4x.x7.units.r090.r}{0.9} & \raggedright \resultnum{v4x.x7.units.r090.delta_auc.i006.a}{0.3} & \raggedright \resultnum{v4x.x7.units.r090.delta_auc.i006.delta}{0.0109932} & \raggedright [\resultnum{v4x.x7.units.r090.delta_auc.i006.interval.i000}{0.0095784}, \resultnum{v4x.x7.units.r090.delta_auc.i006.interval.i001}{0.0126732}] \tabularnewline
\raggedright \resultnum{v4x.x7.units.r090.r}{0.9} & \raggedright \resultnum{v4x.x7.units.r090.delta_auc.i007.a}{0.6} & \raggedright \resultnum{v4x.x7.units.r090.delta_auc.i007.delta}{3.4875e-05} & \raggedright [\resultnum{v4x.x7.units.r090.delta_auc.i007.interval.i000}{1.5e-05}, \resultnum{v4x.x7.units.r090.delta_auc.i007.interval.i001}{5.91984e-05}] \tabularnewline
\raggedright \resultnum{v4x.x7.units.r090.r}{0.9} & \raggedright \resultnum{v4x.x7.units.r090.delta_auc.i008.a}{0.9} & \raggedright \resultnum{v4x.x7.units.r090.delta_auc.i008.delta}{3.75e-07} & \raggedright [\resultnum{v4x.x7.units.r090.delta_auc.i008.interval.i000}{0}, \resultnum{v4x.x7.units.r090.delta_auc.i008.interval.i001}{1.5e-06}] \tabularnewline
\botrule
\end{tabular}
\end{table*}

\clearpage

\subsection{Readings beside the released boundary}\label{sec:x-readings}\phantomsection\label{supp:x1b-release}

\paragraph{Available local readings.}
After the results became available, the cell of the released first boundary against B gave an X3 cell with an available contrast at \resultnum{v4x.comparisons.count.x3.available}{9} signal shares. Both single-panel rates were available at \resultnum{v4x.comparisons.count.x4.available}{3} shares (Tables~\ref{tab:v4x-comparisons} and~\ref{tab:v4x-selected-comparisons}). Selection used the point value of the first boundary and did not interpolate strengths or select by the sustained boundary. A missing or gap-censored first boundary selects no cell, and a finite first boundary with a censored interval end selects its exact cell with that limitation retained. The original X3 label, paired interval, X4 denominators, size flags and quantity labels remain subject to the same local reading limits.

\paragraph{Partial and unavailable readings.}
For X4, \resultnum{v4x.comparisons.count.x4.partial}{6} signal shares had a partial reading (the oracle rate only, outside the declared domain of the fitted-panel procedure), and \resultnum{v4x.comparisons.count.x4.unavailable}{0} had none. For X3, \resultnum{v4x.comparisons.count.x3.unavailable}{0} shares had no reading (Table~\ref{tab:v4x-comparisons}). Local availability does not alter the validity of the main curve, which was valid.

\paragraph{Twin and expanded-feature boundaries.}
The twin's first boundary against B was available at \resultnum{v4x.comparisons.count.x1b.available}{1} signal share; at the other \resultnum{v4x.comparisons.count.x1b.unavailable}{8} it was gap-censored or not attained on the grid (Tables~\ref{tab:v4x-comparisons} and~\ref{tab:v4x-boundary-comparisons}). Table~\ref{tab:v4x-x1b-boundaries} keeps every declared signal share with first and sustained records, intervals and censoring. Table~\ref{tab:v4x-boundary-comparisons} places the twin and both LDAs of the paired X7 run beside the released $a^*(r)$, retaining their own intervals, statuses and censoring in the accompanying tables. These displays add no between-curve test across the different seed sets and do not isolate network dependence.

\paragraph{Two re-runs of the frozen three-feature pipeline (read after the results became available).}
Two exploratory analyses re-ran the frozen three-feature pipeline against B on their own seeds, disjoint from the confirmatory ones, although neither was declared as a replication and no comparison between curves was declared. The oracle single-panel test (Section~\ref{sec:x4}) refit B with the frozen matching in each cell, retrained the frozen discriminant and computed a reference AUC on \resultnum{v4x.x4.sizes.check_prose}{4,000} new panels per arm (\resultnum{v4x.x4a.gaps}{4} fits failed, all at $a\ge0.60$); the frozen-feature arm of the expanded-feature analysis (Section~\ref{sec:x7}) applied the whole frozen curve code, bands included, to its own panels. Both place the strength at which the three features reach the target above 0.03 at every share. No oracle reference AUC at $a\le0.03$ reached 0.70, nor did the upper end of its interval; the largest, at $a=0.03$, was \resultnum{j.x4a.a0300.auc_ref.max}{0.638} at $r=\resultnum{j.x4a.a0300.auc_ref.max.r}{0.90}$. Every frozen-feature first boundary of the expanded-feature run lay, with its whole interval, above 0.03, so its region was large at all nine shares (Table~\ref{tab:v4x-x7-boundaries}). The point boundary is less stable where the AUC at $a=0.10$ lies near 0.70. At $r=0.55$, where the released point AUC at $a=0.10$ was \resultnum{j.b.r055.a1000.auc}{0.702}, the frozen-feature run gave a first boundary of \resultnum{j.x7.frozen_first.r055}{0.30}, inside the released interval [\resultnum{j.b.r055.first.lower}{0.10}, \resultnum{j.b.r055.first.upper}{0.30}], and the oracle reference AUC at $a=0.10$ was \resultnum{j.x4a.r055.a1000.auc_ref}{0.683} (interval \resultnum{j.x4a.r055.a1000.auc_ref.lower}{0.672}--\resultnum{j.x4a.r055.a1000.auc_ref.upper}{0.694}). At $r=0.65$ the frozen-feature interval spanned one grid step where the released interval is a single grid point (Section~\ref{sec:record-gate}). Both re-runs therefore reproduced the region label of the main result on independent seeds, while its point value can move by one grid step where the AUC at $a=0.10$ is close to the target; this is evidence from two exploratory re-runs, not a demonstration that the label is robust.

\paragraph{Misspecification beside the released curve.}
For \resultnum{v4x.comparisons.count.x5.equivalent.available}{35} available X5 condition--signal-share readings, Table~\ref{tab:v4x-stress-comparisons} reports the equivalent strength: the first grid strength whose confirmatory point AUC equals or exceeds the misspecification AUC, without interpolation. A gap before the attainment left \resultnum{v4x.comparisons.count.x5.equivalent.censored}{1} reading censored at the recorded strength; \resultnum{v4x.comparisons.count.x5.equivalent.above-grid}{0} lay above the complete grid. At the released first boundary, \resultnum{v4x.comparisons.count.x5.interval.overlap}{9} misspecification intervals overlapped the confirmatory band (touching an endpoint counts as overlap), \resultnum{v4x.comparisons.count.x5.interval.above}{10} lay above it and \resultnum{v4x.comparisons.count.x5.interval.below}{17} below it. The X5 fit gaps left \resultnum{v4x.comparisons.count.x5.equivalent.unavailable}{36} equivalent strengths and \resultnum{v4x.comparisons.count.x5.interval.unavailable}{36} interval comparisons not estimable. These local descriptive relations neither calibrate the SHP channel nor identify causal influence.

% Generated by v4_exploratory_export; edited during the journal revision
% (agent S, 2026-10-13): captions, headers and categorical labels translated
% into plain English. Every number keeps its \resultnum binding.
\begin{table*}[htbp]
\setlength{\tabcolsep}{2pt}
\caption{Readings beside the released first boundary against B, by $r$: the twin's own first boundary (X1b), the free-share contrast (X3) and the two single-panel rejection rates (X4a, X4b) at the cell of the first boundary, without interpolation. The main curve was available and valid at every $r$. Twin: gap-censored, the twin boundary cannot be declared across a fit gap; not attained on grid, the twin's point AUC stayed below 0.70 on a complete grid. The fitted-panel procedure was declared only at $r=0.50$, 0.70 and 0.90, so elsewhere the X4 reading is partial (oracle only). Machine states, recorded in Portuguese in the archive, are translated here; a missing or censored reading does not imply a null result.}\label{tab:v4x-comparisons}\label{tab:v4x-comparison-states}
\centering\footnotesize
\begin{tabular}{@{}p{0.080\textwidth}p{0.140\textwidth}p{0.220\textwidth}p{0.130\textwidth}p{0.130\textwidth}p{0.220\textwidth}@{}}
\toprule
\raggedright $r$ & \raggedright \mbox{Released} \mbox{first} \mbox{boundary} & \raggedright \mbox{Twin} \mbox{(X1b)} & \raggedright \mbox{Free} \mbox{share} \mbox{(X3)} & \raggedright \mbox{Oracle} \mbox{(X4a)} & \raggedright \mbox{Fitted} \mbox{panel} \mbox{(X4b)} \tabularnewline
\midrule
\raggedright \resultnum{v4x.comparisons.by_r.0-5.r}{0.5} & \raggedright \resultnum{v4x.comparisons.by_r.0-5.main_curve.first.point}{0.3} & \raggedright gap-censored at $a=0.6$ & \raggedright available & \raggedright available & \raggedright available \tabularnewline
\raggedright \resultnum{v4x.comparisons.by_r.0-55.r}{0.55} & \raggedright \resultnum{v4x.comparisons.by_r.0-55.main_curve.first.point}{0.1} & \raggedright gap-censored at $a=0.9$ & \raggedright available & \raggedright available & \raggedright outside the declared domain \tabularnewline
\raggedright \resultnum{v4x.comparisons.by_r.0-6.r}{0.6} & \raggedright \resultnum{v4x.comparisons.by_r.0-6.main_curve.first.point}{0.1} & \raggedright not attained on grid & \raggedright available & \raggedright available & \raggedright outside the declared domain \tabularnewline
\raggedright \resultnum{v4x.comparisons.by_r.0-65.r}{0.65} & \raggedright \resultnum{v4x.comparisons.by_r.0-65.main_curve.first.point}{0.1} & \raggedright not attained on grid & \raggedright available & \raggedright available & \raggedright outside the declared domain \tabularnewline
\raggedright \resultnum{v4x.comparisons.by_r.0-7.r}{0.7} & \raggedright \resultnum{v4x.comparisons.by_r.0-7.main_curve.first.point}{0.1} & \raggedright not attained on grid & \raggedright available & \raggedright available & \raggedright available \tabularnewline
\raggedright \resultnum{v4x.comparisons.by_r.0-75.r}{0.75} & \raggedright \resultnum{v4x.comparisons.by_r.0-75.main_curve.first.point}{0.1} & \raggedright gap-censored at $a=0.9$ & \raggedright available & \raggedright available & \raggedright outside the declared domain \tabularnewline
\raggedright \resultnum{v4x.comparisons.by_r.0-8.r}{0.8} & \raggedright \resultnum{v4x.comparisons.by_r.0-8.main_curve.first.point}{0.1} & \raggedright gap-censored at $a=0.9$ & \raggedright available & \raggedright available & \raggedright outside the declared domain \tabularnewline
\raggedright \resultnum{v4x.comparisons.by_r.0-85.r}{0.85} & \raggedright \resultnum{v4x.comparisons.by_r.0-85.main_curve.first.point}{0.1} & \raggedright not attained on grid & \raggedright available & \raggedright available & \raggedright outside the declared domain \tabularnewline
\raggedright \resultnum{v4x.comparisons.by_r.0-9.r}{0.9} & \raggedright \resultnum{v4x.comparisons.by_r.0-9.main_curve.first.point}{0.1} & \raggedright available & \raggedright available & \raggedright available & \raggedright available \tabularnewline
\botrule
\end{tabular}
\end{table*}

% Generated by v4_exploratory_export; edited during the journal revision
% (agent S, 2026-10-13): captions, headers and categorical labels translated
% into plain English. Every number keeps its \resultnum binding.
\begin{table*}[htbp]
\setlength{\tabcolsep}{2pt}
\caption{Free-share contrast (X3) and single-panel rejection rates (X4a, X4b) at the cell of the released first boundary against B. X3 reading as in Table~\ref{tab:v4x-x3} (reduced: interval entirely below zero; machine label supports-objection). Each rejection rate carries its own procedure's size flag at that $r$ (Table~\ref{tab:v4x-x4-calibration}) and, for the fitted-panel procedure, the frozen quantity label; every entry is reported as a rejection rate (Section~\ref{sec:x4}). Both procedures use the LDA direction trained against the true alternative. The selection neither estimates the confirmatory boundary nor adds a confirmatory contrast.}\label{tab:v4x-selected-comparisons}
\centering\footnotesize
\begin{tabular}{@{}p{0.068\textwidth}p{0.098\textwidth}p{0.127\textwidth}p{0.205\textwidth}p{0.195\textwidth}p{0.244\textwidth}@{}}
\toprule
\raggedright $r$ & \raggedright \mbox{X3} \mbox{reading} & \raggedright $\Delta$AUC & \raggedright \mbox{Interval} & \raggedright \mbox{Oracle} \mbox{(X4a)} \mbox{rate,} \mbox{size} & \raggedright \mbox{Fitted} \mbox{(X4b)} \mbox{rate,} \mbox{size;} \mbox{quantity} \tabularnewline
\midrule
\raggedright \resultnum{v4x.comparisons.by_r.0-5.r}{0.5} & \raggedright reduced & \raggedright $\resultnum{v4x.comparisons.by_r.0-5.x3.delta.point}{-0.27282}$ & \raggedright [$\resultnum{v4x.comparisons.by_r.0-5.x3.delta.interval_95.i000}{-0.280625}$, $\resultnum{v4x.comparisons.by_r.0-5.x3.delta.interval_95.i001}{-0.264818}$] & \raggedright \resultnum{v4x.comparisons.by_r.0-5.x4.P_or.statistics.power_or}{0.447} (compatible) & \raggedright \resultnum{v4x.comparisons.by_r.0-5.x4.P_pu.statistics.power_pu}{0.5} (above: rate) \tabularnewline
\raggedright \resultnum{v4x.comparisons.by_r.0-55.r}{0.55} & \raggedright reduced & \raggedright $\resultnum{v4x.comparisons.by_r.0-55.x3.delta.point}{-0.12614}$ & \raggedright [$\resultnum{v4x.comparisons.by_r.0-55.x3.delta.interval_95.i000}{-0.134898}$, $\resultnum{v4x.comparisons.by_r.0-55.x3.delta.interval_95.i001}{-0.117152}$] & \raggedright \resultnum{v4x.comparisons.by_r.0-55.x4.P_or.statistics.power_or}{0.171} (compatible) & \raggedright outside the declared domain \tabularnewline
\raggedright \resultnum{v4x.comparisons.by_r.0-6.r}{0.6} & \raggedright reduced & \raggedright $\resultnum{v4x.comparisons.by_r.0-6.x3.delta.point}{-0.159901}$ & \raggedright [$\resultnum{v4x.comparisons.by_r.0-6.x3.delta.interval_95.i000}{-0.165875}$, $\resultnum{v4x.comparisons.by_r.0-6.x3.delta.interval_95.i001}{-0.153723}$] & \raggedright \resultnum{v4x.comparisons.by_r.0-6.x4.P_or.statistics.power_or}{0.205} (compatible) & \raggedright outside the declared domain \tabularnewline
\raggedright \resultnum{v4x.comparisons.by_r.0-65.r}{0.65} & \raggedright reduced & \raggedright $\resultnum{v4x.comparisons.by_r.0-65.x3.delta.point}{-0.171633}$ & \raggedright [$\resultnum{v4x.comparisons.by_r.0-65.x3.delta.interval_95.i000}{-0.177105}$, $\resultnum{v4x.comparisons.by_r.0-65.x3.delta.interval_95.i001}{-0.166476}$] & \raggedright \resultnum{v4x.comparisons.by_r.0-65.x4.P_or.statistics.power_or}{0.209} (compatible) & \raggedright outside the declared domain \tabularnewline
\raggedright \resultnum{v4x.comparisons.by_r.0-7.r}{0.7} & \raggedright reduced & \raggedright $\resultnum{v4x.comparisons.by_r.0-7.x3.delta.point}{-0.178288}$ & \raggedright [$\resultnum{v4x.comparisons.by_r.0-7.x3.delta.interval_95.i000}{-0.183328}$, $\resultnum{v4x.comparisons.by_r.0-7.x3.delta.interval_95.i001}{-0.173138}$] & \raggedright \resultnum{v4x.comparisons.by_r.0-7.x4.P_or.statistics.power_or}{0.199} (above) & \raggedright \resultnum{v4x.comparisons.by_r.0-7.x4.P_pu.statistics.power_pu}{0.385} (compatible: power) \tabularnewline
\raggedright \resultnum{v4x.comparisons.by_r.0-75.r}{0.75} & \raggedright reduced & \raggedright $\resultnum{v4x.comparisons.by_r.0-75.x3.delta.point}{-0.172387}$ & \raggedright [$\resultnum{v4x.comparisons.by_r.0-75.x3.delta.interval_95.i000}{-0.177462}$, $\resultnum{v4x.comparisons.by_r.0-75.x3.delta.interval_95.i001}{-0.167671}$] & \raggedright \resultnum{v4x.comparisons.by_r.0-75.x4.P_or.statistics.power_or}{0.282} (compatible) & \raggedright outside the declared domain \tabularnewline
\raggedright \resultnum{v4x.comparisons.by_r.0-8.r}{0.8} & \raggedright reduced & \raggedright $\resultnum{v4x.comparisons.by_r.0-8.x3.delta.point}{-0.173259}$ & \raggedright [$\resultnum{v4x.comparisons.by_r.0-8.x3.delta.interval_95.i000}{-0.17741}$, $\resultnum{v4x.comparisons.by_r.0-8.x3.delta.interval_95.i001}{-0.169218}$] & \raggedright \resultnum{v4x.comparisons.by_r.0-8.x4.P_or.statistics.power_or}{0.301} (compatible) & \raggedright outside the declared domain \tabularnewline
\raggedright \resultnum{v4x.comparisons.by_r.0-85.r}{0.85} & \raggedright reduced & \raggedright $\resultnum{v4x.comparisons.by_r.0-85.x3.delta.point}{-0.194977}$ & \raggedright [$\resultnum{v4x.comparisons.by_r.0-85.x3.delta.interval_95.i000}{-0.200278}$, $\resultnum{v4x.comparisons.by_r.0-85.x3.delta.interval_95.i001}{-0.189724}$] & \raggedright \resultnum{v4x.comparisons.by_r.0-85.x4.P_or.statistics.power_or}{0.321} (compatible) & \raggedright outside the declared domain \tabularnewline
\raggedright \resultnum{v4x.comparisons.by_r.0-9.r}{0.9} & \raggedright reduced & \raggedright $\resultnum{v4x.comparisons.by_r.0-9.x3.delta.point}{-0.201556}$ & \raggedright [$\resultnum{v4x.comparisons.by_r.0-9.x3.delta.interval_95.i000}{-0.207505}$, $\resultnum{v4x.comparisons.by_r.0-9.x3.delta.interval_95.i001}{-0.194782}$] & \raggedright \resultnum{v4x.comparisons.by_r.0-9.x4.P_or.statistics.power_or}{0.368} (compatible) & \raggedright \resultnum{v4x.comparisons.by_r.0-9.x4.P_pu.statistics.power_pu}{0.415} (below: power, lower bound) \tabularnewline
\botrule
\end{tabular}
\end{table*}

% Generated by v4_exploratory_export; edited during the journal revision
% (agent S, 2026-10-13): captions, headers and categorical labels translated
% into plain English. Every number keeps its \resultnum binding.
\begin{table*}[htbp]
\setlength{\tabcolsep}{2pt}
\caption{First boundaries beside the released first boundary against B: the mean-field twin against B (X1b) and the frozen-feature and expanded-feature LDAs of the paired X7 run. Different seed sets are shown side by side without a between-curve test. Gap-censored: not declared across a fit gap; not attained on grid: point AUC below 0.70 on a complete grid. Main curve valid: the frozen validity of the main curve against B at that $r$. Intervals and censoring are in Tables~\ref{tab:v4x-x1b-boundaries} and~\ref{tab:v4x-x7-boundaries}.}\label{tab:v4x-boundary-comparisons}
\centering\footnotesize
\begin{tabular}{@{}p{0.070\textwidth}p{0.100\textwidth}p{0.230\textwidth}p{0.170\textwidth}p{0.170\textwidth}p{0.130\textwidth}@{}}
\toprule
\raggedright $r$ & \raggedright \mbox{Released} & \raggedright \mbox{Twin} \mbox{(X1b)} & \raggedright \mbox{Frozen} \mbox{features} \mbox{(X7} \mbox{run)} & \raggedright \mbox{Expanded} \mbox{features} \mbox{(X7)} & \raggedright \mbox{Main} \mbox{curve} \mbox{valid} \tabularnewline
\midrule
\raggedright \resultnum{v4x.comparisons.by_r.0-5.r}{0.5} & \raggedright \resultnum{v4x.comparisons.by_r.0-5.main_curve.first.point}{0.3} & \raggedright gap-censored at $a=0.6$ & \raggedright \resultnum{v4x.x7.units.r050.x7.first}{0.3} & \raggedright \resultnum{v4x.x7.units.r050.rich.first}{0.03} & \raggedright yes \tabularnewline
\raggedright \resultnum{v4x.comparisons.by_r.0-55.r}{0.55} & \raggedright \resultnum{v4x.comparisons.by_r.0-55.main_curve.first.point}{0.1} & \raggedright gap-censored at $a=0.9$ & \raggedright \resultnum{v4x.x7.units.r055.x7.first}{0.3} & \raggedright \resultnum{v4x.x7.units.r055.rich.first}{0.03} & \raggedright yes \tabularnewline
\raggedright \resultnum{v4x.comparisons.by_r.0-6.r}{0.6} & \raggedright \resultnum{v4x.comparisons.by_r.0-6.main_curve.first.point}{0.1} & \raggedright not attained on grid & \raggedright \resultnum{v4x.x7.units.r060.x7.first}{0.1} & \raggedright \resultnum{v4x.x7.units.r060.rich.first}{0.03} & \raggedright yes \tabularnewline
\raggedright \resultnum{v4x.comparisons.by_r.0-65.r}{0.65} & \raggedright \resultnum{v4x.comparisons.by_r.0-65.main_curve.first.point}{0.1} & \raggedright not attained on grid & \raggedright \resultnum{v4x.x7.units.r065.x7.first}{0.1} & \raggedright \resultnum{v4x.x7.units.r065.rich.first}{0.01} & \raggedright yes \tabularnewline
\raggedright \resultnum{v4x.comparisons.by_r.0-7.r}{0.7} & \raggedright \resultnum{v4x.comparisons.by_r.0-7.main_curve.first.point}{0.1} & \raggedright not attained on grid & \raggedright \resultnum{v4x.x7.units.r070.x7.first}{0.1} & \raggedright \resultnum{v4x.x7.units.r070.rich.first}{0.01} & \raggedright yes \tabularnewline
\raggedright \resultnum{v4x.comparisons.by_r.0-75.r}{0.75} & \raggedright \resultnum{v4x.comparisons.by_r.0-75.main_curve.first.point}{0.1} & \raggedright gap-censored at $a=0.9$ & \raggedright \resultnum{v4x.x7.units.r075.x7.first}{0.1} & \raggedright \resultnum{v4x.x7.units.r075.rich.first}{0.01} & \raggedright yes \tabularnewline
\raggedright \resultnum{v4x.comparisons.by_r.0-8.r}{0.8} & \raggedright \resultnum{v4x.comparisons.by_r.0-8.main_curve.first.point}{0.1} & \raggedright gap-censored at $a=0.9$ & \raggedright \resultnum{v4x.x7.units.r080.x7.first}{0.1} & \raggedright \resultnum{v4x.x7.units.r080.rich.first}{0.01} & \raggedright yes \tabularnewline
\raggedright \resultnum{v4x.comparisons.by_r.0-85.r}{0.85} & \raggedright \resultnum{v4x.comparisons.by_r.0-85.main_curve.first.point}{0.1} & \raggedright not attained on grid & \raggedright \resultnum{v4x.x7.units.r085.x7.first}{0.1} & \raggedright \resultnum{v4x.x7.units.r085.rich.first}{0.01} & \raggedright yes \tabularnewline
\raggedright \resultnum{v4x.comparisons.by_r.0-9.r}{0.9} & \raggedright \resultnum{v4x.comparisons.by_r.0-9.main_curve.first.point}{0.1} & \raggedright \resultnum{v4x.comparisons.by_r.0-9.x1b.first.point}{0.3} & \raggedright \resultnum{v4x.x7.units.r090.x7.first}{0.1} & \raggedright \resultnum{v4x.x7.units.r090.rich.first}{0.003} & \raggedright yes \tabularnewline
\botrule
\end{tabular}
\end{table*}

% Generated by v4_exploratory_export; edited during the journal revision
% (agent S, 2026-10-13): captions, headers and categorical labels translated
% into plain English. Every number keeps its \resultnum binding.
\begin{table*}[htbp]
\setlength{\tabcolsep}{2pt}
\caption{Channel-misspecification stress test (X5) beside the released curve against B. Conditions at zero influence strength: a linear trend in the signal share across waves with amplitude 0.03 or 0.08 (machine labels trend-003, trend-008), a step in the signal share at wave 5 (step-003, step-008), panel conditioning in the panels only (tenure-010, tenure-020) and attrition that depends on the latent state, with retention odds ratio 1.5 or 2 per unit of the household's mean latent state, on the individual latent standard-deviation scale (mnar-ln15, mnar-ln2). Equivalent strength: the first grid strength whose confirmatory point AUC equals or exceeds the misspecification AUC, without interpolation; a gap before it censors the reading (gap column; none means no gap). Relation: position of the misspecification interval relative to the confirmatory band at the released first boundary (overlap includes a shared endpoint). Not estimable: the X5 cell was a fit gap (the step and conditioning conditions at every $r$).}\label{tab:v4x-stress-comparisons}
\centering\footnotesize
\begin{tabular}{@{}p{0.070\textwidth}p{0.270\textwidth}p{0.150\textwidth}p{0.170\textwidth}p{0.110\textwidth}p{0.130\textwidth}@{}}
\toprule
\raggedright $r$ & \raggedright \mbox{Condition} & \raggedright \mbox{Equivalent} \mbox{strength} & \raggedright \mbox{State} & \raggedright \mbox{Gap} \mbox{at} \mbox{$a$} & \raggedright \mbox{Relation} \tabularnewline
\midrule
\raggedright \resultnum{v4x.comparisons.by_r.0-5.r}{0.5} & \raggedright signal-share trend, 0.03 & \raggedright \resultnum{v4x.comparisons.by_r.0-5.x5.r050.a0000.trend-003.equivalent_strength.point}{0.3} & \raggedright available & \raggedright none & \raggedright below \tabularnewline
\raggedright \resultnum{v4x.comparisons.by_r.0-5.r}{0.5} & \raggedright signal-share trend, 0.08 & \raggedright \resultnum{v4x.comparisons.by_r.0-5.x5.r050.a0000.trend-008.equivalent_strength.point}{0.6} & \raggedright available & \raggedright none & \raggedright overlaps \tabularnewline
\raggedright \resultnum{v4x.comparisons.by_r.0-5.r}{0.5} & \raggedright wave-5 step, 0.03 & \raggedright not estimable & \raggedright X5 fit gap & \raggedright  & \raggedright  \tabularnewline
\raggedright \resultnum{v4x.comparisons.by_r.0-5.r}{0.5} & \raggedright wave-5 step, 0.08 & \raggedright not estimable & \raggedright X5 fit gap & \raggedright  & \raggedright  \tabularnewline
\raggedright \resultnum{v4x.comparisons.by_r.0-5.r}{0.5} & \raggedright conditioning, 0.10 & \raggedright not estimable & \raggedright X5 fit gap & \raggedright  & \raggedright  \tabularnewline
\raggedright \resultnum{v4x.comparisons.by_r.0-5.r}{0.5} & \raggedright conditioning, 0.20 & \raggedright not estimable & \raggedright X5 fit gap & \raggedright  & \raggedright  \tabularnewline
\raggedright \resultnum{v4x.comparisons.by_r.0-5.r}{0.5} & \raggedright attrition, odds ratio 1.5 & \raggedright \resultnum{v4x.comparisons.by_r.0-5.x5.r050.a0000.mnar-ln15.equivalent_strength.point}{0.1} & \raggedright available & \raggedright none & \raggedright below \tabularnewline
\raggedright \resultnum{v4x.comparisons.by_r.0-5.r}{0.5} & \raggedright attrition, odds ratio 2 & \raggedright \resultnum{v4x.comparisons.by_r.0-5.x5.r050.a0000.mnar-ln2.equivalent_strength.point}{0.3} & \raggedright available & \raggedright none & \raggedright overlaps \tabularnewline
\raggedright \resultnum{v4x.comparisons.by_r.0-55.r}{0.55} & \raggedright signal-share trend, 0.03 & \raggedright \resultnum{v4x.comparisons.by_r.0-55.x5.r055.a0000.trend-003.equivalent_strength.point}{0.3} & \raggedright available & \raggedright none & \raggedright overlaps \tabularnewline
\raggedright \resultnum{v4x.comparisons.by_r.0-55.r}{0.55} & \raggedright signal-share trend, 0.08 & \raggedright gap-censored & \raggedright censored at gap & \raggedright \resultnum{v4x.comparisons.by_r.0-55.x5.r055.a0000.trend-008.equivalent_strength.censored_at}{0.6} & \raggedright above \tabularnewline
\raggedright \resultnum{v4x.comparisons.by_r.0-55.r}{0.55} & \raggedright wave-5 step, 0.03 & \raggedright not estimable & \raggedright X5 fit gap & \raggedright  & \raggedright  \tabularnewline
\raggedright \resultnum{v4x.comparisons.by_r.0-55.r}{0.55} & \raggedright wave-5 step, 0.08 & \raggedright not estimable & \raggedright X5 fit gap & \raggedright  & \raggedright  \tabularnewline
\raggedright \resultnum{v4x.comparisons.by_r.0-55.r}{0.55} & \raggedright conditioning, 0.10 & \raggedright not estimable & \raggedright X5 fit gap & \raggedright  & \raggedright  \tabularnewline
\raggedright \resultnum{v4x.comparisons.by_r.0-55.r}{0.55} & \raggedright conditioning, 0.20 & \raggedright not estimable & \raggedright X5 fit gap & \raggedright  & \raggedright  \tabularnewline
\raggedright \resultnum{v4x.comparisons.by_r.0-55.r}{0.55} & \raggedright attrition, odds ratio 1.5 & \raggedright \resultnum{v4x.comparisons.by_r.0-55.x5.r055.a0000.mnar-ln15.equivalent_strength.point}{0.1} & \raggedright available & \raggedright none & \raggedright overlaps \tabularnewline
\raggedright \resultnum{v4x.comparisons.by_r.0-55.r}{0.55} & \raggedright attrition, odds ratio 2 & \raggedright \resultnum{v4x.comparisons.by_r.0-55.x5.r055.a0000.mnar-ln2.equivalent_strength.point}{0.3} & \raggedright available & \raggedright none & \raggedright above \tabularnewline
\raggedright \resultnum{v4x.comparisons.by_r.0-6.r}{0.6} & \raggedright signal-share trend, 0.03 & \raggedright \resultnum{v4x.comparisons.by_r.0-6.x5.r060.a0000.trend-003.equivalent_strength.point}{0.1} & \raggedright available & \raggedright none & \raggedright overlaps \tabularnewline
\raggedright \resultnum{v4x.comparisons.by_r.0-6.r}{0.6} & \raggedright signal-share trend, 0.08 & \raggedright \resultnum{v4x.comparisons.by_r.0-6.x5.r060.a0000.trend-008.equivalent_strength.point}{0.3} & \raggedright available & \raggedright none & \raggedright above \tabularnewline
\raggedright \resultnum{v4x.comparisons.by_r.0-6.r}{0.6} & \raggedright wave-5 step, 0.03 & \raggedright not estimable & \raggedright X5 fit gap & \raggedright  & \raggedright  \tabularnewline
\raggedright \resultnum{v4x.comparisons.by_r.0-6.r}{0.6} & \raggedright wave-5 step, 0.08 & \raggedright not estimable & \raggedright X5 fit gap & \raggedright  & \raggedright  \tabularnewline
\raggedright \resultnum{v4x.comparisons.by_r.0-6.r}{0.6} & \raggedright conditioning, 0.10 & \raggedright not estimable & \raggedright X5 fit gap & \raggedright  & \raggedright  \tabularnewline
\raggedright \resultnum{v4x.comparisons.by_r.0-6.r}{0.6} & \raggedright conditioning, 0.20 & \raggedright not estimable & \raggedright X5 fit gap & \raggedright  & \raggedright  \tabularnewline
\raggedright \resultnum{v4x.comparisons.by_r.0-6.r}{0.6} & \raggedright attrition, odds ratio 1.5 & \raggedright \resultnum{v4x.comparisons.by_r.0-6.x5.r060.a0000.mnar-ln15.equivalent_strength.point}{0.1} & \raggedright available & \raggedright none & \raggedright below \tabularnewline
\raggedright \resultnum{v4x.comparisons.by_r.0-6.r}{0.6} & \raggedright attrition, odds ratio 2 & \raggedright \resultnum{v4x.comparisons.by_r.0-6.x5.r060.a0000.mnar-ln2.equivalent_strength.point}{0.3} & \raggedright available & \raggedright none & \raggedright above \tabularnewline
\raggedright \resultnum{v4x.comparisons.by_r.0-65.r}{0.65} & \raggedright signal-share trend, 0.03 & \raggedright \resultnum{v4x.comparisons.by_r.0-65.x5.r065.a0000.trend-003.equivalent_strength.point}{0.1} & \raggedright available & \raggedright none & \raggedright overlaps \tabularnewline
\raggedright \resultnum{v4x.comparisons.by_r.0-65.r}{0.65} & \raggedright signal-share trend, 0.08 & \raggedright \resultnum{v4x.comparisons.by_r.0-65.x5.r065.a0000.trend-008.equivalent_strength.point}{0.3} & \raggedright available & \raggedright none & \raggedright above \tabularnewline
\raggedright \resultnum{v4x.comparisons.by_r.0-65.r}{0.65} & \raggedright wave-5 step, 0.03 & \raggedright not estimable & \raggedright X5 fit gap & \raggedright  & \raggedright  \tabularnewline
\raggedright \resultnum{v4x.comparisons.by_r.0-65.r}{0.65} & \raggedright wave-5 step, 0.08 & \raggedright not estimable & \raggedright X5 fit gap & \raggedright  & \raggedright  \tabularnewline
\raggedright \resultnum{v4x.comparisons.by_r.0-65.r}{0.65} & \raggedright conditioning, 0.10 & \raggedright not estimable & \raggedright X5 fit gap & \raggedright  & \raggedright  \tabularnewline
\raggedright \resultnum{v4x.comparisons.by_r.0-65.r}{0.65} & \raggedright conditioning, 0.20 & \raggedright not estimable & \raggedright X5 fit gap & \raggedright  & \raggedright  \tabularnewline
\raggedright \resultnum{v4x.comparisons.by_r.0-65.r}{0.65} & \raggedright attrition, odds ratio 1.5 & \raggedright \resultnum{v4x.comparisons.by_r.0-65.x5.r065.a0000.mnar-ln15.equivalent_strength.point}{0.1} & \raggedright available & \raggedright none & \raggedright below \tabularnewline
\raggedright \resultnum{v4x.comparisons.by_r.0-65.r}{0.65} & \raggedright attrition, odds ratio 2 & \raggedright \resultnum{v4x.comparisons.by_r.0-65.x5.r065.a0000.mnar-ln2.equivalent_strength.point}{0.3} & \raggedright available & \raggedright none & \raggedright overlaps \tabularnewline
\raggedright \resultnum{v4x.comparisons.by_r.0-7.r}{0.7} & \raggedright signal-share trend, 0.03 & \raggedright \resultnum{v4x.comparisons.by_r.0-7.x5.r070.a0000.trend-003.equivalent_strength.point}{0.1} & \raggedright available & \raggedright none & \raggedright below \tabularnewline
\raggedright \resultnum{v4x.comparisons.by_r.0-7.r}{0.7} & \raggedright signal-share trend, 0.08 & \raggedright \resultnum{v4x.comparisons.by_r.0-7.x5.r070.a0000.trend-008.equivalent_strength.point}{0.3} & \raggedright available & \raggedright none & \raggedright above \tabularnewline
\raggedright \resultnum{v4x.comparisons.by_r.0-7.r}{0.7} & \raggedright wave-5 step, 0.03 & \raggedright not estimable & \raggedright X5 fit gap & \raggedright  & \raggedright  \tabularnewline
\raggedright \resultnum{v4x.comparisons.by_r.0-7.r}{0.7} & \raggedright wave-5 step, 0.08 & \raggedright not estimable & \raggedright X5 fit gap & \raggedright  & \raggedright  \tabularnewline
\raggedright \resultnum{v4x.comparisons.by_r.0-7.r}{0.7} & \raggedright conditioning, 0.10 & \raggedright not estimable & \raggedright X5 fit gap & \raggedright  & \raggedright  \tabularnewline
\raggedright \resultnum{v4x.comparisons.by_r.0-7.r}{0.7} & \raggedright conditioning, 0.20 & \raggedright not estimable & \raggedright X5 fit gap & \raggedright  & \raggedright  \tabularnewline
\raggedright \resultnum{v4x.comparisons.by_r.0-7.r}{0.7} & \raggedright attrition, odds ratio 1.5 & \raggedright \resultnum{v4x.comparisons.by_r.0-7.x5.r070.a0000.mnar-ln15.equivalent_strength.point}{0.03} & \raggedright available & \raggedright none & \raggedright below \tabularnewline
\raggedright \resultnum{v4x.comparisons.by_r.0-7.r}{0.7} & \raggedright attrition, odds ratio 2 & \raggedright \resultnum{v4x.comparisons.by_r.0-7.x5.r070.a0000.mnar-ln2.equivalent_strength.point}{0.1} & \raggedright available & \raggedright none & \raggedright overlaps \tabularnewline
\raggedright \resultnum{v4x.comparisons.by_r.0-75.r}{0.75} & \raggedright signal-share trend, 0.03 & \raggedright \resultnum{v4x.comparisons.by_r.0-75.x5.r075.a0000.trend-003.equivalent_strength.point}{0.1} & \raggedright available & \raggedright none & \raggedright below \tabularnewline
\raggedright \resultnum{v4x.comparisons.by_r.0-75.r}{0.75} & \raggedright signal-share trend, 0.08 & \raggedright \resultnum{v4x.comparisons.by_r.0-75.x5.r075.a0000.trend-008.equivalent_strength.point}{0.3} & \raggedright available & \raggedright none & \raggedright above \tabularnewline
\raggedright \resultnum{v4x.comparisons.by_r.0-75.r}{0.75} & \raggedright wave-5 step, 0.03 & \raggedright not estimable & \raggedright X5 fit gap & \raggedright  & \raggedright  \tabularnewline
\raggedright \resultnum{v4x.comparisons.by_r.0-75.r}{0.75} & \raggedright wave-5 step, 0.08 & \raggedright not estimable & \raggedright X5 fit gap & \raggedright  & \raggedright  \tabularnewline
\raggedright \resultnum{v4x.comparisons.by_r.0-75.r}{0.75} & \raggedright conditioning, 0.10 & \raggedright not estimable & \raggedright X5 fit gap & \raggedright  & \raggedright  \tabularnewline
\raggedright \resultnum{v4x.comparisons.by_r.0-75.r}{0.75} & \raggedright conditioning, 0.20 & \raggedright not estimable & \raggedright X5 fit gap & \raggedright  & \raggedright  \tabularnewline
\raggedright \resultnum{v4x.comparisons.by_r.0-75.r}{0.75} & \raggedright attrition, odds ratio 1.5 & \raggedright \resultnum{v4x.comparisons.by_r.0-75.x5.r075.a0000.mnar-ln15.equivalent_strength.point}{0.03} & \raggedright available & \raggedright none & \raggedright below \tabularnewline
\raggedright \resultnum{v4x.comparisons.by_r.0-75.r}{0.75} & \raggedright attrition, odds ratio 2 & \raggedright \resultnum{v4x.comparisons.by_r.0-75.x5.r075.a0000.mnar-ln2.equivalent_strength.point}{0.1} & \raggedright available & \raggedright none & \raggedright overlaps \tabularnewline
\raggedright \resultnum{v4x.comparisons.by_r.0-8.r}{0.8} & \raggedright signal-share trend, 0.03 & \raggedright \resultnum{v4x.comparisons.by_r.0-8.x5.r080.a0000.trend-003.equivalent_strength.point}{0.1} & \raggedright available & \raggedright none & \raggedright below \tabularnewline
\raggedright \resultnum{v4x.comparisons.by_r.0-8.r}{0.8} & \raggedright signal-share trend, 0.08 & \raggedright \resultnum{v4x.comparisons.by_r.0-8.x5.r080.a0000.trend-008.equivalent_strength.point}{0.3} & \raggedright available & \raggedright none & \raggedright above \tabularnewline
\botrule
\end{tabular}
\end{table*}
\begin{table*}[htbp]
\setlength{\tabcolsep}{2pt}
\addtocounter{table}{-\csname @ne\endcsname}
\providecommand{\theHtable}{\thetable}
\renewcommand{\theHtable}{v4x-stress-comparisons-part1}
\caption{Channel-misspecification stress test (X5) beside the released curve (continued, part 2 of 2)}\label{tab:v4x-stress-comparisons-part1}
\centering\footnotesize
\begin{tabular}{@{}p{0.070\textwidth}p{0.270\textwidth}p{0.150\textwidth}p{0.170\textwidth}p{0.110\textwidth}p{0.130\textwidth}@{}}
\toprule
\raggedright $r$ & \raggedright \mbox{Condition} & \raggedright \mbox{Equivalent} \mbox{strength} & \raggedright \mbox{State} & \raggedright \mbox{Gap} \mbox{at} \mbox{$a$} & \raggedright \mbox{Relation} \tabularnewline
\midrule
\raggedright \resultnum{v4x.comparisons.by_r.0-8.r}{0.8} & \raggedright wave-5 step, 0.03 & \raggedright not estimable & \raggedright X5 fit gap & \raggedright  & \raggedright  \tabularnewline
\raggedright \resultnum{v4x.comparisons.by_r.0-8.r}{0.8} & \raggedright wave-5 step, 0.08 & \raggedright not estimable & \raggedright X5 fit gap & \raggedright  & \raggedright  \tabularnewline
\raggedright \resultnum{v4x.comparisons.by_r.0-8.r}{0.8} & \raggedright conditioning, 0.10 & \raggedright not estimable & \raggedright X5 fit gap & \raggedright  & \raggedright  \tabularnewline
\raggedright \resultnum{v4x.comparisons.by_r.0-8.r}{0.8} & \raggedright conditioning, 0.20 & \raggedright not estimable & \raggedright X5 fit gap & \raggedright  & \raggedright  \tabularnewline
\raggedright \resultnum{v4x.comparisons.by_r.0-8.r}{0.8} & \raggedright attrition, odds ratio 1.5 & \raggedright \resultnum{v4x.comparisons.by_r.0-8.x5.r080.a0000.mnar-ln15.equivalent_strength.point}{0.03} & \raggedright available & \raggedright none & \raggedright below \tabularnewline
\raggedright \resultnum{v4x.comparisons.by_r.0-8.r}{0.8} & \raggedright attrition, odds ratio 2 & \raggedright \resultnum{v4x.comparisons.by_r.0-8.x5.r080.a0000.mnar-ln2.equivalent_strength.point}{0.1} & \raggedright available & \raggedright none & \raggedright below \tabularnewline
\raggedright \resultnum{v4x.comparisons.by_r.0-85.r}{0.85} & \raggedright signal-share trend, 0.03 & \raggedright \resultnum{v4x.comparisons.by_r.0-85.x5.r085.a0000.trend-003.equivalent_strength.point}{0.1} & \raggedright available & \raggedright none & \raggedright below \tabularnewline
\raggedright \resultnum{v4x.comparisons.by_r.0-85.r}{0.85} & \raggedright signal-share trend, 0.08 & \raggedright \resultnum{v4x.comparisons.by_r.0-85.x5.r085.a0000.trend-008.equivalent_strength.point}{0.3} & \raggedright available & \raggedright none & \raggedright above \tabularnewline
\raggedright \resultnum{v4x.comparisons.by_r.0-85.r}{0.85} & \raggedright wave-5 step, 0.03 & \raggedright not estimable & \raggedright X5 fit gap & \raggedright  & \raggedright  \tabularnewline
\raggedright \resultnum{v4x.comparisons.by_r.0-85.r}{0.85} & \raggedright wave-5 step, 0.08 & \raggedright not estimable & \raggedright X5 fit gap & \raggedright  & \raggedright  \tabularnewline
\raggedright \resultnum{v4x.comparisons.by_r.0-85.r}{0.85} & \raggedright conditioning, 0.10 & \raggedright not estimable & \raggedright X5 fit gap & \raggedright  & \raggedright  \tabularnewline
\raggedright \resultnum{v4x.comparisons.by_r.0-85.r}{0.85} & \raggedright conditioning, 0.20 & \raggedright not estimable & \raggedright X5 fit gap & \raggedright  & \raggedright  \tabularnewline
\raggedright \resultnum{v4x.comparisons.by_r.0-85.r}{0.85} & \raggedright attrition, odds ratio 1.5 & \raggedright \resultnum{v4x.comparisons.by_r.0-85.x5.r085.a0000.mnar-ln15.equivalent_strength.point}{0.01} & \raggedright available & \raggedright none & \raggedright below \tabularnewline
\raggedright \resultnum{v4x.comparisons.by_r.0-85.r}{0.85} & \raggedright attrition, odds ratio 2 & \raggedright \resultnum{v4x.comparisons.by_r.0-85.x5.r085.a0000.mnar-ln2.equivalent_strength.point}{0.03} & \raggedright available & \raggedright none & \raggedright below \tabularnewline
\raggedright \resultnum{v4x.comparisons.by_r.0-9.r}{0.9} & \raggedright signal-share trend, 0.03 & \raggedright \resultnum{v4x.comparisons.by_r.0-9.x5.r090.a0000.trend-003.equivalent_strength.point}{0.1} & \raggedright available & \raggedright none & \raggedright below \tabularnewline
\raggedright \resultnum{v4x.comparisons.by_r.0-9.r}{0.9} & \raggedright signal-share trend, 0.08 & \raggedright \resultnum{v4x.comparisons.by_r.0-9.x5.r090.a0000.trend-008.equivalent_strength.point}{0.3} & \raggedright available & \raggedright none & \raggedright above \tabularnewline
\raggedright \resultnum{v4x.comparisons.by_r.0-9.r}{0.9} & \raggedright wave-5 step, 0.03 & \raggedright not estimable & \raggedright X5 fit gap & \raggedright  & \raggedright  \tabularnewline
\raggedright \resultnum{v4x.comparisons.by_r.0-9.r}{0.9} & \raggedright wave-5 step, 0.08 & \raggedright not estimable & \raggedright X5 fit gap & \raggedright  & \raggedright  \tabularnewline
\raggedright \resultnum{v4x.comparisons.by_r.0-9.r}{0.9} & \raggedright conditioning, 0.10 & \raggedright not estimable & \raggedright X5 fit gap & \raggedright  & \raggedright  \tabularnewline
\raggedright \resultnum{v4x.comparisons.by_r.0-9.r}{0.9} & \raggedright conditioning, 0.20 & \raggedright not estimable & \raggedright X5 fit gap & \raggedright  & \raggedright  \tabularnewline
\raggedright \resultnum{v4x.comparisons.by_r.0-9.r}{0.9} & \raggedright attrition, odds ratio 1.5 & \raggedright \resultnum{v4x.comparisons.by_r.0-9.x5.r090.a0000.mnar-ln15.equivalent_strength.point}{0.01} & \raggedright available & \raggedright none & \raggedright below \tabularnewline
\raggedright \resultnum{v4x.comparisons.by_r.0-9.r}{0.9} & \raggedright attrition, odds ratio 2 & \raggedright \resultnum{v4x.comparisons.by_r.0-9.x5.r090.a0000.mnar-ln2.equivalent_strength.point}{0.03} & \raggedright available & \raggedright none & \raggedright below \tabularnewline
\botrule
\end{tabular}
\end{table*}

% Generated by v4_exploratory_export; edited during the journal revision
% (agent S, 2026-10-13): captions, headers and categorical labels translated
% into plain English. Every number keeps its \resultnum binding.
\begin{table*}[htbp]
\setlength{\tabcolsep}{2pt}
\caption{Mean-field twin against B (X1b): first and sustained boundaries of the twin, with boundary intervals, beside the released boundaries against B, for context only and without a between-curve test. All nine twin units and all nine main units were complete and valid, and the main curve was valid. Gap-censored at $a$: the boundary (or an interval end, censored) cannot be declared across the fit gap at that strength; not attained on grid: the point AUC stayed below 0.70 on a complete grid; $>0.9$: the interval end was not reached on the complete grid.}\label{tab:v4x-x1b-boundaries}
\centering\footnotesize
\begin{tabular}{@{}p{0.070\textwidth}p{0.240\textwidth}p{0.240\textwidth}p{0.190\textwidth}p{0.190\textwidth}@{}}
\toprule
\raggedright $r$ & \raggedright \mbox{Twin,} \mbox{first} & \raggedright \mbox{Twin,} \mbox{sustained} & \raggedright \mbox{Released,} \mbox{first} & \raggedright \mbox{Released,} \mbox{sustained} \tabularnewline
\midrule
\raggedright \resultnum{v4x.comparisons.by_r.0-5.r}{0.5} & \raggedright gap-censored at $a=0.6$ & \raggedright gap-censored at $a=0.6$ & \raggedright \resultnum{v4x.comparisons.by_r.0-5.main_curve.first.point}{0.3} [\resultnum{v4x.comparisons.by_r.0-5.main_curve.first.interval.i000}{0.1}, \resultnum{v4x.comparisons.by_r.0-5.main_curve.first.interval.i001}{0.3}] & \raggedright \resultnum{v4x.comparisons.by_r.0-5.main_curve.sustained.point}{0.3} [\resultnum{v4x.comparisons.by_r.0-5.main_curve.sustained.interval.i000}{0.1}, \resultnum{v4x.comparisons.by_r.0-5.main_curve.sustained.interval.i001}{0.3}] \tabularnewline
\raggedright \resultnum{v4x.comparisons.by_r.0-55.r}{0.55} & \raggedright gap-censored at $a=0.9$ & \raggedright gap-censored at $a=0.9$ & \raggedright \resultnum{v4x.comparisons.by_r.0-55.main_curve.first.point}{0.1} [\resultnum{v4x.comparisons.by_r.0-55.main_curve.first.interval.i000}{0.1}, \resultnum{v4x.comparisons.by_r.0-55.main_curve.first.interval.i001}{0.3}] & \raggedright gap-censored at $a=0.9$ \tabularnewline
\raggedright \resultnum{v4x.comparisons.by_r.0-6.r}{0.6} & \raggedright not attained on grid & \raggedright not attained on grid & \raggedright \resultnum{v4x.comparisons.by_r.0-6.main_curve.first.point}{0.1} [\resultnum{v4x.comparisons.by_r.0-6.main_curve.first.interval.i000}{0.1}, \resultnum{v4x.comparisons.by_r.0-6.main_curve.first.interval.i001}{0.3}] & \raggedright \resultnum{v4x.comparisons.by_r.0-6.main_curve.sustained.point}{0.1} [\resultnum{v4x.comparisons.by_r.0-6.main_curve.sustained.interval.i000}{0.1}, \resultnum{v4x.comparisons.by_r.0-6.main_curve.sustained.interval.i001}{0.3}] \tabularnewline
\raggedright \resultnum{v4x.comparisons.by_r.0-65.r}{0.65} & \raggedright not attained on grid & \raggedright not attained on grid & \raggedright \resultnum{v4x.comparisons.by_r.0-65.main_curve.first.point}{0.1} [\resultnum{v4x.comparisons.by_r.0-65.main_curve.first.interval.i000}{0.1}, \resultnum{v4x.comparisons.by_r.0-65.main_curve.first.interval.i001}{0.1}] & \raggedright \resultnum{v4x.comparisons.by_r.0-65.main_curve.sustained.point}{0.1} [\resultnum{v4x.comparisons.by_r.0-65.main_curve.sustained.interval.i000}{0.1}, \resultnum{v4x.comparisons.by_r.0-65.main_curve.sustained.interval.i001}{0.1}] \tabularnewline
\raggedright \resultnum{v4x.comparisons.by_r.0-7.r}{0.7} & \raggedright not attained on grid & \raggedright not attained on grid & \raggedright \resultnum{v4x.comparisons.by_r.0-7.main_curve.first.point}{0.1} [\resultnum{v4x.comparisons.by_r.0-7.main_curve.first.interval.i000}{0.1}, \resultnum{v4x.comparisons.by_r.0-7.main_curve.first.interval.i001}{0.1}] & \raggedright gap-censored at $a=0.9$ \tabularnewline
\raggedright \resultnum{v4x.comparisons.by_r.0-75.r}{0.75} & \raggedright gap-censored at $a=0.9$; interval [\resultnum{v4x.comparisons.by_r.0-75.x1b.first.interval.i000}{0.3}, censored] & \raggedright gap-censored at $a=0.9$ & \raggedright \resultnum{v4x.comparisons.by_r.0-75.main_curve.first.point}{0.1} [\resultnum{v4x.comparisons.by_r.0-75.main_curve.first.interval.i000}{0.1}, \resultnum{v4x.comparisons.by_r.0-75.main_curve.first.interval.i001}{0.1}] & \raggedright \resultnum{v4x.comparisons.by_r.0-75.main_curve.sustained.point}{0.1} [\resultnum{v4x.comparisons.by_r.0-75.main_curve.sustained.interval.i000}{0.1}, \resultnum{v4x.comparisons.by_r.0-75.main_curve.sustained.interval.i001}{0.1}] \tabularnewline
\raggedright \resultnum{v4x.comparisons.by_r.0-8.r}{0.8} & \raggedright gap-censored at $a=0.9$; interval [\resultnum{v4x.comparisons.by_r.0-8.x1b.first.interval.i000}{0.6}, censored] & \raggedright gap-censored at $a=0.9$ & \raggedright \resultnum{v4x.comparisons.by_r.0-8.main_curve.first.point}{0.1} [\resultnum{v4x.comparisons.by_r.0-8.main_curve.first.interval.i000}{0.1}, \resultnum{v4x.comparisons.by_r.0-8.main_curve.first.interval.i001}{0.1}] & \raggedright \resultnum{v4x.comparisons.by_r.0-8.main_curve.sustained.point}{0.1} [\resultnum{v4x.comparisons.by_r.0-8.main_curve.sustained.interval.i000}{0.1}, \resultnum{v4x.comparisons.by_r.0-8.main_curve.sustained.interval.i001}{0.1}] \tabularnewline
\raggedright \resultnum{v4x.comparisons.by_r.0-85.r}{0.85} & \raggedright not attained on grid; interval [\resultnum{v4x.comparisons.by_r.0-85.x1b.first.interval.i000}{0.3}, $>0.9$] & \raggedright not attained on grid; interval [\resultnum{v4x.comparisons.by_r.0-85.x1b.sustained.interval.i000}{0.3}, $>0.9$] & \raggedright \resultnum{v4x.comparisons.by_r.0-85.main_curve.first.point}{0.1} [\resultnum{v4x.comparisons.by_r.0-85.main_curve.first.interval.i000}{0.1}, \resultnum{v4x.comparisons.by_r.0-85.main_curve.first.interval.i001}{0.1}] & \raggedright \resultnum{v4x.comparisons.by_r.0-85.main_curve.sustained.point}{0.1} [\resultnum{v4x.comparisons.by_r.0-85.main_curve.sustained.interval.i000}{0.1}, \resultnum{v4x.comparisons.by_r.0-85.main_curve.sustained.interval.i001}{0.1}] \tabularnewline
\raggedright \resultnum{v4x.comparisons.by_r.0-9.r}{0.9} & \raggedright \resultnum{v4x.comparisons.by_r.0-9.x1b.first.point}{0.3} [\resultnum{v4x.comparisons.by_r.0-9.x1b.first.interval.i000}{0.3}, $>0.9$] & \raggedright \resultnum{v4x.comparisons.by_r.0-9.x1b.sustained.point}{0.3} [\resultnum{v4x.comparisons.by_r.0-9.x1b.sustained.interval.i000}{0.3}, $>0.9$] & \raggedright \resultnum{v4x.comparisons.by_r.0-9.main_curve.first.point}{0.1} [\resultnum{v4x.comparisons.by_r.0-9.main_curve.first.interval.i000}{0.1}, \resultnum{v4x.comparisons.by_r.0-9.main_curve.first.interval.i001}{0.1}] & \raggedright \resultnum{v4x.comparisons.by_r.0-9.main_curve.sustained.point}{0.1} [\resultnum{v4x.comparisons.by_r.0-9.main_curve.sustained.interval.i000}{0.1}, \resultnum{v4x.comparisons.by_r.0-9.main_curve.sustained.interval.i001}{0.1}] \tabularnewline
\botrule
\end{tabular}
\end{table*}


\clearpage
